% Νέα αναθεωρημένη πηγή για τον κανόνα Μ1-J, 31/10/2026.
% Η ιστορική πηγή sources/proposal-v871e-candidate.tex διατηρείται αυτούσια.
% PDF εκ νέου μεταγλώττισης: εκκρεμεί.
% Προγενέστερη αναλυτική πρόταση — μουσικολογική θεμελίωση, τεκμηρίωση πηγών και παραδείγματα ΠΕ0.
\documentclass[11pt,a4paper]{article}

\usepackage[a4paper,left=2.45cm,right=2.45cm,top=2.35cm,bottom=2.45cm,headheight=15pt]{geometry}
\usepackage{fontspec}
\usepackage{polyglossia}
\setdefaultlanguage{greek}
\setotherlanguage{english}
\setmainfont{FreeSerif.otf}[Ligatures=TeX,BoldFont=FreeSerifBold.otf,ItalicFont=FreeSerifItalic.otf,BoldItalicFont=FreeSerifBoldItalic.otf]
\newfontfamily\greekfont{FreeSerif.otf}[Ligatures=TeX,BoldFont=FreeSerifBold.otf,ItalicFont=FreeSerifItalic.otf,BoldItalicFont=FreeSerifBoldItalic.otf]
\newfontfamily\englishfont{FreeSerif.otf}[Ligatures=TeX,BoldFont=FreeSerifBold.otf,ItalicFont=FreeSerifItalic.otf,BoldItalicFont=FreeSerifBoldItalic.otf]
\newfontfamily\coverfont{FreeSerif.otf}[Ligatures=TeX,BoldFont=FreeSerifBold.otf,ItalicFont=FreeSerifItalic.otf,BoldItalicFont=FreeSerifBoldItalic.otf]
\setsansfont{FreeSans.otf}[Scale=MatchLowercase,BoldFont=FreeSansBold.otf,ItalicFont=FreeSansOblique.otf,BoldItalicFont=FreeSansBoldOblique.otf]
\setmonofont{DejaVuSansMono.ttf}[Scale=MatchLowercase,BoldFont=DejaVuSansMono-Bold.ttf,ItalicFont=DejaVuSansMono-Oblique.ttf,BoldItalicFont=DejaVuSansMono-BoldOblique.ttf]

\usepackage{microtype}
\usepackage{setspace}
\usepackage{enumitem}
\usepackage{booktabs,longtable,array,tabularx,xltabular}
\usepackage{amsmath,amssymb}
\usepackage{xcolor}
\usepackage{hyperref}
\usepackage{titlesec}
\usepackage{fancyhdr}
\usepackage{ragged2e}
\usepackage{indentfirst}
\usepackage{xurl}
\usepackage[bottom,hang,flushmargin]{footmisc}
\setlength{\footnotemargin}{1em}
\interfootnotelinepenalty=10000
\usepackage{placeins}
\usepackage{tikz}
\usepackage{graphicx}
\usetikzlibrary{arrows.meta,positioning}
\providecommand{\tightlist}{\setlength{\itemsep}{0pt}\setlength{\parskip}{0pt}}

\definecolor{acc}{HTML}{7A2E2E}
\hypersetup{colorlinks=true,linkcolor=acc,urlcolor=acc,bookmarksdepth=2,
 pdftitle={Η μουσική θέση ως υπολογιστική μονάδα της βυζαντινής μουσικής}}
\newcolumntype{Y}{>{\RaggedRight\arraybackslash}X}

\titleformat{\section}{\large\bfseries\color{acc}}{\thesection}{0.7em}{}
\titlespacing*{\section}{0pt}{1.5em}{0.55em}
\titleformat{\subsection}{\normalsize\bfseries}{\thesubsection}{0.6em}{}
\titlespacing*{\subsection}{0pt}{1.0em}{0.35em}
\titleformat{\subsubsection}{\normalsize\itshape}{\thesubsubsection}{0.6em}{}

\setlist[itemize]{leftmargin=1.5em,itemsep=0.18em,topsep=0.3em}
\setlist[enumerate]{leftmargin=1.7em,itemsep=0.18em,topsep=0.3em}
\setstretch{1.10}
\setlength{\parskip}{0.12em}
\setlength{\parindent}{1.25em}
\setlength{\LTpre}{0.5\baselineskip}
\setlength{\LTpost}{0.5\baselineskip}
\setlength{\emergencystretch}{2em}
\emergencystretch=2em

\pagestyle{fancy}\fancyhf{}
\fancyhead[L]{\small Η μουσική θέση ως υπολογιστική μονάδα}
\fancyhead[R]{\small \thepage}
\renewcommand{\headrulewidth}{0.4pt}

\begin{document}

\hypersetup{pageanchor=false}
\begin{titlepage}
\centering
\vspace*{1.0cm}
{\coverfont\large ΕΛΛΗΝΙΚΟ ΜΕΣΟΓΕΙΑΚΟ ΠΑΝΕΠΙΣΤΗΜΙΟ\par}
{\coverfont Τμήμα Μουσικής Τεχνολογίας και Ακουστικής\par}
\vspace{1.5cm}
{\coverfont\large ΕΡΕΥΝΗΤΙΚΗ ΠΡΟΤΑΣΗ\par}
\vspace{1.3cm}
{\coverfont\LARGE\bfseries Η μουσική θέση ως υπολογιστική μονάδα της βυζαντινής μουσικής\par}
\vspace{0.65cm}
{\coverfont\large Συμβολική αναπαράσταση, ιστορική μετάδοση και γνωσιακή λειτουργία\par}
\vspace{1.1cm}
{\coverfont\large\itshape The Byzantine Melodic Formula as a Computational Unit in New Method Byzantine Chant\par}
\vspace{0.3cm}
{\coverfont\normalsize\itshape Symbolic Representation, Historical Transmission, and Cognitive Function\par}
\vspace{9.8cm}
{\coverfont Γεώργιος Σπυριδάκης\par}
\vspace{0.8cm}
{\coverfont\small Οκτώβριος 2026\par}
\vfill
\end{titlepage}
\hypersetup{pageanchor=true}

\section*{Περίληψη}\addcontentsline{toc}{section}{Περίληψη}

Η βυζαντινή ψαλτική\footnote{Στο πλαίσιο της πρότασης, ο όρος δηλώνει τη μονοφωνική λειτουργική μουσική παράδοση της Ορθόδοξης Εκκλησίας, με γραπτή και προφορική μετάδοση μέσω διδασκαλίας.} αποτελεί τη μονοφωνική λειτουργική μουσική παράδοση της Ορθόδοξης Εκκλησίας και μεταδίδεται επί αιώνες μέσα από γραπτές πηγές και τη διδασκαλία από δάσκαλο σε μαθητή. Από τις αρχές του 19ου αιώνα καταγράφεται με τη Νέα Μέθοδο\footnote{Η Νέα Μέθοδος είναι το αναλυτικό σύστημα σημειογραφίας που καθιερώθηκε στις αρχές του 19ου αιώνα από τους Τρεις Διδασκάλους. Βλ. Χρύσανθος 1832 και τη σχετική ιστοριογραφία της εξήγησης.}, ένα σύστημα σχεσιακής νευματικής σημειογραφίας\footnote{Νευματική σημειογραφία: σύστημα γραφικών χαρακτήρων (νευμάτων) που κωδικοποιούν σχεσιακή μελωδική και συναφή μουσική πληροφορία. Βλ. Raasted 1966· Troelsgård 2011.}, στο οποίο η μελωδική κίνηση αποδίδεται κυρίως ως σχετική άνοδος ή κάθοδος και συνδυάζεται με χρονικά, ποιοτικά και τροπικά στοιχεία μέσα στο σύστημα των οκτώ ήχων\footnote{Ήχος: τροπική κατηγορία της ψαλτικής με χαρακτηριστική οργάνωση βαθμίδων, κινήσεων και καταλήξεων· δεν ταυτίζεται με απλή δυτική κλίμακα ούτε με τον γενικό όρο «ήχος» ως ακουστικό σήμα. Βλ. Χρύσανθος 1832.}.

Βασικό δομικό στοιχείο αυτής της μουσικής είναι οι μουσικές θέσεις\footnote{Η ιστορική θεωρία χρησιμοποιεί τη θέση ως οργανωμένο μελικό σχηματισμό ευρύτερο από το μεμονωμένο σημείο. Ο Μανουήλ Χρυσάφης τη συνδέει με την ένωση σημείων που συγκροτεί μέλος, ενώ ο Γαβριήλ ιερομόναχος την παρομοιάζει με τη λέξη της γραμματικής. Βλ. Conomos 1985b, σσ.~40--41· Hannick--Wolfram 1985, σ.~72.}. Πρόκειται για αναγνωρίσιμους μελικούς σχηματισμούς που επανέρχονται σε διαφορετικά έργα και μπορούν να μεταβάλλονται ουσιαστικά. Μια θέση δεν πραγματώνεται κατ' ανάγκην ως απαράλλακτη ακολουθία φθόγγων\footnote{Φθόγγος: βαθμίδα/τονικό σημείο στο θεωρητικό σύστημα της ψαλτικής. Στην παρούσα πρόταση δεν ταυτίζεται μηχανικά με ένα υπολογιστικό «note token».}. Μπορεί να είναι συντομότερη ή πιο αναπτυγμένη. Μπορεί επίσης να προσαρμόζεται στο τροπικό, κειμενικό και συνθετικό περιβάλλον. Η επιφανειακή ομοιότητα δεν αρκεί για να ορίσει την ταυτότητα. Η έρευνα συγκροτεί \textbf{υποψήφιες οικογένειες θέσεων} με ρητές υπογραφές τύπου. Η οικογένεια χρησιμοποιείται ως ελέγξιμη αναλυτική κατηγορία και όχι ως υπόθεση κοινής γενεαλογικής καταγωγής. Κάθε υπογραφή δηλώνει την αναλυτική κλίμακα, τις διαγνωστικές μουσικές σχέσεις, τον ρόλο τους, τους κανόνες επιτρεπτής πραγμάτωσης και τα κριτήρια αποκλεισμού. Τα κριτήρια είναι ειδικά για κάθε τύπο και μπορούν να εξαρτώνται από ρητά δηλωμένο περιβάλλον. Η επάρκειά τους ελέγχεται σε νέο υλικό από ανεξάρτητους ειδικούς πριν η οικογενειακή ένταξη χρησιμοποιηθεί ως εποπτεία στη μελέτη αναπαραστάσεων (M1) και στη σύγκριση της εξηγητικής εποπτείας (M2), που παρουσιάζονται αμέσως παρακάτω.

Η παράδοση διαθέτει επίσης ένα σπάνιο ιστορικό αρχείο. Με την καθιέρωση της Νέας Μεθόδου, κυρίως ο Γρηγόριος Πρωτοψάλτης και ο Χουρμούζιος Χαρτοφύλαξ\footnote{«Πρωτοψάλτης» και «Χαρτοφύλαξ» είναι ιστορικοί τίτλοι με τους οποίους είναι γνωστοί οι δύο από τους Τρεις Διδασκάλους της Νέας Μεθόδου· εδώ λειτουργούν ως καθιερωμένες ονομασίες προσώπων.} απέδωσαν μεγάλο μέρος του παλαιότερου ρεπερτορίου στη νέα, αναλυτικότερη γραφή. Οι αποδόσεις αυτές ονομάζονται \textbf{εξηγήσεις}\footnote{Εξήγηση: τεκμηριωμένη αναλυτική απόδοση παλαιότερης μουσικής γραφής σε νεότερη σημειογραφία. Η χρήση της ως ιστορικής σχέσης δεν συνεπάγεται από μόνη της ότι ανακατασκευάζεται το παλαιότερο άκουσμα. Βλ. Στάθης 1978· Καραγκούνης 2003· Αρβανίτης 2010· Arvanitis 2022.}. Για ορισμένα μέλη υπάρχουν περισσότερες τεκμηριωμένες εξηγήσεις με γνωστό συντάκτη και πηγή, γεγονός που επιτρέπει να εξεταστεί τι διατηρείται και τι διαφοροποιείται από τη μία απόδοση στην άλλη. Η χρήση αυτών των σχέσεων ως δεδομένων εκπαίδευσης δεν προϋποθέτει ότι οι εξηγήσεις ανακατασκευάζουν αυτούσιο το παλαιότερο άκουσμα.

Το κεντρικό ερώτημα είναι πώς αναπαρίσταται υπολογιστικά μια μουσική μονάδα όταν η ταυτότητά της δεν εξαντλείται στην επανάληψη της ίδιας ακολουθίας φθόγγων. Η ιστορική και μορφολογική έρευνα έχει τεκμηριώσει ονομασμένες θέσεις, επανεμφανίσεις και παραλλαγές τους, αλλά η συστηματικά συγκρίσιμη διερεύνηση της ταυτότητας και της εσωτερικής τους μεταβλητότητας παραμένει αυτοτελές πρόβλημα. Η ίδια θέση μπορεί να αποκτά αισθητά διαφορετική πραγμάτωση, ενώ μουσικολογικά διαφορετικές θέσεις μπορεί να έχουν συγγενή επιφανειακή μορφή. Η βυζαντινή ψαλτική προσφέρει για αυτόν τον έλεγχο έναν σπάνιο συνδυασμό δεδομένων: τεκμηριωμένες πραγματώσεις, ιστορικές σχέσεις μεταξύ συγκεκριμένων αποδόσεων, πλήρη έργα και ακροατές με διαφορετικό βαθμό εμπειρίας στη συγκεκριμένη μουσική γλώσσα.

Βασικό καταλογογραφικό υλικό είναι η συλλογή \emph{Papadic Formulae} της Ιεράς Μονής Αγίου Αντωνίου (St.~Anthony's Greek Orthodox Monastery), η οποία συγκεντρώνει μελωδικούς τύπους του αργού, μελισματικού\footnote{Μελισματικό: μελωδική ανάπτυξη στην οποία μία συλλαβή εκτείνεται σε πολλούς φθόγγους ή μελωδικά γεγονότα.} ρεπερτορίου στο οποίο ανήκουν τα Χερουβικά\footnote{Χερουβικό: μελοποιημένη μορφή του Χερουβικού Ύμνου της Θείας Λειτουργίας και βασικό είδος του παπαδικού ρεπερτορίου. Βλ. Conomos 1974· Καραγκούνης 2003/2005.} και τα Κοινωνικά\footnote{Κοινωνικό: λειτουργικό μέλος που ψάλλεται κατά τη Θεία Κοινωνία, ιδίως κατά την κοινωνία των λειτουργών· στο εξεταζόμενο παπαδικό ρεπερτόριο παρουσιάζει εκτενή μελισματική ανάπτυξη. Βλ. Conomos 1985a και τη σχετική βιβλιογραφία των Κοινωνικών.}. Η συλλογή προέρχεται από 758 συνθέσεις και καταγράφει για κάθε τύπο το έργο και τον συνθέτη από τον οποίο προέρχεται. Οι καταλογογραφημένοι τύποι παρέχουν υποψήφιες σχέσεις, όχι έτοιμη τελική ταξινόμηση. Τα πλήρη έργα, οι τεκμηριωμένες εξηγήσεις και ηχογραφήσεις επαγγελματιών ψαλτών σε κοινό ρεπερτόριο συμπληρώνουν τα δεδομένα.

Η πρώτη μεθοδολογική ενότητα, \textbf{Μελέτη 1 (M1: σύγκριση συμβολικών αναπαραστάσεων, δηλαδή διαφορετικών υπολογιστικών περιγραφών του ίδιου μέλους)}, εξετάζει ποια κωδικοποίηση εισόδου διατηρεί καλύτερα την πληροφορία που απαιτεί η αναγνώριση θέσεων (§~\ref{sec:m1}). Συγκρίνονται το γραφικό σημείο, το ατομικό μελωδικό γεγονός και η σύνθετη νευματική μονάδα, με διαδοχικά εμπλουτισμένα σύνολα σημειογραφικών, χρονικών, τροπικών και κειμενικών γνωρισμάτων. Σε χωριστές δομικές συγκρίσεις, η ρητή αναπαράσταση των θέσεων αντιπαραβάλλεται με μοντέλο σε επίπεδο γεγονότων και με αλγοριθμικά επαγόμενα τμήματα. Η επιλογή βασίζεται στην οριοθέτηση και στην οικογενειακή αναγνώριση σε έργα και πηγές εκτός εκπαίδευσης. Το \textbf{ΠΕ0 (Πακέτο Εργασίας 0)} είναι η αρχική απογραφή και μουσικολογική δοκιμή επάρκειας των διαθέσιμων πηγών, υπογραφών και επισημειώσεων· καθορίζει εκ των προτέρων ποιες συγκρίσεις μπορούν να πραγματοποιηθούν με αξιόπιστα δεδομένα.

Η \textbf{Μελέτη 2 (M2: ελεγχόμενη σύγκριση εξηγητικής εποπτείας, §~\ref{sec:eise})} εξετάζει τι πληροφορία προσφέρει στην εκπαίδευση η τεκμηριωμένη ιστορική σχέση δύο εξηγήσεων του ίδιου χωρίου. Ο \emph{Exegesis-Invariant Segmental Encoder} (\textbf{EISE}, «τμηματικός κωδικοποιητής αναλλοιωτότητας ως προς τις εξηγητικές πραγματώσεις») αναπαριστά τμήματα μεταβλητού μήκους, διακρίνοντας οικογενειακή ταυτότητα και συγκεκριμένη πραγμάτωση. Η σύγκριση οργανώνεται σε τέσσερις συνθήκες: \textbf{B0}, άμεση συμβολική περιγραφή χωρίς εκπαιδευμένη νευρωνική αναπαράσταση· \textbf{B1}, κοινή νευρωνική υποδομή χωρίς τον πρόσθετο σχεσιακό όρο· \textbf{B2}, ίδια υποδομή με θετικά ζεύγη ανεξάρτητα επικυρωμένων μελών της ίδιας οικογένειας· και \textbf{EISE}, θετικά ζεύγη που επιπλέον τεκμηριώνονται από συγκεκριμένη ιστορική συν-εξήγηση. Η κύρια αντίθεση B2--EISE διατηρεί αντιστοιχισμένη αρχιτεκτονική, οικογενειακή πληροφορία, έκθεση και υπολογιστικό προϋπολογισμό· εξετάζει τη διαφορά επίδοσης στο ανεξάρτητο κριτήριο εκτελεστικής αξιολόγησης, όχι την ανακατασκευή ενός μοναδικού ιστορικού ακούσματος. Η κύρια ποσότητα $\Delta_{\mathrm{hist}}=S_{\mathrm{EISE}}-S_{B2}$ είναι η διαφορά συμφωνίας των δύο μοντέλων με το ίδιο εξωτερικό εκτελεστικό κριτήριο. Τα ιστορικά θετικά ζεύγη απαιτούν ανεξάρτητη επαλήθευση του κοινού χωρίου και της οικογενειακής ένταξης, ενώ το ΠΕ0 ελέγχει την επάρκεια του κοινού επιλέξιμου πεδίου. Μικρή ή μη διακριτή διαφορά $\Delta_{\mathrm{hist}}$ δεν ισοδυναμεί με απόδειξη απουσίας ιστορικής πληροφορίας.

Η ανεξάρτητη μουσικολογική επικύρωση πραγματοποιείται στη M3a (§~\ref{sec:m3a}). Οι ανθρώπινες δοκιμασίες περιλαμβάνουν έμπειρους ψάλτες και μουσικούς χωρίς βυζαντινή εκπαίδευση. Στη δοκιμασία κατηγοριοποίησης M3b (§~\ref{sec:m3b}) ο ακροατής ακούει τρία αποσπάσματα και επιλέγει εκείνο που ταιριάζει λιγότερο με τα άλλα δύο, ώστε να εξεταστεί αν οργανώνει τις πραγματώσεις σύμφωνα με τις οικογένειες ακόμη και όταν η επιφανειακή τους μορφή διαφέρει. Στη δοκιμασία προσδοκίας M3c (§~\ref{sec:m3c}) αξιολογεί πόσο απρόσμενη του φαίνεται η συνέχεια ενός αποσπάσματος. Εξετάζεται αν το πληροφοριακό περιεχόμενο της επόμενης οικογένειας εξηγεί πρόσθετη διακύμανση στις ανθρώπινες κρίσεις πέρα από την πρόβλεψη που βασίζεται στα μεμονωμένα μελωδικά γεγονότα. Μία δευτερεύουσα ανάλυση συγκρίνει τη σχέση των κρίσεων με τις δύο υπολογιστικές προβλέψεις όταν το ίδιο απόσπασμα συνοδεύεται από σύντομη ένδειξη λατρευτικού πλαισίου ή από ουδέτερη ένδειξη. Η M4 (§~\ref{sec:m4-broader}) αναλύει τη σύγχρονη εκτελεστική μεταβλητότητα και παρέχει το ανεξάρτητο εκτελεστικό κριτήριο που απαιτεί η κύρια σύγκριση της M2, στο επίπεδο που επιτρέπει η διαθέσιμη κάλυψη. Περαιτέρω αναλύσεις εξετάζουν τη σχέση των θέσεων με τεκμηριωμένες ισοκρατηματικές αποδόσεις, οι οποίες παρέχουν συμπληρωματική μαρτυρία για την τροπική τους λειτουργία\footnote{Ισοκράτημα: συνοδευτική φωνητική στήριξη με φθόγγους που διατηρούνται ή μεταβάλλονται κατά τη μουσική πορεία. Στην πρόταση κάθε ισοκρατηματική απόδοση διατηρεί χωριστή προέλευση και δεν ταυτίζεται αυτομάτως με μοναδική «ορθή» τροπική ετικέτα.} και, όπου υπάρχουν επαρκή δεδομένα, τη διαφοροποίηση των εκτελέσεων μέσα στην ίδια οικογένεια.

Τέλος, η ιστορική εφαρμογή M5 (§~\ref{sec:nikolaos}) περιγράφει πλήρη έργα με χάρτη εμφανίσεων, κάλυψης και τεκμηριωμένων διαδοχών. Για τον Νικόλαο Πρωτοψάλτη Σμύρνης εξετάζονται χωριστά τέσσερις διαστάσεις: ποιες οικογένειες επιλέγονται, πώς διατάσσονται στις τεκμηριωμένες συνεχείς διαδοχές, σε ποιο δομικό και κειμενικό περιβάλλον χρησιμοποιούνται και πώς διαφοροποιούνται οι πραγματώσεις μέσα στην ίδια οικογένεια. Οι διαστάσεις αναφέρονται χωριστά ως προς δηλωμένο σώμα αναφοράς, χωρίς να συγχωνεύονται σε έναν συνολικό βαθμό. Το πλήρες έργο διατηρεί επίσης τις συνδετικές, μη αποδοθείσες ή ανεπίλυτες περιοχές, ώστε η ακολουθιακή ανάλυση να μην δημιουργεί τεχνητές μεταβάσεις πάνω από κενά τεκμηρίωσης. Η ανεξάρτητη μουσικολογική ανάλυση λειτουργεί ως κριτήριο αναφοράς, ενώ ο χάρτης που παράγει το υπολογιστικό σύστημα αποτελεί χωριστή έξοδο· η συμφωνία και η απόκλισή τους αναφέρονται χωρίς να συγχέονται με την ιστορική διαφοροποίηση του έργου ως προς το σώμα αναφοράς.

Η αναμενόμενη συμβολή είναι ένα ελεγχόμενο πλαίσιο αναγνώρισης μουσικών μονάδων μεταβλητής έκτασης και μουσικολογικής ερμηνείας της παραλλαγής τους, με σαφή διάκριση ιστορικής τεκμηρίωσης, υπολογιστικής πρόβλεψης, σύγχρονης εκτέλεσης και ανθρώπινης απόκρισης. Το πλαίσιο δεν προϋποθέτει ότι κάθε έργο διαιρείται πλήρως σε μη επικαλυπτόμενες θέσεις: η πλήρης μουσικολογική αναφορά διατηρεί αβέβαια όρια, συνδετικό υλικό και πιθανές τοπικές επικαλύψεις, ενώ η μορφή της υπολογιστικής εξόδου οριστικοποιείται στην αρχική απογραφή πριν από τη M1. Η μεθοδολογική αρχή μπορεί να εξεταστεί και σε άλλες παραδόσεις που διαθέτουν τεκμηριωμένες μουσικές μονάδες, οργανωμένη παραλλαγή και επαρκή τεκμηρίωση της προέλευσης των πηγών, αλλά κάθε τέτοια γενίκευση απαιτεί ανεξάρτητο έλεγχο.

\clearpage
\setcounter{tocdepth}{1}
\tableofcontents
\clearpage

\section{Εισαγωγή και κεντρική συνεισφορά}

Στη βυζαντινή ψαλτική επανέρχονται μελωδικοί σχηματισμοί που μεταβάλλονται από έργο σε έργο και διατηρούν αναγνωρίσιμα στοιχεία της μουσικής τους ταυτότητας. Η παράδοση περιγράφει τέτοιους σχηματισμούς με την έννοια της \textbf{μουσικής θέσης}. Η παρούσα έρευνα εξετάζει πώς μπορεί ένα υπολογιστικό σύστημα να αναγνωρίζει αυτή την ταυτότητα μέσα στην παραλλαγή: δύο αισθητά διαφορετικές αποδόσεις μπορεί να αποτελούν πραγματώσεις της ίδιας θέσης, ενώ δύο όμοια αποσπάσματα μπορεί να αντιστοιχούν σε διαφορετικές θέσεις. Για να ελεγχθεί αυτό το ερώτημα, η έρευνα συγκρίνει διαφορετικές δομικές κλίμακες αναπαράστασης του μέλους και χρησιμοποιεί μουσικολογικά, ανθρώπινα και εκτελεστικά δεδομένα με διακριτούς ρόλους στην εκπαίδευση και στην αξιολόγηση.

\subsection{Η μουσική θέση και η ταυτότητα μέσα στη μεταβολή}

Το κύριο συμβολικό σώμα χρησιμοποιεί τη \textbf{Νέα Μέθοδο}, το αναλυτικό σύστημα νευματικής σημειογραφίας που καθιερώθηκε το 1814 από τους λεγόμενους \textbf{Τρεις Διδασκάλους}: τον Χρύσανθο τον εκ Μαδύτων, τον Γρηγόριο Πρωτοψάλτη και τον Χουρμούζιο Χαρτοφύλακα. Η γραφή είναι σχεσιακή: τα νεύματα δηλώνουν κυρίως μελωδική κίνηση σε σχέση με το προηγούμενο σημείο και συνδυάζονται με χρονικές, ποιοτικές και τροπικές πληροφορίες.

Η θέση διαθέτει ιστορική και μορφολογική θεμελίωση ως οργανωμένος μελικός σχηματισμός. Ο \textbf{Μανουήλ Χρυσάφης}, λαμπαδάριος και θεωρητικός της ψαλτικής του 15ου αιώνα, συνδέει τη θέση με τη συνένωση σημείων που συγκροτούν μέλος (\emph{Περὶ τῶν ἐνθεωρουμένων τῇ ψαλτικῇ τέχνῃ}, έκδ. D.~E.~Conomos, 1985b, σσ.~40--41). Ο \textbf{Κωνσταντίνος Χ.~Καραγκούνης}, ερευνητής της ιστορίας, της μορφολογίας και της εξήγησης του παπαδικού μέλους, μελέτησε συστηματικά στη διδακτορική του διατριβή \emph{Η Παράδοση και Εξήγηση του Μέλους των Χερουβικών της Βυζαντινής και μετα-Βυζαντινής Μελοποιίας} (2003) τις ονομασμένες θέσεις και τη μελοποιητική τους χρήση (ιδίως σσ.~714--715· πρβλ. Στάθης, 1987, σσ.~343--344). Ο \textbf{Αντώνιος Ε.~Αλυγιζάκης}, στη μελέτη \emph{Η οκταηχία στην ελληνική λειτουργική υμνογραφία} (Θεσσαλονίκη: Πουρναράς, 1985, σσ.~81--82), περιγράφει τις θέσεις ως επαναχρησιμοποιούμενα μελικά πρότυπα που συμμετέχουν στη συγκρότηση φράσεων και ολοκληρωμένων μελών.\footnote{Η βιβλιογραφική ταυτότητα της έντυπης έκδοσης του 1985 έχει επιβεβαιωθεί. Το Εθνικό Αρχείο Διδακτορικών Διατριβών καταγράφει το έργο ως διατριβή του 1986 (ΕΑΔΔ 0578). Η ακριβής διατύπωση των σσ.~81--82 διασταυρώθηκε από τη σχετική δημοσιευμένη παράθεση του Αχιλλέα Χαλδαιάκη και την παρεχόμενη αποδελτίωση· εκκρεμεί αυτοψία των αντίστοιχων έντυπων σελίδων.} Οι ιστορικές περιγραφές τεκμηριώνουν την οργανωμένη και επαναχρησιμοποιούμενη θέση, ενώ ο υπολογιστικός προσδιορισμός των ορίων και των σχέσεων μεταξύ πραγματώσεων αποτελεί αντικείμενο ελέγχου της παρούσας έρευνας. Στο εξεταζόμενο σώμα, κάθε \textbf{εμφάνιση} αντιστοιχεί σε συγκεκριμένο χωρίο έργου και πηγής και κάθε \textbf{πραγμάτωση} στη μορφή που αποκτά εκεί ο τύπος. Οι πραγματώσεις συγκροτούν υποψήφιες \textbf{οικογένειες θέσεων}, των οποίων η συνοχή δοκιμάζεται με τύπο-ειδικές διαγνωστικές σχέσεις και ανεξάρτητες μουσικολογικές κρίσεις.

Χρήσιμες δοκιμασίες των ορίων αυτής της ταυτότητας προσφέρει η διατριβή του Κ.~Χ. Καραγκούνη (2003). Στη σ.~788 συγγενής μετροφωνική επιφάνεια ---δηλαδή συγγενής βασική κίνηση κατά την ανάγνωση των ποσοτικών σημείων--- μπορεί να αντιστοιχεί σε διαφορετικές ιστορικές θέσεις όταν συνεκτιμώνται άλλα γνωρίσματα. Στις σσ.~810--812 καταγράφονται θέσεις με ιδιαίτερα τροπικά συμφραζόμενα· στη σ.~811 η παρακλητική παρουσιάζει ειδικά σταθερότητα στον εξεταζόμενο τρόπο μεταγραφής. Οι περιπτώσεις αιτιολογούν τη διάκριση μελωδικής ομοιότητας, ιστορικού ονόματος και οικογενειακής ένταξης, με χωριστή αξιολόγηση κάθε εμφάνισης.

Ο \textbf{Ιωάννης Αρβανίτης}, στη διδακτορική του διατριβή \emph{Ο ρυθμός των εκκλησιαστικών μελών μέσα από την παλαιογραφική έρευνα και την εξήγηση της παλαιάς σημειογραφίας} (Ιόνιο Πανεπιστήμιο, 2010, τ.~Α΄, §5.2, σσ.~136--142), εξετάζει τη θέση ως αρχή μελικής ομοιότητας παρά τις διαφοροποιήσεις σημειογραφικής γραφής, λεκτικού τονισμού και συλλαβικής κατανομής. Το ΠΕ0 μετατρέπει τις παρατηρήσεις σε υποψήφιες δοκιμασίες συλλαβικής προσαρμογής, κλίμακας, ένθεσης και χρονικής οργάνωσης. Η ειδική ρυθμική του θεώρηση αφορά κυρίως ειρμολογικό και στιχηραρικό ρεπερτόριο· η εφαρμογή της στα παπαδικά Χερουβικά και Κοινωνικά αποτελεί χωριστό εμπειρικό ζήτημα.

Ο \textbf{Nicolae Gheorghiță}, μουσικολόγος με έρευνα στη σημειογραφία και στην ιστορική εξήγηση, μελετά στο άρθρο \emph{Observations on the Technique of Transcription (Exegesis) into the New Method of Analytical Music Notation of the Sunday Koinonikon of the 18th Century} (2010, σσ.~9--26· αναδημοσίευση σε βιβλίο του ίδιου έτους, κεφ.~XIV, σσ.~225--248) σαράντα δύο Κυριακάτικα Κοινωνικά μέσα από επαναλαμβανόμενες θέσεις, τη διαδοχή και τη λειτουργία τους. Πρόκειται για μία εργασία σε δύο σελιδοποιήσεις. Διακρίνει ως παράγοντες της πραγμάτωσης τη βαθμίδα, το γένος, το μελοποιητικό ιδίωμα και, στο συγκεκριμένο παλαιογραφικό περιβάλλον, το χρώμα χειρονομικού σημείου. Αυτές οι παρατηρήσεις τεκμηριώνουν υποψήφιους κανόνες περιβάλλοντος ανά τύπο και δεν μετατρέπονται αυτομάτως σε πλήρες λεξιλόγιο οικογενειών ή υποχρεωτικά διαγνωστικά γνωρίσματα. Ο \textbf{Κωνσταντίνος Ι.~Φωτόπουλος}, στη μελέτη του για τη μελοποιία και τη σύγχρονη προβληματική της (2009), προσθέτει χρήσιμες κειμενικές και συνθετικές παραμέτρους.

Το Σχήμα~\ref{fig:prop-formula-example} δίνει ένα άμεσο οπτικό παράδειγμα του είδους της δομημένης παραλλαγής που καλείται να χειριστεί η αναπαράσταση.

\begin{figure}[htbp]
\centering
\includegraphics[width=0.68\textwidth]{formula_position_example.png}
\caption{\textbf{Μελωδικός τύπος και παραλλαγές σε κοινό προσωδιακό πρότυπο.} Η ακολουθία \textbf{00001} κωδικοποιεί πρότυπο πέντε συλλαβών, όπου \textbf{0} δηλώνει άτονη και \textbf{1} τονισμένη συλλαβή· δεν αποτελεί αναγνωριστικό μουσικής θέσης. Οι έγχρωμες γραφές δηλώνουν καταγεγραμμένες παραλλαγές του ίδιου μελωδικού τύπου. Το παράδειγμα εισάγει οπτικά το πρόβλημα της οικογενειακής ταυτότητας υπό δομημένη παραλλαγή· η τελική οικογενειακή απόδοση παραμένει αντικείμενο μουσικολογικής τεκμηρίωσης και ανεξάρτητου ελέγχου. Πηγή: St.~Anthony's Greek Orthodox Monastery, \emph{Papadic Formulae}.}
\label{fig:prop-formula-example}
\end{figure}
\FloatBarrier

\paragraph{Μικρό παράδειγμα: από την πηγή στην οικογενειακή κρίση.}
Δύο πραγματώσεις μπορεί να έχουν συγγενή μετροφωνική επιφάνεια αλλά να αντιστοιχούν σε διαφορετικές ονομασμένες θέσεις (ενδεικτικό ιστορικό υλικό: Καραγκούνης, 2003, σ.~788). Αντίστροφα, διαφορές συλλαβικής κατανομής ή τονισμού μπορεί να συνυπάρχουν με αναγνωρίσιμη μελική οργάνωση (Αρβανίτης, 2010, σσ.~136--142). Το ΠΕ0 εξετάζει παρόμοιες περιπτώσεις με τον εξής προκαθορισμένο κανόνα:\par\smallskip
\noindent{\small\begin{tabularx}{\linewidth}{@{}p{0.28\linewidth}Y@{}}
\toprule
\textbf{Παρατηρούμενο τεκμήριο} & \textbf{Ελεγχόμενη οικογενειακή κρίση}\\
\midrule
Κοινή βασική μελωδική κίνηση & Έλεγχος των τύπο-ειδικών διαγνωστικών σχέσεων και των ιστορικά μαρτυρημένων διαφορών, όχι αυτόματη απόδοση κοινής οικογένειας.\\
Διαφορετική συλλαβική εφαρμογή & Έλεγχος αν η υπογραφή επιτρέπει την προσαρμογή στο συγκεκριμένο περιβάλλον, χωρίς προαποφασισμένο αποκλεισμό.\\
Ανεπαρκές ή αντιφατικό τεκμήριο & Κρίση \texttt{unresolved} αντί αυθαίρετου \texttt{member} ή \texttt{non\_member}.\\
\bottomrule
\end{tabularx}}\par\smallskip
Ο πίνακας περιγράφει τον \emph{τρόπο κρίσης}, όχι νέες επισημειώσεις των αναφερόμενων ιστορικών παραδειγμάτων. Οι αντίστοιχες υπογραφές οριστικοποιούνται στο υλικό ανάπτυξης και ελέγχονται από διαφορετικούς ειδικούς σε ανεξάρτητες εμφανίσεις.

\subsection{Όροι και κλίμακες ανάλυσης}

Η ίδια μουσική πηγή μπορεί να περιγραφεί σε διαφορετικές κλίμακες. Ένα γραφικό σημείο συμβάλλει στην ανάγνωση ενός μελωδικού γεγονότος· συνδυασμοί νευμάτων σχηματίζουν μεγαλύτερες μονάδες· οι θέσεις οργανώνονται σε σχέση με φραστικές ενότητες και πλήρη έργα. Τα όρια φράσεων και θέσεων σχετίζονται, αλλά δεν συμπίπτουν αναγκαστικά, και μια νέα φραστική ενότητα μπορεί να αρχίζει πριν ολοκληρωθεί μια θέση. Η πρώτη μεθοδολογική ενότητα (M1, §~\ref{sec:m1}) συγκρίνει διαφορετικές κωδικοποιήσεις εισόδου και, σε χωριστές δομικές συγκρίσεις, εξετάζει τη ρητή αναπαράσταση των θέσεων έναντι αναπαραστάσεων σε επίπεδο μελωδικού γεγονότος και αλγοριθμικά επαγόμενων τμημάτων.

Η περιγραφή περιλαμβάνει και τη σχέση \emph{μέρους--όλου}: ένα κοινό μελωδικό υποτμήμα μπορεί να ανήκει σε μεγαλύτερους, διαφορετικά οργανωμένους σχηματισμούς. Η ένθεση τεκμηριώνεται, όπου υπάρχει, σε συγκεκριμένη αναλυτική κλίμακα· η τοπική σύμπτωση δεν οδηγεί από μόνη της σε κοινό \texttt{family\_id}. Πολλαπλές κλίμακες, σύνθετες θέσεις και πραγματικές επικαλύψεις κρίνονται ως διακριτές περιπτώσεις και δεν εξομοιώνονται αυτόματα.

Η συνδετική λειτουργία δεν ταυτίζεται με μη αναγνωρισμένο υλικό ανάμεσα σε θέσεις. Μια τεκμηριωμένη θέση μπορεί να λειτουργεί ως σύνδεσμος ανάμεσα σε μεγαλύτερες μελωδικές μορφές. Στο μητρώο καταγράφεται τότε ως θέση με \texttt{functional\_role=connector}. Μόνο το υλικό που δεν έχει τεκμηριωθεί ως εμφάνιση τύπου μπορεί να λάβει την κατάσταση \texttt{unassigned\_connective\_material}. Η διάκριση αποτρέπει την αφαίρεση πραγματικών formulae από το λεξιλόγιο μόνο επειδή έχουν συνδετική λειτουργία (πρβλ. Gheorghiță, 2010, Πίν.~8--9).

\paragraph{Πραγμάτωση, εκτέλεση και επιτελεστικό πλαίσιο.}
Η \textbf{πραγμάτωση} είναι η συγκεκριμένη μουσική μορφή μιας θέσης σε έργο και πηγή. Η \textbf{εκτέλεση} είναι η πραγματική φωνητική και χρονική εκφορά αυτής της μουσικής μορφής από ψάλτη ή χορό. Το ευρύτερο \textbf{επιτελεστικό πλαίσιο} περιλαμβάνει επιπλέον τη λειτουργική περίσταση, τον χώρο, τους συμμετέχοντες και τις κοινωνικές σχέσεις της ψαλτικής πράξης. Η διάκριση είναι αναλυτική και δεν δημιουργεί νέο πείραμα. Στην πρόταση η οικογενειακή ταυτότητα, η συνθετική πραγμάτωση, η σύγχρονη εκτέλεση και το ευρύτερο λειτουργικό πλαίσιο διατηρούν διαφορετικούς ρόλους. Τα δεδομένα εκτέλεσης χρησιμοποιούνται όπου προβλέπεται στη M2/M4, ενώ το λειτουργικό πλαίσιο καταγράφεται ως μεταδεδομένο ή πειραματική συνθήκη μόνο όταν έχει δηλωθεί ρητά.

Η σχέση της θέσης με τη συλλαβική άρθρωση είναι μη αμφιμονοσήμαντη προς αμφότερες τις κατευθύνσεις. Στις μορφολογικές αναλύσεις του Καραγκούνη, μία εναρκτήρια θέση μπορεί να εκτείνεται στις τρεις πρώτες συλλαβές «Οἱ τὰ Χε-», ενώ σε άλλα σημεία μία συλλαβή αναπτύσσεται σε περισσότερες διαδοχικές θέσεις (Καραγκούνης, 2003, σσ.~207, 209). Η συλλαβή λειτουργεί επομένως ως χωριστό επίπεδο κειμενικής και προσωδιακής πληροφορίας που ευθυγραμμίζεται με τη μουσική δομή, ενώ τα όρια της θέσης προσδιορίζονται μουσικολογικά. Η M1 διατηρεί αυτή τη σχέση στο κοινό ατομικό επίπεδο αναφοράς και εξετάζει τη σύμπτωση ή απόκλιση των ορίων ως μετρήσιμη ιδιότητα της τμηματοποίησης.

Για να διαβάζεται το ερευνητικό σχέδιο χωρίς προϋπάρχουσα γνώση της ορολογίας, οι βασικοί όροι ορίζονται εδώ πριν χρησιμοποιηθούν τεχνικά:

{\small
\begin{xltabular}{\textwidth}{@{}>{\RaggedRight\arraybackslash}p{0.23\textwidth}Y@{}}
\toprule
\textbf{Όρος} & \textbf{Ορισμός στην παρούσα έρευνα}\tabularnewline
\midrule
\endfirsthead
\toprule
\textbf{Όρος} & \textbf{Ορισμός στην παρούσα έρευνα}\tabularnewline
\midrule
\endhead
\midrule
\multicolumn{2}{r}{\footnotesize Συνέχεια στην επόμενη σελίδα}\tabularnewline
\endfoot
\bottomrule
\endlastfoot
\textbf{Μουσική θέση} & Ιστορικός μουσικολογικός όρος για οργανωμένο μελικό σχηματισμό μεταβλητής έκτασης. Η θέση μπορεί να επιτελεί εναρκτήρια, αναπτυκτική, συνδετική ή καταληκτική λειτουργία ανάλογα με το είδος, τον τύπο και τη συγκεκριμένη εμφάνιση. Η ιστορική ονομασία ή προέλευση είναι τεκμήριο. Δεν αποτελεί από μόνη της σύγχρονη οικογενειακή ετικέτα.\tabularnewline
\midrule
\textbf{Μουσική φράση} & Μορφολογική ή μελική ενότητα της μουσικής ροής. Τα όριά της τεκμηριώνονται χωριστά από τα όρια θέσεων. Τα δύο επίπεδα μπορούν να συμπίπτουν ή να αποκλίνουν.\tabularnewline
\midrule
\textbf{Αναλυτική κλίμακα} & Το επίπεδο στο οποίο κρίνεται η μουσική μονάδα. Μπορεί να αφορά αυτοτελή θέση, σύνθετη θέση, ακολουθία θέσεων ή συστατικό μεγαλύτερης μονάδας. Η κλίμακα δηλώνεται πριν από την οικογενειακή κρίση.\tabularnewline
\midrule
\textbf{Οικογένεια θέσεων} & Σύγχρονη αναλυτική κατηγορία εμφανίσεων με κοινή διαγνωστική μουσική οργάνωση. Η ένταξη κρίνεται σε προκαθορισμένη αναλυτική κλίμακα. Χρησιμοποιείται παγωμένη έκδοση υπογραφής τύπου. Η υπογραφή δηλώνει διαγνωστικές σχέσεις με ρόλο και πεδίο εφαρμογής, επιτρεπτές διαμορφώσεις, κριτήρια αποκλεισμού και τεκμηριωμένα θετικά και δύσκολα αντιπαραδείγματα. Η δυνατότητα περισσότερων επαρκών διαμορφώσεων εξετάζεται ανά τύπο, χωρίς να απαιτείται ένας κοινός πυρήνας αναλλοίωτων για όλα τα μέλη. Κοινή ονομασία, λειτουργία, μετροφωνία ή εξηγητική σχέση δεν αρκούν μόνες τους για ένταξη.\tabularnewline
\midrule
\textbf{Συμβατότητα περιβάλλοντος} & Περιγραφική κρίση για τη χρήση μιας θέσης στο συγκεκριμένο τροπικό, κειμενικό και συνθετικό περιβάλλον. Διακρίνεται από την οικογενειακή ταυτότητα. Μόνο συνθήκη που έχει προδηλωθεί ως ταυτοτική διαγνωστική απαίτηση σε μία ενεργή επιτρεπτή διαμόρφωση μπορεί να συμμετέχει στην απόφαση ένταξης· η συνήθης λειτουργία δεν λειτουργεί μόνη της ως ετικέτα.\tabularnewline
\midrule
\textbf{Πραγμάτωση} & Η συγκεκριμένη μουσική μορφή που αποκτά μία θέση σε έργο και πηγή. Περιλαμβάνει μελωδική, χρονική, τροπική και σημειογραφική οργάνωση. Το περιβάλλον καταγράφεται χωριστά.\tabularnewline
\midrule
\textbf{Παραλλαγή πραγματώσεων} & Μελωδική, σημειογραφική, χρονική ή εκτελεστική διαφοροποίηση μιας πραγματώσεως ενός τύπου· δεν ταυτίζεται με την ιστορική ψαλτική παραλλαγή (φθογγική και τροπική ανάγνωση). Η ένταξή της στην οικογένεια εξετάζεται χωριστά.\tabularnewline
\midrule
\textbf{Παραλλαγή (ψαλτική)} & Φθογγική και τροπική ανάγνωση με συλλαβές/ονομασίες του συστήματος των ήχων. Συναφής με τη μετροφωνία, αλλά μη ταυτόσημη σε κάθε ιστορική χρήση· διακρίνεται από την πλήρη μελική εκφορά.\tabularnewline
\midrule
\textbf{Εμφάνιση} & Εντοπισμένη παρουσία μιας πραγμάτωσης σε συγκεκριμένο έργο, πηγή ή εκτέλεση, με σταθερό αναγνωριστικό και τεκμηριωμένα όρια. Όπου περισσότερες οριοθετήσεις είναι μουσικολογικά αποδεκτές, διατηρείται σύνολο επιτρεπτών ζευγών αρχής--τέλους για την ίδια εμφάνιση. Η κοινή συμμετοχή γεγονότων δύο διαφορετικών εμφανίσεων καταγράφεται χωριστά από την αβεβαιότητα των ορίων.\tabularnewline
\midrule
\textbf{Αλυσίδα θέσεων} & Τεκμηριωμένη διαδοχή θέσεων ή οικογενειακών ταυτοτήτων μέσα στη μουσική ροή ενός έργου, όπου η σειρά και η συνέχεια μπορούν να προσδιοριστούν αξιόπιστα· διαφορετικό επίπεδο από την εσωτερική παραλλαγή μιας οικογένειας.\tabularnewline
\midrule
\textbf{Εξηγητική σχέση} & Σύνδεση παλαιότερης μουσικής καταγραφής με συγκεκριμένη μεταγενέστερη αναλυτικότερη απόδοσή της, τεκμηριωμένη από τις ιστορικές πηγές και τη γραμμή προέλευσης των μαρτύρων. Τεκμηριώνει τη σχέση και την προέλευση συγκεκριμένων αποδόσεων, ενώ η κοινή μουσική ταυτότητα και ο τύπος διαφοροποίησης ελέγχονται χωριστά.\tabularnewline
\midrule
\textbf{Μελωδικό γεγονός} & Ελάχιστη μονάδα μελωδικής κίνησης στην υπολογιστική αναπαράσταση, η οποία προκύπτει από την ανάγνωση ενός ή περισσότερων γραφικών σημείων.\tabularnewline
\midrule
\textbf{Ενιαίο συμβολικό επίπεδο αναφοράς} & Κοινή, ενιαία μορφή στην οποία μετατρέπονται οι διαφορετικές πηγές ώστε όλα τα μοντέλα να συγκρίνονται πάνω στην ίδια ακολουθία μουσικών γεγονότων, με διατήρηση της σύνδεσης με την αρχική πηγή και τα σημειογραφικά της στοιχεία.\tabularnewline
\midrule
\textbf{Σημειογραφική χρονική αξία} & Σχετική διάρκεια που προκύπτει από τους χρονικούς χαρακτήρες της Νέας Μεθόδου με δηλωμένο κανόνα ανάγνωσης και κοινή μονάδα. Δεν αποτελεί απόλυτη διάρκεια εκτέλεσης.\tabularnewline
\midrule
\textbf{Αγωγή} & Συνολική ή τοπική ταχύτητα εκφοράς των σημειογραφικών χρονικών αξιών στην εκτέλεση· διακρίνεται από τις ίδιες τις σχετικές χρονικές αναλογίες.\tabularnewline
\midrule
\textbf{Ρυθμική/μετρική οργάνωση} & Ερμηνευτική οργάνωση των χρονικών αξιών σε ισχυρότερες και ασθενέστερες φάσεις ή μεγαλύτερες ομάδες. Χρησιμοποιείται μόνο όπου τεκμηριώνεται και δεν συνάγεται αυτόματα από τις διάρκειες ή τις εκδοτικές διαστολές.\tabularnewline
\midrule
\textbf{Ήχος} & Τροπικό πλαίσιο με χαρακτηριστική οργάνωση βαθμίδων, κινήσεων και καταλήξεων· δεν ταυτίζεται με απλή κλίμακα. Ο γενικός ήχος του έργου διακρίνεται από την τοπική τροπική λειτουργία και τη γραπτή μαρτυρία.\tabularnewline
\midrule
\textbf{Τοπική τροπική λειτουργία} & Αναλυτική κρίση για τον μελικά αναγνωριζόμενο τροπικό χαρακτήρα συγκεκριμένου χωρίου, με τεκμηρίωση βαθμίδων, καταλήξεων, μαρτυριών και αβεβαιότητας. Δεν ταυτίζεται αυτομάτως με οικογενειακή ετικέτα ή με διαστηματική αλλοίωση.\tabularnewline
\midrule
\textbf{Λειτουργικός ρόλος} & Η μουσική λειτουργία μιας εμφάνισης μέσα στο μέλος, όπως εναρκτήρια, καταληκτική ή μη καταληκτική χρήση· λεπτότερες κατηγορίες χρησιμοποιούνται μόνο όπου ορίζονται ρητά για τις ανάγκες της ανάλυσης.\tabularnewline
\end{xltabular}
}\par

\paragraph{Κανόνες οικογενειακής ταυτότητας.}
Η οικογενειακή κρίση δεν βασίζεται σε ένα μόνο επιφανειακό γνώρισμα. Ίδια ιστορική ονομασία δεν συνεπάγεται κατ' ανάγκην ίδια οικογένεια. Ίδια μετροφωνική πορεία δεν συνεπάγεται κατ' ανάγκην ίδιο μέλος. Διαφορετική επιφάνεια δεν συνεπάγεται κατ' ανάγκην διαφορετική οικογένεια. Κοινή λειτουργία δεν αρκεί για κοινή ταυτότητα. Κοινό υποτμήμα δεν αρκεί για κοινή μεγαλύτερη θέση. Τέλος, μια εξηγητική σχέση δεν δημιουργεί αυτομάτως οικογενειακή ένταξη. Κάθε περίπτωση ελέγχεται απέναντι στην παγωμένη υπογραφή του τύπου.

\paragraph{Ιστορική θέση και αναλυτική οικογένεια.}
Ο Χρυσάφης (15ος αιώνας· έκδ. Conomos, 1985b, σσ.~40--41) παρέχει ιστορικό θεωρητικό πλαίσιο για τη θέση ως οργανωμένο μελικό σχηματισμό. Ο \textbf{Σπυρίδων Αντωνόπουλος}, στη διδακτορική του διατριβή \emph{The Life and Works of Manuel Chrysaphes the Lampadarios, and the Figure of Composer in Late Byzantium} (City University London, 2014 [υποβλήθηκε 2013], σσ.~284--286, εικ.~5.36), εξετάζει παραλλαγές αναγνωρίσιμης θέσης. Ο \textbf{Christian Troelsgård}, ερευνητής της μεσοβυζαντινής μουσικής σημειογραφίας, στο εγχειρίδιό του \emph{Byzantine Neumes: A New Introduction to the Middle Byzantine Musical Notation} (2011, §32, σ.~47), συσχετίζει συμβατικούς σχηματισμούς με μεγάλα σημάδια, αλλά και διαφοροποιήσεις των σημείων. Η μελέτη του \textbf{Αντωνίου Μποτωνάκη}, \emph{Η Σοφωτάτη και Αρίστη Παραλλαγή Ιωάννου Ιερέως του Πλουσιαδηνού} (2014, σσ.~337--373), παρουσιάζει τη μέθοδο του Πλουσιαδηνού, όπου η άσκηση σε μετροφωνία και παραλλαγή συνδέεται με τη διδασκαλία ονομασμένων θέσεων και χαρακτηριστικών εναρκτήριων και καταληκτικών φράσεων. Η ιστορική διδασκαλία θεμελιώνει τη μουσική σημασία της θέσης· η παρούσα πρόταση μετατρέπει τις περιγραφές σε ελέγξιμες υπογραφές. Τα παραδείγματα συγκρότησης υπογραφών χαρακτηρίζονται \texttt{definition} και δεν λειτουργούν ως ανεξάρτητη επικύρωση.

Η οντολογία της έρευνας διατηρεί τέσσερα διακριτά επίπεδα: \textbf{ιστορική θέση}, \textbf{αναλυτική οικογένεια}, \textbf{πραγμάτωση} και \textbf{εμφάνιση}. Η οικογένεια δεν κληρονομείται μηχανικά από το ιστορικό όνομα ή από μία σχέση προέλευσης. Για κάθε υποψήφιο τύπο τηρείται εκδοχή υπογραφής (\texttt{type\_signature\_version}), η οποία οριστικοποιείται πριν από τον ανεξάρτητο έλεγχο και επικυρώνεται σε έργα που δεν χρησιμοποιήθηκαν για τη συγκρότησή της. Κάθε μεταγενέστερη επέκταση ή αλλαγή του κανόνα δημιουργεί νέα έκδοση. Εναλλακτικές περιγραφές της οικογενειακής ταυτότητας μπορούν να εξετάζονται διερευνητικά χωρίς να αποτελούν νέες επιβεβαιωτικές υποθέσεις. Ειδικά μια περιγραφή που βασίζεται σε τοπικές σχεσιακές ομοιότητες δεν μετατρέπεται σε οικογενειακή κλάση αν δεν έχει ελεγχθεί ότι ο κανόνας της παράγει συνεκτικές κατηγορίες χωρίς τεχνητή επιβολή μεταβατικότητας.

\paragraph{Παραλλαγή, μέλος και ειδική χρήση των ενδιάμεσων μαρτυριών.}
Η ιστορική ψαλτική \emph{παραλλαγή} αφορά φθογγική και τροπική ανάγνωση, ενώ η \emph{μετροφωνία} ειδικότερα τη σημειογραφικά οριζόμενη διαστηματική κίνηση· οι έννοιες σχετίζονται αλλά δεν συμπίπτουν σε κάθε ιστορική χρήση. Το \emph{μέλος} αναφέρεται στη συγκροτημένη μελωδική εκφορά. Ο \textbf{Ιωάννης Πλουσιαδηνός}, θεωρητικός και ψάλτης του 15ου αιώνα, διακρίνει στην \emph{Ἑρμηνεία τῆς Παραλλαγῆς} σχέσεις «ἀπὸ μέλους» και «ἀπὸ παραλλαγῆς» (δημοσιευμένη μεταγραφή: Μποτωνάκης, 2014, σ.~352). Η εκτενέστερη φράση «ἀπὸ παραλλαγῆς καὶ μέλους ὁμοῦ» εμφανίζεται σε άλλη δημοσιευμένη απόδοση του ίδιου θεωρητικού υλικού (Αλυγιζάκης, 1985, σ.~235, σύμφωνα με την αποδελτίωση· πρβλ. Αχιλλέας Γ.~Χαλδαιάκης, \emph{Βυζαντινομουσικολογικά}, τ.~Γ΄, 2014, σσ.~156--157, υπό επανέλεγχο της συγκεκριμένης σελίδας): η εκδοτική διαφορά καταγράφεται, ενώ εκκρεμεί αυτοψία των κωδίκων Διονυσίου~570 και ΕΒΕ~968. Έτσι ο όρος \emph{παραλλαγή} στην ιστορική θεωρία δεν συγχέεται με την \emph{παραλλαγή} ως διαφοροποιημένη πραγμάτωση ενός τύπου στη σύγχρονη αναλυτική χρήση.

Ο Troelsgård (2011, §67, σ.~69, παράδ.~38) εξετάζει συγκεκριμένα δύο ενδιάμεσες μαρτυρίες: η μαρτυρία πλαγίου του πρώτου αναγνωρίζει \emph{ἀπὸ μέλους} την τυπική καταληκτική συμπεριφορά της προηγούμενης φράσης, ενώ η μαρτυρία τετάρτου \emph{ἀπὸ παραλλαγῆς} προσδιορίζει τη θεωρητική θέση του καταληκτικού και επόμενου αρχικού φθόγγου. Για το συγκεκριμένο χωρίο η μερική τροπική μεταβολή δεν απαιτεί αλλοιώσεις στη μεταγραφή. Η μουσικολόγος \textbf{Μαρία Αλεξάνδρου} και ο Christian Troelsgård, στο άρθρο \emph{The Elements of the Papadikê and Modality Features in Byzantine Chant} (2020, σσ.~178--179), εντάσσουν την ίδια διάκριση στη μεταβαλλόμενη θέση και λειτουργία των ενδιαμέσων μαρτυριών, ενώ ο \textbf{Συμεών Κανάκης}, στο άρθρο \emph{Addressing the Modality of Church Melodies, According to the Variety in the Choice of the Isokratema. Case Studies} (2021, σ.~255) τη χρησιμοποιεί σε διαφορετικό παράδειγμα για την ερμηνευτική επιλογή τριφωνικής αντί τροχικής οργάνωσης. Έτσι το «ἀπὸ μέλους» σημαίνει τροπική αναγνώριση από τη μελωδική συμπεριφορά, η οποία μπορεί, αλλά δεν υποχρεούται, να συνεπάγεται διαστηματικές αλλαγές.

Αναλυτικά διακρίνονται \textbf{(α)} η τεκμηριωμένη σημειογραφική/μετροφωνική κίνηση, \textbf{(β)} η μελική πραγματοποίηση, \textbf{(γ)} η γραπτή θεωρητική αναφορά της μαρτυρίας και \textbf{(δ)} η τοπική τροπική λειτουργία. Το \texttt{local\_modal\_state} αποδίδει την τελευταία με τεκμήριο και αβεβαιότητα, χωριστά από το \texttt{global\_echos} και από το \texttt{family\_id}. Η ετικέτα \emph{ἀπὸ μέλους} δεν αποτελεί αυτοτελή σχέση \texttt{co\_explanation} ούτε απόδειξη διαφορετικής οικογενειακής ταυτότητας. Στην M5 η τροπική λειτουργία ερμηνεύεται μόνο με ανεξάρτητο κριτήριο όταν τροπικές πληροφορίες έχουν χρησιμοποιηθεί προηγουμένως για την ίδια την οικογενειακή απόδοση.

Στην υπολογιστική αλυσίδα, το ενιαίο συμβολικό επίπεδο αναφοράς λειτουργεί ως κοινή γλώσσα ανάμεσα στις διαφορετικές πηγές. Διατηρεί χωριστά όσα παρατηρούνται άμεσα στη σημειογραφία από όσα παράγονται με δηλωμένο υπολογιστικό κανόνα, όπως μια σχετική χρονική αξία ή μια σωρευτική συντεταγμένη. Κάθε τέτοιο πεδίο συνοδεύεται από προέλευση, ώστε να είναι γνωστό από ποια πηγή, ποια σημεία και ποιον κανόνα προέκυψε. Η ίδια σύνδεση επιτρέπει σε ένα αποτέλεσμα του μοντέλου να επιστρέφει στο συγκεκριμένο χωρίο που το παρήγαγε. Οι δοκιμές ανακατασκευής ελέγχουν ότι η μετατροπή προς το κοινό επίπεδο διατηρεί τις πληροφορίες που δηλώνονται ως ανακτήσιμες από την αρχική πηγή.

Η διάκριση των επιπέδων περιγραφής είναι ιδιαίτερα σημαντική και για τον χρόνο, επειδή η ιστορική ρυθμική ορολογία χρησιμοποιεί επίσης τον όρο «θέσις». Στην ιστορική ρυθμική ορολογία οι όροι \emph{θέσις} και \emph{ἄρσις} δηλώνουν διακριτές φάσεις της χρονικής κίνησης. Για να αποφεύγεται η ομωνυμία με τη μελωδική θέση της διατριβής, στο υπόλοιπο κείμενο γίνεται αναφορά στις δύο χρονικές φάσεις· χαρακτηρισμοί ισχυρότερης και ασθενέστερης φάσης χρησιμοποιούνται μόνο όπου στηρίζονται από τη συγκεκριμένη μετρική ερμηνεία. Η σημειογραφική χρονική αξία, η αγωγή και η μετρική ερμηνεία παραμένουν διαφορετικά επίπεδα περιγραφής.

Η θέση εντάσσεται έτσι σε έναν χάρτη δομικών κλιμάκων:
\[
\begin{aligned}
\text{σημειογραφικά στοιχεία}
&\rightarrow \text{ατομικά μελωδικά γεγονότα}
\rightarrow \text{νευματικοί συνδυασμοί}\\[0.25em]
&\rightarrow \boxed{\text{μουσικές θέσεις}}
\leftrightarrow \text{φραστικές/μορφολογικές ενότητες}
\rightarrow \text{έργο}.
\end{aligned}
\]
Η M1 συγκρίνει ρητές θέσεις, αναπαραστάσεις σε επίπεδο μελωδικού γεγονότος και αλγοριθμικά επαγόμενες μονάδες με τις ίδιες δοκιμασίες αξιολόγησης. Η αλγοριθμική ομοιότητα χρησιμεύει στην ανάκτηση και στη γενίκευση· η οικογενειακή ταυτότητα παραμένει χωριστή αναλυτική υπόθεση με δική της τεκμηρίωση.

\subsection{Από την ιστορική μετάδοση στην ελεγχόμενη εποπτεία}

Η \textbf{εξήγηση} αποτελεί ιστορικά τεκμηριωμένη πρακτική αναλυτικότερης σημειογραφικής απόδοσης ενός μέλους, η οποία καλλιεργείται ήδη πριν από τη μεταρρύθμιση του 1814. Ο Απόστολος Κώνστας ο Χίος, σύγχρονος των Τριών Διδασκάλων και ο ίδιος εξηγητής, αποτελεί σημαντικό θεωρητικό μάρτυρα της σχετικής συζήτησης. Το κείμενο του 1820 που σώζεται στον ΕΒΕ 1867 παραπέμπεται μέσω της γ΄ έκδοσης του Στάθη (1993, σ.~90), αλλά η ακριβής έκταση των επιμέρους ισχυρισμών που αποδίδονται στο συγκεκριμένο χωρίο παραμένει αντικείμενο άμεσης αντιπαραβολής της έκδοσης πριν χρησιμοποιηθεί σε ισχυρότερη ιστορική διατύπωση. Στο σώμα της παρούσας έρευνας, οι βασικές αναλυτικές αποδόσεις της Νέας Μεθόδου συνδέονται ιδίως με τον Γρηγόριο Πρωτοψάλτη και τον Χουρμούζιο Χαρτοφύλακα. Η εξηγητική σύνδεση τεκμηριώνει τη σχέση και την προέλευση συγκεκριμένων αποδόσεων, ενώ η κοινή μουσική ταυτότητα, η έκταση της διατήρησης και ο χαρακτήρας της διαφοροποίησης εξετάζονται χωριστά.

Η πρόταση μετατρέπει αυτή την τεκμηριωμένη πληροφορία προέλευσης σε ελεγχόμενο παράγοντα της εκπαίδευσης. Η \textbf{ιστορική μετάδοση} είναι η διαδικασία και η σχέση που τεκμηριώνουν οι πηγές· η \textbf{εποπτεία που βασίζεται σε τεκμηριωμένες ιστορικές σχέσεις} είναι ο τρόπος με τον οποίο το μαθησιακό πείραμα χρησιμοποιεί αυτές τις επαληθευμένες σχέσεις ως θετικά ζεύγη εκπαίδευσης, δηλαδή ως ζεύγη πραγματώσεων που το μοντέλο εκπαιδεύεται να αναπαριστά ως συγγενείς.

Η δεύτερη μεθοδολογική ενότητα (M2, §~\ref{sec:eise}) συγκρίνει δύο τρόπους εκπαίδευσης του Exegesis-Invariant Segmental Encoder (EISE), ενός κωδικοποιητή που μετατρέπει μουσικά τμήματα μεταβλητού μήκους σε αριθμητικές αναπαραστάσεις. Στη συνθήκη EISE, τα θετικά ζεύγη προέρχονται από τεκμηριωμένες εξηγητικές σχέσεις· στην οικογενειακή συνθήκη αναφοράς B2, από αντιστοιχισμένα ζεύγη της ίδιας οικογένειας χωρίς τη συγκεκριμένη ιστορική σύνδεση. Η κοινή αρχιτεκτονική και οι αντιστοιχισμένες συνθήκες εκπαίδευσης επιτρέπουν να εξεταστεί η πρόσθετη αξία της προέλευσης των θετικών σχέσεων.

Το κεντρικό αντικείμενο εξετάζεται με τέσσερις συμπληρωματικές πηγές τεκμηρίωσης, όπως συνοψίζει το Σχήμα~\ref{fig:conceptual-v865}. Οι ρόλοι τους στην εκπαίδευση και στην ανεξάρτητη αξιολόγηση καθορίζονται στις αντίστοιχες μεθοδολογικές ενότητες.

\begin{figure}[htbp]
\centering
\begin{tikzpicture}[
  node distance=12mm and 12mm,
  box/.style={draw=acc, rounded corners, align=center, inner sep=6pt, text width=34mm},
  core/.style={draw=acc, very thick, rounded corners, align=center, inner sep=8pt, text width=42mm},
  arr/.style={Latex-Latex, thick}
]
\node[core] (core) {Ταυτότητα και αναπαράσταση\\οικογένειας θέσεων};
\node[box, above=of core] (hist) {Ιστορικά τεκμηριωμένες\\σχέσεις};
\node[box, left=of core] (mus) {Μουσικολογική\\δομή και λειτουργία};
\node[box, right=of core] (human) {Ανθρώπινη\\κατηγορική οργάνωση};
\node[box, below=of core] (pred) {Ακολουθιακή πρόβλεψη\\και εκτελεστική μεταβλητότητα};
\draw[arr] (hist) -- (core);
\draw[arr] (mus) -- (core);
\draw[arr] (human) -- (core);
\draw[arr] (pred) -- (core);
\end{tikzpicture}
\caption{Κοινό ερευνητικό αντικείμενο της πρότασης. Η αναπαράσταση της οικογενειακής ταυτότητας ελέγχεται με ιστορικά, μουσικολογικά και συμπεριφορικά τεκμήρια, καθώς και με δεδομένα για την οργάνωση των ακολουθιών και για την εκτέλεση. Τα βέλη δηλώνουν σχέση ελέγχου και αξιοποίησης δεδομένων, όχι αιτιώδη ακολουθία.}
\label{fig:conceptual-v865}
\end{figure}
\FloatBarrier

\subsection{Συμβολή και οδηγός ανάγνωσης}

Η έρευνα διατυπώνει την \textbf{ταυτότητα της θέσης μέσα σε τεκμηριωμένη μεταβολή} ως κοινό πρόβλημα μουσικολογίας και Music Information Retrieval (MIR, Ανάκτηση Μουσικής Πληροφορίας). Το MIR αναπτύσσει υπολογιστικές μεθόδους για την αναπαράσταση, αναζήτηση, σύγκριση και ανάλυση μουσικών δεδομένων. Η παρούσα έρευνα ανήκει κυρίως στο \textbf{συμβολικό MIR}, επειδή το βασικό υπολογιστικό υλικό προέρχεται από κωδικοποιημένη μουσική σημειογραφία και οργανώνεται ως ακολουθίες μουσικών γεγονότων. Η ηχογραφημένη εκτέλεση εισέρχεται αργότερα ως χωριστή πηγή εκτελεστικών μετρήσεων, κυρίως στην ενότητα ανάλυσης των σύγχρονων εκτελέσεων (M4, §~\ref{sec:m4-broader}) και, στη M2, ως ανεξάρτητο εκτελεστικό κριτήριο αξιολόγησης.

Η \textbf{M1} εξετάζει ποια κλίμακα αναπαράστασης διατηρεί καλύτερα τις κρίσιμες μουσικολογικές διακρίσεις. Ως ισχυρό μοντέλο αναφοράς χρησιμοποιείται Transformer σε επίπεδο μελωδικού γεγονότος. Πρόκειται για μοντέλο ακολουθιών που χρησιμοποιεί μηχανισμό αυτοπροσοχής (self-attention) για να συνδυάζει πληροφορία από το προηγούμενο μουσικό συμφραζόμενο και να προβλέπει το επόμενο γεγονός. Στην αιτιακή του μορφή χρησιμοποιεί μόνο προηγούμενα γεγονότα. Η σύγκριση εξετάζει επομένως αν η ρητή μουσική θέση προσφέρει χρήσιμη δομική πληροφορία πέρα από όσα μπορεί να μάθει απευθείας ένα μοντέλο από την ακολουθία των γεγονότων (§~\ref{sec:m1}).

Η \textbf{M2} εξετάζει την προέλευση της πληροφορίας με την οποία μαθαίνεται η συγγένεια μεταξύ πραγματώσεων. Ένας κωδικοποιητής (encoder) μετατρέπει ένα μουσικό τμήμα σε διανυσματική αναπαράσταση που προκύπτει από την εκπαίδευση του μοντέλου. Η συνοπτική διανυσματική αναπαράσταση, ή embedding, τοποθετεί κάθε τμήμα σε έναν χώρο όπου οι σχετικές αποστάσεις μπορούν να χρησιμοποιηθούν για ανάκτηση και αναγνώριση συγγένειας. Η εγγύτητα σε αυτόν τον χώρο αποτελεί ιδιότητα που μαθαίνει το μοντέλο σύμφωνα με το πρωτόκολλο εκπαίδευσης· η μουσικολογική οικογενειακή ταυτότητα και η ιστορική σχέση τεκμηριώνονται με τα χωριστά κριτήρια της έρευνας.

Ο EISE παράγει μία συνοπτική διανυσματική αναπαράσταση για την οικογενειακή ταυτότητα και τοπικές καταστάσεις, δηλαδή αριθμητικές αναπαραστάσεις των επιμέρους γεγονότων του τμήματος. Η πρώτη εκπαιδεύεται ώστε να διατηρεί την κοινή ταυτότητα απέναντι σε τεκμηριωμένες εξηγητικές μεταβολές, ενώ οι δεύτερες προορίζονται για τη μελέτη της συγκεκριμένης πραγμάτωσης. Η διατήρηση αυτής της τοπικής πληροφορίας ελέγχεται εμπειρικά πριν από τις αντίστοιχες αναλύσεις (§~\ref{sec:eise} και Παράρτημα~\ref{app:eise-rationale}).

Η \textbf{M3a} παρέχει ανεξάρτητη μουσικολογική επικύρωση των οικογενειών. Η \textbf{M3c} εξετάζει την οικογένεια ως επίπεδο μουσικής πρόβλεψης και η \textbf{M3b} ως κατηγορία ακρόασης. Η \textbf{M4} παρέχει πρώτα το ανεξάρτητο εκτελεστικό κριτήριο αξιολόγησης που χρειάζεται η M2 και, όπου υπάρχει επαρκής κάλυψη, επεκτείνει την ανάλυση της σύγχρονης εκτέλεσης.

Η \textbf{M5} αποτελεί την ιστορική εφαρμογή του επικυρωμένου πλαισίου σε πλήρη έργα. Η σύγκριση με τεκμηριωμένα σώματα αναφοράς εντοπίζει χωριστά τη συγγένεια και τη διαφοροποίηση στην επιλογή οικογενειών θέσεων, στη διαδοχή τους, στη δομική ή κειμενική τοποθέτησή τους και στον τρόπο πραγμάτωσής τους. Η πρώτη πλήρης εφαρμογή αφορά τον Νικόλαο Πρωτοψάλτη Σμύρνης (§~\ref{sec:nikolaos}).

Το \textbf{ΠΕ0} είναι το αρχικό πακέτο εργασίας. Απογράφει πηγές, εμφανίσεις, ιστορικές σχέσεις, δικαιώματα και εκτελεστική κάλυψη. Περιλαμβάνει επίσης μικρό πιλοτικό έλεγχο οικογενειακών υπογραφών. Οι υπογραφές αναθεωρούνται μόνο σε υλικό ανάπτυξης και παγώνουν πριν από ανεξάρτητο έλεγχο αξιοπιστίας. Μόνο μετά τον έλεγχο αυτό παράγονται τα κύρια θετικά ζεύγη B2 και EISE. Το \textbf{ΠΕ1} συγκροτεί το κοινό συμβολικό επίπεδο αναφοράς. Η μεθοδολογική ιεραρχία και τα κριτήρια επάρκειας ακολουθούν αφού έχει οριστεί το μουσικό αντικείμενο και ο ρόλος κάθε ενότητας.

\section{Βιβλιογραφική θεμελίωση και ερευνητικό κενό}
\subsection{Θέση, οικογένεια και ιστορική σχέση}

Η ιστορική θεμελίωση που παρουσιάστηκε στην εισαγωγή συγκροτείται από συμπληρωματικές μαρτυρίες: ο Χρυσάφης (15ος αι., έκδ. Conomos, 1985b, σσ.~40--41) ορίζει τη θέση στο πλαίσιο της μελοποιητικής τέχνης· ο Γαβριήλ ιερομόναχος (15ος αι., έκδ. Hannick--Wolfram, 1985, σ.~72) χρησιμοποιεί τη γραμματική αναλογία της λέξης· ο Χρύσανθος, στο \emph{Θεωρητικόν Μέγα} (1832, §§400, 408), περιγράφει τη μετάδοση και τις διαφορετικές ενέργειες σημείων ανά περιβάλλον. Ο Αλυγιζάκης (1985, σσ.~81--82) συνδέει τη θέση με τη φράση και το πλήρες μέλος· ο Αντωνόπουλος (2013/2014, εικ.~5.36) περιγράφει χαρακτηριστικό τύπο με συνθετικές διαφοροποιήσεις. Ο Troelsgård (2011, §§32, 57, 67) και οι Αλεξάνδρου--Troelsgård (2020, σσ.~175--180) εξηγούν τη σημειογραφική και τροπική λειτουργία αυτών των σχηματισμών. Πρόκειται για μουσικολογικό πλαίσιο πολλών περιόδων, ενώ οι υπογραφές οικογενειών της παρούσας έρευνας προσδιορίζονται ειδικά στο εξεταζόμενο σώμα της Νέας Μεθόδου. Η αναλυτική παρουσίαση και οι εκκρεμείς έλεγχοι πηγών παραμένουν στο Παράρτημα~\ref{app:literature-full}.

Ο Δεσπότης (2007a), στο άρθρο «Ερμηνευτικές προσεγγίσεις στο μουσικό υλικό της ελληνικής ψαλτικής τέχνης», περιγράφει ειδικά ότι η λειτουργία των σημαδοφώνων εξαρτάται από τον συνδυασμό τους με γειτονικά στοιχεία και από τη θέση τους μέσα στη μουσική φράση· αυτή η σχεσιακή παρατήρηση παρέχει μουσικολογικό κίνητρο για τη σύγκριση τοπικής, φραστικής και ευρύτερης δομικής πληροφορίας.

Στη μελέτη του Αρβανίτη για τη θέση ως αρχή ομοιότητας (2010, τ.~Α΄, §5.2, σσ.~136--142), οι διαφορετικές εφαρμογές στο κείμενο, η σχέση μικρότερων με μεγαλύτερους μελικούς σχηματισμούς και η χρονική άρθρωση συνδέονται με το ερώτημα της αναγνωρίσιμης μουσικής μορφής. Αυτή η βιβλιογραφική γραμμή υποστηρίζει μια ανάλυση της θέσης που διακρίνει μελωδική ταυτότητα, σημειογραφική επιφάνεια, συλλαβική εφαρμογή και συγκεκριμένη χρονική πραγμάτωση. Οι παράμετροι μετατρέπονται σε ρητά διαγνωστικά γνωρίσματα μόνο όταν η μουσική πηγή και το υλικό της αντίστοιχης οικογένειας το δικαιολογούν.

Η χρονική οργάνωση αποτελεί χωριστό επίπεδο από τη σημειογραφία και το τροπικό σύστημα. Ο Χρύσανθος (1832, §§144--208) πραγματεύεται τον χρόνο και τον ρυθμό· ο Αρβανίτης (2010, τ.~Α΄, σσ.~136--142 και κεφ.~9) μελετά τη συμμετοχή τους στη μορφή κυρίως ειρμολογικών και στιχηραρικών σχηματισμών. Αυτά αιτιολογούν την εμπειρική σύγκριση του χρονικού στρώματος $T$ στην M1, χωρίς να επιβάλλουν κοινή ρυθμική υπογραφή σε κάθε οικογένεια παπαδικών θέσεων. Σημειογραφική χρονική αξία, μετρική ερμηνεία και εκτελεστικός χρονισμός παραμένουν διακριτά επίπεδα.

Μορφολογικές αναλύσεις των Χερουβικών καταγράφουν ονομασμένες θέσεις, επανεμφανίσεις, παραλλαγές και αλυσίδες θέσεων (Καραγκούνης, 2003, 2005). Δύο ομάδες παρατηρήσεων έχουν άμεση σημασία για τη M1. Στη σ.~788, παρόμοια μετροφωνική επιφάνεια μπορεί να αντιστοιχεί σε διαφορετικές θέσεις στην ιστορική γραφή. Στις σσ.~810--812 συγκρίνονται ονομασμένες θέσεις, εξηγήσεις και τροπικά συμφραζόμενα, με ρητή σταθερότητα της μεταγραφής της παρακλητικής στις περιπτώσεις της σ.~811. Οι παρατηρήσεις αυτές τροφοδοτούν τον σχεδιασμό δύσκολων περιπτώσεων, αλλά μια περίπτωση εισέρχεται ως M1 hard case μόνο αφού ελεγχθεί ότι η σχετική δυσκολία παραμένει στη συγκεκριμένη απόδοση της Νέας Μεθόδου. Στο Χερουβικό που παραδίδεται με το όνομα του Ιωάννη Δαμασκηνού και αναφέρεται στη σχετική βιβλιογραφία ως έργο του ψευδο-Ιωάννου Δαμασκηνού καταγράφονται επίσης αλυσίδες θέσεων που επανέρχονται με περιορισμένες αλλαγές στη σύνθεσή τους (Καραγκούνης, 2005, σσ.~157--163). Η οικογένεια θέσεων αφορά την ταυτότητα συγκεκριμένων πραγματώσεων μιας θέσης· η αλυσίδα αφορά τη διαδοχή θέσεων μέσα σε ένα έργο.

Η θεωρία των μουσικών σχημάτων προσφέρει συγκριτική γλώσσα για την ποικίλη εσωτερική σταθερότητα ενός τύπου. Οι Rabinovitch και Carter-Enyi (2024, §§1.5, 1.13--1.16) διακρίνουν σκελετικούς φθόγγους, περιγράμματα για μελωδική ανάπτυξη (\emph{outlines for diminution}) και περισσότερο σταθερές επιφανειακές φόρμουλες. Η σύγκριση αναδεικνύει ότι οι επιτρεπτές μεταβολές και η διαγνωστική βαρύτητα των γνωρισμάτων χρειάζονται εξειδίκευση ανά τύπο. Στην ψαλτική η εξειδίκευση στηρίζεται στις ίδιες τις μουσικές και σημειογραφικές πηγές, στις τροπικές σχέσεις και στην εφαρμογή του μέλους στο κείμενο.

Η \textbf{οικογένεια θέσεων} ορίζεται ήδη επιχειρησιακά στην εισαγωγή ως ελέγξιμη κατηγορία πραγματώσεων με τύπο-ειδικά διαγνωστικά γνωρίσματα, κανόνες περιβάλλοντος και ρητά όρια. Η αναφορά σε ιστορικό όνομα, κοινή τροπική λειτουργία ή κοινή μετροφωνική επιφάνεια λειτουργεί ως μουσικολογικό τεκμήριο και όχι ως αυτόματος κανόνας ένταξης. Το ΠΕ0 οριστικοποιεί τις υπογραφές πριν από ανεξάρτητη εφαρμογή τους, επιτρέποντας και μουσικολογικά τεκμηριωμένες εναλλακτικές διαμορφώσεις του ίδιου τύπου.

Η μεθοδολογία των tune families προσφέρει άμεσο συγκριτικό προηγούμενο για αυτή τη διαδικασία. Οι Volk και van Kranenburg (2012) χρησιμοποίησαν ρητά κριτήρια, μελωδία αναφοράς, επαναληπτική βελτίωση του πρωτοκόλλου, μέτρηση συμφωνίας και καταγραφή αμφιβολιών. Τα αποτελέσματά τους έδειξαν επίσης ότι η διαγνωστική ισχύς των χαρακτηριστικών μεταβάλλεται μεταξύ οικογενειών και ότι δεν υπάρχει ένα ενιαίο σύνολο χαρακτηριστικών που να λειτουργεί ως καθολικός πυρήνας για όλες. Για τον λόγο αυτό η παρούσα έρευνα αντιμετωπίζει την υπογραφή ως εμπειρικά ελεγχόμενο αναλυτικό εργαλείο και όχι ως εκ των προτέρων γνωσιακό ισχυρισμό. Η εργασία των Finkensiep κ.ά. (2020) ενισχύει την ίδια λογική: η σωστή τοπική μορφή δεν αρκεί πάντοτε για ασφαλή κατηγορική απόφαση και ορισμένες περιπτώσεις παραμένουν αναλυτικά αμφίσημες. Οι κατηγορίες \texttt{member}, \texttt{non\_member} και \texttt{unresolved} διατηρούν επομένως διαφορετικό τεκμηριωτικό ρόλο. Οι εργασίες των Karsdorp--van Kranenburg--Manjavacas (2019) και Nuttall κ.ά. (2021, 2022) παρέχουν συμπληρωματικά προηγούμενα για τη μάθηση συγγένειας και την αξιολόγηση προτύπων με γνώση της αντίστοιχης παράδοσης.


Οι ονομασμένες θέσεις, οι διαφοροποιήσεις της αναλυτικής γραφής και η επανεμφάνιση αλυσίδων παρέχουν συμπληρωματικά, μουσικά συγκεκριμένα κριτήρια. Στο Χερουβικό του ψευδο-Ιωάννου Δαμασκηνού περιγράφεται επαναλαμβανόμενη αλυσίδα θέσεων στο πρώτο τμήμα του έργου και διαφορετική οργάνωση μετά τη λειτουργική τομή της Μεγάλης Εισόδου (Καραγκούνης, 2005, σσ.~158--159). Το παράδειγμα αυτό στηρίζει τη χωριστή ανάλυση της διαδοχής και της προβλεψιμότητας των οικογενειών στο συμβολικό σώμα. Η ανάλυση αφορά τη μορφή του ρεπερτορίου, ενώ η ανθρώπινη M3c μεταβάλλει την πληροφορία λατρευτικού πλαισίου που δίνεται στον ακροατή για τον ίδιο μουσικό στόχο.

Η ειδική επιστημονική συμβολή οργανώνεται σε τρεις μουσικολογικούς άξονες. Πρώτον, εξετάζεται η ταυτότητα της θέσης μέσα στην παραλλαγή και συγκροτείται εκδοχικά ελεγχόμενο λεξιλόγιο τύπων με ανεξάρτητη επισημείωση και δύσκολα αντιπαραδείγματα. Δεύτερον, η M2 εξετάζει τι πρόσθετη πληροφορία προσφέρει η τεκμηριωμένη συν-εξήγηση για την οργάνωση της ενδοοικογενειακής παραλλαγής μέσα σε αντιστοιχισμένο πεδίο. Τρίτον, η M5 εξετάζει πώς οι επικυρωμένες οικογένειες επιλέγονται, διατάσσονται, τοποθετούνται λειτουργικά και πραγματώνονται μέσα σε πλήρη έργα. Οι τεχνικές δικλίδες της πρότασης υπηρετούν την αξιοπιστία αυτών των τριών ερωτημάτων.

Η ιστορική συγγένεια, η οικογενειακή ένταξη και η μελωδική ομοιότητα έχουν διαφορετική προέλευση. Οι θέσεις σε συγγενή ορθόδοξα μονοφωνικά ρεπερτόρια και η παραλλαγή μελωδικών τύπων παρέχουν συγκριτικό πλαίσιο, με διαφορετικά όμως συστήματα πηγών και μουσικής λειτουργίας (Amargianakis, 1977· Grigorjeva, 2023· Ryabinets, 2025· Tofan, 2024). Η εργασία των Janssen κ.ά. (2017) δείχνει ότι η παραλλαγή σε μεταδιδόμενα τραγούδια μπορεί να εξεταστεί συστηματικά σε κλίμακα σώματος. Στην παρούσα πρόταση οι επαληθευμένες εξηγητικές σχέσεις επιτρέπουν ειδικότερη σύγκριση: η B2 και ο EISE μοιράζονται την ίδια εποπτεία ως προς την οικογενειακή ένταξη και αντιστοιχισμένες συνθήκες εκπαίδευσης, ενώ διαφέρει η προέλευση των θετικών ζευγών. Η σύγκριση κρίνεται με ένα ανεξάρτητο εκτελεστικό κριτήριο αξιολόγησης που ορίζεται πριν από την τελική αξιολόγηση.

Στη διδακτορική του \textbf{Δημοσθένη Σπανουδάκη}, \emph{Σύγχρονα μοντέλα ανάλυσης Βυζαντινής Μουσικής: Το φαινόμενο της μελισματικότητας στα κεκραγάρια} (ΑΠΘ, 2017), η συγκριτική εξέταση κεκραγαρίων συνδυάζει παλαιογραφική, δομική και ποσοτική μουσικολογική ανάλυση. Η αναλυτική αποδελτίωση του έργου αναδεικνύει τη μετροφωνική δομή, τις θέσεις και τα μεγάλα σημάδια, το τροπικό περιβάλλον και την εξηγητική πρωτοβουλία ως κρίσιμες όψεις της εξήγησης (πρβλ. τ.~Α΄, σσ.~166--176· οι ακριβείς διατυπώσεις απαιτούν αυτοψία πριν από την τελική υποβολή). Οι όψεις αυτές προσφέρουν μουσικολογικό πλαίσιο για τη διάκριση διατηρούμενων σχέσεων, αναλυτικής ανάπτυξης και εξαρτήσεων από συγκεκριμένο εξηγητή. Η κύρια M2 εξετάζει την αξία τεκμηριωμένων σχέσεων συγκεκριμένων χωρίων μέσα σε ήδη ανεξάρτητα προσδιορισμένες οικογένειες.

Η εξήγηση προσφέρει συγκεκριμένες ιστορικές σχέσεις ανάμεσα σε αποδόσεις. Η ιστοριογραφική συζήτηση των Στάθη (1978), Καραγκούνη (2003), Αρβανίτη (2010, 2022), Gheorghiță (2010) και Lingas (2003) προσδιορίζει πώς τα τεκμήρια αυτά χρησιμοποιούνται και ποια ερωτήματα παραμένουν ανοιχτά για το παλαιότερο άκουσμα. Ο Gheorghiță χρησιμοποιεί αναδρομική συγκριτική αντιπαραβολή παλαιότερων και αναλυτικότερων γραφών. Η διαδικασία αυτή είναι μουσικολογική μέθοδος σύγκρισης και δεν ταυτίζεται με τον EISE. Στην παρούσα έρευνα προηγείται η επαλήθευση των ιστορικών σχέσεων και έπειτα οι σχέσεις αυτές χρησιμοποιούνται ως σχεσιακή εποπτεία. Η M2 δεν επιχειρεί ανακατασκευή ενός μοναδικού «αυθεντικού» ιστορικού ακούσματος. Η μουσική διατήρηση και η διαφοροποίηση εκτιμώνται χωριστά. Πλήρεις ιστορικές και βιβλιογραφικές διευκρινίσεις παρέχονται στο Παράρτημα~\ref{app:literature-full}.

\subsection{Σώματα, υπολογιστικά προηγούμενα και ακρόαση}

Τα Χερουβικά και τα Κοινωνικά, δύο κατηγορίες μελών της Θείας Λειτουργίας, παρέχουν πλήρη έργα, κοινά λειτουργικά κείμενα και ιστορικά διαφοροποιημένες αποδόσεις (Conomos, 1974, 1985a· Καραγκούνης, 2003, 2005· Gheorghiță, 2010). Η χωριστή παπαδική συλλογή της Ιεράς Μονής Αγίου Αντωνίου καταγράφει τύπους από 758 συνθέσεις και λειτουργεί ως καταλογογραφική αφετηρία και ασθενής εποπτεία. Η άλλη διαδικτυακή συλλογή με περισσότερες από 10.000 δηλωμένες φόρμουλες αφορά διαφορετικό ρεπερτοριακό πεδίο. Στο ΠΕ0 ελέγχεται αν κάθε εγγραφή μπορεί πράγματι να ενταχθεί στο σώμα και αν η διαδρομή της προς τις πηγές είναι επαληθεύσιμη.

Παλαιότερες και νεότερες υπολογιστικές προσεγγίσεις της βυζαντινής μουσικής καλύπτουν την ανάλυση ύμνων, την κωδικοποίηση σημειογραφίας και τη μοντελοποίηση μελωδικής σύνθεσης (Schiødt, 1971· Bouris κ.ά., 2018· Anton--Trăușan-Matu, 2018). Η ευρύτερη υπολογιστική βιβλιογραφία εξετάζει τη μελωδική τμηματοποίηση και την ανακάλυψη τμηματικών προτύπων (Cambouropoulos, 2006· Conklin--Anagnostopoulou, 2006), την επιλογή μονάδων σε μονοφωνικά σώματα (Cornelissen κ.ά., 2020a, 2020b· Lanz--Hajič, 2023, 2025), την κωδικοποίηση συμβολικής μουσικής (Fradet κ.ά., 2023· Kaliakatsos-Papakostas κ.ά., 2025) και την ανάκτηση μελωδικής συγγένειας. Ο Cambouropoulos (2006) εξετάζει ειδικότερα πώς οι επαναλαμβανόμενες μελωδικές δομές μπορούν να συνεισφέρουν στον εντοπισμό πιθανών ορίων, ενώ οι Conklin και Anagnostopoulou (2006) αναζητούν επαναλαμβανόμενα πρότυπα σε μελωδίες που έχουν ήδη αναπαρασταθεί ως ακολουθίες τμημάτων. Οι προσεγγίσεις αυτές παρέχουν υπολογιστικό υπόβαθρο για την οριοθέτηση και την ανακάλυψη επαναλαμβανόμενης τμηματικής δομής, χωρίς να ταυτίζονται με την αναγνώριση μουσικολογικά τεκμηριωμένων οικογενειών θέσεων. Η M1 συγκρίνει αυτές τις σχετικές δυνατότητες σε κοινές μουσικές εργασίες: γραφικά σημεία, γεγονότα και σύνθετες νευματικές μονάδες είναι εναλλακτικές κωδικοποιήσεις εισόδου, ενώ το μοντέλο γεγονότων, η επαγόμενη τμηματοποίηση και η ρητή θέση ορίζουν χωριστές συγκρίσεις δομικής κλίμακας.

Η διδακτορική του \textbf{André Holzapfel}, \emph{Similarity Methods for Computational Ethnomusicology} (Πανεπιστήμιο Κρήτης, 2010), παρέχει μεθοδολογικό προηγούμενο για αναζήτηση συγγενών μελωδικών φράσεων και αξιολόγηση διαφορετικών αναπαραστάσεων σε παραδοσιακά ρεπερτόρια. Οι συγκρίσεις ρυθμικών και μελωδικών γνωρισμάτων αναδεικνύουν τη σημασία της εμπειρικής επιλογής αναπαράστασης. Οι περιπτώσεις στις οποίες δύο ανακτώμενα αποσπάσματα μοιράζονται μόνο ένα μικρό μελωδικό υποτμήμα υποδεικνύουν τη σημασία του σωστού επιπέδου ανάλυσης: η ανάκτηση πραγματικής τοπικής ομοιότητας και η απόδοση κοινής ταυτότητας στη μεγαλύτερη μονάδα είναι δύο χωριστές κρίσεις. Η διάκριση αυτή υποστηρίζει το σχεδιασμό δύσκολων περιπτώσεων στις M1 και M3a.

Στη δυτική ιστορική και υπολογιστική μουσικολογία, ο Byros (2012) συνδέει ιστορικά σχήματα και τεκμήρια πρόσληψης στην ανάλυση της \emph{Eroica}, ενώ οι Park, Nam και Park (2020) μελετούν έργα ως προς ιστορικές αναφορές. Οι \textbf{Gilad Rabinovitch και Aaron Carter-Enyi}, στο άρθρο \emph{Melodic Organization and Sequential Ordering of Galant Schemata} (2024), αναλύουν την εσωτερική οργάνωση των galant schemata και τη διαδοχή τους. Στο επιλεγμένο σώμα 27 εκθέσεων μείζονος τρόπου από τη δεκαετία του 1740, η αρχική αναλυτική επισημείωση των σκελετικών φθόγγων και σχημάτων έγινε χειροκίνητα· ακολούθησαν υπολογιστική κατωφλίωση και οπτικοποίηση. Εξετάζονται και μη ειδικές σκελετικές ακολουθίες που δεν εντάσσονται σε επώνυμο σχήμα. Οι εμφανίσεις οργανώνονται αναλυτικά σε πέντε μορφικά στάδια και συγκρίνονται κοινές διαδρομές. Στη συνέχεια οι συγγραφείς εξετάζουν συγκεκριμένα έργα του C.~P.~E.~Bach. Η μετάβαση από εσωτερική οργάνωση τύπου σε διαδοχή, μορφική θέση και ερμηνεία συγκεκριμένου έργου αποτελεί ουσιαστικό πλαίσιο για τα μουσικολογικά ερωτήματα της M5. Οι ίδιοι αξιοποιούν τα μοντέλα φορμουλών της ελληνικής ψαλτικής του Mavromatis, συμμετέχοντας σε αμφίδρομο θεωρητικό διάλογο. Τα κριτήρια οικογενειακής ταυτότητας, οι κανόνες αναφοράς και οι μετρήσεις της παρούσας έρευνας ορίζονται ειδικά για το εξεταζόμενο παπαδικό σώμα.

Στη βυζαντινή μουσική, ο Panayotis Mavromatis, στο άρθρο \emph{A Hidden Markov Model of Melody Production in Greek Church Chant} (2005a) και στη διδακτορική του \emph{The Echoi of Modern Greek Church Chant in Written and Oral Transmission: A Computational Model and Its Cognitive Implications} (Eastman School of Music, University of Rochester, 2005b), θεμελιώνει πιθανοτική μοντελοποίηση της μελωδικής παραγωγής και της μουσικής δομής· στις εργασίες του 2009 αναπτύσσει περαιτέρω την επαγωγή μουσικής δομής από δεδομένα. Στην ανακοίνωση \emph{Hierarchical Structural Patterns in Improvised Music: Implications for Cognition} (SMT, 2019), για την οποία είναι διαθέσιμη η επίσημη δημοσιευμένη περίληψη, περιγράφει ιεραρχικά μοντέλα Markov όπου οι νότες συνδυάζονται σε μελωδικές φόρμουλες, οι φόρμουλες σε \emph{phrase prototypes} και αυτά σε μεγαλύτερες μορφικές ενότητες. Η ακολουθία αυτή παρέχει άμεσο θεωρητικό πλαίσιο για ερωτήματα διαδοχής και πρόβλεψης σε επίπεδο φόρμουλας, χωρίς να προδικάζει τον τρόπο με τον οποίο ορίζονται οι οικογένειες θέσεων της παρούσας έρευνας. Ο \textbf{Πολύκαρπος Γ. Πολυκαρπίδης}, στη διδακτορική του \emph{Υπολογιστική προσέγγιση της Βυζαντινής Μουσικής: Η περίπτωση του Ειρμολογίου του Θεοφάνη Καρύκη} (Εθνικό και Καποδιστριακό Πανεπιστήμιο Αθηνών, 2023), αναπτύσσει ιεραρχική υπολογιστική αναπαράσταση της Μεσοβυζαντινής σημειογραφίας, με σύνδεση μουσικού και κειμενικού υλικού, φρασεοποίηση και εφαρμογές συγκριτικής ανάλυσης. Σε συναφή εργασία, οι Polykarpidis, Kalofonos, Balageorgos και Anagnostopoulou (2022) συγκρίνουν τρία σώματα ειρμολογικού μέλους με κοινή ψηφιακή αναπαράσταση. Οι Panteli--Purwins (2013) συγκρίνουν ποσοτικά θεωρία και εκτελεστική πρακτική σε επίπεδο διαστημάτων και τροπικών βαθμίδων. Η παρούσα M5 εστιάζει σε διαφορετικό συνδυασμό: πλήρες έργο της Νέας Μεθόδου, ανάλυση σε μουσικολογικά τεκμηριωμένες οικογένειες θέσεων, ρητό σώμα αναφοράς με προέλευση και χρονολογικό έλεγχο, και χωριστό προφίλ επιλογής, διαδοχής, λειτουργικής τοποθέτησης και ενδοοικογενειακής πραγμάτωσης. Στη βιβλιογραφία που εξετάστηκε δεν εντοπίστηκε άμεσο ισοδύναμο αυτού του πλήρους συνδυασμού σε βυζαντινό μέλος· η διατύπωση της συμβολής περιορίζεται στον συγκεκριμένο σχεδιασμό, στις τέσσερις διαστάσεις χαρακτηρισμού και στα εξεταζόμενα δεδομένα.


Η επιλογή της μονάδας έχει συνέπειες πέρα από την επίδοση σε ένα έργο ανάκτησης. Ο έλεγχος πρέπει να καλύπτει όρια εμφανίσεων, συγγενείς πραγματώσεις, ανοικτή αναγνώριση οικογενειών και γενίκευση σε νέα έργα και πηγές. Η ρητή θέση, τα επαγόμενα τμήματα και το μοντέλο μεμονωμένων γεγονότων αξιολογούνται με κοινούς στόχους. Η ανεξάρτητη μουσικολογική δοκιμή ελέγχει συγχρόνως αν διατηρούνται οι ιστορικά τεκμηριωμένες διακρίσεις ανάμεσα σε δύσκολα θετικά και δύσκολα αρνητικά παραδείγματα.

Η ανθρώπινη μελέτη στηρίζεται σε δύο βιβλιογραφικές γραμμές. Η πρώτη αφορά την οργάνωση μελωδικών οικογενειών στην ακρόαση (Volk--van Kranenburg, 2012) και τη μελέτη της μελωδικής ταυτότητας υπό μεταβολές της επιφάνειας (Schindler κ.ά., 2013). Η δεύτερη συνδέει την πιθανοτική μουσική προσδοκία με την αντίληψη ορίων, τον εκφραστικό χρονισμό και την επεξεργασία μελωδικών ακολουθιών (Pearce κ.ά., 2010· Gingras κ.ά., 2016· Hansen κ.ά., 2021· Koelsch κ.ά., 2013· Herff κ.ά., 2021· Di Liberto κ.ά., 2020). Οι Kern κ.ά. (2022) παρέχουν επιπλέον λόγο να διακρίνονται η μακροχρόνια μάθηση του μοντέλου και το μήκος του άμεσου μουσικού συμφραζομένου. Η M3c εξετάζει ειδικά αν το πληροφοριακό περιεχόμενο της επόμενης οικογένειας εξηγεί ανθρώπινες κρίσεις πέρα από το πληροφοριακό περιεχόμενο των μεμονωμένων γεγονότων.

\subsection{Το ακριβές ερευνητικό κενό}

Οι σχετικές εργασίες έχουν προσφέρει μορφολογικές περιγραφές θέσεων, μεθόδους μελωδικής τμηματοποίησης και συγγένειας, υπολογιστική μελέτη της μετάδοσης, μοντέλα μουσικής προσδοκίας και, σε άλλα ρεπερτόρια, ποσοτικό χαρακτηρισμό πλήρων έργων ως προς ιστορικά σώματα. Το ακριβές κενό της πρότασης έχει δύο διαδοχικά μέρη. Πρώτον, εξετάζεται στην \emph{ίδια μουσική μονάδα} πώς συνδέονται η αναπαράσταση, η προέλευση της εποπτείας, η μουσικολογική συνοχή και οι ανεξάρτητοι ανθρώπινοι και εκτελεστικοί έλεγχοι. Δεύτερον, αφού η μονάδα επικυρωθεί, εξετάζεται αν μπορεί να στηρίξει ερμηνεύσιμο χαρακτηρισμό πλήρους βυζαντινού έργου ως προς τεκμηριωμένο περιβάλλον αναφοράς, με χωριστή αποτίμηση του λεξιλογίου οικογενειών, της διαδοχής τους, της λειτουργικής τοποθέτησης και της πραγμάτωσης μέσα στην ίδια οικογένεια.

\paragraph{Μονάδα αναπαράστασης.} Οι έρευνες τμηματοποίησης σε μονοφωνικά σώματα και η μάθηση συγγένειας σε οικογένειες μελωδιών προσφέρουν ισχυρά εναλλακτικά μοντέλα (Cornelissen κ.ά., 2020a, 2020b· Lanz--Hajič, 2025· Karsdorp--van Kranenburg--Manjavacas, 2019). Η αναγνώριση θεωρητικά ορισμένων σχημάτων στη δυτική πολυφωνία παρέχει ένα επιπλέον συναφές μεθοδολογικό προηγούμενο (Gjerdingen, 2007· Katsiavalos κ.ά., 2019· Finkensiep κ.ά., 2020· Neuwirth κ.ά., 2023). Η M1 εξετάζει τη θέση μεταβλητής έκτασης απέναντι σε μονάδες γεγονότος και επαγόμενα τμήματα, με τις ίδιες δοκιμασίες αξιολόγησης και δύσκολες περιπτώσεις που ορίζει η μουσικολογική τεκμηρίωση: συγγενή επιφάνεια με διαφορετική ταυτότητα και διαφορετική μορφή με κοινή ταυτότητα.

\paragraph{Προέλευση της εποπτείας.} Η υπολογιστική μελέτη της ιστορικής μετάδοσης έχει χρησιμοποιήσει συγγενείς μελωδίες, στοίχιση και φυλογενετικές προσεγγίσεις για να περιγράψει μεταβολές και σχέσεις πηγών (Savage κ.ά., 2022· Hajič κ.ά., 2023, 2025· Ballen κ.ά., 2026). Πειράματα αλυσιδωτής προφορικής μετάδοσης εξετάζουν χωριστά μηχανισμούς μεταβολής της μελωδίας (Anglada-Tort κ.ά., 2023· Popescu--Rohrmeier, 2024). Η M2 θέτει διαφορετικό ελεγχόμενο ερώτημα: τι μεταβάλλεται στην ανεξάρτητη εξωτερική αξιολόγηση όταν τα θετικά ζεύγη εκπαίδευσης προέρχονται από επαληθευμένες εξηγητικές σχέσεις, με ίδια εποπτεία ως προς την οικογενειακή ένταξη και αντιστοιχισμένους τους άλλους όρους σε σχέση με τη B2;

\paragraph{Μουσικολογική γνώση, εκτέλεση και πλήρες έργο.} Η μορφολογική βιβλιογραφία των Χερουβικών τεκμηριώνει ονομασμένες θέσεις, επανεμφανίσεις και αλυσίδες θέσεων (Καραγκούνης, 2003, 2005). Η σύγκριση θεωρίας και σύγχρονης εκτέλεσης έχει επίσης εξεταστεί ποσοτικά σε επίπεδο διαστημάτων και τροπικών βαθμίδων (Panteli--Purwins, 2013). Η προφορική μετάδοση, η διδασκαλία και η ερμηνεία του μέλους έχουν δική τους μουσικολογική βιβλιογραφία (Karas, 1982· Lind, 2011· Χατζάρα, 2015· Κανάκης, 2020). Η παρούσα πρόταση συγκροτεί λεξικό με εμφανίσεις και αμφίσημες περιπτώσεις, συγκρίνει τη μορφολογία της παραλλαγής σε κλίμακα ρεπερτορίου, εξετάζει τι διατηρούν συγκεκριμένες εξηγήσεις σε επίπεδο θέσης και συνδέει την ιστορική οργάνωση με χωριστά μετρημένη εκτελεστική μεταβλητότητα. Η M5 προσθέτει επίπεδο πλήρους έργου: το έργο περιγράφεται με χάρτη επικυρωμένων εμφανίσεων και τεκμηριωμένων διαδοχών και χαρακτηρίζεται ως προς ρητά δηλωμένο σώμα αναφοράς, με σύγκριση σε επίπεδο ανεξάρτητων μελοποιών ως προς το δηλωμένο σώμα αναφοράς. Η πρώτη εφαρμογή αφορά τον Νικόλαο (Τόμα, 2018) και αξιοποιεί, όπου επαληθεύονται και παραμένουν ανεξάρτητες από την ανάπτυξη του συστήματος, δημοσιευμένες μουσικολογικές αντιστοιχίσεις συγκεκριμένων έργων.

\paragraph{Ανθρώπινη ακρόαση.} Η κατηγοριοποίηση συγγενών μελωδιών και η πιθανοτική μουσική προσδοκία έχουν δικά τους ερευνητικά προηγούμενα (Volk--van Kranenburg, 2012· Pearce κ.ά., 2010· Koelsch κ.ά., 2013· Herff κ.ά., 2021· Di Liberto κ.ά., 2020). Οι δύο ανθρώπινες υποθέσεις ελέγχουν αν οι τεκμηριωμένες οικογένειες θέσεων οργανώνουν τις κρίσεις των ακροατών πέρα από την επιφανειακή ομοιότητα και αν το πληροφοριακό περιεχόμενο σε επίπεδο οικογένειας εξηγεί πόσο απρόσμενη ακούγεται η συνέχεια πέρα από το επίπεδο γεγονότος. Ο διαχωρισμός των τεκμηρίων συγκρότησης από τα σύνολα αξιολόγησης καθιστά τις τέσσερις πλευρές του προβλήματος συγκρίσιμες χωρίς να συγχέει τους συμπερασματικούς τους ρόλους.

\paragraph{Θέση στο MIR και στη μηχανική μάθηση.} Το MIR στρέφεται συστηματικότερα σε μουσικές παραδόσεις εκτός του δυτικού κανόνα. Η αξιολόγηση μοντέλων θεμελίωσης σε σώματα από διάφορες παραδόσεις δείχνει ότι αναπαραστάσεις που μαθαίνονται κυρίως από δυτική μουσική μεταφέρονται άνισα (Papaioannou κ.ά., 2025). Στο λατινικό λειτουργικό μέλος, ενιαία σώματα και εργαλεία κάνουν πλέον εφικτή την επαναλήψιμη υπολογιστική έρευνα (Dvořáková κ.ά., 2026· Lanz--Hajič, 2025). Η ψαλτική της Νέας Μεθόδου απουσιάζει από αυτές τις αξιολογήσεις και δεν διαθέτει, εξ όσων εντοπίστηκε, σώμα αξιολόγησης σε επίπεδο θέσης.

Στη μηχανική μάθηση, το ερώτημα της πηγής των θετικών ζευγών έχει ήδη τεθεί. Η αντιστοίχιση παρατηρήσεων του ίδιου υποκείμενου αντικειμένου έχει μελετηθεί ως εναλλακτική της αναλλοιωτότητας υπό κοινή κατηγορία (Mahajan κ.ά., 2021), και η μεταφορά τρόπων ενδοκατηγορικής μεταβολής σε νέες κατηγορίες έχει χρησιμοποιηθεί όταν τα παραδείγματα είναι λίγα (Hariharan--Girshick, 2017). Και στις δύο περιπτώσεις η σχέση «ίδιο αντικείμενο» δίνεται από τη συλλογή ή υποτίθεται. Η εξήγηση της Νέας Μεθόδου προσφέρει κάτι σπανιότερο: ιστορικά τεκμηριωμένες παράλληλες αποδόσεις του ίδιου χωρίου, καταγεγραμμένες ανεξάρτητα από την οικογενειακή ταξινόμηση, σε μουσικές ακολουθίες μεταβλητού μήκους. Η διατριβή αξιοποιεί αυτή την ιδιότητα σε δύο ελεγχόμενες συγκρίσεις. Η πρώτη εξετάζει αν η τεκμηριωμένη σχέση δίνει αναπαράσταση με μεγαλύτερη εξωτερική εγκυρότητα από την κοινή οικογένεια (M2). Η δεύτερη εξετάζει αν οι μεταβολές που μαθαίνονται από αυτήν μεταφέρονται σε οικογένειες χωρίς αντίστοιχα ζεύγη (§~\ref{sec:atl}). Η τμηματική μάθηση από μερικές και εναλλακτικές επισημειώσεις συγγενεύει με τα CRF υπό περιορισμούς (Papay κ.ά., 2022). Εδώ προσαρμόζεται σε εναλλακτικά αποδεκτά όρια μουσικών μονάδων με ανοικτό σύνολο οικογενειών.

\section{Ερευνητικό ερώτημα και υποθέσεις}

\paragraph{Οδηγός αναγνώστη για τα στάδια και τις συνθήκες.}
Οι συμβολισμοί \textbf{M1--M5} δηλώνουν τις πέντε ενότητες της έρευνας: επιλογή αναπαράστασης (M1), σύγκριση εκπαίδευσης με ιστορικά εξηγητικά ζεύγη (M2), ανεξάρτητη μουσικολογική και ανθρώπινη επικύρωση (M3a--M3c), ανάλυση σύγχρονης εκτελεστικής μεταβλητότητας (M4) και χαρακτηρισμό πλήρων έργων (M5). Το \textbf{ΠΕ0} είναι η αρχική απογραφή και δοκιμή των τεκμηρίων και των κανόνων οικογενειακής ένταξης, πριν από τη χρήση τους ως δεδομένων εκπαίδευσης. Για την M2, \textbf{B0} σημαίνει άμεση συμβολική περιγραφή χωρίς μάθηση νευρωνικής αναπαράστασης, \textbf{B1} κοινή νευρωνική υποδομή χωρίς πρόσθετη σχεσιακή εποπτεία, \textbf{B2} μάθηση με ζεύγη επικυρωμένων μελών της ίδιας οικογένειας και \textbf{EISE} τον ίδιο κωδικοποιητή με θετικά ζεύγη από επαληθευμένες ιστορικές συν-εξηγήσεις. Η κύρια ερευνητική σύγκριση αφορά B2 έναντι EISE: ίδια οικογενειακή πληροφορία, διαφορετική τεκμηριωμένη προέλευση των σχέσεων εκπαίδευσης. Οι λεπτομέρειες και οι συνθήκες αντιστοίχισης ορίζονται στην §~\ref{sec:eise}.


Το κεντρικό ερώτημα είναι ποια αναπαράσταση διαχωρίζει διαφορετικές θέσεις παρά την επιφανειακή τους συγγένεια και αναγνωρίζει κοινή ταυτότητα παρά τις τεκμηριωμένες μεταβολές των πραγματώσεών τους. Στο επίπεδο εφαρμογής τίθεται ένα δεύτερο ερώτημα: πώς χαρακτηρίζεται ένα πλήρες έργο ως προς ένα τεκμηριωμένο μουσικό περιβάλλον αναφοράς όταν η ανάλυση γίνεται στο επίπεδο αυτών των οικογενειών και όχι μόνο των μεμονωμένων μελωδικών γεγονότων; Οι επόμενες συγκρίσεις χρησιμοποιούν λίγα σταθερά σύμβολα, που συνοψίζονται εδώ πριν από τις υποθέσεις.

Κάθε κύρια σύγκριση συνδέεται με συγκεκριμένο \textbf{εκτιμώμενο μέγεθος (estimand)}: την ποσότητα που το αντίστοιχο πείραμα έχει σχεδιαστεί να εκτιμήσει. Η σαφής διάκριση των εκτιμώμενων μεγεθών περιορίζει το εύρος κάθε συμπεράσματος στο ερώτημα και στα δεδομένα που το υποστηρίζουν.

Για τις αναλύσεις μουσικής πρόβλεψης χρησιμοποιείται το πληροφοριακό περιεχόμενο ή surprisal μιας παρατηρημένης επιλογής,
\[
I=-\log_2 P.
\]
Μια επιλογή στην οποία το μοντέλο αποδίδει μεγάλη πιθανότητα έχει μικρό πληροφοριακό περιεχόμενο, ενώ μια λιγότερο πιθανή επιλογή έχει μεγαλύτερο. Στη συνέχεια το $I_F$ εφαρμόζει την έννοια στο επίπεδο οικογένειας θέσεων, το $I_{\mathrm{event}}$ στο επίπεδο μελωδικού γεγονότος και το $I_{\Omega}$ στη M5 ως προς το δηλωμένο περιβάλλον αναφοράς $\Omega$.

\par\noindent\textbf{Βασικά σύμβολα για την ανάγνωση.}\par\smallskip
\noindent{\small
\begin{tabularx}{\linewidth}{@{}>{\RaggedRight\arraybackslash}p{0.22\linewidth}Y@{}}
\toprule
\textbf{Σύμβολο} & \textbf{Τι δηλώνει}\tabularnewline
\midrule

$R^\star$ & Αναπαράσταση που επιλέγεται στα δεδομένα ανάπτυξης.\tabularnewline
$S_m$ & Συμφωνία του ιστορικού προφίλ του μοντέλου $m$ με το ανεξάρτητο εκτελεστικό προφίλ.\tabularnewline
$\Delta_{\mathrm{hist}}$ & Διαφορά $S_{\mathrm{EISE}}-S_{B2}$: αξία της προέλευσης της εποπτείας.\tabularnewline
$\Delta_{\mathrm{learn}}$ & Διαφορά $S_{\mathrm{EISE}}-S_{B0}$: σύγκριση με άμεση συμβολική περιγραφή.\tabularnewline
$H_{\mathrm{pred}},H_{\mathrm{cat}}$ & Ανθρώπινες υποθέσεις για μουσική προσδοκία και κατηγορική οργάνωση.\tabularnewline
$I_F,I_{\mathrm{event}}$ & Πληροφοριακό περιεχόμενο σε επίπεδο οικογένειας και γεγονότος.\tabularnewline
$R$ στη M3c & Μη μουσική ένδειξη λατρευτικού πλαισίου· χωριστή μεταβλητή από την $R^\star$.\tabularnewline
$\Omega$ στη M5 & Τεκμηριωμένο περιβάλλον αναφοράς: δηλωμένο σώμα έργων, οριστικοποιημένη αναπαράσταση, μοντέλο πιθανότητας και κανόνες ένταξης στο σώμα.\tabularnewline
$I_{\Omega}$ & Πληροφοριακό περιεχόμενο μιας επιλογής οικογένειας ως προς το συγκεκριμένο τεκμηριωμένο περιβάλλον αναφοράς $\Omega$.\tabularnewline
$\Delta_{\mathrm{ctx}}$ & Διαφορά των δύο συντελεστών αλληλεπίδρασης με την ένδειξη $R$.\tabularnewline
\end{tabularx}
}\par

Η M1 επιλέγει την αναπαράσταση $R^\star$ μόνο στα δεδομένα ανάπτυξης. Η M2 συγκρίνει την EISE με την αντιστοιχισμένη οικογενειακή συνθήκη B2 στο ανεξάρτητο εκτελεστικό κριτήριο αξιολόγησης. Η κύρια αντίθεση είναι $\Delta_{\mathrm{hist}}=S_{\mathrm{EISE}}-S_{B2}$· η $\Delta_{\mathrm{learn}}=S_{\mathrm{EISE}}-S_{B0}$ ανήκει στη χωριστή οικογένεια συγκρίσεων της εκπαίδευσης.

Το $\Delta_{\mathrm{hist}}$ αφορά αποκλειστικά το προκαθορισμένο ανεξάρτητο κριτήριο αξιολόγησης της M2. Θετική τιμή αντιστοιχεί σε μεγαλύτερη συμφωνία της EISE συνθήκης με αυτό το κριτήριο σε σχέση με την B2, τιμή κοντά στο μηδέν σε μικρή ή μη διακριτή πρόσθετη διαφορά υπό το συγκεκριμένο πρωτόκολλο, και αρνητική τιμή σε μεγαλύτερη συμφωνία της B2. Όταν η κάλυψη ή η αξιοπιστία των απαιτούμενων δεδομένων δεν ικανοποιεί το προκαθορισμένο κριτήριο επάρκειας, δηλώνεται ότι το μέγεθος δεν μπορεί να εκτιμηθεί αξιόπιστα. Η ερμηνεία αφορά την πληροφοριακή αξία της προέλευσης των θετικών σχέσεων υπό αντιστοιχισμένη εκπαίδευση και το συγκεκριμένο ανεξάρτητο κριτήριο αξιολόγησης.



Η M3a εξετάζει την ταυτότητα, τα όρια και την παραλλαγή των οικογενειών με ανεξάρτητες μουσικολογικές κρίσεις. Ο συντελεστής $\beta_F^{\mathrm{cat}}$ αποτυπώνει τη συμβολή της οικογενειακής ένταξης στην τριαδική επιλογή, με συνεκτίμηση της επιφανειακής απόστασης. Ο $\beta_F^{\mathrm{pred}}$ αποτυπώνει τη σχέση του $I_F$ με την ανθρώπινη κρίση για το πόσο απρόσμενη είναι η συνέχεια, αφού ληφθεί υπόψη το $I_{\mathrm{event}}$. Οι δύο ανθρώπινες επιβεβαιωτικές υποθέσεις διατυπώνονται ως
\[
H_{\mathrm{cat}}:\ \beta_F^{\mathrm{cat}}>0,
\qquad
H_{\mathrm{pred}}:\ \beta_F^{\mathrm{pred}}>0.
\]
Η M3b ελέγχει αν οι ακροατές οργανώνουν διαφορετικές πραγματώσεις σύμφωνα με τις οικογένειες πέρα από την επιφανειακή ομοιότητα. Η M3c ελέγχει αν το πληροφοριακό περιεχόμενο σε επίπεδο οικογένειας προσθέτει εξηγητική πληροφορία για το πόσο απρόσμενη ακούγεται η συνέχεια πέρα από το επίπεδο γεγονότος. Στους δύο στατιστικούς ελέγχους των ανθρώπινων υποθέσεων εφαρμόζεται η διόρθωση Holm, ώστε να ελέγχεται η πιθανότητα ψευδούς θετικού συμπεράσματος στην κοινή επιβεβαιωτική οικογένεια (§~\ref{app:stats-multiplicity}).

Η ένδειξη λατρευτικού πλαισίου $R$ επιτρέπει την κύρια δευτερεύουσα αντίθεση $\Delta_{\mathrm{ctx}}$, η οποία συγκρίνει τη μεταβολή της σχέσης της ανθρώπινης κρίσης με τις δύο τυποποιημένες μεταβλητές πληροφοριακού περιεχομένου. Ο έλεγχος γενίκευσης της οικογενειακής διάκρισης σε νέες συνθήκες (Cross-Condition Generalization Performance, CCGP· Bernardi κ.ά., 2020· πρωτόκολλο στην §~\ref{app:stats-m3b}), η ευθυγράμμιση των αναπαραστάσεων μέσα στην ίδια οικογένεια και οι αλληλεπιδράσεις εμπειρίας είναι κύρια δευτερεύοντα αποτελέσματα. Η βιωματική απόκριση είναι διερευνητική. Οι ευρύτερες M4 και M5, καθώς και η υπό προϋποθέσεις επέκταση της M3c σε μετρήσεις εγκεφαλικής δραστηριότητας (M3ν, §~\ref{sec:m3nu}), ακολουθούν την ιεραρχία της §~\ref{sec:core-hierarchy}. Ο χάρτης υποθέσεων, μεγεθών και καθεστώτων βρίσκεται στο Παράρτημα~\ref{app:estimand-map}.

\section{Κεντρική μεθοδολογική αρχιτεκτονική}

\subsection{Από τις πηγές στη γνωσιακή λειτουργία}

Η μεθοδολογία οργανώνεται σε μία ακολουθία εξαρτήσεων. Το αρχικό πακέτο εργασίας ΠΕ0 απογράφει τις πραγματικές πηγές, τα δικαιώματα, τις υποψήφιες οικογενειακές και εξηγητικές σχέσεις, την εκτελεστική κάλυψη και τη δυνατότητα συλλογής δεδομένων από συμμετέχοντες και καθορίζει τους διαχωρισμούς ανάπτυξης και ανεξάρτητης αξιολόγησης. Πριν από τις M1--M2 οριστικοποιείται αποκλειστικά στο υλικό ανάπτυξης το λεξιλόγιο που απαιτείται για την εκπαίδευση και την επιλογή αναπαράστασης. Η M1 οριστικοποιεί την κλίμακα αναπαράστασης και η M2 εκπαιδεύει και συγκρίνει B2/EISE πάνω στο ίδιο λεξιλόγιο ανάπτυξης. Η M3a παρέχει την ανεξάρτητη μουσικολογική επικύρωση, με τις κρίσεις των ειδικών διαχωρισμένες από την ανάπτυξη και χωρίς ανατροφοδότηση προς τα επιβεβαιωτικά αποτελέσματα των M1--M2. Η M3c προηγείται της M3b στη συλλογή από τους ίδιους συμμετέχοντες. Η σειρά διατηρεί αυθόρμητη την πρώτη κρίση μουσικής προσδοκίας.

Η ακολουθία αυτή δεν είναι μόνο χρονική· αντιστοιχεί και σε διαφορετικούς ρόλους των δεδομένων. Τα δεδομένα \textbf{ανάπτυξης} χρησιμοποιούνται για επιλογές όπως η αναπαράσταση και οι υπερπαράμετροι. Μετά την οριστικοποίηση αυτών των επιλογών, η αξιολόγηση πραγματοποιείται σε έργα, πηγές ή άλλες μονάδες που έχουν διατηρηθεί εκτός της αντίστοιχης διαδικασίας ανάπτυξης. Η ανεξάρτητη M3a προσθέτει μουσικολογική αξιολόγηση χωρίς αναδρομική αλλαγή των επιβεβαιωτικών αποτελεσμάτων M1--M2. Με αυτόν τον διαχωρισμό, η κατασκευή ενός μοντέλου, η επιλογή του και η τελική αξιολόγησή του βασίζονται σε διακριτούς ρόλους δεδομένων.

Οι κύριες εξαρτήσεις και οι υπό προϋποθέσεις επεκτάσεις συνοψίζονται στο Σχήμα~\ref{fig:experimental-architecture}. Στο σχήμα, EEG/MEG σημαίνει ηλεκτροεγκεφαλογραφία/μαγνητοεγκεφαλογραφία, δηλαδή καταγραφή νευρικής δραστηριότητας, και EGG σημαίνει ηλεκτρογλωττογραφία, δηλαδή καταγραφή της φωνητικής επαφής. Πρόκειται για διακριτές, υπό προϋποθέσεις επεκτάσεις της M3c και της M4 αντιστοίχως.

\begin{figure}[htbp]
\centering
\begin{tikzpicture}[
box/.style={draw=acc,rounded corners,align=center,font=\small,inner sep=4pt,text width=25mm,minimum height=9mm},
smallbox/.style={draw=acc,rounded corners,align=center,font=\scriptsize,inner sep=3pt,text width=25mm,minimum height=9mm},
arr/.style={-{Latex[length=2mm]},thick},
gatearr/.style={-{Latex[length=2mm]},thick,dashed}
]
\node[box] (pe) at (0,0) {ΠΕ0/ΠΕ1\\Πηγές και λεξιλόγιο};
\node[box] (m1) at (3.25,0) {M1\\Αναπαράσταση};
\node[box] (m2) at (6.5,0) {M2\\B2--EISE};
\node[smallbox] (minm4) at (9.75,0) {Ανεξάρτητο εκτελεστικό\\κριτήριο M4};
\draw[arr] (pe)--(m1);
\draw[arr] (m1)--(m2);
\draw[arr] (minm4)--(m2);

\node[box] (m3a) at (3.25,-1.55) {M3a\\Ανεξάρτητη επικύρωση};
\node[box] (m3c) at (6.5,-1.55) {M3c\\$H_{\mathrm{pred}}$};
\node[box] (m3b) at (9.75,-1.55) {M3b\\$H_{\mathrm{cat}}$};
\draw[arr] (m1)--(m3a);
\draw[arr] (m3a)--(m3c);
\draw[arr] (m3c)--(m3b);

\node[box] (validated) at (3.25,-3.1) {Επικυρωμένο\\πλαίσιο};
\node[box] (broadm4) at (6.5,-3.1) {Ευρύτερη M4\\Εκτέλεση};
\node[box] (m5) at (9.75,-3.1) {M5\\Νικόλαος};
\draw[arr] (m3a)--(validated);
\draw[arr] (validated)--(broadm4);
\draw[arr] (validated.south) -- ++(0,-10mm) -| (m5.south);

\node[smallbox] (neural) at (13,-1.55) {M3ν\\EEG/MEG};
\node[smallbox] (egg) at (13,-3.1) {EGG\\διερευνητικό M4};
\draw[gatearr] (m3c.north) -- ++(0,5mm) -| (neural.north);
\draw[gatearr] (broadm4.south) -- ++(0,-5mm) -| (egg.south);
\end{tikzpicture}
\caption{Ιεραρχία και εξαρτήσεις του προγράμματος. Η M2 απαιτεί το αναγκαίο ανεξάρτητο εκτελεστικό κριτήριο αξιολόγησης από τη M4. Η M3a αξιολογεί ανεξάρτητα το μουσικολογικό λεξιλόγιο. Το βέλος M3c προς M3b δηλώνει τη σειρά συλλογής των ανθρώπινων δεδομένων· οι διακεκομμένες συνδέσεις δηλώνουν ενεργοποίηση υπό προϋποθέσεις. Η ευρύτερη M4 και η M5 αξιοποιούν το επικυρωμένο πλαίσιο.}
\label{fig:experimental-architecture}
\end{figure}

\FloatBarrier

\subsection{Υπολογιστικά εργαλεία: τι κάνει και τι λύνει το καθένα}\label{sec:tool-rationale}

Τα υπολογιστικά εργαλεία επιλέγονται επειδή απαντούν σε διαφορετικό μέρος του ίδιου προβλήματος. Ο ρόλος τους συνοψίζεται εδώ, ώστε οι τεχνικές λεπτομέρειες των επόμενων ενοτήτων να διαβάζονται με γνωστό σκοπό.

{\footnotesize
\begin{xltabular}{\linewidth}{@{}>{\RaggedRight\arraybackslash}p{0.19\linewidth}>{\RaggedRight\arraybackslash}p{0.25\linewidth}Y@{}}
\toprule
\textbf{Εργαλείο} & \textbf{Τι κάνει} & \textbf{Γιατί χρησιμοποιείται εδώ}\tabularnewline
\midrule
\endfirsthead
\toprule
\textbf{Εργαλείο} & \textbf{Τι κάνει} & \textbf{Γιατί χρησιμοποιείται εδώ}\tabularnewline
\midrule
\endhead
\bottomrule
\endlastfoot
\textbf{Ενιαίο συμβολικό επίπεδο αναφοράς} & Μετατρέπει διαφορετικές πηγές σε κοινή ακολουθία μουσικών γεγονότων, διατηρώντας σύνδεση με τα αρχικά σημεία. & Επιτρέπει σε όλες τις αναπαραστάσεις να συγκρίνονται στα ίδια δεδομένα και σε κάθε υπολογιστικό αποτέλεσμα να επιστρέφει στην πηγή.\tabularnewline
\midrule
\textbf{Transformer γεγονότων} & Νευρωνικό μοντέλο ακολουθιών που προβλέπει κάθε γεγονός από το προηγούμενο συμφραζόμενο και μπορεί να μάθει μακρύτερες εξαρτήσεις χωρίς ρητή ετικέτα «θέσης». & Αποτελεί βασικό μοντέλο αναφοράς της M1: αν η ρητή θέση δεν προσθέτει πληροφορία πέρα από ένα μοντέλο αναφοράς σε επίπεδο μελωδικού γεγονότος, αυτό είναι ουσιαστικό αποτέλεσμα.\tabularnewline
\midrule
\textbf{Επαγόμενη τμηματοποίηση} & μη παραμετρικός πιθανοτικός μηχανισμός που αναζητεί επαναλαμβανόμενα τμήματα χωρίς να του δίνονται εκ των προτέρων τα όρια των θέσεων. Η κύρια υλοποίηση βασίζεται σε διαδικασία Pitman--Yor, δηλαδή μοντέλο που επιτρέπει το λεξιλόγιο τμημάτων να αυξάνεται όσο αυξάνονται τα δεδομένα. & Ελέγχει αν οι μονάδες που αναγνωρίζει η παράδοση συμπίπτουν με μονάδες που προκύπτουν μόνο από τη στατιστική επανάληψη. Συγγενής μη παραμετρική πιθανοτική τμηματοποίηση έχει χρησιμοποιηθεί στο γρηγοριανό μέλος (Lanz--Hajič, 2025).\tabularnewline
\midrule
\textbf{EISE} & Νευρωνικός κωδικοποιητής τμημάτων μεταβλητού μήκους με χωριστή προβολή οικογενειακής ταυτότητας και τοπικές καταστάσεις για τη συγκεκριμένη πραγμάτωση. & Σχεδιάζεται για την αναγνώριση της οικογενειακής ταυτότητας και τη διατήρηση της πληροφορίας για τη συγκεκριμένη πραγμάτωση. Η διατήρηση της τοπικής πληροφορίας ελέγχεται χωριστά πριν από τη χρήση της στις M2, M4 και M5.\tabularnewline
\midrule
\textbf{Ημι-Μαρκοβιανό CRF} & Στη μη επικαλυπτόμενη διαδρομή, μοντέλο τμημάτων που αποφασίζει από κοινού πού αρχίζει και τελειώνει ένα τμήμα μεταβλητού μήκους και σε ποια οικογένεια αποδίδεται. & Τα όρια της θέσης είναι μέρος του προβλήματος και η επισημείωση είναι μερική. Η καταλληλότητα της μη επικαλυπτόμενης μορφής εξόδου ελέγχεται στο ΠΕ0 πριν από τη M1.\tabularnewline
\midrule
\textbf{Αντιθετική/μετρική μάθηση} & Εκπαίδευση με ζεύγη ή ομάδες όπου το μοντέλο μαθαίνει ποιες πραγματώσεις πρέπει να βρίσκονται κοντά και ποιες να διαχωρίζονται. & Επιτρέπει να αλλάζει μόνο η προέλευση των θετικών ζευγών ανάμεσα σε B2 και EISE, κρατώντας την υπόλοιπη εκπαίδευση αντιστοιχισμένη.\tabularnewline
\midrule
\textbf{B0/B1/B2} & Συνθήκες ελέγχου: άμεση συμβολική περιγραφή χωρίς μάθηση αναπαράστασης (B0), κοινή νευρωνική υποδομή χωρίς τον πρόσθετο σχεσιακό όρο της B2 (B1), και αντιστοιχισμένη εποπτεία με ζεύγη ίδιας οικογένειας (B2). & Διαχωρίζουν τι προέρχεται από την άμεση περιγραφή, τι από την κοινή μαθησιακή υποδομή και τι από την οικογενειακή/ιστορική προέλευση των θετικών σχέσεων. Οι ακριβείς απώλειες και οι αντιστοιχίσεις ορίζονται στην §~\ref{sec:eise}.\tabularnewline
\midrule
\textbf{PPM} & Το Prediction by Partial Matching είναι μοντέλο μεταβλητού μήκους συμφραζομένου: προβλέπει την επόμενη οικογένεια από όσο μακρύ προηγούμενο υποστηρίζουν τα δεδομένα και, όταν το ειδικό συμφραζόμενο είναι αραιό, επιστρέφει σε συντομότερο και καλύτερα υποστηριζόμενο συμφραζόμενο (backoff). & Είναι ερμηνεύσιμο και κατάλληλο για σχετικά μικρά σώματα. Χρησιμοποιείται στη M3c για το πληροφοριακό περιεχόμενο οικογένειας και στη M5 για την προβλεψιμότητα ως προς δηλωμένο σώμα αναφοράς. Η επιλογή συνδέεται με τη βιβλιογραφία πιθανοτικής μουσικής προσδοκίας και IDyOM (Pearce, 2018).\tabularnewline
\end{xltabular}
}\par

\subsection{Ιεραρχία και ελάχιστος πυρήνας}\label{sec:core-hierarchy}

Η \textbf{προεγγραφή (preregistration)} καταγράφει πριν από την κύρια αξιολόγηση τα βασικά εκτιμώμενα μεγέθη, τους κανόνες ανάλυσης, τις απαιτήσεις δεδομένων και τα κριτήρια ενεργοποίησης. Ένα \textbf{κριτήριο ενεργοποίησης} ελέγχει αν τα διαθέσιμα δεδομένα έχουν την κάλυψη, την ανεξαρτησία και την αξιοπιστία των μετρήσεων που απαιτεί μια προκαθορισμένη ανάλυση. Το αποτέλεσμα αυτού του ελέγχου καθορίζει ποιο επίπεδο ανάλυσης είναι επιστημονικά υποστηρίξιμο· εφαρμόζεται ανεξάρτητα από το πρόσημο ή την επιθυμητή κατεύθυνση του αποτελέσματος.

Οι τρεις καταστάσεις διαθεσιμότητας του ανεξάρτητου εκτελεστικού κριτηρίου αξιολόγησης ορίζονται πριν από την τελική αξιολόγηση: GREEN σημαίνει επαρκή λεπτομερή αντιστοίχιση ιστορικών και εκτελεστικών δεδομένων, AMBER επαρκή κάλυψη μόνο σε προκαθορισμένη αδρότερη κλίμακα και RED ανεπαρκή κάλυψη ή αξιοπιστία για επιβεβαιωτική εκτελεστική σύγκριση (Παράρτημα~\ref{app:corpora-gate}).

Ο \textbf{βασικός υπολογιστικός και μουσικολογικός πυρήνας} συγκροτείται από τη σύγκριση αναπαραστάσεων της M1, την αντιστοιχισμένη σύγκριση B2--EISE της M2, την ανεξάρτητη μουσικολογική επικύρωση της M3a και την εξωτερική εκτελεστική τεκμηρίωση που χρειάζεται για τον προκαθορισμένο στόχο της M2. Η τελευταία χρησιμοποιεί το αναγκαίο μέρος των δεδομένων και μετρήσεων της M4. Οι ευρύτερες αναλύσεις της M4 αναπτύσσουν τη μελέτη της εκτελεστικής παραλλαγής στο επίπεδο που επιτρέπει η πραγματική κάλυψη.

Οι \textbf{ανθρώπινες μελέτες} παρέχουν το κύριο γνωσιακό επίπεδο επικύρωσης με δύο επιβεβαιωτικές υποθέσεις. Η M3c εξετάζει την πρόσθετη σχέση του πληροφοριακού περιεχομένου σε επίπεδο οικογένειας με την ανθρώπινη μουσική προσδοκία, ενώ η M3b εξετάζει την οργάνωση των πραγματώσεων σε οικογένειες κατά την ακρόαση. Η M3c συλλέγεται πριν από τη M3b, ώστε η τριαδική σύγκριση οικογενειών να μην κατευθύνει την πρώτη κρίση προσδοκίας.

Το $\Delta_{\mathrm{ctx}}$, το CCGP, η ευθυγράμμιση των αναπαραστάσεων μέσα στην ίδια οικογένεια, οι αλληλεπιδράσεις εμπειρίας, η βιωματική απόκριση, οι ευρύτερες αναλύσεις της M4 και η M5 εμπλουτίζουν τη μουσική, γνωσιακή και ιστορική τεκμηρίωση με το ιδιαίτερο καθεστώς κάθε ανάλυσης. Η M5 είναι προγραμματισμένη ιστορική εφαρμογή του επικυρωμένου πλαισίου στον Νικόλαο Πρωτοψάλτη Σμύρνης, με ελάχιστη και πλήρη εκδοχή ανάλογα με την επάρκεια των δεδομένων. Η ελάχιστη εκδοχή \textbf{M5-min} παράγει, σε συγκρίσιμο σώμα πλήρων έργων της Νέας Μεθόδου, τουλάχιστον τα $D_{\mathrm{fam}}$ (κατανομή οικογενειών) και $D_{\mathrm{within}}$ (ενδοοικογενειακή διαφοροποίηση) με αβεβαιότητα σε επίπεδο έργου. Η πλήρης εκδοχή προσθέτει τη διαδοχή οικογενειών ($D_{\mathrm{seq}}$) και τη λειτουργική τοποθέτησή τους ($D_{\mathrm{func}}$), όταν πληρούνται τα αντίστοιχα κριτήρια επάρκειας. Η κατανομή αναφοράς σε επίπεδο μελοποιού ενεργοποιείται μόνο όταν το ΠΕ0 τεκμηριώσει επαρκή αριθμό ανεξάρτητων μελοποιών σύμφωνα με το προκαθορισμένο κατώφλι $K_{\min}$. Η δημοσιευμένη περίπτωση Ιωαννίδη (2022) λειτουργεί ως χωριστός εξωτερικός έλεγχος εφαρμογής και δεν αποτελεί προϋπόθεση για την εκτέλεση της M5-min.

{\footnotesize
\begin{xltabular}{\linewidth}{@{}>{\RaggedRight\arraybackslash}p{0.105\linewidth}>{\RaggedRight\arraybackslash}p{0.23\linewidth}>{\RaggedRight\arraybackslash}p{0.29\linewidth}Y@{}}
\toprule
\textbf{Βαθμίδα} & \textbf{Ενότητες} & \textbf{Συμπερασματικό καθεστώς} & \textbf{Όταν λείπει η απαιτούμενη κάλυψη}\tabularnewline
\midrule
\endfirsthead
\toprule
\textbf{Βαθμίδα} & \textbf{Ενότητες} & \textbf{Συμπερασματικό καθεστώς} & \textbf{Όταν λείπει η απαιτούμενη κάλυψη}\tabularnewline
\midrule
\endhead
\midrule
\multicolumn{4}{r}{\footnotesize Συνέχεια στην επόμενη σελίδα}\tabularnewline
\endfoot
\bottomrule
\endlastfoot

1. Πυρήνας & M1, M2, M3a και ελάχιστο ανεξάρτητο εκτελεστικό κριτήριο της M4 & $R^\star$· $\Delta_{\mathrm{hist}}$ επιβεβαιωτικό σε GREEN/AMBER· χωριστή $\Delta_{\mathrm{learn}}$· ανεξάρτητη μουσικολογική κρίση & Το ΠΕ0 ορίζει την κλίμακα στην οποία τα δεδομένα επιτρέπουν αξιόπιστη εκτίμηση· αν δεν υπάρχει ανεξάρτητο κριτήριο αξιολόγησης, δηλώνεται ότι η $\Delta_{\mathrm{hist}}$ δεν μπορεί να εκτιμηθεί αξιόπιστα.\tabularnewline
2. Ανθρώπινη επικύρωση & M3c, έπειτα M3b & $H_{\mathrm{pred}}$, $H_{\mathrm{cat}}$· οικογένεια δύο επιβεβαιωτικών ελέγχων με Holm & Για κάθε υπόθεση δηλώνεται χωριστά αν τα διαθέσιμα δεδομένα επιτρέπουν αξιόπιστη εκτίμηση, σύμφωνα με τα κριτήρια της §~\ref{sec:estimability}.\tabularnewline
3. M5-min & M5a/M5b: επιλογή οικογενειών και ενδοοικογενειακή πραγμάτωση & προκαθορισμένα περιγραφικά μεγέθη· η κατανομή αναφοράς ανά μελοποιό ενεργοποιείται μόνο μετά τον έλεγχο επάρκειας του ΠΕ0 & Με μικρότερο αριθμό επιλέξιμων μελοποιών αναφέρονται μόνο οι διαθέσιμες εκτιμήσεις, η κάλυψη και τα διαστήματα αβεβαιότητας, χωρίς εκατοστημόριο ή θέση σε κατανομή συνθετών.\tabularnewline
4. Ενίσχυση & $\Delta_{\mathrm{ctx}}$, CCGP, ευθυγράμμιση, εμπειρία, ευρύτερη M4, πλήρης M5· βασική μελέτη τεκμηριωμένων μετασχηματισμών (§~\ref{sec:atl}) & δευτερεύοντα ή διερευνητικά, κατά περίπτωση· η μελέτη μετασχηματισμών δεσμεύεται εφόσον περάσει την πύλη επάρκειας & $D_{\mathrm{seq}}$, $D_{\mathrm{func}}$ και ιστορικές υποομάδες ενεργοποιούνται μόνο με επαρκή κάλυψη και ανεξαρτησία πηγών.\tabularnewline
5. Υπό προϋποθέσεις & M3ν· δείκτες EGG μέσα στη M4· υποψήφιες παραλλαγές από τη μελέτη μετασχηματισμών και δεύτερο ρεπερτόριο (§~\ref{sec:atl}) & νευρική επέκταση, διερευνητική εκτελεστική ανάλυση, χρήση υποψήφιων παραλλαγών στη μάθηση αναπαραστάσεων και μεταφορά στο ζναμένιο μέλος & Η επέκταση ενεργοποιείται όταν πληρούνται τα κριτήριά της.\tabularnewline
\end{xltabular}
}\par

Η \textbf{υλοποίηση της βαθμίδας 1} ακολουθεί την κλίμακα που τεκμηριώνει το ΠΕ0 και περιλαμβάνει τις συγκρίσεις της M1, την ανεξάρτητη μουσικολογική αξιολόγηση της M3a και την αντιστοιχισμένη σύγκριση της M2. Η διαφορά $\Delta_{\mathrm{hist}}$ αξιολογείται επιβεβαιωτικά σε GREEN ή AMBER, ανεξάρτητα από το πρόσημο του αποτελέσματος. Σε RED η διαθέσιμη εκτελεστική σύγκριση αναφέρεται διερευνητικά· όταν λείπει αξιόπιστο ανεξάρτητο κριτήριο αξιολόγησης, δηλώνεται ότι η $\Delta_{\mathrm{hist}}$ δεν μπορεί να εκτιμηθεί αξιόπιστα. Το ΠΕ0 καταγράφει έγκαιρα την πραγματική κάλυψη και κατευθύνει τον διαθέσιμο χρόνο στις υποστηριζόμενες συγκρίσεις. Η τεκμηρίωση της διατριβής συνθέτει τα αποτελέσματα που μπορούν να εκτιμηθούν, τις ανεξάρτητες μουσικολογικές κρίσεις και τα σαφώς δηλωμένα όρια των υπόλοιπων συμπερασμάτων. Η βαθμίδα 2 παρέχει τις προγραμματισμένες κύριες γνωσιακές δοκιμές $H_{\mathrm{pred}}$ και $H_{\mathrm{cat}}$· για κάθε αποτέλεσμα δηλώνεται αν τα διαθέσιμα δεδομένα επιτρέπουν αξιόπιστη εκτίμηση.

Η νευρική M3ν ενεργοποιείται μόνο με γραπτή επιβεβαίωση υποδομής EEG/MEG, κατάλληλης τεχνογνωσίας και χωριστής δεοντολογικής έγκρισης. Η διερευνητική ανάλυση δεικτών φωνητικής επαφής μέσα στη M4 εξαρτάται από την πρόσβαση και την ποιότητα του συγχρονισμένου ηλεκτρογλωττογραφήματος. Η ενεργοποίηση αυτών των επεκτάσεων δεν αποτελεί προϋπόθεση για την εκτίμηση των ήδη ορισμένων κύριων υποθέσεων.



\subsection{Διαχωρισμός ορισμού, εκπαίδευσης και αξιολόγησης}

Ως \textbf{διαρροή πληροφορίας (data leakage)} νοείται κάθε περίπτωση στην οποία πληροφορία από υλικό που προορίζεται για ανεξάρτητη αξιολόγηση επηρεάζει την εκπαίδευση, την επιλογή αναπαράστασης, τη ρύθμιση του μοντέλου ή την κατασκευή των ίδιων ετικετών που αργότερα αξιολογούνται. Στο παρόν υλικό η ανεξαρτησία πρέπει να προστατεύεται σε επίπεδο ολόκληρου έργου, πηγής, ομάδας συγγενών πηγών και τεκμηριωμένης σχέσης. Για παράδειγμα, τμήματα του ίδιου Χερουβικού σε εκπαίδευση και τελική αξιολόγηση θα επέτρεπαν στο μοντέλο να αξιοποιήσει ειδικά γνωρίσματα του έργου αντί να ελεγχθεί η γενίκευσή του σε νέο έργο.

Οι οικογένειες, οι τεκμηριωμένες ιστορικές σχέσεις και τα σύνολα αξιολόγησης διατηρούν σαφή τεκμηριωμένη προέλευση. Τα τρία είδη συνδέσεων μεταξύ πραγματώσεων διατηρούνται χωριστά: το $E_{\mathrm{sim}}$ προκύπτει από μελωδική ομοιότητα που υπολογίζει το μοντέλο, το $E_{\mathrm{fam}}$ δηλώνει ένταξη στην ίδια οικογένεια και το $E_{\mathrm{ex}}$ καταγράφει τεκμηριωμένη ιστορική εξηγητική σχέση (πλήρεις κανόνες στο Παράρτημα~\ref{app:splits}). Η τελική ανεξάρτητη επικύρωση από ειδικούς πραγματοποιείται σε υλικό και υποψήφιες ετικέτες που έχουν διαχωριστεί από τη διαδικασία κατασκευής του συνόλου αναφοράς. Το ιστορικό σύνολο ελέγχου της §~\ref{sec:hist-discriminability} οριστικοποιείται στο ΠΕ0 και παραμένει εκτός εκπαίδευσης, επιλογής υπερπαραμέτρων και επιλογής αναπαράστασης. Η ανθρώπινη τράπεζα ερεθισμάτων κατασκευάζεται μετά την οριστικοποίηση του επικυρωμένου λεξιλογίου οικογενειών και αναγνωριστικά με ελεγχόμενες εκδόσεις συνδέουν κάθε ερέθισμα με έργο, πηγή, οικογένεια, $I_F$, $I_{\mathrm{event}}$, επιφανειακές αποστάσεις και ανάθεση συνθήκης.

Η επιλογή μοντέλου της M1 πραγματοποιείται στα δεδομένα ανάπτυξης. Οι αξιολογήσεις της M2 σε υλικό εκτός εκπαίδευσης, το ιστορικό σύνολο ελέγχου και τα ανθρώπινα ερεθίσματα αξιολόγησης δεν χρησιμοποιούνται για επιλογή υπερπαραμέτρων ή αναπαράστασης. Περιπτώσεις του ιστορικού συνόλου μπορούν αργότερα να αποτελέσουν προκαθορισμένο σύνολο υποψήφιων ερεθισμάτων της M3b μόνο βάσει μουσικολογικών ετικετών και κριτηρίων επιφανειακής απόστασης, ποτέ βάσει της επιτυχίας ή αποτυχίας της M1 στις συγκεκριμένες περιπτώσεις. Ο χώρος των ανθρώπινων αναπαραστάσεων της M3b προσαρμόζεται χωρίς ετικέτες οικογένειας και οι υπερπαράμετροί του επιλέγονται με πρόβλεψη τριάδων σε υλικό ανάπτυξης που έχει κρατηθεί εκτός της προσαρμογής του μοντέλου. Το $I_F$ της M3c οριστικοποιείται μετά την οριστικοποίηση του λεξιλογίου οικογενειών και ο έλεγχος επάρκειας του μοντέλου διαδοχής οικογενειών (PPM, §~\ref{sec:m3c}) αξιολογείται σε έργα εκτός εκπαίδευσης πριν από την επιλογή των κύριων ανθρώπινων δοκιμών.

\subsection{Κοινές αρχές συγκρισιμότητας και αναπαραγωγιμότητας}

Κάθε κύρια αντίθεση αλλάζει μία σαφώς ορισμένη πηγή πληροφορίας. Η M1 συγκρίνει κλίμακες αναπαράστασης με κοινό στόχο. Η M2 συγκρίνει προέλευση θετικών σχέσεων με αντιστοιχισμένο περιβάλλον εκπαίδευσης. Η M3b συγκρίνει οικογενειακή ένταξη μέσα σε κοινή περιοχή επιφανειακών αποστάσεων. Η M3c διαχωρίζει τα δύο επίπεδα πρόβλεψης με ερεθίσματα που συνδυάζουν υψηλό $I_F$ και χαμηλό $I_{\mathrm{event}}$, ή το αντίστροφο. Οι δύο αυτοί αντιδιαγώνιοι συνδυασμοί υποστηρίζουν τη χωριστή εκτίμηση της συμβολής των δύο μεταβλητών, οι οποίες παραμένουν συνεχείς στα στατιστικά μοντέλα. Το $R$ αλλάζει μόνο τη μη μουσική ένδειξη λατρευτικού πλαισίου, ενώ ο στόχος και το άμεσο μουσικό πρόθεμα παραμένουν ίδια.

Οι λεπτομερείς κανόνες διαχωρισμών δεδομένων, αντιστοίχισης, ελέγχου πολλαπλότητας, κατασκευής ερεθισμάτων, αντιστάθμισης, προσομοίωσης ισχύος και προεγγραφής δίνονται στα Παραρτήματα~\ref{app:splits}--\ref{app:human}.

\section{Σύγκριση αναπαραστάσεων}\label{sec:m1}

Η M1 εξετάζει πώς πρέπει να αναπαραστήσουμε το μουσικό υλικό ώστε να εντοπίζουμε θέσεις και να αναγνωρίζουμε τις οικογένειές τους σε νέα έργα και πηγές. Το μουσικολογικό ερώτημα είναι ποια κλίμακα και ποια πληροφοριακά στρώματα διατηρούν καλύτερα τις τεκμηριωμένες συγγένειες και διακρίσεις των θέσεων, των οποίων η έκταση μπορεί να διασχίζει συλλαβικά και σημειογραφικά όρια. Η M1 διακρίνει την επιλογή κωδικοποίησης της εισόδου από τη σύγκριση μεταξύ γεγονότων, επαγόμενων τμημάτων και ρητών θέσεων.

\subsection{Ενιαίο συμβολικό επίπεδο αναφοράς και \texorpdfstring{$U\times F$}{U x F}}

Κάθε πηγή μετατρέπεται σε ακολουθία \textbf{ατομικών μελωδικών γεγονότων}, διατηρώντας χωριστά το επίπεδο του γραφικού σημείου και το επίπεδο του μελωδικού γεγονότος. Το \emph{γραφικό σημείο} είναι το σύμβολο όπως εμφανίζεται στην πηγή· το \emph{ατομικό γεγονός} είναι η ελάχιστη μονάδα μελωδικής κίνησης που προκύπτει από την ανάγνωση ενός ή περισσότερων σημείων. Η αντιστοίχιση διατηρείται αμφίδρομα: κανένα γεγονός δεν αποκόπτεται από τα σημεία που το παρήγαγαν.

Μια \textbf{σωρευτική σημειογραφική συντεταγμένη} ορίζει κοινό άξονα πάνω στον οποίο ευθυγραμμίζονται θέσεις, συλλαβές και ηχητικές παρατηρήσεις. Λειτουργεί ως \emph{κοινός χάρακας}: επιτρέπει να αντιστοιχιστούν γραφικά σημεία, ατομικά γεγονότα, συλλαβές, θέσεις και χρονικά σημεία της ηχογράφησης χωρίς να απαιτείται να έχουν το ίδιο μήκος. Η αντιστοίχιση αυτή είναι αναγκαία στο παπαδικό μέλος, όπου μία εκτεταμένη συλλαβή μπορεί να περιλαμβάνει πολλαπλές εσωτερικές δομικές μεταβάσεις.

\paragraph{Διαλειτουργικότητα και εισαγωγή δεδομένων.}
Το ενιαίο ερευνητικό επίπεδο αναφοράς παραμένει ανεξάρτητο από τη μορφή οποιουδήποτε εξωτερικού λογισμικού. Υφιστάμενες δομημένες μορφές της Νέας Μεθόδου, όπως τα αρχεία του ανοικτού προγράμματος μουσικής σημειογραφίας \textbf{Neanes}, μπορούν να χρησιμοποιηθούν ως πηγές εισαγωγής μέσω ντετερμινιστικών προσαρμογέων που διατηρούν το αρχικό γραφικό σημείο, τις τροπικές και χρονικές ενδείξεις, το λεκτικό υλικό και την προέλευση. Η εξωτερική μορφή δεν ταυτίζεται με την οντολογία της ερευνητικής πρότασης: κάθε μετατροπή ελέγχεται σε προκαθορισμένο δείγμα με ανακατασκευή ή ισοδύναμο έλεγχο πεδίων πριν τροφοδοτήσει το ενιαίο επίπεδο αναφοράς. Για έντυπες εκδόσεις της Νέας Μεθόδου μπορεί να δοκιμαστεί, ως \textbf{προαιρετική διαδρομή απόκτησης δεδομένων}, υφιστάμενο ανοικτό εργαλείο οπτικής αναγνώρισης χαρακτήρων (OCR) βυζαντινής σημειογραφίας συμβατό με το οικοσύστημα Neanes. Η αναγνώριση αυτής της γραφής έχει ειδικά τεχνικά προηγούμενα (Gezerlis--Theodoridis, 2002· Dalitz κ.ά., 2008· Kokkinidis κ.ά., 2016). Η έξοδος OCR δεν εισέρχεται αυτόματα στο σώμα: το ΠΕ0 μετρά σφάλμα μετατροπής σε μικρό προκαθορισμένο δείγμα, ελέγχει τα δικαιώματα της συγκεκριμένης ψηφιοποίησης και απαιτεί χειροκίνητη επαλήθευση. Αν η διαδρομή δεν περάσει το κριτήριο ποιότητας, απενεργοποιείται χωρίς καμία μεταβολή του βασικού σώματος ή των κύριων υποθέσεων.

Η σύγκριση των $u_1$--$u_3$ και η διαδοχική προσθήκη χρονικής, τροπικής και κειμενοπροσωδιακής πληροφορίας έχουν και μουσικολογικό κίνητρο. Η ανάλυση του Δεσπότη (2007a) δείχνει ότι η μελωδικο-ρυθμική λειτουργία μπορεί να εξαρτάται από σχέσεις μεταξύ σημείων, από την τοποθέτησή τους στη φράση και από ευρύτερες μορφολογικές και χρονικές οργανώσεις. Το πείραμα ελέγχει αν αυτή η πληροφορία από το συμφραζόμενο προσφέρει μετρήσιμο πρόσθετο όφελος για τα προκαθορισμένα υπολογιστικά έργα, χωρίς να προδικάζει ποιο επίπεδο θα αποδειχθεί διακριτικό.

Το πλέγμα $U\times F$ (\(U\): μονάδα εισόδου, \(F\): διαθέσιμα γνωρίσματα· ορισμοί στον πίνακα που ακολουθεί) συγκρίνει εναλλακτικές κωδικοποιήσεις εισόδου για τη μοντελοποίηση και οριοθέτηση της μουσικής θέσης· δεν αποτελεί κατάλογο εναλλακτικών μουσικολογικών μονάδων. Ο $T_{\mathrm{evt}}$ και η επαγόμενη τμηματοποίηση αποτελούν χωριστές συγκρίσεις δομικής κλίμακας έναντι της ρητής θέσης.

Το πλέγμα αναλύει δύο σχεδιαστικές διαστάσεις: την κωδικοποίηση εισόδου και το σύνολο γνωρισμάτων. Η M1 εξετάζει αν η διαφορά επίδοσης συνδέεται με τον τρόπο κωδικοποίησης ή με την πρόσθετη πληροφορία των γνωρισμάτων. Οι συγκρίσεις δομικής κλίμακας εξετάζονται χωριστά, με τον Transformer γεγονότων, την επαγόμενη τμηματοποίηση και τη ρητή θέση.

Μουσικολογικά, οι δύο διαστάσεις απαντούν σε διαφορετικά ερωτήματα: ο παράγοντας \(U\) ελέγχει σε ποια τοπική κλίμακα κωδικοποιείται αποτελεσματικότερα η σχεσιακή πληροφορία, ενώ ο \(F\) ελέγχει ποια είδη πληροφορίας απαιτούνται για τη διάκριση. Η σχεσιακή και χρονική διάσταση έχει ειδικό αναλυτικό κίνητρο στη νεότερη μελέτη της Νέας Μεθόδου (Δεσπότης, 2007a), η τροπική πληροφορία στη θεωρητική και σημειογραφική βιβλιογραφία (Χρύσανθος, 1832· Raasted, 1966· Troelsgård, 2011), ενώ η κειμενική και προσωδιακή διάσταση συνδέεται με τη φραστική λειτουργία και την προσαρμογή του παπαδικού μέλους στο κείμενο, όπως τεκμηριώνεται στις μορφολογικές αναλύσεις που χρησιμοποιούνται στην παρούσα έρευνα (Καραγκούνης, 2003, 2005).

Οι τροπικοί όροι του πεδίου $M$ χρησιμοποιούνται με την ψαλτική τους σημασία: η μαρτυρία\footnote{Μαρτυρία: σημειογραφική ένδειξη που δηλώνει ή επιβεβαιώνει τροπική και φθογγική κατάσταση· διακρίνεται από τη «μαρτυρία» με τη συνήθη ιστοριογραφική έννοια. Βλ. Χρύσανθος 1832.} δηλώνει τροπική και φθογγική κατάσταση, ενώ η φθορά\footnote{Φθορά: σημειογραφική ένδειξη που μεταβάλλει ή εξειδικεύει την τροπική και διαστηματική κατάσταση. Βλ. Χρύσανθος 1832.} δηλώνει μεταβολή ή εξειδίκευση της διαστηματικής/τροπικής οργάνωσης.\par

Το πείραμα έχει επομένως \textbf{διασταυρωμένη διάταξη \(3\times4\) με διατεταγμένη σωρευτική κλίμακα γνωρισμάτων}: οι τρεις κωδικοποιήσεις εισόδου συγκρίνονται σε όλα τα επίπεδα γνωρισμάτων, ενώ το \(F\) είναι ένας παράγοντας τεσσάρων διαδοχικά εμπλουτισμένων επιπέδων. Η διάταξη εξετάζει τη μονάδα εισόδου, τις διαδοχικές προσθήκες πληροφορίας και την αλληλεπίδραση των δύο παραγόντων. Οι ομάδες \(T\), \(M\) και \(L\) δεν διασταυρώνονται ως ανεξάρτητοι δυαδικοί παράγοντες· οι μεταβάσεις από \(N\) σε \(N{+}T\), \(N{+}T{+}M\) και \(N{+}T{+}M{+}L\) ερμηνεύονται ως πρόσθετη πληροφοριακή συνεισφορά κάθε νέου συνόλου υπό τα ήδη διαθέσιμα γνωρίσματα και όχι ως ανεξάρτητες κύριες επιδράσεις των \(T\), \(M\) και \(L\).

\[
U \times F,\qquad
U=\{u_1,u_2,u_3\},\qquad
F=\{N,\ N{+}T,\ N{+}T{+}M,\ N{+}T{+}M{+}L\}
\]

\begin{center}\small
\begin{tabularx}{\textwidth}{@{}p{0.075\textwidth}p{0.245\textwidth}Y@{}}
\toprule
\multicolumn{3}{@{}l}{\textbf{Παράγοντας \(U\): κωδικοποίηση εισόδου}}\tabularnewline
\midrule
\(u_1\) & Γραφικό σύμβολο & Το σύμβολο όπως εμφανίζεται στην πηγή· σημείο αναφοράς\tabularnewline
\(u_2\) & Ατομικό μελωδικό γεγονός & Ελάχιστη μονάδα μελωδικής κίνησης\tabularnewline
\(u_3\) & Σύνθετη νευματική μονάδα & Ομαδοποίηση με ντετερμινιστικό κανόνα\tabularnewline
\midrule
\multicolumn{3}{@{}l}{\textbf{Παράγοντας \(F\): σύνολο γνωρισμάτων}}\tabularnewline
\midrule
\(N\) & Σημειογραφικά & Σχετική μελωδική κίνηση και πλήρης ταυτότητα των φωνητικών/ποσοτικών και ποιοτικών γραφικών στοιχείων της Νέας Μεθόδου που συνεισφέρουν στο γεγονός αναφοράς\tabularnewline
\(T\) & Χρονικά & Χρονικοί χαρακτήρες και σχετικές σημειογραφικές χρονικές αξίες σε κοινή μονάδα ανάγνωσης, με δηλωμένο κανόνα παραγωγής και διατήρηση της προέλευσης/αβεβαιότητας\tabularnewline
\(M\) & Τροπικά και κατάστασης & Ήχος σε ισχύ, μαρτυρία, φθορά\tabularnewline
\(L\) & Κειμενικά και προσωδιακά & Συλλαβική μονάδα, λέξη/κοινό κειμενικό χωρίο, λεκτικός τόνος όπου εφαρμόζεται, όρια λέξης/φράσης, μελισματική πυκνότητα\tabularnewline
\bottomrule
\end{tabularx}
\end{center}

Το $L$ καταγράφει τη συντακτική, προσωδιακή και λειτουργική οργάνωση του κειμένου των Χερουβικών και Κοινωνικών, χωρίς να προϋποθέτει σταθερό ποιητικό μέτρο. Τυχόν μετρική ερμηνεία ή οργάνωση της μουσικής χρονικής κίνησης καταγράφεται χωριστά, με τη συγκεκριμένη τεκμηρίωσή της.

Το σύνολο χρονικών γνωρισμάτων $T$ διατηρεί τους παρατηρούμενους χρονικούς χαρακτήρες χωριστά από τις παράγωγες σχετικές χρονικές αξίες. Οι τιμές αυτές εκφράζονται σε κοινή μονάδα ανάγνωσης και δεν κωδικοποιούν απόλυτη αγωγή εκτέλεσης. Εκδοτικές διαστολές ή άλλες μεταγενέστερες μετρικές ομαδοποιήσεις καταγράφονται ως ιδιότητες της συγκεκριμένης πηγής και δεν χρησιμοποιούνται ως αλήθεια αναφοράς για ιστορικό μέτρο. Ποσότητες που απαιτούν γνωστά όρια ολόκληρης θέσης ---όπως η συνολική ανάπτυξη $Q_x$ ή το κανονικοποιημένο χρονικό προφίλ $p_x$ της M5--- υπολογίζονται μόνο \emph{μετά} την οριοθέτηση και δεν παρέχονται ως είσοδοι σε μοντέλο που καλείται να εντοπίσει τα ίδια όρια.

Η διατήρηση των ποιοτικών ενδείξεων στο $N$ δεν προϋποθέτει άμεση ιστορική αντιστοιχία προς τις μεγάλες υποστάσεις της παλαιάς γραφής\footnote{Μεγάλες υποστάσεις: ιστορικά σημειογραφικά στοιχεία της παλαιότερης γραφής των οποίων η λειτουργία δεν εξαντλείται στην ποσοτική μελωδική κίνηση και συνδέεται στη βιβλιογραφία με ποιοτική, εκτελεστική και μελική πληροφορία· βλ. Στάθης 1978· Καραγκούνης 2003.}· αποτρέπει την εκ των προτέρων απώλεια σημειογραφικής πληροφορίας που μπορεί να αποδειχθεί διακριτική.

Στα μη λεξικά χωρία η συλλαβική άρθρωση και η μελισματική πυκνότητα διατηρούνται όπου τεκμηριώνονται, ενώ λέξη, κειμενικό χωρίο και λεκτικός τόνος δηλώνονται ως μη εφαρμοστέα, διακριτά από άγνωστες τιμές. Οι συγκρινόμενες συνθήκες χρησιμοποιούν κοινό επιλέξιμο υλικό. Η μετάβαση προς \(N{+}T{+}M{+}L\) ερμηνεύεται ως όφελος του πλήρους συνόλου γνωρισμάτων \(L\), όχι ως καθαρή επίδραση μιας μεμονωμένης λεξικής ή προσωδιακής μεταβλητής. Τα όρια της μουσικής θέσης ορίζονται ανεξάρτητα· η σύμπτωσή τους με αλλαγή συλλαβής ή λέξης μετράται χωριστά ως δευτερεύον διαγνωστικό τμηματοποίησης.


Το πλέγμα δίνει δώδεκα συνθήκες με \emph{κοινό πρωτόκολλο μοντελοποίησης}, \emph{κοινό χώρο εξόδου}, \emph{κοινές μετρικές} και ίδιο προϋπολογισμό ρύθμισης. Η χωρητικότητα της κοινής αρχιτεκτονικής, η διάσταση των διανυσματικών ενσωματώσεων και ο προϋπολογισμός ρύθμισης υπερπαραμέτρων οριστικοποιούνται από κοινού, ενώ οι διαφορές στον αριθμό παραμέτρων που προκαλεί το λεξιλόγιο καταγράφονται και ελέγχονται με ανάλυση ευαισθησίας με αντιστοιχισμένο υπολογιστικό κόστος. Παράλληλα γίνεται ρητός έλεγχος ότι τα αναγνωριστικά των διακριτών μονάδων δεν επανεισάγουν έμμεσα πεδία που έχουν αφαιρεθεί από κάποια διαμόρφωση γνωρισμάτων. Οι έλεγχοι αυτοί περιορίζουν το ενδεχόμενο οι διαφορές επίδοσης να προκύπτουν από πρόσβαση σε μη προβλεπόμενη πληροφορία ή από διαφορετική χωρητικότητα των μοντέλων. Το πλέγμα επιτρέπει να εξεταστούν χωριστά:

\begin{enumerate}
\item η \textbf{επίδραση της κωδικοποίησης εισόδου} --- ποια κωδικοποίηση είναι καταλληλότερη, με σταθερή πληροφορία·
\item η \textbf{κύρια επίδραση της διαμόρφωσης γνωρισμάτων} --- ποια πρόσθετη πληροφορία συνεισφέρει κάθε διαδοχικός εμπλουτισμός, με σταθερή κωδικοποίηση εισόδου·
\item η \textbf{αλληλεπίδρασή τους} --- αν το όφελος μιας κωδικοποίησης εξαρτάται από τα διαθέσιμα γνωρίσματα.
\end{enumerate}

Με αυτή τη διάκριση, το ερώτημα «ποια αναπαράσταση είναι καταλληλότερη» αποκτά ερμηνεύσιμη και αναπαραγώγιμη απάντηση.

Ως εξωτερικός έλεγχος της δομικής κλίμακας, ο $T_{\mathrm{evt}}$ εκπαιδεύεται απευθείας στην ακολουθία ατομικών μελωδικών γεγονότων με την ίδια επιλέξιμη πληροφορία και το ίδιο διαθέσιμο παρελθόν. Η σύγκριση γίνεται σε κοινούς στόχους αξιολόγησης και σε έργα/ομάδες συγγενών πηγών που έχουν παραμείνει εκτός εκπαίδευσης. Με αυτόν τον τρόπο ελέγχεται αν ένα ισχυρό μοντέλο ακολουθίας μπορεί να ανακτήσει τη σχετική δομή χωρίς ρητή αναπαράσταση της θέσης. Οι τεχνικές προδιαγραφές του ελέγχου δίνονται στο Παράρτημα~\ref{app:event-transformer}.

Η πλήρης αμφίδρομη χαρτογράφηση γραφικών σημείων και ατομικών γεγονότων, ο ντετερμινιστικός ορισμός της σύνθετης νευματικής μονάδας, τα λεξιλόγια πεδίων και οι έλεγχοι ανακατασκευής δίνονται στο Παράρτημα~\ref{app:repr}.

\subsection{Έλεγχος ιστορικής διακριτικής ικανότητας σε ανεξάρτητες δύσκολες περιπτώσεις}\label{sec:hist-discriminability}

Ως προκαθορισμένη δευτερεύουσα ανάλυση της M1 συγκροτείται στο ΠΕ0 ένα \textbf{ιστορικό σύνολο ελέγχου υψηλής δυσκολίας}. Περιλαμβάνει δύο χωριστές κατηγορίες περιπτώσεων: (α) μουσικολογικά τεκμηριωμένα \emph{δύσκολα αρνητικά}, δηλαδή διαφορετικές θέσεις με ισχυρή μετροφωνική ή επιφανειακή συγγένεια, με αφετηρία περιπτώσεις όπως εκείνες της σ.~788 του Καραγκούνη· και (β) \emph{δύσκολα θετικά}, δηλαδή τεκμηριωμένες πραγματώσεις της ίδιας θέσης με αισθητή αναλυτική ή τροπική διαφοροποίηση, με αφετηρία τις αντιπαραβολές των σσ.~810--812. Η παλαιά σημειογραφία δεν κωδικοποιείται ως νέο υπολογιστικό σώμα: από την παλαιότερη πλευρά διατηρούνται μόνο τα μουσικολογικά τεκμήρια που ορίζουν τη σχέση, ενώ η υπολογιστική μέτρηση εφαρμόζεται στις τεκμηριωμένες αποδόσεις της Νέας Μεθόδου.

Το σύνολο προκαθορίζεται και \textbf{οριστικοποιείται στο ΠΕ0 πριν από την επιλογή αναπαράστασης}. Οι συγκεκριμένες εμφανίσεις, σχέσεις και αποφάσεις ετικέτας που το συγκροτούν δεν χρησιμοποιούνται για εκπαίδευση, επιλογή υπερπαραμέτρων, επιλογή των $R^\star$, $\mathcal R^\star$ και $R_{\mathrm{sens}}$, ούτε για κατασκευή δύσκολων αρνητικών παραδειγμάτων της M1 ή για την κατασκευή των ίδιων ετικετών που καλείται να αξιολογήσει. Η γενική οντολογία μπορεί να στηρίζεται στην ίδια θεωρητική βιβλιογραφία, με καταγραφή προέλευσης. Όπου η ανεξαρτησία συγκεκριμένης περίπτωσης δεν μπορεί να τεκμηριωθεί, η περίπτωση παραμένει περιγραφική και δεν εισέρχεται στην εκτίμηση AUROC της κύριας στατιστικής σύγκρισης (εμβαδόν κάτω από την καμπύλη λειτουργικού χαρακτηριστικού δέκτη, μέτρο του πόσο καλά διακρίνονται τα δύο είδη ζευγών). Μετά την οριστικοποίηση της κύριας διαδικασίας επιλογής, οι ήδη εκπαιδευμένες αναπαραστάσεις αξιολογούνται στο ιστορικό σύνολο μία φορά ως ανεξάρτητος έλεγχος ανθεκτικότητας. Το αποτέλεσμα μπορεί να ερμηνευθεί στη M3a, αλλά δεν επιτρέπεται να αλλάξει αναδρομικά την επιλεγμένη αναπαράσταση.

Η κύρια στατιστική αξιολόγηση εφαρμόζεται μία φορά στην ήδη οριστικοποιημένη $R^\star$. Για την $R^\star$ υπολογίζεται η ίδια απόσταση αναπαράστασης $d(z_i,z_j)$ και ορίζεται
\[
s_{ij}=-d(z_i,z_j),\qquad y_{ij}\in\{0,1\}.
\]
Εδώ $y_{ij}=1$ δηλώνει προκαθορισμένο δύσκολο θετικό ζεύγος κοινής θέσης, ενώ $y_{ij}=0$ προκαθορισμένο δύσκολο αρνητικό ζεύγος διαφορετικών θέσεων.
Εφόσον το κριτήριο επάρκειας δεδομένων τεκμηριώσει επαρκή αριθμό \textbf{ανεξάρτητων θέσεων}, επαρκή κάλυψη και των δύο τύπων αντίθεσης και αξιόπιστες ιστορικές αντιστοιχίσεις, η κύρια συνοπτική μετρική είναι το AUROC του $s_{ij}$ στην $R^\star$. Η αβεβαιότητα εκτιμάται με επαναδειγματοληψία στο επίπεδο της υποκείμενης θέσης/ιστορικής μονάδας, διατηρώντας από κοινού τα εξαρτώμενα ζεύγη. Οι διαμορφώσεις $N$, $N{+}T$, $N{+}T{+}M$ και $N{+}T{+}M{+}L$ αναφέρονται ως προκαθορισμένη δευτερεύουσα περιγραφική ανάλυση του επιπέδου στο οποίο εμφανίζεται ή χάνεται η διακριτική πληροφορία και δεν δημιουργούν χωριστές επιβεβαιωτικές δοκιμές. Η ακριβής ομαδοποιημένη υλοποίηση οριστικοποιείται στο ΠΕ0 μετά την απογραφή της πραγματικής δομής εξάρτησης.

Αν το κριτήριο επάρκειας δεδομένων δεν ικανοποιείται, δεν αποδίδεται στο AUROC συμπερασματική στατιστική ερμηνεία. Το προκαθορισμένο παραδοτέο είναι περιγραφικός πίνακας ανά ιστορική περίπτωση με τις αποστάσεις στα επίπεδα $N$, $N{+}T$, $N{+}T{+}M$ και $N{+}T{+}M{+}L$, ώστε να φαίνεται σε ποιο επίπεδο πληροφορίας εμφανίζεται ή χάνεται η απαιτούμενη διάκριση ή συγγένεια. Ο έλεγχος αφορά τη \textbf{διατήρηση μουσικολογικά τεκμηριωμένων διακρίσεων και συγγενειών στην αναλυτική γραφή}, όχι την ιστορική πιστότητα της εξήγησης προς το παλαιό άκουσμα.

\subsection{Επαγόμενη τμηματοποίηση και οικογένειες}\label{sec:induced-segmentation}

Ως μεθοδολογικά διαφορετικό μοντέλο αναφοράς χρησιμοποιείται μη επιβλεπόμενη τμηματοποίηση μεταβλητού μήκους της οικογένειας Pitman--Yor, δηλαδή πιθανοτικού μοντέλου που προτείνει όρια τμημάτων από τις επαναλήψεις του υλικού χωρίς ετικέτες θέσεων, εκπαιδευμένη αποκλειστικά στο υλικό εκπαίδευσης και αξιολογημένη στα ίδια έργα και στις ίδιες ομάδες συγγενών πηγών που έχουν διατηρηθεί εκτός εκπαίδευσης. Τα επαγόμενα τμήματα οργανώνονται σε λεξιλόγιο οικογενειών χωρίς ετικέτες θέσεων και αντιπαραβάλλονται με το επαληθευμένο οικογενειακό υποσύνολο του καταλόγου και τις δημοσιευμένες οικογένειες. Το μοντέλο αναφοράς παραμένει \textbf{εκτός} του συμμετρικού πλέγματος των δώδεκα θεωρητικά προκαθορισμένων συνθηκών, ώστε η διασταυρωμένη διάταξη \(U\times F\) να διατηρεί καθαρή ερμηνεία.

Η επαγόμενη συνθήκη αξιολογείται στα ίδια τέσσερα επίπεδα γνωρισμάτων και δεν αποκτά πρόσθετη πρόσβαση σε μελλοντική πληροφορία στην ανάλυση της μουσικής διαδοχής. Το κρίσιμο ερώτημα είναι αν μια μονάδα που βελτιστοποιείται από τα δεδομένα διατηρεί επίσης τη μουσικολογική ταυτότητα, την τροπική λειτουργία και το εκτελεστικό αποτύπωμα της θέσης. Οι πλήρεις κανόνες επαγωγής, ομαδοποίησης, αντιστοίχισης και γεωμετρικών διαγνωστικών ελέγχων δίνονται στις §§~\ref{app:repr-induced} και \ref{app:robust-repr}.

\paragraph{Μουσικολογική αιτιολόγηση των γνωρισμάτων.} Οι παρατηρήσεις του Αρβανίτη (2010, §5.2) για τη συλλαβική εφαρμογή και τη σημειογραφική μεταβολή αιτιολογούν τη χωριστή διατήρηση των επιπέδων $N$, $T$ και $L$. Οι συγκρίσεις αναπαραστάσεων του Holzapfel (2010, κεφ.~4, 7) υποστηρίζουν την εμπειρική επιλογή και την ανεξάρτητη αξιολόγηση της μελωδικής συγγένειας. Η M1 συγκρίνει τις δώδεκα κωδικοποιήσεις εισόδου $U\times F$ με κοινό πρωτόκολλο. Η σύγκριση μεμονωμένου γεγονότος, επαγόμενου τμήματος και ρητής θέσης αποτελεί χωριστό ερώτημα δομικής κλίμακας. Η διαγνωστική σημασία των χρονικών και κειμενικών γνωρισμάτων κρίνεται στα συγκεκριμένα δεδομένα, με τα όρια καθορισμένα ανεξάρτητα από τα προς πρόβλεψη κειμενικά ή χρονικά παράγωγα.

\subsection{Επιλογή στα σύνολα ανάπτυξης και οριστικοποίηση της εισόδου του EISE}

Η επιλογή γίνεται μόνο στα δεδομένα ανάπτυξης, με διαχωρισμό κατά ολόκληρες ομάδες συγγενών πηγών. Η $R^\star$ επιλέγεται με μέγιστο $F_{1,\mathrm{seg+fam}}$ για πλήρεις εμφανίσεις με σωστή οικογένεια στην προκαθορισμένη αρχιτεκτονική A και με $G_A\ge g_{\min}$ σε γνωστά όρια. Το $g_{\min}$ καθορίζεται πριν από την εκπαίδευση του πλέγματος. Αποτελεί τη λειτουργική είσοδο της M2, ανεξάρτητα από το αν τεκμηριώνεται μοναδική στατιστική υπεροχή. Το σύνολο $\mathcal R^\star$ αναφέρει χωριστά την αβεβαιότητα της κατάταξης και της οικογενειακής επάρκειας· δεν ορίζει ισοδύναμες αναπαραστάσεις ούτε επιβάλλει κύρια M2 για κάθε μέλος. Αναφέρεται υποχρεωτικά η οικογενειακή επίδοση που ενδεχομένως θυσιάζεται έναντι της καλύτερης οικογενειακής συνθήκης. Η $R^\star$ και ο κανόνας μιας ενδεχόμενης πρόσθετης αναπαράστασης ευαισθησίας οριστικοποιούνται πριν από την κύρια εκπαίδευση της M2. Η τελική αξιολόγηση πραγματοποιείται αργότερα σε χωριστές ομάδες συγγενών πηγών που δεν χρησιμοποιήθηκαν για ανάπτυξη. Το ιστορικό σύνολο ελέγχου της §~\ref{sec:hist-discriminability} και τα ανεξάρτητα εκτελεστικά και ανθρώπινα αποτελέσματα δεν συμμετέχουν στον κανόνα επιλογής.

Ο ακριβής κανόνας επιλογής, τα διαστήματα αβεβαιότητας, ο χειρισμός ισοπαλιών και η πολιτική πολλαπλότητας δίνονται στην §~\ref{app:stats-selection}. Η γεωμετρία αξιολόγησης ελέγχεται χωριστά ώστε τυχόν πλεονέκτημα της μονάδας γεγονότος να μην προκύπτει απλώς επειδή αποτελεί το κοινό επίπεδο αναφοράς.

\section{EISE και κύρια ελεγχόμενη σύγκριση B2--EISE}\label{sec:eise}

Η M2 συγκρίνει δύο εκδοχές του ίδιου μοντέλου με αντιστοιχισμένες συνθήκες εκπαίδευσης. Η B2 χρησιμοποιεί επιλέξιμα ζεύγη πραγματώσεων της ίδιας οικογένειας χωρίς τεκμηριωμένη άμεση ιστορική εξηγητική ή αντιγραφική εξάρτηση, ενώ το EISE χρησιμοποιεί ζεύγη με επαληθευμένη εξηγητική σχέση, αντιστοιχία του συγκεκριμένου χωρίου και ανεξάρτητα τεκμηριωμένη κοινή οικογενειακή ένταξη. Η αξιολόγηση εξετάζει ποια εκδοχή αποδίδει καλύτερα τη δομή της παραλλαγής που μετριέται ανεξάρτητα στις σύγχρονες εκτελέσεις. Τα κριτήρια επιλεξιμότητας και η επάρκεια του κοινού πεδίου οριστικοποιούνται πριν από την εξέταση της $\Delta_{\mathrm{hist}}$.

\paragraph{Το μουσικολογικό νόημα της σύγκρισης.} Η συγκριτική έρευνα του Σπανουδάκη (2017) για τα κεκραγάρια συνδέει τη μετροφωνική δομή, το τροπικό πλαίσιο και την εκάστοτε εξηγητική ανάπτυξη. Η σύνδεση αυτή προσφέρει συγκεκριμένο μουσικό υπόβαθρο στο ερώτημα της M2: ποια πρόσθετη πληροφορία παρέχει μια επαληθευμένη εξηγητική σχέση για τη διαφοροποίηση πραγματώσεων που αναγνωρίζονται ήδη ως μέλη της ίδιας οικογένειας; Οι συνθήκες B2 και EISE μοιράζονται την οικογενειακή πληροφορία και τον υπολογιστικό προϋπολογισμό. Το $\Delta_{\mathrm{hist}}$ εκτιμά τη διαφορά της συμφωνίας τους με το προκαθορισμένο ανεξάρτητο εκτελεστικό κριτήριο στο επιλέξιμο κοινό πεδίο. Ο ρόλος του εξηγητή, η μετροφωνική οργάνωση και η προέλευση των πηγών τεκμηριώνουν την ανάγκη των υφιστάμενων ελέγχων αντιστοίχισης και εξάρτησης.

\subsection{Μαθησιακό πρόβλημα και βασική αρχιτεκτονική}

Η $R^\star$ είναι η λειτουργική επιλογή της M1 και αποτελεί την κοινή είσοδο των B1, B2 και EISE στην κύρια M2. Το σύνολο $\mathcal R^\star$ χρησιμοποιείται για την αναφορά της αβεβαιότητας και για την προκαθορισμένη επιλογή έως μίας πρόσθετης αναπαράστασης ευαισθησίας, σύμφωνα με την §~\ref{app:stats-selection}. Η M2 δεν ανατροφοδοτεί την επιλογή της M1.

\paragraph{Βιβλιογραφική και αρχιτεκτονική θεμελίωση.}
Η αρχιτεκτονική \textbf{Exegesis-Invariant Segmental Encoder (EISE)} αποτελεί τεκμηριωμένη βασική επιλογή για το συγκεκριμένο ερευνητικό πρόβλημα. Η επιλογή της στηρίζεται σε κριτική σύνθεση τεσσάρων γραμμών έρευνας που αντιστοιχούν άμεσα στις απαιτήσεις της διατριβής. Τα υπό συνθήκη τυχαία πεδία (conditional random fields, CRF), στην ημι-Μαρκοβιανή τους μορφή, επιτρέπουν να βαθμολογούνται ολόκληρα τμήματα μεταβλητού μήκους και να εκτιμώνται από κοινού όρια και ετικέτες (Sarawagi--Cohen, 2004), ενώ τα νευρωνικά τμηματικά μοντέλα δείχνουν πώς η αναπαράσταση ενός τμήματος μπορεί να μαθαίνεται από τα γεγονότα που το συγκροτούν (Kong κ.ά., 2016). Η βιβλιογραφία των ελλιπών επισημειώσεων δείχνει ότι η απουσία ετικέτας δεν πρέπει να μετατρέπεται μηχανικά σε αρνητικό παράδειγμα και ότι η αβεβαιότητα μπορεί να ενσωματώνεται στην πιθανοφάνεια ή σε λανθάνουσες δομές (Tsuboi κ.ά., 2008· Effland--Collins, 2021). Στη μουσική, η άμεση πρόβλεψη γεγονότων ως ολοκληρωμένων χρονικών τμημάτων έχει ήδη υποστηριχθεί με νευρωνικά ημι-Μαρκοβιανά CRF (Yan κ.ά., 2021). Η μελωδική τμηματοποίηση και η ανακάλυψη επαναλαμβανόμενων προτύπων δείχνουν τα όρια της καθαρά αλγοριθμικής τμηματοποίησης, ενώ η μάθηση μετρικής (metric learning) δείχνει πώς οι αποστάσεις μεταξύ αναπαραστάσεων μπορούν να εκπαιδευτούν ώστε να αποτυπώνουν σχέσεις μεταξύ μουσικών τμημάτων (Conklin--Anagnostopoulou, 2006· Karsdorp--van Kranenburg--Manjavacas, 2019· Lanz--Hajič, 2025).

Ο EISE επιλέγεται ως \textbf{βασική αρχιτεκτονική} επειδή συνδυάζει τις απαιτήσεις του συγκεκριμένου προβλήματος: θέσεις μεταβλητής έκτασης, αβέβαια όρια και μερική επισημείωση, κοινή οικογενειακή ταυτότητα με διατήρηση της διαφοροποίησης της συγκεκριμένης πραγμάτωσης και ελεγχόμενη αλλαγή μόνο της προέλευσης των σχεσιακών ζευγών μεταξύ B2 και EISE. Η πρόσθετη αξία της ημι-Μαρκοβιανής αποκωδικοποίησης ελέγχεται στα συμβολικά δεδομένα ανάπτυξης. Συγκρίσεις σε άλλα πεδία δείχνουν ότι απλούστερα CRF μπορούν να επιτυγχάνουν συγκρίσιμη επίδοση όταν η τμηματική δομή δεν προσθέτει επαρκή πληροφορία (Zaratiana κ.ά., 2023). Η τελική διαδικασία οριστικοποιείται με τον περιορισμένο έλεγχο που προβλέπεται πριν από την κύρια M2. Η αναλυτική βιβλιογραφική και τεχνική σύγκριση δίνεται στο Παράρτημα~\ref{app:eise-rationale}.

Ο EISE δέχεται την οριστικοποιημένη συμβολική ακολουθία και υπολογίζει μία τοπική κατάσταση για κάθε μελωδικό γεγονός. Για κάθε υποψήφιο τμήμα, οι αντίστοιχες καταστάσεις συνοψίζονται χωρίς να προϋποτίθεται γνωστή οικογένεια. Από τη σύνοψη προκύπτει η αναπαράσταση οικογενειακής ταυτότητας $z_{\mathrm{id}}$, ενώ οι τοπικές καταστάσεις διατηρούνται για τη μελέτη της συγκεκριμένης πραγμάτωσης.

Η αρχική βασική υλοποίηση χρησιμοποιεί \textbf{ημι-Μαρκοβιανό CRF μηδενικής τάξης}. Το μοντέλο βαθμολογεί ολόκληρα τμήματα μεταβλητού μήκους και συνδυάζει την οριοθέτηση με την οικογενειακή απόδοση. Στις μερικώς επισημειωμένες ακολουθίες, οι μη διαθέσιμες πληροφορίες παραμένουν λανθάνουσες και η πιθανοφάνεια αθροίζει πάνω στις αναλύσεις που είναι συμβατές με την πραγματικά διαθέσιμη τεκμηρίωση, αντί να αντιμετωπίζει την απουσία επισημείωσης ως αρνητική ετικέτα.

Τη δομημένη εκπαίδευση συμπληρώνει η σχεσιακή εποπτεία, η οποία εφαρμόζεται στον χώρο της οικογενειακής ταυτότητας. Τα θετικά ζεύγη δηλώνουν πραγματώσεις που η συγκεκριμένη συνθήκη εκπαίδευσης ζητεί από το μοντέλο να αναπαραστήσει ως συγγενείς, ενώ η οικογενειακή βαθμολόγηση και τα πρότυπα αναφοράς διατηρούν τις διακρίσεις μεταξύ διαφορετικών οικογενειών. Παράλληλος τοπικός στόχος ελέγχει ότι η κοινή οικογενειακή ταυτότητα δεν επιτυγχάνεται με απώλεια της πληροφορίας που χρειάζεται για την εσωτερική παραλλαγή.

Η M1 επιλέγει την $R^\star$ με την προκαθορισμένη κοινή τμηματική αρχιτεκτονική A και την κύρια μετρική $F_{1,\mathrm{seg+fam}}(R,A)$, υπό το κατώφλι $G_A(R)\ge g_{\min}$ σε γνωστά όρια (Παράρτημα~\ref{app:stats-selection}). Η αυτοτελής σύγκριση A/B/C εκτελείται και στις δώδεκα αναπαραστάσεις, με κύριο μέγεθος τη μέση αντίθεση A--C. Η ίδια A αποτελεί την τμηματική υποδομή της Α2 και της M2, με B0 ίση με $A(R^\star)$. Το $\mathcal R^\star$ καταγράφει τη χωριστή στατιστική αβεβαιότητα.

Οι τεκμηριωμένες ιστορικές ακμές $E_{\mathrm{ex}}^{\mathrm{train}}$ οργανώνουν τα θετικά ζεύγη του EISE. Η οικογενειακή αναγνώριση χρησιμοποιεί πρότυπα αναφοράς και διατηρεί χωριστά το τεκμηριωμένο υπόβαθρο, την περίπτωση που δεν αποδίδεται με ασφάλεια σε γνωστή οικογένεια και την απόφαση αποχής. Η πλήρης μαθηματική προδιαγραφή, οι συναρτήσεις απώλειας, η αναγνώριση ανοικτού συνόλου, ο περιορισμένος έλεγχος επιλογής της τμηματικής διαδικασίας και η λανθάνουσα μετρική δίνονται στο Παράρτημα~\ref{app:eise}.

\paragraph{Διάκριση ιστορικών τροπικών μαρτυριών και εξηγητικών ζευγών.} Η σχέση «ἀπὸ παραλλαγῆς»/«ἀπὸ μέλους» τεκμηριώνει ιστορικά διαφορετικές τροπικές αναγνώσεις της ίδιας μελωδικής περιοχής (Troelsgård, 2011, σ.~69· Αλεξάνδρου--Troelsgård, 2020, σσ.~178--179· Πλουσιαδηνός, έκδ. Μποτωνάκης, 2014, σ.~352). Η σχέση αυτή αποθηκεύεται ως θεωρητικό/τροπικό τεκμήριο και δεν αρκεί αυτοτελώς για τη δημιουργία θετικού \texttt{co\_explanation} ζεύγους. Η κύρια EISE απαιτεί τις ήδη προκαθορισμένες ανεξάρτητα επαληθευμένες συν-εξηγήσεις συγκεκριμένου κοινού χωρίου και ανεξάρτητη οικογενειακή κρίση.

\subsection{Τέσσερις συγκρίσιμες συνθήκες: \texorpdfstring{$B0/B1/B2/\mathrm{EISE}$}{B0/B1/B2/EISE}}

Οι τέσσερις συνθήκες οργανώνουν μια σειρά ελεγχόμενων συγκρίσεων. Η B0 ταυτίζεται με το $A(R^\star)$ και δείχνει τι παρέχει η συμβολική περιγραφή υπό την προκαθορισμένη τμηματική υποδομή, η B1 τι προσθέτει η κοινή μαθησιακή υποδομή, η B2 τι προσθέτει η εποπτεία με επιμελημένες σχέσεις κοινής οικογένειας και η EISE συνθήκη εξετάζει την αντικατάσταση αυτών των θετικών σχέσεων από τεκμηριωμένες εξηγητικές σχέσεις. Η κύρια B2--EISE αντίθεση διατηρεί κοινή αρχιτεκτονική, επιλέξιμα δεδομένα, εποπτεία οικογενειακής ένταξης και υπολογιστικό προϋπολογισμό, ώστε η προέλευση των θετικών ζευγών να αποτελεί την κεντρική συστηματική διαφορά.

Η ειδική αξία της εξηγητικής εποπτείας ελέγχεται με τέσσερις συνθήκες πάνω στα ίδια επιλέξιμα δεδομένα:

{\footnotesize
\begin{xltabular}{0.98\linewidth}{@{}p{0.09\linewidth}p{0.25\linewidth}Y@{}}
\toprule
\textbf{Συνθήκη} & \textbf{Εποπτεία / εκτίμηση} & \textbf{Ερευνητικό ερώτημα}\tabularnewline
\midrule
\endfirsthead
\toprule
\textbf{Συνθήκη} & \textbf{Εποπτεία / εκτίμηση} & \textbf{Ερευνητικό ερώτημα}\tabularnewline
\midrule
\endhead
\midrule
\multicolumn{3}{r}{\footnotesize Συνέχεια στην επόμενη σελίδα}\tabularnewline
\endfoot
\bottomrule
\endlastfoot

B0 & Τμηματική αναφορά $A(R^\star)$ χωρίς σχεσιακή μάθηση αναπαράστασης & Πόσο από το σύγχρονο προφίλ εκτελεστικής μεταβλητότητας εξηγεί ήδη η καλύτερη άμεση συμβολική περιγραφή της ιστορικής παραλλαγής;\tabularnewline
\midrule
B1 & Ίδιος τμηματικός κωδικοποιητής, ίδια μερική τμηματοποίηση και οικογενειακή εποπτεία, ίδια πρότυπα αναφοράς και ίδιος βοηθητικός έλεγχος τοπικής πληροφορίας, \textbf{χωρίς πρόσθετη σχεσιακή εποπτεία} & Πόση συμφωνία με το ανεξάρτητο εκτελεστικό προφίλ προσφέρει η κοινή νευρωνική υποδομή χωρίς πρόσθετη σχεσιακή εποπτεία;\tabularnewline
\midrule
B2 & Ίδιος κωδικοποιητής με θετικά ζεύγη που σχηματίζονται εκ νέου εντός της ίδιας οικογένειας & Ποια συμφωνία με το ίδιο ανεξάρτητο εκτελεστικό κριτήριο αξιολόγησης επιτυγχάνει η αντιστοιχισμένη οικογενειακή εποπτεία;\tabularnewline
\midrule
EISE & Ίδιος κωδικοποιητής με εξηγητικές ακμές τεκμηριωμένες από τις ιστορικές πηγές & Πώς μεταβάλλεται η συμφωνία με το προκαθορισμένο ανεξάρτητο εκτελεστικό κριτήριο αξιολόγησης όταν τα θετικά ζεύγη οργανώνονται μέσω τεκμηριωμένων εξηγητικών σχέσεων αντί μέσω αντιστοιχισμένων σχέσεων μεταξύ πραγματώσεων της ίδιας οικογένειας;\tabularnewline
\end{xltabular}
}\par

Μετά την οριστικοποίηση της αρχιτεκτονικής, οι B1, B2 και EISE χρησιμοποιούν κοινή προδιαγραφή, κοινούς κανόνες προεπεξεργασίας και εκπαίδευσης, κοινή πολιτική υποψήφιων τμημάτων, κοινά επιλέξιμα δεδομένα και αντιστοιχισμένο υπολογιστικό προϋπολογισμό. Κάθε συνθήκη εκπαιδεύεται χωριστά και αποκτά τις δικές της παραμέτρους, πρότυπα αναφοράς και λανθάνουσα μετρική σύμφωνα με την ίδια διαδικασία. Στην κρίσιμη σύγκριση B2--EISE μεταβάλλεται η οργάνωση των θετικών σχέσεων: ο B2 χρησιμοποιεί αντιστοιχισμένα ζεύγη κοινής οικογένειας και ο EISE εξηγητικές ακμές τεκμηριωμένες από τις ιστορικές πηγές. Ο $B2_{\mathrm{sim}}$, όταν η απογραφή του ΠΕ0 παρέχει επαρκή κάλυψη, προσθέτει ως αυστηρότερο μη ιστορικό σημείο αναφοράς τις ισχυρότερες ενδοοικογενειακές αντιστοιχίσεις που επάγει η οριστικοποιημένη B1. Η πλήρης αντιστοίχιση έκθεσης, προέλευσης, δυσκολίας και υπολογιστικού κόστους δίνεται στο Παράρτημα~\ref{app:match-b2}.

Ο $T_{\mathrm{evt}}$ αποτελεί το κύριο μοντέλο αναφοράς ακολουθίας στο επίπεδο ατομικών γεγονότων. Συμπληρωματικά συμβολικά και νευρωνικά μοντέλα αναφοράς περιγράφονται στο Παράρτημα~\ref{app:eise-baselines}, ενώ η συνθήκη Pitman--Yor παραμένει το χωριστό μοντέλο αναφοράς τμηματοποίησης της §~\ref{sec:induced-segmentation}. Ο $T_{\mathrm{evt}}$ αξιολογείται σε κοινούς στόχους και ίδιους διαχωρισμούς δεδομένων· οι ακατέργαστες τιμές perplexity διαφορετικών λεξιλογίων χρησιμοποιούνται μόνο διαγνωστικά. Οι τεχνικές προδιαγραφές και οι αναλύσεις αφαίρεσης συγκεντρώνονται στα Παραρτήματα~\ref{app:event-transformer} και \ref{app:eise-baselines}. Η τεκμηρίωση των κριτηρίων αξιολόγησης και της προέλευσης των δομικών επισημειώσεων ακολουθεί τη σχετική βιβλιογραφία MIR (Urbano κ.ά., 2013· Smith κ.ά., 2011).

\subsection{Σύγκριση ιστορικού και ανεξάρτητου εκτελεστικού προφίλ}\label{sec:eise-transfer}

\subsubsection{Αναγκαίο ανεξάρτητο εκτελεστικό κριτήριο για τη M2}\label{sec:m4-min}

Η M2 εξετάζει αν οι περιοχές μιας θέσης που παρουσιάζουν μεγαλύτερη ιστορική παραλλαγή αντιστοιχούν σε περιοχές με μεγαλύτερη μεταβλητότητα μεταξύ σύγχρονων εκτελεστών. Κάθε μοντέλο παράγει το δικό του ιστορικό προφίλ παραλλαγής. Η M2 συγκρίνει πόσο συμφωνεί καθένα από αυτά τα προφίλ με την ανεξάρτητη εκτελεστική μέτρηση.

Η κύρια αξιολόγηση B2--EISE χρησιμοποιεί το αναγκαίο μέρος της M4: επιλέξιμες ηχογραφήσεις του ίδιου υλικού από πολλαπλούς εκτελεστές, αξιόπιστη ευθυγράμμιση με τη γραπτή πηγή και εκτίμηση της διασποράς μεταξύ εκτελεστών στην ίδια μουσική συντεταγμένη. Το ΠΕ0 επαληθεύει την πρόσβαση, την ανεξαρτησία, την κάλυψη και την επάρκεια του κοινού πεδίου αντιστοίχισης ιστορικών πραγματώσεων και σύγχρονων εκτελέσεων πριν από την τελική σύγκριση. Το ανεξάρτητο εκτελεστικό κριτήριο υπολογίζεται στο κοινό συγκρίσιμο υλικό που τεκμηριώνει το ΠΕ0, με τους ίδιους κανόνες για B0, B2 και EISE. Το λεκτικό καθεστώς και ο τύπος μουσικού τμήματος καταγράφονται χωριστά και μπορούν να χρησιμοποιούνται σε περιγραφές ανά προκαθορισμένη υποομάδα και σε αναλύσεις ευαισθησίας όταν υπάρχει επαρκής κοινή κάλυψη· δεν εισάγεται προκαταβολικός καθολικός αποκλεισμός κρατημάτων από το κύριο πεδίο. Οι ευρύτερες αναλύσεις εκτελεστικής παραλλαγής οργανώνονται χωριστά στη M4 σύμφωνα με την πραγματική κάλυψη.

Η προκαθορισμένη εκτελεστική μεταβλητή αξιολόγησης του κεντρικού εξωτερικού ελέγχου της M2 είναι η διασπορά μεταξύ των εκτελεστών του κανονικοποιημένου λογαριθμικού χρονισμού, εφόσον το ΠΕ0 επιβεβαιώσει επαρκή αξιοπιστία μέτρησης· το τονικό ύψος παραμένει προκαθορισμένη δευτερεύουσα μεταβλητή αξιολόγησης. Η αξιολόγηση στα ανεξάρτητα εκτελεστικά δεδομένα γίνεται \textbf{με σταθερές τις παραμέτρους του κωδικοποιητή}: μετά την οριστικοποίηση της εκπαίδευσης δεν επιτρέπεται περαιτέρω προσαρμογή του κωδικοποιητή, της μετρικής ή των κατωφλίων από τα εκτελεστικά αποτελέσματα. Η πλήρης προδιαγραφή των μεταβλητών δίνεται στα Παραρτήματα~\ref{app:stats} και \ref{app:impl}.

Η απαιτούμενη κάλυψη ορίζεται από τις ανεξάρτητες οικογένειες, τα έργα, τις ομάδες συγγενών πηγών και τους εκτελεστές, μαζί με την αξιοπιστία της μέτρησης και τις προσομοιώσεις ακρίβειας του ΠΕ0. Σε GREEN χρησιμοποιείται η λεπτομερής εσωτερική συντεταγμένη· σε AMBER οι προκαθορισμένες μεγαλύτερες μουσικές περιοχές. Σε RED τα διαθέσιμα προφίλ αναφέρονται διερευνητικά, ενώ η πλήρης απουσία μετρήσιμου ανεξάρτητου εκτελεστικού κριτηρίου αξιολόγησης καταγράφεται ως αδυναμία αξιόπιστης εκτίμησης. Η επιλογή άλλης ήδη προβλεπόμενης μεταβλητής επιτρέπεται μόνο με τον κανόνα αξιοπιστίας των μετρήσεων του ΠΕ0, πριν από την επιβεβαιωτική αξιολόγηση (§~\ref{app:stats-eise-transfer}).

\subsubsection{Μέτρο συμφωνίας και συγκριτικές αντιθέσεις}

Ο ανεξάρτητος έλεγχος συγκρίνει δύο προφίλ πάνω σε αντίστοιχες εσωτερικές περιοχές $z$ της ίδιας οικογένειας. Το $A(z)$ περιγράφει την ιστορικά τεκμηριωμένη μεταβλητότητα της περιοχής και το $\Gamma(z)$ τη μεταβλητότητα που μετράται στις σύγχρονες εκτελέσεις. Για κάθε συνθήκη $m$, το $S_m$ συνοψίζει τη συσχέτιση κατάταξης του ιστορικού και του ανεξάρτητου εκτελεστικού προφίλ στο κοινό επιλέξιμο πεδίο, με ίσο βάρος ανά επιλέξιμη οικογένεια και, για τις εκπαιδευμένες συνθήκες, μέσο όρο στις προκαθορισμένες επαναλήψεις. Η σύγχρονη εκτελεστική μεταβλητότητα λειτουργεί εδώ ως εξωτερικό κριτήριο αξιολόγησης της οργανωμένης παραλλαγής. \textbf{Θετική συμφωνία με το ανεξάρτητο εκτελεστικό προφίλ} τεκμηριώνεται όταν
\[
S_{\mathrm{EISE}}>0
\]
με υποστήριξη από το προκαθορισμένο διάστημα. Το $S_{B0}>0$ αναφέρεται χωριστά ως τεκμήριο διατήρησης της γεωμετρίας παραλλαγής χωρίς εκπαίδευση αναπαράστασης και δεν προστίθεται ως τρίτη κύρια EISE αντίθεση.

Οι δύο κύριες προκαθορισμένες αντιθέσεις είναι
\[
\Delta_{\mathrm{hist}}=S_{\mathrm{EISE}}-S_{B2},
\qquad
\Delta_{\mathrm{learn}}=S_{\mathrm{EISE}}-S_{B0}.
\]

Η τιμή του $\Delta_{\mathrm{hist}}$ διαβάζεται ως διαφορά δύο επιδόσεων στο ίδιο ανεξάρτητο εκτελεστικό κριτήριο. Θετικό πρόσημο δείχνει μεγαλύτερο $S$ για την EISE συνθήκη, αρνητικό μεγαλύτερο $S$ για την B2 και τιμή κοντά στο μηδέν μικρή ή μη διακριτή διαφορά υπό το συγκεκριμένο πρωτόκολλο και την εκτιμώμενη αβεβαιότητα. Η ερμηνεία συνοδεύεται πάντοτε από διάστημα αβεβαιότητας και από το αποτέλεσμα του ελέγχου επάρκειας δεδομένων. Το μέγεθος εξετάζει αν η τεκμηριωμένη προέλευση των θετικών ζευγών προσφέρει πρόσθετη προβλεπτική οργάνωση για αυτό το ανεξάρτητο κριτήριο αξιολόγησης. Η ιστορική ερμηνεία των εξηγητικών σχέσεων παραμένει έργο της ανάλυσης των πηγών και της μουσικολογικής ερμηνείας.
Η $\Delta_{\mathrm{learn}}$ αποτιμά χωριστά τη διαφορά της γεωμετρίας που προκύπτει από την εκπαίδευση έναντι της ισχυρής άμεσης συμβολικής περιγραφής. Η αιτιώδης ιστορική ερμηνεία, η δυνατότητα ανάκτησης της ίδιας πληροφορίας από άλλες αναπαραστάσεις και το είδος του μετασχηματισμού εξετάζονται χωριστά. Οι δύο κατηγορίες ιστορικών σχέσεων διακρίνουν την αναλυτική διαφοροποίηση της ίδιας μουσικής ταυτότητας ($E_{\mathrm{an}}$) από τη μελωδική παραλλαγή μέσα σε τεκμηριωμένη ιστορική σχέση ($E_{\mathrm{tr}}$· κριτήρια στο Παράρτημα~\ref{app:corpora-supervision}). Η στρωμάτωσή τους εξετάζει ειδικά τη διαφοροποίηση του αποτελέσματος ανά είδος τεκμηριωμένης σχέσης. Η πλήρης διαχρονική--συγχρονική ερμηνεία απαιτεί θετική συμφωνία με το ανεξάρτητο εκτελεστικό προφίλ, θετικές συγκριτικές αντιθέσεις και διατήρηση του αποτελέσματος στον προκαθορισμένο έλεγχο με προσαρμογή για προκαθορισμένους συγχυτικούς παράγοντες.

Η πύλη επάρκειας του κοινού πεδίου B2--EISE συμβολίζεται ως $G_{\mathrm{pair}}$, ενώ η ήδη ορισμένη πύλη επάρκειας του ανεξάρτητου εκτελεστικού κριτηρίου συμβολίζεται ως $G_{\mathrm{perf}}$. Η $\Delta_{\mathrm{hist}}$ διατηρεί επιβεβαιωτικό καθεστώς μόνο στο πεδίο που υποστηρίζουν και οι δύο πύλες: GREEN/GREEN επιτρέπει την πλήρη προκαθορισμένη σύγκριση, οποιοσδήποτε συνδυασμός GREEN/AMBER ή AMBER/AMBER περιορίζει το συμπέρασμα στο αυστηρότερο κοινό επιλέξιμο πεδίο, ενώ RED σε κρίσιμη πύλη δεν επιτρέπει πλήρη επιβεβαιωτικό ισχυρισμό για το αντίστοιχο estimand. Οι συσχετίσεις που υπολογίζονται χωριστά ανά οικογένεια, οι ομάδες ιστορικών και εκτελεστικών δεδομένων $K\times L$ (ορισμός στο Παράρτημα~\ref{app:stats-eise-transfer}), η κοινή αναδειγματοληψία, τα ταυτόχρονα διαστήματα εμπιστοσύνης, η διόρθωση Holm και το προσαρμοσμένο ιεραρχικό μοντέλο κεντρικής τάσης και κλίμακας (location--scale) ορίζονται στις §§~\ref{app:stats-eise-transfer} και \ref{app:match-resampling}. Οι αντίστοιχες εκτελεστικές υποθέσεις παραμένουν δευτερεύουσες και αναφέρονται με αυτοτελή μουσικολογική ερμηνεία.

\subsection{Σύνδεση του EISE με την ανθρώπινη γνωσιακή αξιολόγηση}\label{sec:eise-human-link}

Μετά την οριστικοποίηση των B2/EISE μοντέλων, οι αναπαραστάσεις τους μεταφέρονται στη M3b ως εναλλακτικές υποθέσεις για την οργάνωση της παραλλαγής μέσα στην ίδια οικογένεια. Η σύγκριση πραγματοποιείται σε κοινές πραγματώσεις εκτός εκπαίδευσης και αξιολογεί την υπολειμματική ευθυγράμμιση των αναπαραστάσεων αφού ληφθούν υπόψη οι προκαθορισμένες συμβολικές και ακουστικές μήτρες αποστάσεων μεταξύ των ίδιων ερεθισμάτων (representational dissimilarity matrices, RDM). Η ανθρώπινη αξιολόγηση δεν ανατροφοδοτεί την εκπαίδευση των μοντέλων και δεν χρησιμοποιείται για επιλογή οικογενειακών ετικετών ή υπερπαραμέτρων. Το πλήρες ανθρώπινο πρωτόκολλο δίνεται στο Παράρτημα~\ref{app:human}.

\subsection{Μάθηση τεκμηριωμένων μετασχηματισμών}\label{sec:atl}

Η σύγκριση B2--EISE εξετάζει αν η συν-εξήγηση του ίδιου χωρίου βοηθά την αναπαράσταση εκεί όπου υπάρχει. Ένα συγγενές ερώτημα αφορά τις οικογένειες για τις οποίες σώζεται μία μόνο εξήγηση: μπορούν οι κανονικότητες που εμφανίζονται ανάμεσα σε παράλληλες εξηγήσεις να μαθευτούν και να μεταφερθούν σε αυτές; Για τη διερεύνησή του προβλέπεται πιθανοτικό μοντέλο στοίχισης και μεταβολής, στη γραμμή των μουσικών μέτρων απόστασης επεξεργασίας, που ορίζουν την απόσταση δύο μελωδιών ως το ελάχιστο συνολικό κόστος εισαγωγών, διαγραφών, αντικαταστάσεων, συγχωνεύσεων και κατατμήσεων φθόγγων (Mongeau \& Sankoff, 1990), και της εκτίμησης των κοστών αυτών από παραδείγματα αντί για τον εκ των προτέρων καθορισμό τους (Ristad \& Yianilos, 1998). Οι πιθανότητες ανάπτυξης, σύμπτυξης και τροποποίησης εκτιμώνται από εξηγήσεις του ίδιου χωρίου και εξαρτώνται από τη ζώνη μέσα στη θέση, τον ήχο και τον λειτουργικό ρόλο. Επειδή οι δύο εξηγήσεις μοιράζονται το κείμενο, η στοίχιση γίνεται μέσα στη συλλαβή. Το μοντέλο περιγράφει αντιστοιχίσεις ανάμεσα σε δύο αποδόσεις ενός κοινού χωρίου και δεν υποθέτει ότι ο ένας εξηγητής μετασχημάτισε τη γραφή του άλλου.

Η κύρια δοκιμασία είναι προβλεπτική. Οικογένειες με παράλληλες εξηγήσεις αποκρύπτονται πλήρως από την εκτίμηση του μοντέλου, και ελέγχεται αν οι πραγματικές εξηγήσεις τους προβλέπονται καλύτερα από ό,τι με μετασχηματισμούς που ορίζει εκ των προτέρων η θεωρία και από το ίδιο μοντέλο εκτιμημένο σε ζεύγη της ίδιας οικογένειας στο ίδιο κειμενικό σημείο χωρίς εξηγητική σχέση. Η δεύτερη σύγκριση, με ίδιο μοντέλο, όγκο δεδομένων και έκθεση, απομονώνει την πληροφορία που προσθέτει η ιστορική τεκμηρίωση. Παράλληλα μετριέται ποιες διαγνωστικές σχέσεις των υπογραφών διατηρούνται όταν αλλάζει η σημειογραφική επιφάνεια: η συνολική διαστηματική πορεία ή η βαθμίδα κατάληξης μπορεί να μένει σταθερή, ενώ ο αριθμός και η διάρκεια των σημείων διαφέρουν. Η μέτρηση γίνεται σε όλα τα ζεύγη του ίδιου χωρίου και περιγράφει πόσο συχνά η εξήγηση μεταβάλλει μια σχέση· οι περιπτώσεις στις οποίες η μεταβολή μετακινεί την οικογενειακή απόδοση αναφέρονται χωριστά. Το ερώτημα συνδέεται με τη συζήτηση για τη σχέση ομοιότητας και κατηγοριοποίησης στις μελωδικές παραλλαγές (Cambouropoulos, 2009· Volk \& van Kranenburg, 2012). Εδώ όμως οι αποδεκτές παραλλαγές προέρχονται από ιστορικά τεκμήρια και όχι από υπόθεση του ερευνητή.

Αν η απογραφή το επιτρέψει, το μοντέλο μπορεί να προτείνει υποψήφιες παραλλαγές για οικογένειες χωρίς δεύτερη εξήγηση. Οι παραλλαγές αυτές ελέγχονται ως προς την ορθότητα της γραφής και τις υπογραφές και κρίνονται τυφλά από ειδικούς. Δεν αποκτούν καθεστώς ιστορικού τεκμηρίου και δεν εισέρχονται στη M2. Η υποδομή συνθετικών δεδομένων (EISE-Synth) καταγράφει την προέλευση κάθε παραγόμενου αντικειμένου, και η χρήση τους ως θετικών ζευγών εξετάζεται χωριστά, μετά την οριστικοποίηση του κύριου μοντέλου EISE. Η εκτίμηση του μοντέλου, ο έλεγχος των διαγνωστικών σχέσεων και η προβλεπτική δοκιμασία αποτελούν δεσμευμένη δευτερεύουσα μελέτη (βαθμίδα~4), εφόσον το ΠΕ0 δείξει επαρκή κάλυψη οικογενειών και ανεξάρτητων πηγών. Η μελέτη έχει δικό της χρονικό όριο στο χρονοδιάγραμμα και δεν αντλεί χρόνο από τις βαθμίδες 1--3· αν το μοντέλο δεν περάσει τους ελέγχους εκτίμησης έως το προκαθορισμένο ορόσημο, περιορίζεται στην άμεση μέτρηση των διαγνωστικών σχέσεων στα ιστορικά ζεύγη. Ένα επαρκώς τεκμηριωμένο αρνητικό αποτέλεσμα θα προσδιορίσει τα όρια της μεταφοράς. Η παραγωγή υποψήφιων παραλλαγών και η χρήση τους στη μάθηση αναπαραστάσεων μένουν επέκταση υπό προϋποθέσεις (βαθμίδα~5). Η μη ενεργοποίηση της μελέτης δεν επηρεάζει τις υποθέσεις και τα συμπεράσματα των βαθμίδων 1--3.

\paragraph{Δεύτερο ρεπερτόριο.} Ο έλεγχος μεταφοράς σε άλλη παράδοση προβλέπεται κατά προτεραιότητα στο ζναμένιο μέλος, όπως διασώζεται στις κοινότητες των Παλαιοπίστων της περιοχής της λίμνης Πείψι στην Εσθονία, καθώς και σε κοινότητες της Λετονίας. Η γραφή του κατάγεται από τη βυζαντινή νευματική σημειογραφία, και οι ποπέβκι λειτουργούν ως μελωδικοί τύποι συγγενείς προς τις θέσεις (Grigorjeva, 2023). Οι φίτες είναι μακροί μελισματικοί τύποι γραμμένοι συνοπτικά, που οι πηγές καταγράφουν και αναλυτικά με απλά σημάδια (ραζβόντ). Η σχέση τους προσφέρει λειτουργική αναλογία προς τη σχέση μεγάλης υπόστασης και εξήγησης, που θα διερευνηθεί χωρίς να προϋποτίθεται κοινός ιστορικός μηχανισμός. Όπου διαφορετικές πηγές δίνουν διαφορετικό ραζβόντ της ίδιας φίτας, η ίδια σύγκριση ---τεκμηριωμένες αναπτύξεις, κανόνες των θεωρητικών εγχειριδίων και ζεύγη χωρίς τεκμηριωμένη σχέση--- μπορεί να επαναληφθεί σε δεύτερη παράδοση με κοινή καταγωγή γραφής. Τα χειρόγραφα με παράλληλη γραμμική σημειογραφία έχουν ήδη χρησιμοποιηθεί για υπολογιστική ανάγνωση της γραφής (Gusev κ.ά., 2020). Η μελέτη ενεργοποιείται εφόσον επιβεβαιωθούν πρόσβαση σε ψηφιοποιημένες πηγές, δικαιώματα και ειδικός ανάγνωσης της γραφής, σε συνεργασία με ερευνητικούς φορείς και κοινότητες στην Εσθονία και στη Λετονία. Το γρηγοριανό μέλος, που έχει μελετηθεί εκτενώς υπολογιστικά, χρησιμοποιείται ως πλαίσιο σύγκρισης και όχι ως πεδίο νέας ανάλυσης.

\section{Μουσικολογική και γνωσιακή επικύρωση της μονάδας}\label{sec:m3-human}

\subsection{M3a: μουσικολογική θεμελίωση και ανεξάρτητη επικύρωση}\label{sec:m3a}

Το μέρος της M3a που συγκροτεί τον βασικό πυρήνα είναι η ανεξάρτητη μουσικολογική αξιολόγηση των ορίων, της οικογενειακής ταυτότητας και των τεκμηριωμένων σχέσεων μεταξύ πραγματώσεων. Οι πρόσθετες αναλύσεις της εσωτερικής οργάνωσης και της λειτουργίας των θέσεων αναπτύσσονται ως ενισχυτική τεκμηρίωση σύμφωνα με την κάλυψη. Η M3a συνδέει έτσι το οικογενειακό λεξιλόγιο με τη δομή και την πράξη της παράδοσης. Για κάθε οικογένεια οργανώνει τεκμηριωμένες εμφανίσεις, χαρακτηριστικά μελωδικής ταυτότητας, τροπική λειτουργία, κειμενική και προσωδιακή τοποθέτηση και εύρος εκτελεστικής παραλλαγής. Η ανάλυση εξετάζει ποια χαρακτηριστικά διατηρούνται μεταξύ εκδοχών και αν οι μεταβολές τους συνδέονται με διαφορετικές μουσικές λειτουργίες.

Το λεξιλόγιο ανάπτυξης και οι υποψήφιες ετικέτες που χρησιμοποιούνται στην εκπαίδευση οριστικοποιούνται πριν από τις M1--M2. Το $C_{\mathrm{StA}}$ χρησιμοποιείται ως καταλογογραφική αφετηρία και το $C_{\mathrm{occ}}$ καταγράφει τεκμηριωμένες εμφανίσεις. Η τελική μουσικολογική αξιολόγηση της M3a πραγματοποιείται χωριστά, από ανεξάρτητους ειδικούς που επιλέγονται με προκαθορισμένα κριτήρια εμπειρίας· οι κρίσεις τους συγκροτούν το $C_{\mathrm{exp}}$. Η συμφωνία, η διαφωνία και η αποχή καταγράφονται ως αποτελέσματα με ρητή προέλευση. Οι αποκλίνουσες κρίσεις τροφοδοτούν τη μουσικολογική ερμηνεία των ορίων και της εσωτερικής διαφοροποίησης των οικογενειών. Παραδοτέα είναι το τεκμηριωμένο λεξιλόγιο, οι περιγραφές της παραλλαγής και ο χάρτης αντιστοιχιών ανάμεσα στις πηγές, στις κρίσεις ειδικών και στις υπολογιστικές αναπαραστάσεις.

Η χρονική οργάνωση εξετάζεται εδώ ως χαρακτηριστικό της πραγμάτωσης και όχι ως προϋπόθεση οικογενειακής ταυτότητας. Σε μικρό διαγνωστικό υποσύνολο των ήδη επιλεγμένων συγκρίσεων, μετά την αρχική κρίση οικογενειακής ένταξης, οι ειδικοί απαντούν στο ερώτημα: \emph{«Οι διαφορές στις σημειογραφικές χρονικές αξίες και στην εσωτερική ανάπτυξη είναι συμβατές με την ίδια μουσική θέση; Αν όχι, σε ποια περιοχή εντοπίζεται η διαφορά και ποιο τεκμήριο στηρίζει την κρίση;»} Οι απαντήσεις κωδικοποιούνται ως συμβατές, μερικώς συμβατές, ασύμβατες ή ανεπαρκή στοιχεία. Όπου είναι εφικτό, σημειώνονται περιοχές σχετικής χρονικής σταθερότητας ή ελαστικότητας. Η διερευνητική αυτή καταγραφή δεν δημιουργεί νέα οικογένεια εκτιμώμενων μεγεθών και δεν επιστρέφει στην εκπαίδευση ή στον επανακαθορισμό των οικογενειών.

Οι ανεξάρτητοι ειδικοί εφαρμόζουν τις υπογραφές και σε περιπτώσεις αλλαγής συλλαβικής κατανομής, τονισμού ή αναλυτικής κλίμακας, όταν τέτοιες περιπτώσεις έχουν προαπογραφεί. Η απόφαση \texttt{member}/\texttt{non\_member}/\texttt{unresolved} συνοδεύεται από το ενεργό διαγνωστικό γνώρισμα και την πηγή του. Για τις περιπτώσεις κοινής μελωδικής υποακολουθίας αποδίδεται χωριστά η ταυτότητα στην επιλεγμένη μεγαλύτερη κλίμακα. Η διαδικασία εξετάζει αν τα κριτήρια παραμένουν εφαρμόσιμα σε νέο έργο χωρίς προσαρμογή στο αποτέλεσμα.

\paragraph{Έλεγχος διακριτικής επάρκειας των υπογραφών.}
Η M3a ελέγχει τις παγωμένες υπογραφές σε νέο υλικό. Εξετάζει δύσκολα θετικά, δύσκολα αρνητικά και περιπτώσεις αμφίβολης αναλυτικής κλίμακας. Ελέγχει επίσης αν μια υπογραφή είναι υπερβολικά ευρεία ή υπερβολικά στενή. Η αποτυχία αυτού του ελέγχου αποτελεί αποτέλεσμα της M3a. Δεν χρησιμοποιείται για εκ των υστέρων επαναετικετοποίηση της κύριας M2. Οι νέες πηγές του Gheorghiță και του Καραγκούνη παρέχουν υποψήφιο υλικό για αυτούς τους ελέγχους, όχι έτοιμο gold standard.

\subsubsection{Μουσικολογική αναγνώριση και ιστορική οργάνωση}
Η ανάλυση εξετάζει όρια, οικογενειακή ένταξη, λειτουργικό ρόλο και τεκμηριωμένες σχέσεις μεταξύ συγκεκριμένων αποδόσεων. Περιπτώσεις από το στιχηραρικό υλικό του Κανάκη μπορούν να καταγραφούν ως υποψήφια δύσκολα θετικά μόνο μετά από ανεξάρτητη τεκμηρίωση της κοινής θέσης και φέρουν \texttt{cross\_repertoire=true}; δεν χρησιμοποιούνται για εκπαίδευση ή επιλογή της $R^\star$. Η ιστορική ανάλυση $A(z)$ συγκρίνει αντίστοιχες εσωτερικές περιοχές δύο τεκμηριωμένων πραγματώσεων της ίδιας θέσης. Σε κάθε περιοχή καταγράφει διατήρηση, μουσική μεταβολή, αβεβαιότητα ή αδυναμία αξιολόγησης. Η μονάδα παρατήρησης είναι το συγκρίσιμο ιστορικό ζεύγος στην αντίστοιχη εσωτερική περιοχή. Η εξαγωγή συμπερασμάτων σε επίπεδο οικογένειας χρησιμοποιεί ιεραρχική μερική συνένωση (partial pooling) και αναδειγματοληψία στο επίπεδο της οικογένειας.

\subsubsection{Πληροφοριακή και τροπική οργάνωση}
Η πληροφοριακή ανάλυση εξετάζει την τοπική οργάνωση της εντροπίας και του πληροφοριακού περιεχομένου με αποκλειστικά προγενέστερο συμφραζόμενο. Η τροπική λειτουργία μιας πραγμάτωσης τεκμηριώνεται συνδυαστικά από τη γραπτή σημειογραφία, τις μαρτυρίες και φθορές, τις καταλήξεις, τη μελωδική πορεία, τα τοπικά κέντρα και τη συγκριτική μουσικολογική ανάλυση. Το ισοκράτημα παρέχει χωριστή πηγή μαρτυρίας: οι καταγεγραμμένες ή εκτελεστικές ισοκρατηματικές αποδόσεις διατηρούν τη δική τους προέλευση και μπορεί να περιλαμβάνουν περισσότερες από μία τεκμηριωμένες λύσεις για το ίδιο χωρίο.\footnote{Η ισοφωνία ως θεωρητική σχέση ή κατάσταση δεν ταυτίζεται με την ισοκρατηματική τροχιά, δηλαδή τη συγκεκριμένη ακολουθία φθόγγων που καταγράφεται ή εκτελείται ως ίσον.} Η σχετική ανάλυση στηρίζεται σε ελεγμένες περιπτώσεις της βιβλιογραφίας και της διατριβής του Κανάκη για τη διάκριση ορίων θέσεων και φράσεων και για ισοκρατηματικές αναγνώσεις· το άρθρο Kanakis (2021) διατηρείται ως βιβλιογραφική αναφορά χωρίς να του αποδίδεται ειδικός ισχυρισμός πέρα από όσα έχουν ελεγχθεί απευθείας.

\paragraph{Τροπικό--λειτουργικό σχήμα επισημείωσης.} Για κάθε επιλέξιμη πραγμάτωση μπορεί να αποθηκεύεται το προαιρετικό σχήμα $\Phi_x$ με πεδία ενεργής τροπικής κατάστασης, οκταηχικής οργάνωσης, τοπικού κέντρου, καταληκτικής/λειτουργικής κατάστασης και συνόλου τεκμηριωμένων ισοκρατηματικών λύσεων. Το $\Phi_x$ είναι σχήμα επισημείωσης και όχι νέο μοντέλο ή εκτιμώμενο μέγεθος· επιτρέπει τιμές \texttt{unknown}, \texttt{ambiguous} και πολλαπλές τεκμηριωμένες λύσεις. Η φραστική οριοθέτηση αποθηκεύεται χωριστά. Κάθε στοιχείο του \texttt{isokratema\_solution\_set} συνοδεύεται από τύπο πηγής, αναγνωριστικό πηγής, χρονολογία, πρόσωπο ή εκδότη, έργο/ρεπερτόριο, σχολική ή εκδοτική γραμμή όπου τεκμηριώνεται και επιστημικό καθεστώς. Θεωρητική πρόταση, έντυπη ή χειρόγραφη καταγραφή, ηχογραφημένη εκτέλεση και σύγχρονη ανακατασκευή ειδικού διατηρούνται ως διαφορετικές κατηγορίες.

\paragraph{Αποτροπή κυκλικότητας.} Αν η οριστικοποιημένη $R^\star$ περιέχει το στρώμα $M$, πρόβλεψη τροπικής κατάστασης από την ίδια αναπαράσταση αποτελεί έλεγχο εσωτερικής συνέπειας και όχι ανεξάρτητη επικύρωση της τροπικής πληροφορίας. Αντίστοιχα, μια ισοκρατηματική λύση που παράγεται από θεωρητικό κανόνα μπορεί να χρησιμοποιείται ως rule-based baseline, αλλά δεν λειτουργεί ταυτόχρονα ως target της ίδιας περίπτωσης. Επιβεβαιωτική πρόβλεψη απαιτεί ανεξάρτητη καταγραφή ή παρατήρηση.

\paragraph{I1 -- οικογενειακή συνοχή ισοκρατηματικών προφίλ.} Ως προκαθορισμένο δευτερεύον διαγνωστικό εξετάζεται, όπου η κάλυψη το επιτρέπει, αν πραγματώσεις της ίδιας οικογένειας εμφανίζουν μεγαλύτερη ομοιότητα ισοκρατηματικής τροχιάς από κατάλληλα αντιστοιχισμένες πραγματώσεις διαφορετικών οικογενειών, μετά από έλεγχο ήχου, καταληκτικής λειτουργίας, δομικής θέσης και προέλευσης της ισοκρατηματικής μαρτυρίας. Δεν εισάγεται νέο $H$ ή νέο κύριο estimand.

\paragraph{I2 -- φραστική άρθρωση και όρια θέσεων.} Ορίζονται ανεξάρτητα δύο σύνολα ορίων, $B_{\mathrm{phrase}}$ και $B_{\mathrm{position}}$. Η φραστική οριοθέτηση παράγεται μόνο στα έργα με τεκμηριωμένη ισοκρατηματική πληροφορία και στο πλήρως ελεγμένο υποσύνολο ανάπτυξης, πριν και ανεξάρτητα από την ανάγνωση του ισοκρατήματος. Για κάθε αλλαγή ισοκρατηματικής βάσης $c$ υπολογίζονται χωριστά οι αποστάσεις προς το πλησιέστερο επιλέξιμο φραστικό όριο και προς το πλησιέστερο όριο θέσης. Οι περιπτώσεις ταξινομούνται σε κοινά όρια, μόνο φράσης και μόνο θέσης. Για να μην ευνοείται το πυκνότερο σύνολο ορίων, χρησιμοποιείται μηδενική κατανομή με αναδιάταξη ή κυκλική μετατόπιση των αλλαγών του ισοκρατήματος μέσα στην ίδια προκαθορισμένη δομική ζώνη. Το τυποποιημένο διαγνωστικό γράφεται
\[
Z_b=\frac{\mathbb E_0[d_b]-d_b^{\mathrm{obs}}}{\operatorname{SD}_0(d_b)},
\]
και η αντίθεση $Z_{\mathrm{phrase-only}}-Z_{\mathrm{position-only}}$ αναφέρεται ως δευτερεύουσα ανάλυση της M3a, όχι ως νέο κεντρικό estimand. Οι I1 και I2 ενεργοποιούνται μόνο στο ήδη διαθέσιμο επιλέξιμο υποσύνολο με τις απαιτούμενες ανεξάρτητες επισημειώσεις· η μη επάρκεια αυτού του υποσυνόλου δεν εμποδίζει την ολοκλήρωση του βασικού πυρήνα της M3a.

Η ευρύτερη εκτελεστική ανάλυση της M4 συμπληρώνει τη μουσικολογική τεκμηρίωση με ανεξάρτητες μετρήσεις παραλλαγής (§~\ref{sec:m4-broader}).

\subsection{M3c: ιεραρχική μουσική πρόβλεψη και ένδειξη λατρευτικού πλαισίου}\label{sec:m3c}

Στη M3c ο ακροατής ακούει ένα μουσικό απόσπασμα και τη συνέχειά του και αξιολογεί πόσο απρόσμενη του φαίνεται η συνέχεια. Η $H_{\mathrm{pred}}$ εξετάζει αν το πληροφοριακό περιεχόμενο της οικογένειας στην οποία ανήκει η συνέχεια εξηγεί πρόσθετη διακύμανση αυτής της κρίσης, αφού ληφθεί υπόψη το πληροφοριακό περιεχόμενο των επιμέρους μελωδικών γεγονότων. Η M3c πραγματοποιείται στην πρώτη ανθρώπινη συνεδρία. Το μοντέλο πρόβλεψης από μερική αντιστοίχιση (Prediction by Partial Matching, PPM) οικογενειών υπολογίζει $I_F$ μόνο από το προηγούμενο μουσικό συμφραζόμενο και η αλυσίδα σε επίπεδο μελωδικού γεγονότος παράγει το $I_{\mathrm{event}}$ για την ίδια συνέχεια. Το μοντέλο διαδοχής των οικογενειών συγκρίνεται σε έργα εκτός εκπαίδευσης με $P(F_t\mid M,L)$, ώστε η προβλεπτική συμβολή της προηγούμενης ακολουθίας οικογενειών να μετράται χωριστά από τις κατανομές που οφείλονται στην τροπική κατάσταση και στο δομικό, κειμενικό και προσωδιακό περιβάλλον.

Το PPM εκτιμά την πιθανότητα της επόμενης οικογένειας χρησιμοποιώντας όσο μεγαλύτερο προηγούμενο συμφραζόμενο υποστηρίζουν τα διαθέσιμα δεδομένα και επιστρέφει σε μικρότερο συμφραζόμενο όταν η μεγαλύτερη ακολουθία είναι πολύ σπάνια. Από αυτή την πιθανότητα προκύπτει το $I_F=-\log_2P$: μεγαλύτερο $I_F$ δηλώνει λιγότερο προβλέψιμη οικογένεια υπό το συγκεκριμένο μοντέλο. Το $I_{\mathrm{event}}$ υπολογίζει την αντίστοιχη πληροφοριακή ποσότητα στο επίπεδο των μεμονωμένων μελωδικών γεγονότων.

Το PPM ορίζει προβλεψιμότητα ως προς το σώμα στο οποίο εκπαιδεύεται. Στη M3c το σώμα χρησιμοποιείται ως προκαθορισμένη προσεγγιστική αναφορά για τη μουσική έκθεση των σημερινών ακροατών. Στη M5 ο ίδιος μηχανισμός επαναχρησιμοποιείται με διαφορετικό, ρητά δηλωμένο σώμα $D_{\Omega}$ και παράγει $I_{\Omega}$ για τον χαρακτηρισμό έργου ως προς το συγκεκριμένο περιβάλλον αναφοράς. Η αλλαγή σώματος μεταβάλλει την κατανομή αναφοράς, ενώ η αρχιτεκτονική και οι κανόνες επιστροφής σε συντομότερο συμφραζόμενο παραμένουν κοινοί. Με αυτόν τον τρόπο η ανθρώπινη προσδοκία της M3c και η ιστορικά προσδιορισμένη μουσική προβλεψιμότητα της M5 παραμένουν διαφορετικά εκτιμώμενα μεγέθη με κοινό υπολογιστικό μηχανισμό.

Η προκαθορισμένη ανάλυση της επανάληψης και προβλεψιμότητας οικογενειών στα τμήματα πλήρων Χερουβικών είναι δευτερεύουσα μελέτη του συμβολικού σώματος (§~\ref{app:human-functional}); παραμένει χωριστή από την ανθρώπινη συνθήκη $R$.


Το πληροφοριακό περιεχόμενο σε επίπεδο οικογένειας δεν ταυτίζεται με τη μουσικολογική έννοια της έντασης--λύσης που συζητά ο Δεσπότης (2007a)· ενδεχόμενη σχέση μεταξύ των δύο παραμένει ερμηνευτικό ερώτημα πέρα από τα προκαθορισμένα μεγέθη της M3c.

Κάθε δοκιμή παρουσιάζει μία συνέχεια. Η τράπεζα ερεθισμάτων εξισορροπεί τους δύο αντιδιαγώνιους συνδυασμούς του $I_F\times I_{\mathrm{event}}$ ---υψηλό $I_F$ με χαμηλό $I_{\mathrm{event}}$, ή το αντίστροφο--- και περιλαμβάνει τυφλό έλεγχο μουσικής φυσικότητας για κάθε κατασκευασμένη συνέχεια. Τα $I_F$ και $I_{\mathrm{event}}$ τυποποιούνται πριν από τη σύγκριση των συντελεστών αλληλεπίδρασης.

Η συνθήκη $R$ χρησιμοποιεί μη μουσική ένδειξη λατρευτικού συμφραζομένου και ουδέτερη ένδειξη αντίστοιχης διάρκειας. Κάθε μουσική συνέχεια παρουσιάζεται μία φορά ανά συμμετέχοντα. Κάθε συμμετέχων ακούει διαφορετικές συνέχειες και στις δύο τιμές του $R$, ενώ κάθε συγκεκριμένη συνέχεια παρουσιάζεται και στις δύο συνθήκες μεταξύ συμμετεχόντων μέσω σχήματος λατινικού τετραγώνου (Latin square) και αντισταθμισμένων λιστών. Το κύριο δευτερεύον μέγεθος $\Delta_{\mathrm{ctx}}$ συγκρίνει πώς η ένδειξη μεταβάλλει τη στάθμιση του πληροφοριακού περιεχομένου σε επίπεδο οικογένειας και γεγονότος. Η συνεχής εμπειρία $E_{\mathrm{Byz}}$ επιτρέπει να εξεταστεί πώς η εξοικείωση με τη συγκεκριμένη μουσική γλώσσα σχετίζεται με αυτή τη στάθμιση.

Η ιεραρχική περιγραφή του Mavromatis (2019, επίσημη περίληψη συνεδρίου, σσ.~56--57) συνδέει νότες, μελωδικές φόρμουλες, \emph{phrase prototypes} και μεγαλύτερες μορφικές μονάδες. Η έρευνα των Rabinovitch--Carter-Enyi (2024) εξετάζει αντίστοιχα τη διαδοχή μουσικών σχημάτων και τις μορφικές τους διαδρομές. Τα προηγούμενα αυτά αιτιολογούν μουσικοθεωρητικά την πιθανότητα προσδοκίας σε περισσότερες από μία κλίμακες. Η M3c εξετάζει εμπειρικά αν το $I_F$ εξηγεί τις ανθρώπινες κρίσεις πέρα από το $I_{\mathrm{event}}$, με τους προβλεπόμενους ελέγχους συμφραζομένου και ανεξαρτησίας.

\subsection{M3b: κατηγορική αναπαράσταση της οικογένειας θέσεων}\label{sec:m3b}

Η M3b χρησιμοποιεί τριάδες επιλογής του διαφορετικού. Ο ακροατής ακούει τρεις αποδόσεις και επιλέγει εκείνη που ταιριάζει λιγότερο με τις άλλες δύο. Οι \textbf{μεικτές τριάδες} περιλαμβάνουν σχέσεις μέσα στην ίδια οικογένεια και μεταξύ διαφορετικών οικογενειών και επιλέγονται από κοινή περιοχή συμβολικών και ακουστικών αποστάσεων, στην οποία εκπροσωπούνται και οι δύο τύποι σχέσεων σύμφωνα με τα προκαθορισμένα κριτήρια αντιστοίχισης. Οι \textbf{ενδοοικογενειακές τριάδες} περιλαμβάνουν τρεις αποδόσεις της ίδιας οικογένειας και εξετάζουν τη σχετική οργάνωση των διαφορών μέσα στην οικογένεια.

\begin{center}\small
\begin{tabularx}{0.9\textwidth}{@{}p{0.27\textwidth}YY@{}}
\toprule
& \textbf{Μικρή επιφανειακή απόσταση} & \textbf{Μεγάλη επιφανειακή απόσταση}\tabularnewline
\midrule
\textbf{Ίδια οικογένεια} & εύκολα θετικά & \textbf{δύσκολα θετικά}\tabularnewline
\midrule
\textbf{Διαφορετική οικογένεια} & \textbf{δύσκολα αρνητικά} & εύκολα αρνητικά\tabularnewline
\bottomrule
\end{tabularx}
\end{center}

Η υπόθεση $H_{\mathrm{cat}}$ ελέγχεται με τις μεικτές τριάδες και μοντέλο τριαδικής επιλογής και διασταυρωμένη δομή τυχαίων επιδράσεων για ακροατή, ερέθισμα, τριάδα και οικογένεια. Η ανάλυση εκτιμά τη γεωμετρία των ανθρώπινων αναπαραστάσεων από τις τριάδες, χωρίς ετικέτες οικογένειας. Η αξιοπιστία του ανθρώπινου χώρου, το ανώτατο όριο αξιοπιστίας και η πρόβλεψη τριάδων σε υλικό ανάπτυξης που έχει κρατηθεί εκτός της προσαρμογής του ανθρώπινου χώρου προηγούνται του CCGP και της ευθυγράμμισης μοντέλου--ανθρώπου.

Το $\mathrm{CCGP}_F$ εξετάζει αν η διάκριση των οικογενειών που μαθαίνεται σε ορισμένες συνθήκες γενικεύει σε νέες συνθήκες. Η ανάλυση μέσα σε κάθε οικογένεια κατασκευάζει RDM ανά οικογένεια και συγκρίνει την υπολειμματική συμφωνία του ανθρώπινου χώρου με EISE και B2, αφού ληφθούν υπόψη οι προκαθορισμένες συμβολικές και ακουστικές μήτρες αποστάσεων μεταξύ των ίδιων ερεθισμάτων (representational dissimilarity matrices, RDM). Η συνάθροιση των αποτελεσμάτων ανά οικογένεια αποτρέπει την ψευδοεπανάληψη που θα προέκυπτε αν οι πολλές ζευγαρωτές παρατηρήσεις αντιμετωπίζονταν ως ανεξάρτητες.

\subsection{Ενισχυτικές ανθρώπινες αναλύσεις}\label{sec:human-secondary}

Το κύριο δευτερεύον $\Delta_{\mathrm{ctx}}$ εξετάζει τη σχετική μεταβολή των δύο συντελεστών υπό την ένδειξη λατρευτικού πλαισίου. Η M3b εξετάζει επιπλέον το $\mathrm{CCGP}_F$, την υπολειμματική συμφωνία των ανθρώπινων αναπαραστάσεων με B2/EISE μέσα στην ίδια οικογένεια και τη σχέση αυτών των αποτελεσμάτων με τη συνεχή εμπειρία στην ψαλτική. Οι σύντομες αξιολογήσεις βιωμένης απόκρισης έχουν διερευνητικό καθεστώς και αναλύονται με τις προϋποθέσεις της §~\ref{app:human-affect}.

\subsection{Ευρύτερη εκτελεστική ανάλυση M4}\label{sec:m4-broader}

Πέρα από το ανεξάρτητο εκτελεστικό κριτήριο αξιολόγησης που απαιτεί η M2, η M4 εξετάζει την εκτελεστική παραλλαγή των θέσεων στο επίπεδο που επιτρέπει η κάλυψη: χρονισμό, τονικό ύψος, επιλεγμένα ακουστικά γνωρίσματα και το χωριστό εκτελεστικό προφίλ $\Gamma(z)$, δηλαδή τη μετρημένη μεταβλητότητα των εκτελέσεων στην εσωτερική περιοχή $z$ της θέσης (§~\ref{app:stats-performance}). Στον χρονικό άξονα διακρίνονται ρητά η συνολική αγωγή της ηχογράφησης, ο κανονικοποιημένος χρονισμός μιας περιοχής, η καταληκτική επιμήκυνση, η τοπική μεταβολή αγωγής και η πυκνότητα μελωδικών γεγονότων. Η κύρια ανάλυση χρησιμοποιεί τον όρο αγωγής και τη διαδικασία εκτίμησης της §~\ref{app:stats-performance}. Η προσαρμογή ως προς τον τύπο τμήματος εξετάζεται χωριστά ως προκαθορισμένη ανάλυση ευαισθησίας και συγκρίνεται με την κύρια εκτίμηση. Η μετρική έμφαση χρησιμοποιείται μόνο όπου υπάρχει ανεξάρτητα επικυρωμένη ερμηνεία. Η ιστορική ανάλυση $A(z)$ και η εκτελεστική $\Gamma(z)$ διατηρούν διαφορετική μονάδα παρατήρησης και δηλώνουν ρητά τη μερική εξάρτηση όταν μοιράζονται δεδομένα. Οι λεπτομερείς μετρήσεις, έλεγχοι και αποφάσεις GREEN/AMBER/RED βρίσκονται στα Παραρτήματα~\ref{app:stats-performance}, \ref{app:corpora-gate} και \ref{app:impl}.

Εφόσον το ΠΕ0 επιβεβαιώσει πρόσβαση, συγχρονισμό και ποιότητα του ηλεκτρογλωττογραφήματος, η M4 μπορεί να εξετάσει διερευνητικά δείκτες φωνητικής επαφής· το τεχνικό πρωτόκολλο βρίσκεται στην §~\ref{app:impl-vars}.

\subsection{M3ν: νευρική επέκταση}\label{sec:m3nu}

Η M3ν χρησιμοποιεί τα ερεθίσματα και τις μεταβλητές πρόβλεψης της M3c για να ελέγξει με EEG/MEG αν το $I_F$ προσθέτει νευρική εξηγητική πληροφορία πέρα από το $I_{\mathrm{event}}$ και τα ακουστικά χαρακτηριστικά. Το $R\times I_F$ εξετάζεται χωριστά. Η επέκταση ενεργοποιείται με γραπτή επιβεβαίωση υποδομής, εμπειρογνωμοσύνης ανάλυσης και ανεξάρτητης δεοντολογικής έγκρισης. Η μη ενεργοποίησή της διατηρεί τις υποθέσεις και τα συμπεράσματα του βασικού ερευνητικού προγράμματος. Η τεχνική προδιαγραφή παραμένει στην §~\ref{app:human-neural}.

\subsection{Σύνθεση ανεξάρτητων μεθόδων}

Η τελική σύνθεση διατηρεί διακριτές τις ερωτήσεις κάθε πηγής τεκμηρίωσης. Η μουσικολογική θεμελίωση απαντά πώς ορίζεται και τεκμηριώνεται η οικογένεια. Η M3b απαντά πώς οργανώνεται η συγγένεια στην ανθρώπινη ακρόαση. Η M3c απαντά πώς η οικογένεια συμμετέχει στην πρόβλεψη. Η εκτελεστική ανάλυση απαντά πώς πραγματώνεται η παραλλαγή. Η σύγκλιση αυτών των αποτελεσμάτων προσφέρει ισχυρή τεκμηρίωση επειδή τα επιμέρους εκτιμώμενα μεγέθη έχουν διαφορετικό αντικείμενο και διαφορετικούς μηχανισμούς σφάλματος.

\section{Δεδομένα και δυνατότητα υλοποίησης}

\subsection{Συμβολικά σώματα, Άγιος Αντώνιος και ΔΑΜΑΣΚΗΝΟΣ}\label{sec:training-pool}

Η παπαδική συλλογή \emph{Papadic Formulae} της Ιεράς Μονής Αγίου Αντωνίου Αριζόνας λειτουργεί ως βασική καταλογογραφική βάση και υποδομή τεκμηρίωσης των πηγών υποψηφίων τύπων, εμφανίσεων και προέλευσης και ως διαδρομή προς τα πλήρη έργα και τις πηγές. Οι 758 συνθέσεις δηλώνουν την έκταση της συλλογής προέλευσης και δεν ταυτίζονται με τον αριθμό των επιλέξιμων παραδειγμάτων εκπαίδευσης. Η κάλυψη καταγράφεται χωριστά για καταληκτικές, εναρκτήριες, συνδετικές και λοιπές λειτουργίες, ώστε η ισχυρή καταληκτική εκπροσώπηση του καταλόγου να μην γενικεύεται σε ολόκληρο το λεξιλόγιο θέσεων. Οι 10.000+ δηλωμένες φόρμουλες της χωριστής ειρμολογικής/στιχηραρικής\footnote{Ειρμολογικό και στιχηραρικό δηλώνουν διαφορετικά ιστορικά μελοποιητικά και ρεπερτοριακά πεδία από το παπαδικό.} συλλογής καταγράφονται ως διαφορετική συλλογή και δεν συγχέονται με την παπαδική.

Το ΠΕ0 επαληθεύει τη σύνδεση κάθε επιλέξιμης εγγραφής με συγκεκριμένο έργο, πηγή και ομάδα συγγενών πηγών, διαχωρίζει οικογενειακή από εξηγητική εποπτεία και καταγράφει τη διαθέσιμη κάλυψη ανά οικογένεια και ήχο. Για το ανεξάρτητο εκτελεστικό κριτήριο χρειάζονται επίσης ηχογραφήσεις κοινού ρεπερτορίου από πολλαπλούς εκτελεστές. Τέτοιο υλικό, σε τυποποιημένες συνθήκες καταγραφής, περιγράφεται για το σώμα ΔΑΜΑΣΚΗΝΟΣ από τους Dimitrios Delviniotis και Georgios Kouroupetroglou (2011). Η σχετική μεθοδολογία επισημείωσης των εκκλησιαστικών ακουστικών σωμάτων έχει περιγραφεί από τους Chryssochoidis, Delviniotis και Kouroupetroglou (2007) με χρήση του λογισμικού ανάλυσης ήχου PRAAT και σύνδεση συλλαβικών ορίων, σημειογραφίας και μουσικών/διαστηματικών μεταδεδομένων. Το ΠΕ0 επαληθεύει αν τα αντίστοιχα αρχεία επισημείωσης και το σχήμα κωδικοποίησης είναι πράγματι διαθέσιμα στην παρούσα πρόσβαση και σε ποια έκταση μπορούν να μειώσουν το κόστος της ευθυγράμμισης χωρίς να θεωρούνται αυτομάτως πρότυπο για την επαλήθευση των εκτελεστικών μετρήσεων. Η απογραφή πρόσβασης, δικαιωμάτων χρήσης, ιστορικής--εκτελεστικής επικάλυψης και ποιότητας ευθυγράμμισης προσδιορίζει το τεκμηριωμένο πεδίο της εκτελεστικής ανάλυσης.

\subsection{ΠΕ0: απογραφή, επάρκεια και οριστικοποίηση κλίμακας}\label{sec:pilot-decision}

Η πρώτη φάση δημιουργεί μητρώο ανά υποψήφια/επικυρωμένη οικογένεια×ήχο με έργα, εμφανίσεις, ατομικά γεγονότα, ανεξάρτητους μάρτυρες και ομάδες συγγενών πηγών, επαληθευμένες ιστορικές ακμές, εφικτές αντιστοιχισμένες ζευγοποιήσεις B2 και, όπου απαιτείται, εκτελεστική κάλυψη. Ανά εμφάνιση καταγράφονται ο λειτουργικός ρόλος ---καταληκτικός, εναρκτήριος, συνδετικός ή άλλος---, το λεκτικό καθεστώς ---λεξικό, συλλαβές χωρίς λεξικό περιεχόμενο, μικτό ή άγνωστο---, αν τα πεδία $L$ εφαρμόζονται στο συγκεκριμένο χωρίο, η πληρότητα εξαγωγής από την πηγή και ο ρόλος κάθε προσώπου στη συγκεκριμένη μαρτυρία. Για σταθερά λειτουργικά κείμενα καταγράφονται επιπλέον αναγνωριστικό κειμένου, κειμενικό χωρίο και βαθμός μελισματικής ανάπτυξης· στα Κοινωνικά η διασταύρωση γίνεται μόνο μεταξύ μελοποιήσεων του ίδιου κοινωνικού στίχου. Για τα κρατήματα και άλλα εκτενή μη λεξικά τμήματα η απογραφή διαχωρίζει την παρουσία στη συγκεκριμένη πηγή, την πληρότητα ή έκταση του παραδιδόμενου τμήματος, τεκμηριωμένη σύντμηση ή αντικατάσταση και απόδοση σε μελοποιό ή μεταγενέστερη παρέμβαση. Η απουσία στην πηγή δεν ταυτίζεται με τεκμηριωμένη παράλειψη, ενώ φθορά ή ατελής διατήρηση του μάρτυρα καταγράφεται χωριστά από τη σύντμηση. Με βάση αυτή την απογραφή οριστικοποιείται η υψηλότερη μουσικά και μετρητικά υποστηριζόμενη κλίμακα ανάλυσης πριν από την τελική αξιολόγηση.

\paragraph{Στοχευμένες περιπτώσεις παραλλαγής.} Οι κάρτες τύπου στο ΠΕ0 αξιοποιούν, όπου υπάρχουν κατάλληλες πηγές, τεκμηριωμένες περιπτώσεις ανωμοτονίας, ανισοσυλλαβίας, διαφορετικής σημειογραφικής απόδοσης και ένθεσης μικρότερου σχηματισμού σε μεγαλύτερη μουσική μονάδα (Αρβανίτης, 2010, σσ.~136--142). Οι παρατηρήσεις καταγράφονται στα υπάρχοντα \texttt{diagnostic\_relations}, \texttt{diagnostic\_scope}, \texttt{realization\_rules} και \texttt{analysis\_scale}, με σαφή διάκριση παραδείγματος ανάπτυξης από ανεξάρτητο δείγμα αξιοπιστίας. Ένα κοινό υποτμήμα μπορεί να αιτιολογεί υψηλή μελωδική ομοιότητα χωρίς να καθορίζει κοινή ταυτότητα της σύνθετης θέσης (πρβλ. Holzapfel, 2010, κεφ.~7). Η χρονική οργάνωση καταγράφεται με ρόλο \texttt{required}, \texttt{supportive} ή ως στοιχείο πραγμάτωσης μόνο όταν αυτό υποστηρίζεται για τον συγκεκριμένο τύπο. Δεν δημιουργείται νέο υποχρεωτικό υποπείραμα.

\paragraph{Ελεγχόμενη ετερογένεια ταυτότητας και τροπικού πλαισίου.} Οι παρατηρήσεις των Αντωνόπουλου (2013/2014, εικ.~5.36, σσ.~284--285), Troelsgård (2011, §§32, 57, 67, σσ.~47, 61, 69) και Αλεξάνδρου--Troelsgård (2020, σσ.~178--180, πίν.~9.1--9.4) χρησιμοποιούνται ως ενδεικτικό μουσικολογικό υλικό \texttt{definition} για ενδεχόμενες τροπικές μεταφορές, διαφορετικά μεγάλα σημάδια και αναγνωρίσιμες μελικές σχέσεις. Ο χαρακτηρισμός \emph{ἀπὸ παραλλαγῆς} ή \emph{ἀπὸ μέλους} τεκμηριώνεται χωριστά ανά χειρογραφικό μάρτυρα, και αμφισβητούμενες ενδιάμεσες μαρτυρίες παραμένουν \texttt{unresolved}. Μία οικογένεια μπορεί στο υλικό ανάπτυξης να αποκτήσει περισσότερες προδηλωμένες επαρκείς διαγνωστικές διαμορφώσεις, μόνο με μουσικολογική τεκμηρίωση και νέο ανεξάρτητο έλεγχο. Η εναλλαγή διαμορφώσεων δεν παράγεται από τον αλγόριθμο εκ των υστέρων στο τελικό δείγμα.

\paragraph{Πιλοτική συγκρότηση οικογενειακών υπογραφών.}
Το ΠΕ0 ξεκινά με μικρό σύνολο πραγματικών υποψήφιων τύπων. Η παρακλητική, το κύλισμα και το πίασμα χρησιμοποιούνται ως πρώτες περιπτώσεις ανάπτυξης, με βάση τους μουσικούς πίνακες του Gheorghiță (2010, σσ.~228--245). Δεν θεωρούνται εκ των προτέρων επικυρωμένες οικογένειες. Το αντικένωμα, το λύγισμα και το έτερο παρακάλεσμα της σ.~788 του Καραγκούνη χρησιμοποιούνται ως υποψήφιες ιστορικές δύσκολες περιπτώσεις. Μετατρέπονται σε M1 hard negatives μόνο αν η αυτοψία των αντίστοιχων εξηγήσεων της Νέας Μεθόδου επιβεβαιώσει ότι η απαιτούμενη δυσκολία ομοιότητας παραμένει. Το τρομικό εντάσσεται μόνο μετά από χωριστή άμεση επαλήθευση των σχετικών σελίδων και πηγών. Επιλέγεται επίσης μία πλήρως ελέγξιμη εξηγητική περίπτωση με σαφές κοινό χωρίο και τεκμηριωμένη γραμμή πηγών.

Για κάθε υποψήφιο τύπο κατασκευάζεται αρχική κάρτα. Δηλώνεται ένα αντιπροσωπευτικό τεκμηριωμένο παράδειγμα ανάπτυξης, χωρίς να λειτουργεί ως gold standard. Δηλώνονται επίσης η αναλυτική κλίμακα, οι διαγνωστικές μουσικές σχέσεις και ο ρόλος τους, οι επιτρεπτές μεταβολές, οι συνθήκες πλαισίου όπου τεκμηριώνονται, τα θετικά παραδείγματα και τα δύσκολα αρνητικά. Δηλώνονται επίσης οι περιπτώσεις στις οποίες η κρίση πρέπει να παραμένει \texttt{unresolved}. Η αρχική κάρτα αναθεωρείται μόνο στο υλικό ανάπτυξης.

Δύο ανεξάρτητοι ειδικοί εφαρμόζουν την παγωμένη υπογραφή σε νέο δείγμα. Πρώτα κρίνουν τα όρια. Έπειτα ελέγχουν κάθε ενεργό διαγνωστικό γνώρισμα ως \texttt{satisfied}, \texttt{violated} ή \texttt{not\_assessable}, με ορατό τον προδηλωμένο ρόλο και το πεδίο εφαρμογής του. Τέλος αποδίδουν $J(x,T;r)$. Οι διαφωνίες ταξινομούνται ως \texttt{boundary}, \texttt{scale}, \texttt{modal}, \texttt{diagnostic\_feature}, \texttt{source\_insufficient}, \texttt{historical\_naming} ή \texttt{other}. Η αρχική ανεξάρτητη κρίση διατηρείται στο μητρώο. Η συμφωνία αναφέρεται συνολικά, ανά οικογένεια και ανά διαγνωστικό γνώρισμα, ώστε να διακρίνεται η αξιοπιστία της τελικής κατηγορίας από την αξιοπιστία των επιμέρους κριτηρίων. Τα ακριβή κατώφλια και η μετρική συμφωνίας οριστικοποιούνται στο ΠΕ0 πριν από την κύρια χρήση των ετικετών.

Η ακολουθία είναι δεσμευτική:
\[
\begin{aligned}
\text{candidate signature} &\rightarrow \text{development refinement}
\rightarrow \text{signature freeze}\\
&\rightarrow \text{independent reliability test}
\rightarrow \text{pair generation}.
\end{aligned}
\]
Η αποτυχία αξιοπιστίας ενός τύπου δεν διορθώνεται με εκ των υστέρων χρήση αποτελεσμάτων M1 ή M2. Ο τύπος παραμένει \texttt{candidate\_only}, αναθεωρείται σε νέα δεδομένα ανάπτυξης ή αποσύρεται από το κύριο λεξιλόγιο.

\paragraph{Έλεγχος επάρκειας της συμβολικής αναπαράστασης.}
Πριν από τη M1 ελέγχεται αν τα παρατηρήσιμα γνωρίσματα που απαιτούν οι υπογραφές μπορούν να αναπαρασταθούν στα υπάρχοντα στρώματα $N/T/M/L$ ή στα μεταδεδομένα. Ελέγχονται ιδίως βαθμίδα, φθορά, γένος, σχετική μελωδική κίνηση, χρονική πληροφορία, συλλαβική σύνδεση και κειμενική ευθυγράμμιση. Το αποτέλεσμα δεν δημιουργεί νέα αρχιτεκτονική ή νέο $U$. Πληροφορία που υπάρχει μόνο σε παλαιότερη γραφή δεν συμπληρώνεται τεχνητά στη Νέα Μέθοδο. Καταγράφεται ως \texttt{observed}, \texttt{inferred} ή \texttt{unknown}, με σαφή προέλευση.

Για τον διαχρονικό έλεγχο που ενεργοποιείται μετά την ικανοποίηση του αντίστοιχου κριτηρίου ενεργοποίησης καταγράφονται επιπλέον διαθέσιμοι ανεξάρτητοι ιστορικοί εκτελεστικοί μάρτυρες, κοινό ρεπερτόριο, ποιότητα της διαθέσιμης μεταφοράς και αξιοπιστία χρονικής ευθυγράμμισης. Στο ίδιο στάδιο ελέγχονται η πραγματική διαθεσιμότητα των δημοσιευμένων επισημειώσεων του ΔΑΜΑΣΚΗΝΟΣ, ένας προσαρμογέας από δομημένη μορφή Νέας Μεθόδου προς το ενιαίο επίπεδο αναφοράς και, σε μικρό προκαθορισμένο δείγμα έντυπης έκδοσης, η δυνατότητα εφαρμογής OCR ως προαιρετικής οδού εισαγωγής. Οι δύο τελευταίες διαδρομές συμπληρώνουν την απόκτηση του βασικού σώματος, με ενεργοποίηση βάσει της μετρημένης πιστότητας μετατροπής και των όρων χρήσης.

Στο ίδιο προπαρασκευαστικό στάδιο ολοκληρώνεται ο πιλοτικός έλεγχος της §~\ref{app:corpora-annotation-pilot}. Η δοκιμή χωρίζει ρητά την ανάπτυξη της υπογραφής από τον νέο έλεγχο αξιοπιστίας. Οι εμφανίσεις που χρησιμοποιούνται για να αλλάξει ένας κανόνας δεν επανεμφανίζονται στο δείγμα αξιοπιστίας της ίδιας έκδοσης. Οι πιλοτικές επισημειώσεις δεν εισέρχονται σε επιβεβαιωτικές αναλύσεις. Παραδοτέα είναι το παγωμένο εγχειρίδιο, οι εκδόσεις των υπογραφών, οι μετρήσεις συμφωνίας και ο κατάλογος τύπων που επιτρέπεται να τροφοδοτήσουν τις M1--M2.

Στο πλήρως ελεγμένο υποσύνολο ανάπτυξης απογράφονται χωριστά η αβεβαιότητα των ορίων, οι τεκμηριωμένες επικαλύψεις εμφανίσεων και οι συνδετικές ή ανεπίλυτες περιοχές. Πριν από τη M1 οριστικοποιείται η μορφή της κύριας δομημένης εξόδου. Η μη επικαλυπτόμενη έξοδος διατηρείται όταν η απώλεια τεκμηριωμένων εμφανίσεων παραμένει μέσα στα προκαθορισμένα πρακτικά περιθώρια, δεν αποκλείει συστηματικά οικογένειες ή λειτουργίες που είναι αναγκαίες για τους κύριους ισχυρισμούς και διατηρεί το κοινό πεδίο που απαιτεί η M2. Τα περιθώρια ορίζονται πριν από οποιαδήποτε σύγκριση μοντέλων από την ακρίβεια της επισημείωσης, την κάλυψη των κρίσιμων υποομάδων που έχουν προκαθοριστεί ανά οικογένεια, ήχο, λειτουργικό ρόλο και ιστορική προέλευση και τις ήδη προβλεπόμενες προσομοιώσεις ισχύος και ακρίβειας. Αν ο περιορισμός μη επικάλυψης προκαλεί ουσιώδη απώλεια του πληθυσμού αναφοράς, ενεργοποιείται περιορισμένος έλεγχος εφικτότητας για ανίχνευση υποψήφιων τμημάτων που επιτρέπει τοπικές επικαλύψεις, με την ίδια αρχιτεκτονική κωδικοποιητή, τις ίδιες οικογενειακές κεφαλές και την ίδια λογική B2--EISE. Ο έλεγχος περιλαμβάνει μερική επισημείωση, αναγνώριση οικογενειών εκτός λεξιλογίου και υπολογιστικό κόστος. Αν επιλεγεί αυτή η μορφή εξόδου, τα μεταγενέστερα μοντέλα και οι συνθήκες σύγκρισης προσαρμόζονται σε αυτήν και περιορίζονται πρώτα οι προαιρετικές συγκρίσεις ώστε να διατηρηθεί το χρονοδιάγραμμα. Αν δεν είναι εφικτή στην προβλεπόμενη κλίμακα, δηλώνεται ρητά το στενότερο πεδίο στο οποίο μπορεί να στηριχθεί η αυτόματη ανάλυση. Η απόφαση λαμβάνεται πριν από τη M1 και χωρίς πρόσβαση στα τελικά αποτελέσματα M1/M2 ή στο ανεξάρτητο εκτελεστικό κριτήριο.

Η επιλογή επιπέδου βασίζεται στην κάλυψη, στην ανεξαρτησία πηγών, στην ποιότητα αντιστοίχισης, στην αξιοπιστία των μετρήσεων, στον χρόνο κωδικοποίησης και στις προκαθορισμένες προσομοιώσεις --- όχι στα παρατηρημένα \(S,\rho,\Delta\). Στις καταστάσεις GREEN και AMBER η σύγκριση της EISE με το ανεξάρτητο εκτελεστικό κριτήριο έχει επιβεβαιωτικό καθεστώς· στη RED αναφέρεται διερευνητικά. Τα ακριβή κατώφλια, οι προσομοιώσεις ακρίβειας και οι κανόνες μετάβασης σε αδρότερο επίπεδο δίνονται στις §§~\ref{app:corpora-gate} και \ref{app:stats-eise-transfer}.

Το ΠΕ0 χρονομετρά αντιπροσωπευτικά έργα ανά τύπο πηγής, υπολογίζει ώρες κωδικοποίησης και ελέγχου ανά πλήρες έργο και καταγράφει ποιο υλικό διατίθεται ήδη σε μηχανικά αναγνώσιμη μορφή. Ζητείται από τη Μονή Αγίου Αντωνίου η μορφή των πρωτογενών αρχείων και οι όροι πρόσβασης· οι διαδρομές Neanes και OCR αξιολογούνται με τους ήδη ορισμένους ελέγχους πιστότητας. Με βάση διαθέσιμο χρόνο και προβλέψεις ακρίβειας, η κλίμακα οριστικοποιείται κατά σειρά: σύνολα ανάπτυξης και ανεξάρτητης αξιολόγησης M1--M2, ιστορικό σύνολο υψηλής δυσκολίας και σώμα M5. Όταν το κόστος υπερβαίνει τον προϋπολογισμό, περιορίζονται πρώτα οι πρόσθετες ιστορικές ομάδες αναφοράς, το πλήρες $D_{\mathrm{seq}}/D_{\mathrm{func}}$ και οι ευρύτερες εφαρμογές της M5. Διατηρείται κατά προτεραιότητα η \textbf{M5-min}, δηλαδή $D_{\mathrm{fam}}$ και $D_{\mathrm{within}}$ σε διαθέσιμο συγκρίσιμο σώμα της Νέας Μεθόδου· αν $K<K_{\min}$, η M5-min παραμένει περιγραφική χωρίς κατάταξη σε επίπεδο μελοποιού. Έπειτα ορίζεται η μικρότερη κλίμακα ελέγχου που ικανοποιεί τις προσομοιώσεις και τους διαχωρισμούς κατά έργο και πηγή. Αν ένα αναγκαίο σύνολο υπολείπεται ακόμη αυτών των κριτηρίων, δηλώνεται ότι το αντίστοιχο μέγεθος δεν μπορεί να εκτιμηθεί αξιόπιστα στην προβλεπόμενη κλίμακα.\par

Η δυνατότητα εκτέλεσης των συμβολικών και τροπικών αναλύσεων παραμένει ανεξάρτητη από την κάλυψη πολλαπλών εκτελέσεων. Η κάλυψη από πολλαπλές εκτελέσεις στηρίζει την εκτελεστική ανάλυση και τη σύγκριση ιστορικού και ανεξάρτητου εκτελεστικού προφίλ, ενώ οι υπόλοιπες αναλύσεις αξιοποιούν το συμβολικό σώμα με τα δικά τους δεδομένα και εκτιμώμενα μεγέθη.

\subsection{Κριτήρια επάρκειας για αξιόπιστη εκτίμηση}\label{sec:estimability}

Το ΠΕ0 συνδέει κάθε συμπερασματικό αποτέλεσμα με τα δεδομένα που απαιτούνται για την εκτίμησή του. Στη M1 η αναπαράσταση $R^\star$ προκύπτει από το $F_{1,\mathrm{seg+fam}}$ πλήρους αναγνώρισης στην A υπό το οικογενειακό ελάχιστο $g_{\min}$, το οποίο οριστικοποιείται πριν από την εκπαίδευση των δώδεκα συνθηκών με τη διαδικασία της §~\ref{app:stats-selection}· το σύνολο $\mathcal R^\star$ αναφέρει την αβεβαιότητα της επιλογής ανεξάρτητα από τη λειτουργική απόφαση. Η στατιστική αξιολόγηση του AUROC στο ιστορικό σύνολο υψηλής δυσκολίας προϋποθέτει επαρκείς ανεξάρτητες περιπτώσεις· διαφορετικά αναφέρεται μόνο η προκαθορισμένη περιγραφική αντιπαραβολή.

Για τις αντιθέσεις της M2 χρειάζεται κοινό και αξιόπιστο πεδίο ιστορικής--εκτελεστικής αντιστοίχισης. Η κατάσταση GREEN επιτρέπει επιβεβαιωτική λεπτομερή κλίμακα και η AMBER την προκαθορισμένη αδρότερη κλίμακα. Η RED οδηγεί σε διερευνητική αναφορά· όπου το ανεξάρτητο κριτήριο αξιολόγησης δεν μπορεί να μετρηθεί, δηλώνεται ότι η αντίστοιχη αντίθεση δεν μπορεί να εκτιμηθεί αξιόπιστα. Η αξιολόγηση των συμβολικών και μουσικολογικών αποτελεσμάτων διατηρεί τις δικές της προϋποθέσεις και το δικό της πεδίο τεκμηρίωσης.

Η ανθρώπινη M3c απαιτεί επικυρωμένο λεξιλόγιο οικογενειών, επιτυχία του προκαθορισμένου ελέγχου του PPM, κάλυψη των αντιδιαγώνιων ερεθισμάτων και επαρκή μουσική φυσικότητα. Η M3b απαιτεί κοινό πεδίο τιμών αποστάσεων και αξιοπιστία της δοκιμασίας τριαδικής επιλογής. Η ανεπαρκής κάλυψη καταγράφεται για το αντίστοιχο εκτιμώμενο μέγεθος, σύμφωνα με το Παράρτημα~\ref{app:human-gates}. Η ακολουθιακή συνιστώσα της M5 απαιτεί επαρκείς τεκμηριωμένες διαδοχές οικογενειών και συγκρίσιμο σώμα αναφοράς. Η M5-min μπορεί να εφαρμοστεί με $D_{\mathrm{fam}}$ και $D_{\mathrm{within}}$ ακόμη όταν ο αριθμός ανεξάρτητων μελοποιών δεν επαρκεί για κατανομή αναφοράς σε επίπεδο μελοποιού. Η κατανομή αναφοράς και κάθε εμπειρική κατάταξη ενεργοποιούνται μόνο με $K\ge K_{\min}$, όπου το $K_{\min}$ οριστικοποιείται στο ΠΕ0 με προσομοίωση. Όπου η κάλυψη δεν στηρίζει τις συνθήκες ήχου, είδους ή τόπου, ενεργοποιείται προκαθορισμένο απλούστερο μοντέλο και το αποτέλεσμα αφορά το επίπεδο ανάλυσης στο οποίο τα διαθέσιμα δεδομένα παραμένουν αξιόπιστα συγκρίσιμα. Πρόσθετες ιστορικές ομάδες αναφοράς ενεργοποιούνται μόνο μετά τον έλεγχο χρονολογίας, εξάρτησης πηγών, κοινού εξηγητή και κοινής κάλυψης των ίδιων οικογενειών.

\subsection{Λίγα τεκμηριωμένα ζεύγη: στρατηγική εκτίμησης}\label{sec:small-data}

Ο αριθμός των χωρίων με παράλληλες εξηγήσεις και των στοιχισμένων ιστορικών πηγών αναμένεται μικρός. Η πρόταση οργανώνεται ώστε το μέγεθος αυτό να περιορίζει την ακρίβεια και όχι την εγκυρότητα των συμπερασμάτων. Η μονάδα ανεξαρτησίας είναι το έργο μαζί με την ομάδα συγγενών πηγών, και η ισχύς των ελέγχων υπολογίζεται σε επίπεδο οικογένειας με προσομοίωση πριν από την πρόσβαση στα αποτελέσματα· οι πύλες του ΠΕ0 προσαρμόζουν την κλίμακα. Κάθε μαθημένη αναπαράσταση συγκρίνεται με ισχυρούς απλούς συγκριτές, όπως το άμεσο συμβολικό προφίλ, οι κανόνες των υπογραφών και ο Transformer σε επίπεδο γεγονότος. Αν δεν τους ξεπερνά, το αποτέλεσμα αναφέρεται ως εύρημα. Όπου τα ζεύγη είναι λίγα, προτιμώνται μοντέλα με ιεραρχική συρρίκνωση, και η προεκπαίδευση γίνεται χωρίς ετικέτες.

Το ίδιο το μέγεθος του δείγματος γίνεται ερευνητικό ερώτημα. Οι κύριες συγκρίσεις της M2 και της §~\ref{sec:atl} επαναλαμβάνονται με αυξανόμενο ποσοστό των διαθέσιμων τεκμηριωμένων ζευγών, ώστε να φανεί πόση ιστορική εποπτεία χρειάζεται, πού αποδίδει περισσότερο και πότε παύει να προσφέρει. Παράλληλα, ένα διαγνωστικό σύνολο κρίνεται ανεξάρτητα από ειδικούς. Περιλαμβάνει μεγάλες επιφανειακές μεταβολές που διατηρούν την ταυτότητα και μικρές μεταβολές που αλλάζουν κρίσιμη σχέση. Το σπάνιο ιστορικό υλικό χρησιμοποιείται κατά προτεραιότητα για αξιολόγηση. Τα συνθετικά δεδομένα χρησιμεύουν μόνο στον έλεγχο των μεθόδων και της ισχύος. Η γενίκευση εξετάζεται πρώτα μέσα στο σώμα, με δύο διακριτούς ελέγχους: σε οικογένειες που αποκρύπτονται από την εκτίμηση και σε νέα έργα και ομάδες πηγών. Ένα δεύτερο ρεπερτόριο προστίθεται μόνο όταν επιτρέπει συγκρίσιμο ερώτημα.

\subsection{Πρόσβαση, διάθεση δεδομένων και εφικτότητα της έρευνας}\label{sec:data-rights}

Το σύνολο εκπαίδευσης, οι όροι χρονολογίας και προέλευσης, οι άδειες χρήσης και η δυνατότητα αναδιανομής οριστικοποιούνται πριν από την επιβεβαιωτική αξιολόγηση. Για υλικό που δεν μπορεί να αναδιανεμηθεί δημοσιοποιούνται το σχήμα δεδομένων, τα μεταδεδομένα, τα αναγνωριστικά προέλευσης, ο κώδικας και οι οριστικοποιημένες ρυθμίσεις στον βαθμό που επιτρέπεται.

Η τεχνική ετοιμότητα αφορά πλήρως καθορισμένη προδιαγραφή: ενιαίο συμβολικό επίπεδο αναφοράς με τεκμηριωμένη προέλευση, τμηματικό κωδικοποιητή με κοινό υπολογιστικό προϋπολογισμό, αλυσίδα ευθυγράμμισης παρτιτούρας--εκτέλεσης και καταγεγραμμένη διαδικασία αναπαραγωγής των διαχωρισμών, των αρχικοποιήσεων και των ρυθμίσεων. Η πρώτη λειτουργική υλοποίηση του EISE προβλέπεται στους μήνες 4--12 και θα αποτιμήσει στην πράξη μνήμη, χρόνο και διατήρηση τοπικής πληροφορίας. Ο κωδικοποιητής σχεδιάζεται για έναν επιταχυντή γραφικών (GPU) μεσαίας κατηγορίας, με κοινό προϋπολογισμό εκπαίδευσης που θα προσδιοριστεί από τον πιλοτικό έλεγχο.

Για την εφαρμογή στον Νικόλαο, το ΠΕ0 απογράφει πλήρη έργα και όχι μόνο απομονωμένες φόρμουλες: αριθμό και σειρά οικογενειών, ήχο, είδος, δομικό και κειμενικό τόπο, συνθέτη, εξηγητή, πηγή, ανεξαρτησία μάρτυρα και όλες τις διαθέσιμες χρονολογήσεις. Οι στόχοι των 15--25 έργων Νικολάου και 30--40 έργων αναφοράς είναι όρια απογραφής· η πραγματική κλίμακα ακολουθεί τον μετρημένο χρόνο κωδικοποίησης, τον αριθμό ανεξάρτητων μελοποιών και τις επιλέξιμες συγκρίσεις. Αν ο μετρημένος διαθέσιμος χρόνος δεν επαρκεί, περιορίζεται πρώτα η H.7 στις συγκρίσεις ακριβών κοινών χωρίων, έπειτα το σώμα της M5 και τελευταία τα σύνολα των M1/M2· η απουσία κατηγορικού αποτελέσματος στην H.7 δεν ερμηνεύεται ως αρνητικό εύρημα. Πρώτο παραδοτέο της M5 είναι το προφίλ ως προς το δηλωμένο περιβάλλον αναφοράς σε συγκρίσιμο σώμα της Νέας Μεθόδου. Η λεπτομερής δημοσιευμένη ανάλυση του Ιωαννίδη (2022, \emph{Θεολογία} 93(4), 173--193) έχει επαληθευθεί στην επίσημη έκδοση του περιοδικού και χρησιμοποιείται ως προκαθορισμένη ανεξάρτητη δημοσιευμένη μελέτη περίπτωσης· οι συγκεκριμένες μουσικές αντιστοιχίσεις ελέγχονται πάντοτε στο πρωτογενές άρθρο και στις αντίστοιχες μουσικές πηγές πριν από την υπολογιστική σύγκριση. Η πρόσθετη εργασία αφορά μεταδεδομένα ρόλων, έλεγχο σημειογραφικής ισοδυναμίας, κατανομή αναφοράς σε επίπεδο μελοποιού και μουσικολογική ανάλυση των χωρίων με τη μεγαλύτερη τεκμηριωμένη διαφοροποίηση.

Η υπολογιστική συνιστώσα απαιτεί εξειδίκευση στη σύγκριση αναπαραστάσεων, στη σχεσιακή μάθηση και στην αξιολόγηση της γενίκευσης. Η ερμηνεία των πηγών, η τροπική επικύρωση και η αξιοπιστία των εκτελεστικών μετρήσεων απαιτούν συμπληρωματική τεχνογνωσία στη βυζαντινή μουσικολογία και στην ακουστική.

Η πλήρης απογραφή των σωμάτων και των κριτηρίων ένταξης βρίσκεται στο Παράρτημα~\ref{app:corpora}, οι διαχωρισμοί στο Παράρτημα~\ref{app:splits} και η υλοποίηση/αναπαραγωγιμότητα στο Παράρτημα~\ref{app:impl}.

\subsection{Συμπληρωματική ανάλυση κοινού κειμενικού χωρίου}

Σε πλήρη Χερουβικά του ίδιου ήχου και σε Κοινωνικά με ακριβώς κοινό στίχο συγκρίνεται η μουσική μεταχείριση προκαθορισμένων κοινών χωρίων μεταξύ ανεξάρτητων μελοποιήσεων. Η ανάλυση περιγράφει χωριστά τις κοινές τάσεις και την ετερογένεια μεταξύ μελοποιών και εξετάζει ειδικά τη διαφοροποίηση του Νικολάου στα ίδια χωρία. Όπου η απογραφή τεκμηριώνει διαφορετικά χωρία της ίδιας ανεξάρτητα ορισμένης κειμενικής κατηγορίας, επαρκή κοινή κάλυψη των ουσιωδών δομικών συνθηκών και δυνατότητα διαχωρισμού από τον συγκεκριμένο λειτουργικό τόπο, προστίθεται δευτερεύουσα σύνοψη ανά κατηγορία. Η οικογενειακή απόδοση παραμένει ανεξάρτητη από τις νέες κειμενικές ετικέτες. Η εργασία είναι προκαθορισμένη και δευτερεύουσα, χρησιμοποιεί τα δεδομένα και τους κανόνες σύγκρισης της M5 και δεν προσθέτει νέα επιβεβαιωτική υπόθεση· το αναλυτικό πρωτόκολλο βρίσκεται στην §~\ref{sec:text-locus}.

\section{Ιστορική εφαρμογή και χαρακτηρισμός πλήρους έργου: Νικόλαος Πρωτοψάλτης Σμύρνης}\label{sec:nikolaos}

\paragraph{Από το λεξιλόγιο στη μορφή του έργου.} Οι Rabinovitch και Carter-Enyi (2024, §§2--4) περιγράφουν διαδοχές galant schemata σε προσδιορισμένο σώμα, συνοπτικές μορφικές διαδρομές και τη χρήση τους για την ερμηνεία συγκεκριμένων μεταγενέστερων έργων. Το προηγούμενο συνδέει τοπικούς μουσικούς τύπους, σειρά εμφάνισης και μορφική τοποθέτηση. Για τη M5 η ανάλογη μουσικολογική ερώτηση αφορά την επιλογή, διαδοχή, λειτουργική χρήση και ανάπτυξη επικυρωμένων οικογενειών θέσεων σε συγκεκριμένα Χερουβικά και Κοινωνικά. Η βιβλιογραφική αυτή σχέση λειτουργεί ως συγκριτικό πλαίσιο και οι ποσοτικές διαστάσεις της M5 παραμένουν διακριτές.

\subsection{M5a: χαρακτηρισμός έργου ως προς τεκμηριωμένο περιβάλλον αναφοράς}\label{sec:m5-reference}

Στην περιγραφική σύνοψη των ακολουθιών καταγράφονται χωριστά οι αναγνωρισμένες οικογένειες και τα δομημένα, αλλά μη αποδοθέντα σε γνωστή οικογένεια χωρία. Η καταγραφή αυτών των περιοχών, ανάλογη προς τις \emph{non-specific} σκελετικές ακολουθίες των Rabinovitch--Carter-Enyi (2024, §2.5), προστατεύει την ανοικτή αναγνώριση. Ένα χωρίο μπορεί να έχει μουσική οργάνωση χωρίς ασφαλή οικογενειακή απόδοση. Τυχόν κατώφλια συνοπτικής απεικόνισης συχνότητας/μεταβάσεων δηλώνονται πριν από την κύρια σύγκριση ή συνοδεύονται από ευαισθησία, επειδή η αναλυτική μορφή ενός διαγράμματος εξαρτάται από τις επιλογές κατωφλίωσης (Rabinovitch--Carter-Enyi, 2024, §2.9). Τα όρια επάρκειας και οι κύριες ποσότητες της M5 παραμένουν τα προκαθορισμένα.

Η μελοτροπική θεμελίωση των Αλεξάνδρου--Troelsgård (2020, σσ.~179--180) και του Troelsgård (2011, §57, σ.~61) στηρίζει την εξέταση της επιλογής και διαδοχής οικογενειών σε τοπικά τροπικά περιβάλλοντα. Η οικογενειακή ετικέτα και το τροπικό αποτέλεσμα διατηρούν χωριστά πεδία· όπου το \texttt{local\_modal\_state} αποτελεί είσοδο σε ταυτοτικό κριτήριο, δεν χρησιμοποιείται ως ανεξάρτητο εξωτερικό αποτέλεσμα χωρίς προκαθορισμένη απόκρυψη του αντίστοιχου γνωρίσματος και αξιολόγηση ευαισθησίας. Η τροπική αυτή ανάγνωση παραμένει ερμηνευτική διερεύνηση μέσα στους υπάρχοντες άξονες M5, χωρίς νέο estimand.

Η M5 αποτελεί την ιστορική εφαρμογή του επικυρωμένου πλαισίου σε πλήρη έργα. Προϋποθέτει την οριστικοποιημένη αναπαράσταση $R^\star$, το επικυρωμένο λεξιλόγιο οικογενειών και τους κανόνες αποχής που έχουν προκύψει από τις προηγούμενες ενότητες. Το νέο αντικείμενο είναι το \textbf{πλήρες έργο}: η σύνθεση περιγράφεται ως χάρτης τεκμηριωμένων εμφανίσεων μέσα στη μουσική ροή, με καταγραφή της σειράς, της κάλυψης, του ήχου, του δομικού και κειμενικού τόπου, του είδους και της προέλευσης κάθε εμφάνισης, καθώς και των περιοχών που παραμένουν συνδετικές, μη αποδοθείσες ή ανεπίλυτες. Ο σχεδιασμός τυποποιεί σε συγκρίσιμη μορφή μια καθιερωμένη μορφολογική πρακτική των αναλύσεων Χερουβικών: περιγραφή έργων μέσω του λεξιλογίου θέσεων, των ακολουθιών τους, της καταληκτικής και τροπικής οργάνωσης και του τρόπου ανάπτυξης των επιμέρους σχηματισμών (Καραγκούνης, 2003, 2005).

Η υπολογιστική ροή της M5 είναι
\[
\begin{aligned}
\text{πλήρες έργο}
&\rightarrow \text{τμηματοποίηση}
\rightarrow \text{οικογενειακή απόδοση ή αποχή}\\
&\rightarrow \text{χάρτης έργου}
\rightarrow \Omega
\rightarrow \text{πολυδιάστατος χαρακτηρισμός}.
\end{aligned}
\]
Κάθε αναγνωρισμένη εμφάνιση διατηρεί τη θέση της μέσα στο έργο και σύνδεση με τη συγκεκριμένη πηγή. Η τελική έξοδος συνδυάζει συνοπτικά μεγέθη με ανακτήσιμα τοπικά χωρία, ώστε κάθε υπολογιστική διαφοροποίηση να μπορεί να εξεταστεί μουσικολογικά στο υλικό που την παρήγαγε.

Η M5 διατηρεί δύο διακριτές διαδικασίες. Η \textbf{ανεξάρτητη μουσικολογική ανάλυση} συγκροτεί το σύνολο $\mathcal H_{\mathrm{acc}}=\{H_1,\ldots,H_q\}$ των αποδεκτών αναλύσεων με βάση προκαθορισμένους κανόνες, χωρίς γνώση των τιμών των μεγεθών της M5. Η \textbf{υπολογιστική ανάλυση} $\widehat H_{\mathrm{model}}$ είναι ο χάρτης που παράγει το σύστημα χωρίς να του παρέχεται η πλήρης οριοθέτηση και οικογενειακή απόδοση των ειδικών. Μία προκαθορισμένη $H_0\in\mathcal H_{\mathrm{acc}}$ μπορεί να χρησιμοποιείται για παρουσίαση και αναπαραγωγιμότητα, αλλά δεν θεωρείται μοναδική αλήθεια.

Για κάθε διαθέσιμη διάσταση $k\in\{\mathrm{fam},\mathrm{seq},\mathrm{func},\mathrm{within}\}$ αναφέρονται χωριστά τα $D_k(H;\Omega)$ για τις αποδεκτές αναλύσεις και το $D_k(\widehat H_{\mathrm{model}};\Omega)$. Για βαθμωτό μέγεθος ορίζεται το εύρος ευαισθησίας
\[
\mathcal I_k^{\mathrm{acc}}=\left[\min_{H\in\mathcal H_{\mathrm{acc}}}D_k(H;\Omega),\;\max_{H\in\mathcal H_{\mathrm{acc}}}D_k(H;\Omega)\right],
\]
και η απόσταση της υπολογιστικής τιμής από αυτό το εύρος. Τιμή του μοντέλου μέσα στο $\mathcal I_k^{\mathrm{acc}}$ δηλώνει συμβατότητα της συγκεκριμένης σύνοψης με τις αποδεκτές μουσικολογικές αναλύσεις, όχι τέλεια συμφωνία τμηματοποίησης ή οικογενειακής απόδοσης. Οι δομικές διαφορές συγκρίνονται επιπλέον με boundary/family metrics ως προς το σύνολο $\mathcal H_{\mathrm{acc}}$. Το $\mathcal I_k^{\mathrm{acc}}$ είναι \textbf{εύρος ευαισθησίας στην ανάλυση} και όχι διάστημα εμπιστοσύνης. Όταν το $\mathcal H_{\mathrm{acc}}$ προσεγγίζεται από πεπερασμένο διαθέσιμο υποσύνολο αποδεκτών αναλύσεων, το αναφερόμενο εύρος χαρακτηρίζεται προσεγγιστικό και μπορεί να υποεκτιμά το πλήρες εύρος αναλυτικής αβεβαιότητας. Αν διαφορετικές αποδεκτές αναλύσεις οδηγούν σε αντίθετες ουσιαστικές ερμηνείες για την ίδια διάσταση, το αντίστοιχο συμπέρασμα χαρακτηρίζεται αναλυτικά ασταθές.


Ως \textbf{τεκμηριωμένο περιβάλλον αναφοράς} ορίζεται
\[
\Omega=(D_{\Omega},\rho,\Pi,\mathcal K_{\Omega}),
\]
όπου $D_{\Omega}$ είναι το ρητό σύνολο έργων και πηγών που χρησιμοποιούνται ως αναφορά, $\rho$ η οριστικοποιημένη αναπαράσταση και η ανάθεση οικογενειών από M1--M3a, $\Pi$ το κοινό πιθανοτικό μοντέλο ακολουθιών και $\mathcal K_{\Omega}$ οι κανόνες ένταξης στο σώμα, χρονολογίας και ανεξαρτησίας. Κάθε $D_{\Omega}$ συνοδεύεται από συνθέτη, έργο, πηγή ή έκδοση, εξηγητή όπου υπάρχει, χρονολογία ή terminus ante quem, σημειογραφικό καθεστώς, είδος, ήχο και σχέση εξάρτησης από άλλους μάρτυρες. Οι συγκρίσεις αφορούν επομένως συγκεκριμένο τεκμηριωμένο σώμα και όχι αφηρημένο ή ενιαίο πρότυπο της παράδοσης.

Για έργο-στόχο $w$, ο κύριος κανόνας χρονολογικού ελέγχου επιτρέπει στο σώμα αναφοράς έργα που τεκμηριωμένα προηγούνται ή είναι σύγχρονα και ανεξάρτητα από το έργο-στόχο. Τα έργα του ίδιου μελοποιού και οι δικές του εξηγήσεις εξαιρούνται από την κύρια αναφορά όταν εξετάζεται η ιδιοτυπία του συγκεκριμένου μελοποιού. Αβέβαιες χρονολογίες, μεταγενέστερες εκδόσεις παλαιότερου υλικού και εξαρτημένοι μάρτυρες καταγράφονται χωριστά και οδηγούν σε αντίστοιχα περιορισμένο ιστορικό ισχυρισμό.

Το PPM της M3c επαναχρησιμοποιείται εδώ ως μοντέλο της ακολουθίας οικογενειών:
\[
P_{\Omega}(F_t\mid F_{<t},M_t,L_t,G_t),
\qquad
I_{\Omega}(t)=-\log_2 P_{\Omega}(F_t\mid F_{<t},M_t,L_t,G_t),
\]
Εδώ το $F_{<t}$ δηλώνει τις προηγούμενες οικογένειες της τεκμηριωμένης συνεχούς διαδοχής, το $M_t$ την τροπική κατάσταση, το $L_t$ το δομικό, κειμενικό και προσωδιακό περιβάλλον και το $G_t$ το είδος του μέλους, όπως καταγράφεται στα μεταδεδομένα του έργου. Όταν οι συνδυασμένες συνθήκες είναι αραιές, εφαρμόζεται ο προκαθορισμένος κανόνας επιστροφής σε λιγότερο ειδικό συμφραζόμενο. Μεγαλύτερο $I_{\Omega}$ σημαίνει ότι η παρατηρημένη επιλογή οικογένειας είναι λιγότερο προβλέψιμη υπό το συγκεκριμένο $D_{\Omega}$· το μέγεθος δεν αποτελεί από μόνο του μέτρο μουσικολογικής απόστασης, ποιότητας ή ιστορικής καινοτομίας. Η ιστορική του ερμηνεία αφορά τη στατιστική οργάνωση των τεκμηριωμένων μουσικών πηγών του σώματος αναφοράς. Η σχέση με ανθρώπινες κρίσεις εξετάζεται χωριστά στη M3c.

Η M5 παράγει \textbf{πολυδιάστατο προφίλ διαφοροποίησης} αντί για έναν συνολικό βαθμό. Οι τέσσερις διαστάσεις απαντούν σε διαφορετικά ερωτήματα για το ίδιο έργο: \textbf{ποιες οικογένειες επιλέγονται} ($D_{\mathrm{fam}}$), \textbf{πώς διατάσσονται στη μουσική ροή} ($D_{\mathrm{seq}}$), \textbf{σε ποιο δομικό, κειμενικό και τροπικό περιβάλλον χρησιμοποιούνται} ($D_{\mathrm{func}}$) και \textbf{πώς πραγματώνονται εσωτερικά οι εμφανίσεις της ίδιας οικογένειας} ($D_{\mathrm{within}}$). Η χωριστή αναφορά τους διατηρεί διακριτές μορφές μουσικής διαφοροποίησης που μπορούν να συνυπάρχουν με διαφορετικούς συνδυασμούς.

Η M5 διατηρεί τον χάρτη του πλήρους έργου και αναφέρει ρητά το υλικό στο οποίο στηρίζεται κάθε διάσταση του προφίλ.

Η \texttt{context\_compatibility} μπορεί να συνοδεύει τις εμφανίσεις ως περιγραφικό στρώμα. Δηλώνει αν μια χρήση είναι τεκμηριωμένα συνήθης, τεκμηριωμένα ασυνήθιστη ή ανεπαρκώς γνωστή. Δεν μετατρέπεται σε βαθμό ποιότητας του έργου. Η κύρια οικογενειακή απόδοση δεν επιτρέπεται να χρησιμοποιεί τη μορφική ή λειτουργική θέση ως υποκατάστατο της ταυτότητας όταν η ίδια πληροφορία πρόκειται αργότερα να αξιολογηθεί μέσω του $D_{\mathrm{func}}$. Μόνο συνθήκη που έχει παγώσει ως \texttt{required} για συγκεκριμένο τύπο μπορεί να συμμετέχει στην απόφαση, και τότε αναφέρεται χωριστή ανάλυση ευαισθησίας του $D_{\mathrm{func}}$ χωρίς αυτή την πληροφορία ως κριτήριο ένταξης, εφόσον το υλικό το επιτρέπει. Έτσι η M5 διακρίνει την ταυτότητα μιας θέσης από τη συνθετική της τοποθέτηση και περιορίζει την κυκλικότητα ανάμεσα στην αναγνώριση και στην επακόλουθη λειτουργική ανάλυση. Το ιστορικό τροπικό--λειτουργικό προφίλ ενός έργου συγκροτείται από τις πηγές του ίδιου του έργου ---μαρτυρίες, φθορές, καταλήξεις, μελωδική πορεία και τεκμηριωμένες ιστορικές αποδόσεις--- ενώ μεταγενέστερες ισοκρατηματικές εκδόσεις ή εκτελέσεις αποθηκεύονται ως χωριστή ερμηνευτική μαρτυρία με χρονολογία και προέλευση και δεν χρησιμοποιούνται αναδρομικά ως ιστορική ετικέτα της σύνθεσης. Το λεκτικό καθεστώς και ο τύπος μουσικού τμήματος καταγράφονται χωριστά και οι αντίστοιχες συγκρίσεις παρουσιάζονται ως στρωματοποιημένες περιγραφές ή αναλύσεις ευαισθησίας όπου υπάρχει επαρκής κοινή κάλυψη και τεκμηριωμένη προέλευση. Τα κρατήματα παραμένουν συνδεδεμένα με τη συγκεκριμένη πηγή και την απόδοσή τους, ώστε μεταγενέστερη προσθήκη, αβέβαιη καταγωγή ή μερική διατήρηση να διακρίνεται από τεκμηριωμένη μελοποιητική πρακτική. Ο πλήρης χάρτης μπορεί να διατηρεί αβέβαια ή μερικά τμήματα με τη σωστή κατάσταση, χωρίς να τα προσμετρά ως πλήρως συγκρίσιμα εκεί όπου αυτό δεν τεκμηριώνεται. Όταν η παρουσία ή η πληρότητα κρατήματος διαφέρει μεταξύ των συγκρινόμενων πηγών, οι συγκριτικές τιμές των μεγεθών της M5 υπολογίζονται πρώτα στο κοινό υλικό που παραδίδουν όλες οι συγκρινόμενες πηγές· οι αντίστοιχες τιμές στο πλήρες διαθέσιμο έργο αναφέρονται χωριστά ως ανάλυση ευαισθησίας. Οι κειμενικές κατηγορίες της §~\ref{sec:text-locus} δεν προστίθενται αυτομάτως στους υπάρχοντες κανόνες του συγκρίσιμου κελιού $c$, δηλαδή του συνδυασμού προκαθορισμένων μουσικών συνθηκών και συνθηκών προέλευσης των πηγών μέσα στον οποίο γίνεται η σύγκριση.

Οι τέσσερις κύριες διαστάσεις είναι:
\begin{enumerate}[label=\textbf{\arabic*.},leftmargin=2.2em]
\item \textbf{Επιλογή οικογενειών ($D_{\mathrm{fam}}$).} Κύριο μέτρο είναι η απόκλιση Jensen--Shannon μεταξύ της κατανομής οικογενειών του έργου και της αντίστοιχης κατανομής του $D_{\Omega}$ σε κοινό πεδίο τιμών, υπό τις επιλέξιμες τροπικές και ειδολογικές συνθήκες. Στην κύρια μεταξύ μελοποιών αναφορά κάθε επιλέξιμος μελοποιός λαμβάνει ίσο συνολικό βάρος· στάθμιση κατά τη συχνότητα στο σώμα αναφέρεται ως ανάλυση ευαισθησίας, ώστε η μεγάλη εκπροσώπηση ενός μελοποιού να μην μετατρέπεται σιωπηρά σε κυρίαρχη αναφορά.
\item \textbf{Διαδοχή οικογενειών ($D_{\mathrm{seq}}$).} Το μέσο $I_{\Omega}$ του έργου συγκρίνεται με την αντίστοιχη κατανομή αναφοράς σε επίπεδο μελοποιού, εξισωμένη ως προς μήκος και επιλέξιμες συνθήκες, όταν υπάρχουν επαρκή δεδομένα για αξιόπιστη εκτίμησή της. Η διαφορά του μέσου πληροφοριακού περιεχομένου από την αναφορά δείχνει πόσο λιγότερο ή περισσότερο προβλέψιμη είναι η ακολουθία ως προς αυτή την κατανομή αναφοράς. Χωριστό υπο-μέτρο καταγράφει την επανεμφάνιση ακολουθιών οικογενειών και τον ανασυνδυασμό γνωστών υποακολουθιών σε διαφορετικό συνθετικό περιβάλλον, ακολουθώντας τη μουσικολογική διάκριση του Καραγκούνη (2005, σσ.~175--176).
\item \textbf{Λειτουργική τοποθέτηση ($D_{\mathrm{func}}$).} Εδώ το $G$ δηλώνει το είδος του μέλους, όπως καταγράφεται στα μεταδεδομένα του έργου. Συγκρίνονται οι υπό συνθήκη κατανομές $P_w(F\mid L,M,G)$ και $P_{\Omega}(F\mid L,M,G)$ στις τομές με επαρκή υποστήριξη. Η κύρια σύνοψη χρησιμοποιεί σταθμισμένη Jensen--Shannon απόκλιση μεταξύ των επιλέξιμων τομών, ενώ οι τοπικοί λογαριθμικοί λόγοι ανά οικογένεια και τόπο παραμένουν ερμηνεύσιμες δευτερεύουσες έξοδοι.
\item \textbf{Πραγμάτωση μέσα στην οικογένεια ($D_{\mathrm{within}}$).} Η σύγκριση χρησιμοποιεί τις τοπικές καταστάσεις του EISE που διατηρούν την εσωτερική παραλλαγή και όχι μόνο την προβολή ταυτότητας $z_{\mathrm{id}}$, η οποία έχει εκπαιδευτεί να συμπτύσσει ενδοοικογενειακές διαφορές. Η πραγμάτωση συγκρίνεται με το κέντρο, τη διασπορά και το τεκμηριωμένο εύρος της ίδιας οικογένειας στο $D_{\Omega}$, με τις ακριβείς ποσότητες να ορίζονται στο Παράρτημα~\ref{app:match-nikolaos}. Το $D_{\mathrm{within}}$ παραμένει ενιαίο μέγεθος· όπου υπάρχουν αξιόπιστες κοινές εσωτερικές περιοχές, συνοδεύεται από χωριστό ερμηνεύσιμο χρονικό διαγνωστικό $(Q_x,p_x,d_{\mathrm{time}})$ που περιγράφει τη συνολική σημειογραφική ανάπτυξη και την κατανομή του χρόνου χωρίς να προστίθεται πέμπτη διάσταση M5.

Το χρονικό διαγνωστικό υπολογίζεται μόνο μετά την οριοθέτηση. Το $Q_x$ μετρά τη συνολική σημειογραφική χρονική ανάπτυξη και μπορεί να γραφεί ως $Q_x=n_x\bar v_x$, ώστε να διακρίνεται η συμβολή του αριθμού γεγονότων από τη μέση χρονική αξία τους. Το $p_x$ περιγράφει πώς κατανέμεται αυτή η ανάπτυξη στις αντίστοιχες εσωτερικές περιοχές της θέσης. Το $d_{\mathrm{time}}$ συνοψίζει την απόσταση αυτού του σχετικού προφίλ από το προφίλ αναφοράς της ίδιας οικογένειας. Οι περιοχές ορίζονται χωρίς χρονική πληροφορία, ώστε η ποσότητα που εξετάζεται να μην καθορίζει προηγουμένως την ίδια την αντιστοίχιση.
\end{enumerate}

Οι τέσσερις διαστάσεις διατηρούνται χωριστές. Ένα έργο μπορεί, για παράδειγμα, να χρησιμοποιεί συνηθισμένο λεξιλόγιο οικογενειών και συγχρόνως να παρουσιάζει ασυνήθιστες διαδοχές ή διαφορετική ενδοοικογενειακή πραγμάτωση. Για κάθε κύρια ποσότητα κατασκευάζεται \textbf{κατανομή αναφοράς σε επίπεδο μελοποιού}: κάθε άλλος επιλέξιμος μελοποιός αντιμετωπίζεται διαδοχικά ως στόχος, αφαιρείται από το αντίστοιχο σώμα αναφοράς και το προφίλ του υπολογίζεται με τους ίδιους κανόνες. Η εμπειρική θέση ή το εκατοστημόριο του Νικολάου αναφέρεται μόνο όταν ο αριθμός ανεξάρτητων επιλέξιμων μελοποιών είναι τουλάχιστον $K_{\min}$. Το $K_{\min}$ οριστικοποιείται στο ΠΕ0 με προσομοίωση της διακριτικής ικανότητας και του σφάλματος της εμπειρικής κατάταξης, πριν εξεταστούν αποτελέσματα M5. Με $K<K_{\min}$, το προφίλ αναφέρεται περιγραφικά με διαστήματα αβεβαιότητας σε επίπεδο έργου, τον πραγματικό αριθμό έργων και μελοποιών και χωρίς ισχυρισμό θέσης σε πληθυσμιακή κατανομή συνθετών. Για ένα μεμονωμένο έργο-στόχο, το ίδιο το έργο παραμένει σταθερό. Η αβεβαιότητα εκτιμάται με αναδειγματοληψία ολόκληρων έργων ή/και μελοποιών του σώματος αναφοράς, διάδοση της βαθμονομημένης αβεβαιότητας οικογενειακής απόδοσης και ανάλυση ευαισθησίας σε εναλλακτικά αποδεκτά όρια θέσεων. Οι θέσεις του ίδιου έργου δεν αντιμετωπίζονται ως ανεξάρτητα δείγματα.

Η ελάχιστη εκδοχή της M5 χρησιμοποιεί συγκρίσιμο υλικό της Νέας Μεθόδου. Πρόσθετη σύγκριση της προγενέστερης και της σύγχρονης ομάδας αναφοράς ενεργοποιείται μόνο όταν το ΠΕ0 τεκμηριώσει επαρκείς συνεχείς διαδοχές, ανεξάρτητες πηγές και κοινή κάλυψη των ίδιων οικογενειών. Σε αυτή την περίπτωση η μεταβολή του $I_{\Omega}$ περιγράφει μετατόπιση της προβλεψιμότητας μεταξύ δύο δηλωμένων σωμάτων. Η μεταγενέστερη ομάδα αναφοράς παραμένει διερευνητική επέκταση. Παλαιότερες σημειογραφίες εντάσσονται μόνο μέσω ελεγμένων αντιστοιχίσεων πηγής--εξήγησης και με χωριστή δήλωση του επιπέδου στο οποίο παραμένουν συγκρίσιμες.

\paragraph{Προδιαγραφές εισόδου--εξόδου και ελάχιστη/πλήρης M5.}
Η είσοδος της M5 είναι πλήρες κωδικοποιημένο έργο από συγκεκριμένη τεκμηριωμένη πηγή. Η έξοδος περιλαμβάνει πάντοτε \textbf{χάρτη έργου}: τεκμηριωμένες ή αναγνωρισμένες εμφανίσεις με όρια, οικογένεια ή αποχή, βαθμονομημένη βεβαιότητα, ήχο, δομικό και κειμενικό τόπο, πλησιέστερους τεκμηριωμένους συγγενείς και προέλευση, μαζί με τις μεταξύ τους τεκμηριωμένες σχέσεις σειράς και τις ενδιάμεσες συνδετικές, μη αποδοθείσες ή ανεπίλυτες περιοχές. Αναφέρεται υποχρεωτικά η \textbf{κάλυψη} του έργου, τόσο ως ποσοστό επιλέξιμων μουσικών γεγονότων που ανήκουν σε αναγνωρισμένες θέσεις όσο και ως πλήθος αναγνωρισμένων/μη αναγνωρισμένων περιοχών. Τα ποσοτικά προφίλ ερμηνεύονται μόνο για το καλυπτόμενο υλικό.

Η \textbf{M5-min} χαρακτηρίζει ολόκληρο το έργο ως προς το λεξιλόγιο και τις πραγματώσεις και δεσμεύεται για $D_{\mathrm{fam}}$ και $D_{\mathrm{within}}$. Ο χάρτης διατηρεί πάντοτε τη σειρά και τον τόπο των θέσεων, ακόμη όταν οι βαθμονομημένες συγκρίσεις $D_{\mathrm{seq}}$ και $D_{\mathrm{func}}$ δεν μπορούν να εκτιμηθούν αξιόπιστα. Η \textbf{M5-full} προσθέτει τις δύο αυτές διαστάσεις όταν η κάλυψη μεταβάσεων, λειτουργικών τόπων και ανεξάρτητων μελοποιών ικανοποιεί τα προκαθορισμένα κριτήρια επάρκειας. Η κατάταξη σε κατανομή αναφοράς σε επίπεδο μελοποιού αποτελεί επιπλέον έξοδο και ενεργοποιείται μόνο με $K\ge K_{\min}$.

Η ακολουθιακή ανάλυση εφαρμόζεται μόνο όπου τεκμηριώνεται σαφής διαδοχή οικογενειών. Μη αναγνωρισμένη εμφάνιση, χωριστό συνδετικό υλικό, ανεπίλυτη περιοχή ή ανεπίλυτη επικάλυψη \textbf{διακόπτει} την κύρια ανάλυση της διαδοχής. Αν η ακολουθία είναι $F_1\rightarrow\mathrm{UNKNOWN}\rightarrow F_2$, δεν δημιουργείται τεχνητή μετάβαση $F_1\rightarrow F_2$, ούτε επιβάλλεται αυθαίρετη σειρά σε επικαλυπτόμενες εμφανίσεις χωρίς τεκμηριωμένη διάταξη. Μετά από εσωτερική διακοπή το ιστορικό οικογενειακού συμφραζομένου μηδενίζεται, αλλά η πρώτη επόμενη εμφάνιση σημειώνεται ως επανεκκίνηση και όχι ως πραγματική αρχή έργου. Η πολιτική συμμετοχής αυτών των εμφανίσεων στο $I_{\Omega}$ οριστικοποιείται πριν από την M5 και εφαρμόζεται απαράλλακτα στον Νικόλαο και στο σώμα αναφοράς. Έτσι η μερική κάλυψη δεν παράγει ψευδείς μεταβάσεις ή ψευδές ιστορικό συμφραζόμενο.


\paragraph{Ανεξαρτησία αναγνώρισης και λειτουργικής εξήγησης.} Οι Rabinovitch και Carter-Enyi (2024, §2.6, §4.3 και υποσημ.~37) επισημαίνουν ότι η αναγνώριση ενός σχήματος μπορεί να επηρεάζεται από την ήδη γνωστή μορφική θέση του. Για τη M5, η κύρια οικογενειακή απόδοση προηγείται της αξιολόγησης $D_{\mathrm{func}}$ και καταγράφεται ποια λειτουργική πληροφορία ήταν ορατή στην αντίστοιχη κρίση. Όταν ο τύπος διαθέτει εκ των προτέρων παγωμένο, ενεργό \texttt{required} γνώρισμα λειτουργικού περιβάλλοντος, η επίδρασή του εξετάζεται με την προβλεπόμενη ανάλυση ευαισθησίας. Η προδιαγραφή λειτουργεί ως προστασία της ερμηνείας και όχι ως νέα οικογένεια πειραμάτων.

\subsection{M5b: εφαρμογή στον Νικόλαο Πρωτοψάλτη Σμύρνης}\label{sec:m5-nikolaos}

Η πρώτη πλήρης εφαρμογή εξετάζει τη μελοποιητική πρακτική του \textbf{Νικολάου Γεωργίου, Πρωτοψάλτη Σμύρνης} του 19ου αιώνα (Τόμα, 2018). Το μουσικολογικό ερώτημα είναι συγκεκριμένο: \textbf{πού ακριβώς βρίσκεται η διαφοροποίηση των έργων του ως προς συγκρίσιμα, χρονολογικά ελεγχόμενα έργα άλλων μελοποιών με τεκμηριωμένη και συγκρίσιμη προέλευση των πηγών;} Η ανάλυση εξετάζει αν η διαφοροποίηση εμφανίζεται κυρίως στην επιλογή οικογενειών, στη σειριακή τους οργάνωση, στη λειτουργική τους τοποθέτηση, στην πραγμάτωση γνωστών οικογενειών ή σε συνδυασμό αυτών των διαστάσεων.

Το ΠΕ0 καταγράφει τους μελοποιούς και τα έργα που μπορούν να ενταχθούν στις ομάδες αναφοράς, μαζί με την ιστορική προέλευση κάθε πηγής. Η προγενέστερη ομάδα αναφοράς $C_{\mathrm{pre}}$ ενεργοποιείται μόνο όταν περνά τα κριτήρια κοινής κάλυψης, ανεξαρτησίας και ελέγχου εξηγητών· τότε περιλαμβάνει προγενέστερους μελοποιούς, όπως τον Δανιήλ Πρωτοψάλτη, τον Πέτρο Λαμπαδάριο Πελοποννήσιο και τον Πέτρο Βυζάντιο, όταν το υλικό τους είναι διαθέσιμο σε τεκμηριωμένη εξήγηση και πληροί τα κριτήρια συγκρισιμότητας. Η κύρια σύγχρονη ομάδα αναφοράς $C_{\mathrm{contemp}}$ περιλαμβάνει επιλέξιμους μελοποιούς του 19ου αιώνα, όπως ο Χουρμούζιος, ο Γρηγόριος, ο Θεόδωρος Φωκαέας και ο Ιωάννης Πρωτοψάλτης, με τον Νικόλαο πάντοτε εκτός της αναφοράς: $C_{\mathrm{contemp}}\setminus \mathrm{Nik}$. Κάθε πρόσωπο καταγράφεται με τον πραγματικό ρόλο του στη συγκεκριμένη σχέση ---μελοποιός, εξηγητής, εκδότης ή φορέας πηγής--- και οι χρονολογίες σύνθεσης ή απόδοσης, σωζόμενου μάρτυρα, εξήγησης και έκδοσης παραμένουν χωριστές. Επειδή ο Γρηγόριος και ο Χουρμούζιος μπορούν να εμφανίζονται ταυτόχρονα ως εξηγητές του $C_{\mathrm{pre}}$ και ως μελοποιοί του $C_{\mathrm{contemp}}$, η ιστορική σύγκριση συνοδεύεται από προκαθορισμένη \textbf{ανάλυση ευαισθησίας κοινού εξηγητή}: επαναλαμβάνεται μετά την αφαίρεση από το $C_{\mathrm{pre}}$ των πραγματώσεων που εξηγήθηκαν από πρόσωπο το οποίο συμμετέχει ως μελοποιός στο $C_{\mathrm{contemp}}$, και επιπλέον με leave-one-exegete-out όπου το πλήθος πηγών το επιτρέπει. Η κύρια σύγκριση απαιτεί την κάλυψη που ορίζει ο προκαθορισμένος έλεγχος και επαρκείς ανεξάρτητους μελοποιούς· το $C_{\mathrm{pre}}$ λειτουργεί ως πρόσθετη ιστορική αναφορά μόνο όταν οι εξηγητές, οι εξαρτήσεις και η κοινή κάλυψη των ίδιων οικογενειών μπορούν να ελεγχθούν.

Η ανάλυση πλήρους Χερουβικού ή Κοινωνικού επιστρέφει πάντοτε τον χάρτη του έργου, την κάλυψη, την οικογενειακή απόδοση ή αποχή, τους πλησιέστερους τεκμηριωμένους συγγενείς, τον ήχο και τον δομικό και κειμενικό τόπο. Η M5-min συνοψίζει τα $D_{\mathrm{fam}}$ και $D_{\mathrm{within}}$ που μπορούν να εκτιμηθούν αξιόπιστα. Η M5-full προσθέτει $I_{\Omega}$/$D_{\mathrm{seq}}$ και $D_{\mathrm{func}}$ μόνο όταν ικανοποιούνται τα αντίστοιχα κριτήρια επάρκειας, ενώ η εμπειρική θέση στην κατανομή αναφοράς σε επίπεδο μελοποιού αναφέρεται μόνο με $K\ge K_{\min}$. Το αποτέλεσμα είναι ερμηνεύσιμο μουσικολογικό προφίλ που συνδέει κάθε διαθέσιμη διάσταση με συγκεκριμένα μουσικά χωρία, χωρίς να παρουσιάζει αριθμητικά αποτελέσματα για διαστάσεις που δεν μπορούν να εκτιμηθούν αξιόπιστα.

Η υπάρχουσα ιστορική βιβλιογραφία προσφέρει το πλαίσιο για την επιλογή των ερωτημάτων και των μελοποιών και έργων αναφοράς. Οι αναφορές σε διαφοροποίηση ύφους ή σε επιφυλάξεις συγχρόνων λειτουργούν ως ιστορικό κίνητρο προς έλεγχο και διασταυρώνονται με τις αντίστοιχες πηγές πριν χρησιμοποιηθούν σε ισχυρότερη ερμηνεία. Οι υπολογιστικές ποσότητες χαρακτηρίζουν τη σχέση του έργου με το δηλωμένο $D_{\Omega}$. Η απόδοση ιστορικής καινοτομίας, επιρροής, πατρότητας ή αιτιακής μετάδοσης στηρίζεται σε πρόσθετη τεκμηρίωση από τις ιστορικές πηγές.

\subsection{Ανεξάρτητη δημοσιευμένη μελέτη περίπτωσης}

Ως ανεξάρτητη δημοσιευμένη μελέτη περίπτωσης χρησιμοποιείται η μελέτη του Χρυσοβαλάντη Ιωαννίδη, «Η σύνθεση του Νικολάου Πρωτοψάλτου Σμύρνης στο Κοινωνικό ``Εἰς πᾶσαν τὴν γῆν'' σε ήχο πλ. Δ΄ και ο βαθμός εξάρτησής της από την προγενέστερη παράδοση της εκκλησιαστικής μελοποιίας», \emph{Θεολογία} 93(4), 2022, σσ.~173--193. Η επίσημη δημοσίευση έχει επαληθευθεί και περιέχει ρητή δομική συνεξέταση με το ομόηχο Κοινωνικό του Πέτρου Λαμπαδαρίου, καθώς και τοπικές σχέσεις με τον Χουρμούζιο. Οι ακριβείς μουσικές αντιστοιχίσεις της χρησιμοποιούνται ως προκαθορισμένος εξωτερικός έλεγχος μόνο με τα αντίστοιχα χωρία και τις πηγές τους. Τα σχετικά έργα, μάρτυρες και δημοσιευμένες κρίσεις παραμένουν εκτός ανάπτυξης, επιλογής αναπαράστασης, ορισμού οικογενειών και βαθμονόμησης.

Η M5 διατηρεί περιγραφικό και προκαθορισμένο χαρακτήρα και δεν προσθέτει νέα οικογένεια επιβεβαιωτικών υποθέσεων. Η επιτυχία της κρίνεται χωριστά για κάθε διάσταση που μπορεί να εκτιμηθεί αξιόπιστα, από την επάρκεια και την κάλυψη του σώματος, τη σταθερότητα έναντι εναλλακτικών επιλέξιμων σωμάτων αναφοράς και την ικανότητα του χάρτη και των διαθέσιμων προφίλ να οδηγούν σε συγκεκριμένες μουσικολογικές αντιπαραβολές των πηγών. Η αδυναμία αξιόπιστης εκτίμησης μιας πρόσθετης διάστασης δεν ακυρώνει το προκαθορισμένο M5-min. Το πλήρες τεχνικό πρωτόκολλο, τα κριτήρια επάρκειας για τις αντίστοιχες εκτιμήσεις και οι αναλύσεις ευαισθησίας δίνονται στην §~\ref{app:match-nikolaos}.

\section{Κεντρική συνεισφορά, πρόγραμμα και συμπέρασμα}

\subsection{Αναμενόμενη υπολογιστική, γνωσιακή και μουσικολογική γνώση}

\paragraph{Υπολογιστική συνεισφορά.}
Σε υπολογιστικό επίπεδο, η διατριβή αναμένεται να απαντήσει σε τέσσερα διακριτά ζητήματα και, εφόσον τα δεδομένα επαρκούν, σε ένα πέμπτο για το συμβολικό MIR, τη μάθηση αναπαραστάσεων και την υπολογιστική μουσικολογία.
\begin{enumerate}[label=\textbf{\arabic*.},leftmargin=2.2em]
\item \textbf{Αξία της μονάδας αναπαράστασης.} Προσδιορίζεται πότε η ρητή μουσική θέση υποστηρίζει τμηματοποίηση, ανάκτηση συγγενών πραγματώσεων, αναγνώριση οικογενειών με δυνατότητα αποχής και γενίκευση σε νέα έργα και ομάδες συγγενών πηγών διαφορετικά από ένα ισχυρό Transformer σε επίπεδο μελωδικού γεγονότος και από αλγοριθμικά επαγόμενες μονάδες.
\item \textbf{Πληροφορία της εξηγητικής εποπτείας.} Η αντιστοιχισμένη B2--EISE σύγκριση εξετάζει τι πρόσθετη πληροφορία προσφέρει η τεκμηριωμένη συν-εξήγηση του ίδιου ιστορικού χωρίου για την οργάνωση της ενδοοικογενειακής παραλλαγής. Οι δύο συνθήκες χρησιμοποιούν την ίδια οικογενειακή πληροφορία και αντιστοιχίζονται ως προς την έκθεση στα δεδομένα, τη δυσκολία των ζευγών, την αρχιτεκτονική και τον υπολογιστικό προϋπολογισμό. Το $\Delta_{\mathrm{hist}}$ ερμηνεύεται μέσα στο κοινό επιλέξιμο πεδίο της δοκιμής και σε σχέση με το ανεξάρτητο εκτελεστικό κριτήριο.
\item \textbf{Εξωτερική εγκυρότητα της αναπαράστασης.} Οι οριστικοποιημένες αναπαραστάσεις αξιολογούνται με μουσικολογικές κρίσεις, τροπική και εκτελεστική οργάνωση, ανθρώπινη κατηγοριοποίηση και πρόβλεψη της μελωδικής διαδοχής. Έτσι η υπολογιστική επίδοση συνδέεται με ανεξάρτητα κριτήρια μουσικής σημασίας.
\item \textbf{Χαρακτηρισμός πλήρους έργου ως προς αναφορά.} Η M5 συνδυάζει το επικυρωμένο λεξιλόγιο, την αναπαράσταση EISE και το PPM σε πρωτόκολλο με ρητά δηλωμένο περιβάλλον αναφοράς, στο οποίο περιγράφονται χωριστά η επιλογή οικογενειών, οι τεκμηριωμένες διαδοχές τους, η λειτουργική τους τοποθέτηση και η ενδοοικογενειακή πραγμάτωση. Το χρονολογικό όριο και η κατανομή αναφοράς σε επίπεδο μελοποιού μετατρέπουν τον χαρακτηρισμό ενός μελοποιού σε ελέγξιμη σύγκριση με δηλωμένο σύνολο έργων και πηγών.
\item \textbf{Μεταφερσιμότητα της τεκμηριωμένης παραλλαγής (υπό πύλη επάρκειας).} Η επέκταση της §~\ref{sec:atl} εξετάζει αν οι αντιστοιχίσεις ανάμεσα σε παράλληλες εξηγήσεις του ίδιου χωρίου μπορούν να μαθευτούν από λίγα ιστορικά ζεύγη και να προβλέψουν τις πραγματικές εξηγήσεις οικογενειών που το μοντέλο δεν έχει δει. Για το συμβολικό MIR προσφέρει μοντέλο μελωδικής μεταβολής εκτιμημένο από ιστορικά τεκμήρια· για τη μουσικολογία, μετρήσιμη περιγραφή του τι διατηρεί η εξήγηση πέρα από τη σημειογραφική επιφάνεια.
\end{enumerate}
Η υπολογιστική αρχή είναι δυνητικά μεταφέρσιμη σε πολιτισμικά εδραιωμένες μονάδες μεταβλητής έκτασης των οποίων η κοινή ταυτότητα συνυπάρχει με οργανωμένη παραλλαγή. Η εφαρμογή της σε άλλο ρεπερτόριο απαιτεί ανεξάρτητη επικύρωση των αντίστοιχων μονάδων και σχέσεων.

\paragraph{Γνωσιακή συνεισφορά.}
Η ανθρώπινη μελέτη διαχωρίζει τρία επίπεδα. Η M3b ελέγχει αν η οικογενειακή ένταξη αντιστοιχεί σε ανθρώπινη κατηγορική οργάνωση όταν η επιφανειακή απόσταση ελέγχεται ρητά και αν η εσωτερική παραλλαγή οργανώνεται με γεωμετρία που συμφωνεί με τις υπολογιστικές αναπαραστάσεις. Η M3c εξετάζει αν το πληροφοριακό περιεχόμενο σε επίπεδο οικογένειας προσθέτει εξηγητική ισχύ πέρα από το πληροφοριακό περιεχόμενο των επιμέρους μελωδικών γεγονότων. Η συνθήκη λατρευτικού συμφραζομένου ελέγχει ως κύριο δευτερεύον αποτέλεσμα αν η μεταβολή της σχέσης με την ανθρώπινη κρίση για το πόσο απρόσμενη ακούγεται η συνέχεια είναι περισσότερο θετική ---ή λιγότερο αρνητική--- για το επίπεδο της οικογένειας από ό,τι για το επίπεδο του γεγονότος. Η βιωματική απόκριση παραμένει διερευνητική και η νευρική μελέτη με χρονική ανάλυση ενεργοποιείται μόνο εφόσον πληρωθούν τα προβλεπόμενα κριτήρια υποδομής και δεοντολογίας.

\paragraph{Μουσικολογική συνεισφορά.}
Η μουσικολογική συμβολή της έρευνας έγκειται στη συστηματική τυποποίηση και συγκριτική εξέταση της σχέσης θέσης, πραγματώσεων, παραλλαγής και ιστορικής μετάδοσης. Η προηγούμενη ιστορική και μορφολογική έρευνα έχει ήδη τεκμηριώσει σε μεγάλη κλίμακα την επανεμφάνιση και τις παραλλαγές των θέσεων· η παρούσα ερευνητική πρόταση μετατρέπει αυτή τη γνώση σε ρητό, επαναλήψιμο πλαίσιο στο οποίο η κοινή ταυτότητα, οι συγκεκριμένες πραγματώσεις, οι εμφανίσεις και οι ιστορικές σχέσεις εξετάζονται χωριστά και συγκριτικά. Η συμβολή οργανώνεται σε πέντε είδη γνώσης.
\begin{enumerate}[label=\textbf{\arabic*.},leftmargin=2.2em]
\item \textbf{Τυπολογία των θέσεων.} Παράγεται τεκμηριωμένο και εκδοχικά ελεγχόμενο λεξικό οικογενειών θέσεων. Κάθε οικογένεια συνδέεται με διαγνωστικές σχέσεις, κανόνες πλαισίου, πραγματώσεις, αμφίσημες περιπτώσεις και τεκμήρια προέλευσης. Η αντιπαραβολή καταλόγων, πηγών και ανεξάρτητων κρίσεων επιτρέπει να ελεγχθεί αν ένας προτεινόμενος τύπος έχει συνεκτικό πυρήνα, πολυθετική δομή ---δηλαδή διαφορετικούς συνδυασμούς γνωρισμάτων χωρίς ένα μοναδικό γνώρισμα κοινό σε όλα τα μέλη---, ανάγκη διάσπασης ή ασαφή όρια. Ο ήχος, η λειτουργική θέση, το έργο, η πηγή, ο μελοποιός και ο εξηγητής καταγράφονται με τον ρόλο που τεκμηριώνεται σε κάθε περίπτωση και δεν μετατρέπονται αυτομάτως σε κριτήρια ταυτότητας.
\item \textbf{Μορφολογία της παραλλαγής.} Για τις τεκμηριωμένες οικογένειες ελέγχεται ποια χαρακτηριστικά στηρίζουν την αναγνώριση και πώς κατανέμεται η παραλλαγή στο εσωτερικό τους: ανάπτυξη, σύντμηση, διαποίκιλση, προσαρμογή και, όπου τεκμηριώνεται στο κατάλληλο επίπεδο, ανασυνδυασμός. Η σύγκριση με σύγχρονες εκτελέσεις περιγράφει ποιες περιοχές εμφανίζουν μεγαλύτερη ή μικρότερη παρατηρούμενη εκτελεστική μεταβλητότητα.
\item \textbf{Δομική ανάλυση και προφίλ πλήρους έργου.} Τα επιλέξιμα πλήρη έργα περιγράφονται με χάρτη εμφανίσεων, λειτουργικών ρόλων και σχέσεων προς τις οικογένειές τους· η ακολουθιακή ανάλυση εφαρμόζεται στις τεκμηριωμένες συνεχείς διαδοχές. Η M5 προσθέτει ρητό περιβάλλον αναφοράς και πολυδιάστατο προφίλ $D_{\mathrm{fam}}$, $D_{\mathrm{seq}}$, $D_{\mathrm{func}}$ και $D_{\mathrm{within}}$, ώστε η μελοποιητική οργάνωση Χερουβικών και Κοινωνικών να γίνεται συγκρίσιμη ως προς λεξιλόγιο, διαδοχή, λειτουργική τοποθέτηση και ενδοοικογενειακή πραγμάτωση. Η αντιστοίχιση πολλαπλών μελοποιήσεων του ίδιου λειτουργικού κειμένου επιτρέπει επιπλέον να περιγραφεί η μουσική μεταχείριση προκαθορισμένων κοινών χωρίων και, στο επίπεδο που επιτρέπουν τα δεδομένα, να διακριθούν κοινές συμβάσεις του εξεταζόμενου ρεπερτορίου από τη διαφοροποίηση συγκεκριμένου μελοποιού.
\item \textbf{Εξήγηση και ιστορική μετάδοση.} Οι τεκμηριωμένες εξηγητικές σχέσεις επιτρέπουν συγκριτική μελέτη συγκεκριμένων αποδόσεων και, όπου οι πηγές το επιτρέπουν, διαδοχικών γραμμών μετάδοσης. Η ανάλυση εξετάζει σε επίπεδο θέσης τι διατηρείται, τι διαφοροποιείται και πώς οργανώνεται η μορφολογική μεταβολή, με τη μουσική συγγένεια καταγεγραμμένη χωριστά από τη συγκεκριμένη ιστορική σχέση.
\item \textbf{Θέση, ήχος και μελοποιητική τεχνική.} Προσδιορίζεται πώς οι οικογένειες κατανέμονται στους ήχους και στους λειτουργικούς ρόλους και πώς η χρήση τους συνδέεται με την τροπική οργάνωση. Στον Νικόλαο Πρωτοψάλτη Σμύρνης, το ίδιο πλαίσιο εξετάζει σε ποια από τις τέσσερις διαστάσεις του προφίλ ως προς το σώμα αναφοράς συγκεντρώνεται η διαφοροποίηση και ποια συγκεκριμένα χωρία την παράγουν, επιτρέποντας την άμεση επιστροφή από το ποσοτικό αποτέλεσμα στη μουσική πηγή.
\end{enumerate}
Το αποτέλεσμα είναι ένα αναλυτικό πλαίσιο μεταφέρσιμο σε άλλα έργα του παπαδικού ρεπερτορίου και ένα σώμα γνώσης για το πώς η θέση συγκροτεί, μεταδίδει και ανανεώνει τη μελοποιία.

\subsection{Πακέτα εργασίας, χρονοδιάγραμμα και παραδοτέα}

{\footnotesize
\begin{xltabular}{\textwidth}{@{}p{0.11\textwidth}p{0.15\textwidth}Y@{}}
\toprule
\textbf{ΠΕ} & \textbf{Μήνες} & \textbf{Κύρια εργασία}\tabularnewline
\midrule
\endfirsthead
\toprule
\textbf{ΠΕ} & \textbf{Μήνες} & \textbf{Κύρια εργασία}\tabularnewline
\midrule
\endhead
\midrule
\multicolumn{3}{r}{\footnotesize Συνέχεια στην επόμενη σελίδα}\tabularnewline
\endfoot
\bottomrule
\endlastfoot

ΠΕ0 & 1--6 & Απογραφή σωμάτων, δικαιώματα, προέλευση, κριτήρια ένταξης, εκτελεστική κάλυψη και έναρξη δεοντολογικής διαδικασίας M3b/M3c\tabularnewline
\midrule
ΠΕ1 & 4--12 & Ενιαίο συμβολικό επίπεδο αναφοράς, μοντέλο δεδομένων, πρώτη λειτουργική αλυσίδα EISE, υλοποίηση της μορφής εξόδου που έχει οριστικοποιηθεί στο ΠΕ0 και οριστικοποίηση του πρωτοκόλλου M1\tabularnewline
\midrule
ΠΕ2 / M1 & 10--24 & Σύγκριση αναπαραστάσεων, Transformer σε επίπεδο μελωδικού γεγονότος, επαγόμενη τμηματοποίηση και επιλογή της κύριας κλίμακας αναπαράστασης\tabularnewline
\midrule
ΠΕ3 / M3a & 12--24 & Ανεξάρτητη μουσικολογική επικύρωση, κρίσεις ειδικών και συγκρότηση του επικυρωμένου λεξιλογίου οικογενειών για M3b/M3c\tabularnewline
\midrule
M3c μοντέλο & 20--28 & Πιλοτικό PPM με το λεξιλόγιο ανάπτυξης, τελικό PPM/$I_F$ μετά την οριστικοποίηση του επικυρωμένου λεξιλογίου οικογενειών, τράπεζα ερεθισμάτων και κριτήρια επάρκειας του μοντέλου\tabularnewline
\midrule
ΠΕ3β / M2 & 24--36 & Τελικός περιορισμένος έλεγχος της τμηματικής διαδικασίας μετά την $R^\star$ και πλήρης αντιστοιχισμένη σύγκριση B2--EISE με $\Delta_{\mathrm{hist}}$\tabularnewline
\midrule
M3 πιλοτικοί έλεγχοι & 28--30 & Πιλοτικοί έλεγχοι M3c/M3b, μουσική φυσικότητα, αξιοπιστία, κοινό πεδίο τιμών, κάλυψη των αντιδιαγώνιων συνδυασμών και προσομοίωση στατιστικής ισχύος\tabularnewline
\midrule
M3 συλλογή ανθρώπινων δεδομένων & 30--40 & Συνεδρία 1: M3c και Συνεδρία 2: M3b στους ίδιους συμμετέχοντες, με κοινά μεταδεδομένα εμπειρίας και τρόπο συλλογής ορισμένο στον πιλοτικό έλεγχο\tabularnewline
\midrule
M4 & 30--42 & Αναγκαίο ανεξάρτητο κριτήριο αξιολόγησης έως την τελική αξιολόγηση M2 (μήνας 36)· συνέχεια ευρύτερων αναλύσεων σύμφωνα με GREEN/AMBER/RED\tabularnewline
\midrule
M5 / M3ν & 36--44 & Χαρακτηρισμός πλήρων έργων ως προς δηλωμένο περιβάλλον αναφοράς και εφαρμογή στον Νικόλαο· εφόσον ενεργοποιηθεί, νευρική επέκταση με τις ίδιες μεταβλητές πρόβλεψης\tabularnewline
\midrule
ΠΕ7 & 42--48 & Σύνθεση, δημοσιεύσεις, τεκμηρίωση και αναπαραγώγιμη διάθεση\tabularnewline
\end{xltabular}
}\par

Το αναγκαίο ανεξάρτητο εκτελεστικό κριτήριο αξιολόγησης της M2 τεκμηριώνεται πριν από την τελική αντιστοιχισμένη αξιολόγηση B2--EISE, μέσα στο προβλεπόμενο διάστημα της M4. Η κρίσιμη εξάρτηση για τις ανθρώπινες δοκιμασίες είναι η οριστικοποίηση του επικυρωμένου λεξιλογίου οικογενειών μέχρι τον μήνα 24· το λεξιλόγιο ανάπτυξης των M1--M2 έχει οριστικοποιηθεί νωρίτερα στο ΠΕ1. Το τελικό $I_F$ υπολογίζεται μετά από αυτό το σημείο, οι κύριοι πιλοτικοί έλεγχοι πραγματοποιούνται στους μήνες 28--30 και η ανθρώπινη συλλογή ξεκινά με σταθερή τράπεζα ερεθισμάτων και οριστικοποιημένες μεταβλητές πρόβλεψης.

\paragraph{Δημοσιευτικός σχεδιασμός.} Η σειρά των δημοσιεύσεων ακολουθεί την ιεραρχία των βαθμίδων και ξεκινά από ό,τι δεν εξαρτάται από την έκβαση των υποθέσεων. Πρώτη στοχεύεται η δημοσίευση του σώματος αξιολόγησης σε επίπεδο θέσης, με το μητρώο προέλευσης και τους συγκριτές (ISMIR ή TISMIR). Ακολουθεί η μελέτη των τεκμηριωμένων μετασχηματισμών με την πρόβλεψη σε αποκρυμμένες οικογένειες. Η μελέτη αυτή έχει μικρό υπολογιστικό κόστος και μπορεί να ολοκληρωθεί με περιορισμένο αριθμό ζευγών, πριν από την πλήρη M2. Η σύγκριση B2--EISE, μαζί με τις καμπύλες ως προς το πλήθος τεκμηριωμένων ζευγών και την εξωτερική επικύρωση, αποτελεί το κύριο άρθρο της διατριβής. Τα μουσικολογικά αποτελέσματα για τη θέση και τον Νικόλαο Πρωτοψάλτη Σμύρνης απευθύνονται σε μουσικολογικά περιοδικά.


\subsection{Μελλοντικές κατευθύνσεις}

Η οριστικοποιημένη μονάδα και τα εργαλεία σύγκρισης μπορούν να στηρίξουν μεταγενέστερη μελέτη νέων πληθυσμών ακροατών, περισσότερων τρόπων καταγραφής της φωνής και πρόσθετων ιστορικών σωμάτων. Κάθε επέκταση θα χρειαστεί δικό της σχεδιασμό, δεδομένα και δεοντολογική τεκμηρίωση· η παρούσα πρόταση συγκεντρώνει τον ερευνητικό της φόρτο στις πέντε βαθμίδες της §~\ref{sec:core-hierarchy}.

\subsection{Συμπέρασμα}

Η ερευνητική πρόταση αντιμετωπίζει ένα κοινό πρόβλημα της συμβολικής Ανάκτησης Μουσικής Πληροφορίας και της βυζαντινής μουσικολογίας: πώς διατηρείται και πώς μπορεί να αναπαρασταθεί η ταυτότητα μιας μουσικής μονάδας όταν οι πραγματώσεις της μεταβάλλονται ουσιαστικά. Η θέση παρέχει ιστορικά και αναλυτικά τεκμηριωμένο αντικείμενο για αυτό το ερώτημα, ενώ η διάκριση ανάμεσα σε επιφανειακή ομοιότητα, οικογενειακή ένταξη και συγκεκριμένη ιστορική σχέση επιτρέπει να μετατραπεί σε σειρά ελεγχόμενων συγκρίσεων.

Η M1 εξετάζει την αξία της μονάδας αναπαράστασης, η M2 την πληροφοριακή αξία της ιστορικής προέλευσης, η M3a τη μουσικολογική συνοχή των οικογενειών, η M3b την κατηγορική οργάνωσή τους στην ανθρώπινη ακρόαση και η M3c τη σχέση της οικογενειακής πρόβλεψης με τη μουσική προσδοκία, καθώς και τη μεταβολή αυτής της σχέσης υπό ένδειξη λατρευτικού πλαισίου. Η M4 μετρά την εκτελεστική παραλλαγή και παρέχει το ανεξάρτητο εκτελεστικό κριτήριο της M2 και η M5 αξιοποιεί το επικυρωμένο πλαίσιο για ιστορική ανάλυση της μελοποιητικής πρακτικής του Νικολάου Πρωτοψάλτη Σμύρνης.

Με αυτή τη διάρθρωση, η υπολογιστική σύγκριση και η μουσικολογική τεκμηρίωση απαντούν από κοινού στο ερώτημα της ταυτότητας της θέσης μέσα στη μεταβολή. Η ευρύτερη μεθοδολογική συμβολή αφορά τη μελέτη πολιτισμικά εδραιωμένων μουσικών μονάδων μεταβλητής έκτασης όταν η παραλλαγή αποτελεί οργανωμένο χαρακτηριστικό της μουσικής παράδοσης. Η εφαρμογή του πλαισίου σε άλλα ρεπερτόρια προϋποθέτει ανεξάρτητα τεκμήρια για τις μονάδες, τις σχέσεις συγγένειας και την ιστορική ή πολιτισμική τους οργάνωση.

\appendix
\newcommand{\LatinAppendixLetter}[1]{\ifcase#1\or A\or B\or C\or D\or E\or F\or G\or H\or I\or J\fi}
\renewcommand{\thesection}{\LatinAppendixLetter{\value{section}}}
\renewcommand{\thesubsection}{\thesection.\arabic{subsection}}
\renewcommand{\thesubsubsection}{\thesubsection.\arabic{subsubsection}}

\section{Χάρτης ενοτήτων, υποθέσεων και εκτιμώμενων μεγεθών}\label{app:estimand-map}

Ο χάρτης αντιστοιχίζει κάθε υπόθεση και εκτιμώμενο μέγεθος στην ενότητα M1--M5, στη βαθμίδα τεκμηρίωσης και στο καθεστώς ανάλυσής του. Οι ετικέτες Υ χρησιμοποιούνται για την αναγνώριση των αντίστοιχων στατιστικών ελέγχων και συνδέονται ρητά με τις μεθοδολογικές ενότητες.

{\footnotesize
\begin{xltabular}{0.98\linewidth}{@{}>{\RaggedRight\arraybackslash}p{0.13\linewidth}>{\RaggedRight\arraybackslash}p{0.075\linewidth}>{\RaggedRight\arraybackslash}p{0.115\linewidth}>{\RaggedRight\arraybackslash}p{0.21\linewidth}Y@{}}
\toprule
\textbf{Μέγεθος} & \textbf{Βαθμίδα} & \textbf{Ενότητα} & \textbf{Καθεστώς} & \textbf{Ρόλος}\tabularnewline
\midrule
\endfirsthead
\toprule
\textbf{Μέγεθος} & \textbf{Βαθμίδα} & \textbf{Ενότητα} & \textbf{Καθεστώς} & \textbf{Ρόλος}\tabularnewline
\midrule
\endhead
\midrule
\multicolumn{5}{r}{\footnotesize Συνέχεια στην επόμενη σελίδα}\tabularnewline
\endfoot
\bottomrule
\endlastfoot

$R^\star$ & 1 & M1 & επιλογή στην ανάπτυξη & αναπαράσταση με μέγιστο $F_{1,\mathrm{seg+fam}}$ στην A και οικογενειακό κριτήριο $g_{\min}$\tabularnewline
$\Delta_{\mathrm{hist}}$ & 1 & M2 & κύριο· επιβεβαιωτικό σε GREEN/AMBER & B2--EISE στο ανεξάρτητο εκτελεστικό κριτήριο του ελάχιστου μέρους της M4, σε κοινό εξωτερικά τεκμηριωμένο μητρώο εμφανίσεων\tabularnewline
$\Delta_{\mathrm{learn}}$ & 1 & M2 & δεύτερη κύρια αντίθεση EISE & EISE--B0 στο ίδιο ανεξάρτητο εκτελεστικό κριτήριο\tabularnewline
$H_{\mathrm{pred}}$ & 2 & M3c & ανθρώπινη επιβεβαιωτική & πληροφοριακό περιεχόμενο οικογένειας πέρα από το γεγονός\tabularnewline
$H_{\mathrm{cat}}$ & 2 & M3b & ανθρώπινη επιβεβαιωτική & κατηγορική οργάνωση στην ακρόαση\tabularnewline
$\Delta_{\mathrm{ctx}}$ & 4 & M3c & κύριο δευτερεύον & διαφορά συντελεστών αλληλεπίδρασης με την ένδειξη $R$\tabularnewline
M3ν & 5 & κλάδος M3c & νευρική επέκταση υπό προϋποθέσεις & πρόσθετη συμβολή $I_F$ σε EEG/MEG· χωριστό $R\times I_F$\tabularnewline
EGG & 5 & M4 & διερευνητικό υπό προϋποθέσεις & δείκτες φωνητικής επαφής κατά θέση και εσωτερική περιοχή\tabularnewline
Χάρτης έργου / κάλυψη & 3 & M5a/M5b & υποχρεωτικό παραδοτέο· M5-min & χάρτης εμφανίσεων, σχέσεων σειράς, αποχών, συνδετικών ή ανεπίλυτων περιοχών και προέλευσης, μαζί με την πραγματική κάλυψη του έργου\tabularnewline
$D_{\mathrm{fam}}$ & 3 & M5a/M5b & προκαθορισμένο περιγραφικό· M5-min & διαφοροποίηση της κατανομής οικογενειών ως προς $D_{\Omega}$\tabularnewline
$D_{\mathrm{within}}$ & 3 & M5a/M5b & προκαθορισμένο περιγραφικό· M5-min & ενδοοικογενειακή απόσταση στις τοπικές καταστάσεις EISE που διατηρούν την παραλλαγή\tabularnewline
$D_{\mathrm{seq}}$ & 4 & M5a/M5b & προκαθορισμένο περιγραφικό, εφόσον υπάρχει επαρκής κάλυψη των ακολουθιών & διαφορά του μέσου πληροφοριακού περιεχομένου από την αναφορά, επανεμφάνιση και ανασυνδυασμός ακολουθιών οικογενειών\tabularnewline
$D_{\mathrm{func}}$ & 4 & M5a/M5b & προκαθορισμένο περιγραφικό, εφόσον υπάρχει επαρκής κάλυψη τόπων & διαφοροποίηση $P(F\mid L,M,G)$ ως προς το σώμα αναφοράς\tabularnewline
$I_{\Omega}$ & 4 & M5a & προκαθορισμένο περιγραφικό & πληροφοριακό περιεχόμενο ως προς το τεκμηριωμένο περιβάλλον $\Omega$\tabularnewline
Προφίλ μουσικής έμφασης & 4 & §~\ref{sec:text-locus}/M5 & προκαθορισμένη δευτερεύουσα & αντιστοιχισμένη περιγραφή κοινών κειμενικών χωρίων μεταξύ μελοποιών· κατηγορική σύνοψη μόνο όπου το επιτρέπουν η κάλυψη και ο διαχωρισμός από τον δομικό τόπο\tabularnewline
\midrule
Υ2\textsubscript{E} & 4 & M1/M3a & προκαθορισμένη δευτερεύουσα & αναλλοιότητα ταυτότητας σε εξηγητικές σχέσεις\tabularnewline
Υ3 & 4 & M3a & προκαθορισμένη δευτερεύουσα & πληροφοριακή τροχιά της θέσης\tabularnewline
Υ5 & 1/4 & M4/M2 & στόχος M2· ευρύτερη δευτερεύουσα & ελάχιστη εξωτερική τεκμηρίωση και ευρύτερη παραλλαγή\tabularnewline
Υ8 & 4 & M1/M3a & δευτερεύουσα / διερευνητική & διακρίνουσα εγκυρότητα\tabularnewline
\midrule
Άξονες Α--Ε & παλαιά αντιστ. & M1--M5 & ιστορικές ετικέτες Παραρτημάτων & Α: M1/M3a· Β: M3a· Γ: M4/M2· Δ: M3a· Ε: σύνθεση\tabularnewline
\end{xltabular}
}\par

\section{Αναλυτική προδιαγραφή αναπαραστάσεων}\label{app:repr}

\subsection{Γεγονός αναφοράς και γραφικά σημεία}\label{app:repr-canonical}

Το παρόν Παράρτημα χρησιμοποιεί τη διάκριση γραφικού σημείου--ατομικού γεγονότος και την κοινή σημειογραφική συντεταγμένη της §~\ref{sec:m1}. Εξειδικεύονται εδώ οι απαιτήσεις αμφίδρομης αντιστοίχισης, προέλευσης και ελέγχου ανακατασκευής.

Τα πεδία του ενιαίου επιπέδου αναφοράς διακρίνονται ως προς την προέλευσή τους. \textbf{Παρατηρούμενο πεδίο} είναι πληροφορία που διαβάζεται άμεσα από τη συγκεκριμένη πηγή, ενώ \textbf{παράγωγο πεδίο} υπολογίζεται από παρατηρούμενες πληροφορίες με δηλωμένο και αναπαραγώγιμο κανόνα. Για κάθε γεγονός διατηρούνται αναγνωριστικά της πηγής και των γραφικών σημείων από τα οποία προέκυψε, καθώς και ο κανόνας παραγωγής κάθε παράγωγης ποσότητας. Ο έλεγχος ανακατασκευής εφαρμόζει την αντίστροφη διαδρομή στα πεδία που έχουν δηλωθεί ως ανακτήσιμα και ελέγχει ότι η κανονικοποίηση δεν έχει εισαγάγει μη καταγεγραμμένη απώλεια ή μεταβολή της πληροφορίας.

\subsection{\texorpdfstring{Μονάδες και πεδία \(N/T/M/L\)}{Μονάδες και πεδία N/T/M/L}}\label{app:repr-fields}

Το πλέγμα $U\times F$, οι μονάδες $u_1$--$u_3$ και οι τέσσερις διαδοχικά εμπλουτισμένες διαμορφώσεις γνωρισμάτων ορίζονται στην §~\ref{sec:m1} και δεν επαναλαμβάνονται εδώ. Το παρόν παράρτημα εξειδικεύει μόνο την κωδικοποίηση των πεδίων, τις καταστάσεις των πεδίων ως παρατηρημένων, μη εφαρμοστέων ή άγνωστων και τους ελέγχους ανακατασκευής που απαιτούνται για την υλοποίηση.

Για κάθε χωρίο αποθηκεύεται χωριστά το λεκτικό καθεστώς \texttt{lexical}, \texttt{nonlexical\_vocables}, \texttt{mixed} ή \texttt{unknown}, χωρίς να ταυτίζεται αυτομάτως με την ιστορική ονομασία του χωρίου. Κάθε πεδίο του $L$ φέρει κατάσταση \emph{παρατηρημένο}, \emph{μη εφαρμοστέο} ή \emph{άγνωστο/ελλιπές}. Στα λεξικά χωρία καταγράφονται χωριστά η συλλαβική μονάδα, η λέξη, το κοινό κειμενικό χωρίο μέσα στο λειτουργικό κείμενο, ο λεκτικός τόνος όπου είναι μουσικά ερμηνεύσιμος, τα όρια λέξης/φράσης και η μελισματική πυκνότητα. Για κάθε μελωδικό γεγονός με λεκτικό υπόστρωμα διατηρείται επιπλέον η σύνδεση με τη συλλαβή του σταθερού κειμένου, το αναγνωριστικό του κειμενικού χωρίου και η προέλευση ή κατάσταση της αντιστοίχισης. Οι μελοποιητικές επαναλήψεις συνδέονται με την ίδια κανονική συλλαβή, ενώ οι επιμέρους εμφανίσεις των γεγονότων παραμένουν διακριτές. Άγνωστη αντιστοίχιση διατηρεί ρητή άγνωστη τιμή και αιτία. Στα μη λεξικά χωρία η συλλαβική άρθρωση και η σχέση γεγονότων προς τεκμηριωμένες φωνητικές μονάδες μπορούν να διατηρούνται, ενώ η αναφορά σε κανονική συλλαβή λειτουργικού κειμένου, η λέξη, το κειμενικό χωρίο και ο λεκτικός τόνος είναι μη εφαρμοστέα όταν δεν υπάρχει λεξικό υπόστρωμα. Παύσεις ή χρονικά γεγονότα χωρίς εκφωνούμενη συλλαβή ακολουθούν δηλωμένο κανόνα απόδοσης και δεν αποκτούν τεχνητή κειμενική συλλαβή. Η μελισματική πυκνότητα υπολογίζεται μόνο όταν ο παρονομαστής είναι έγκυρος. Η ένδειξη «μη εφαρμοστέο» μπορεί από μόνη της να αποκαλύπτει το λεκτικό καθεστώς· συνεπώς το όφελος του $L$ ερμηνεύεται ως όφελος του πλήρους συνόλου γνωρισμάτων $L$ και όχι ως καθαρή σημασιολογική ή λεξική επίδραση. Το Αλληλούια καταγράφεται ως \texttt{lexical} με χωριστή ένδειξη λειτουργικής επευφημίας· ο τύπος μουσικού τμήματος καταγράφεται ανεξάρτητα από το λεκτικό καθεστώς. Τα αναγνωριστικά προέλευσης αποτελούν υποδομή ελέγχου και η χρήση τους ως μαθησιακής πληροφορίας ακολουθεί αποκλειστικά τις υπάρχουσες διαμορφώσεις $N/T/M/L$ και τους ελέγχους της M1 για έμμεση επανεισαγωγή αφαιρεμένων πεδίων. Τα μουσικά όρια θέσης παραμένουν ανεξάρτητα από τα κειμενικά όρια· η σύμπτωσή τους με αλλαγή συλλαβής, λέξης ή φράσης αναφέρεται ως δευτερεύον διαγνωστικό. Αν απαιτηθεί περιορισμός ως προς το ποια χωρία μπορούν να ενταχθούν στην ανάλυση, εφαρμόζεται στις ίδιες εμφανίσεις για όλες τις συνθήκες της συγκεκριμένης σύγκρισης.



\subsection{Επαγόμενη τμηματοποίηση και οικογένειες}\label{app:repr-induced}

Το διασταυρωμένο πλέγμα \(U\times F\) διατηρεί ως παράγοντα μονάδας τις τρεις \textbf{θεωρητικά προκαθορισμένες δομικές κλίμακες} \(U=\{u_1,u_2,u_3\}\), οι οποίες συγκροτούν τη συμμετρική βάση για την εκτίμηση κύριων επιδράσεων και αλληλεπιδράσεων. Παράλληλα εισάγεται μία μεθοδολογικά διαφορετική συνθήκη: \textbf{μη επιβλεπόμενη, αλγοριθμικά επαγόμενη τμηματοποίηση μεταβλητού μήκους}, η οποία προκύπτει από τη στατιστική δομή του σώματος χωρίς ετικέτες θέσεων. Η συνθήκη αυτή λειτουργεί ως \textbf{εναλλακτική συνθήκη αναφοράς} απέναντι στις θεωρητικά προκαθορισμένες μονάδες και διευρύνει το ερώτημα από την επιλογή μεταξύ προκαθορισμένων κλιμάκων στην αντιπαραβολή θεωρητικά προκαθορισμένης και επαγόμενης από τα δεδομένα δομικής οργάνωσης.

Η κύρια υλοποίηση βασίζεται σε μη παραμετρική πιθανοτική τμηματοποίηση της ίδιας μεθοδολογικής οικογένειας με τις πρόσφατες εφαρμογές \textbf{ιεραρχικά εμφωλευμένων μοντέλων Pitman--Yor} στο γρηγοριανό μέλος. Η ακριβής οικογένεια μοντέλου οριστικοποιείται στο ΠΕ0 μετά από πιλοτική επιβεβαίωση ότι μπορεί να εφαρμοστεί στο ενιαίο επίπεδο αναφοράς της βυζαντινής μουσικής. Η αλγοριθμικά επαγόμενη τμηματοποίηση εκπαιδεύεται αποκλειστικά στο σώμα εκπαίδευσης και αξιολογείται στα ίδια έργα που έχουν διατηρηθεί εκτός εκπαίδευσης και στις ίδιες ομάδες συγγενών πηγών με τις υπόλοιπες συνθήκες.

Τα επαγόμενα τμήματα οργανώνονται επιπλέον σε \textbf{αλγοριθμικά επαγόμενο λεξιλόγιο οικογενειών}. Η ομαδοποίηση πραγματοποιείται χωρίς ετικέτες οικογένειας, πάνω στην ίδια γεωμετρία αναπαράστασης και με τα ίδια μέτρα ομοιότητας που χρησιμοποιεί η μουσικολογική δοκιμή. Κάθε επαγόμενη συστάδα αποκτά αναγνωριστικό, προέλευση και αντιπροσωπευτικό μέλος της συστάδας (medoid) και αντιστοιχίζεται στις τεκμηριωμένες οικογένειες μέσω του ήδη ορισμένου κλειστού συνόλου σχέσεων \emph{ισοδύναμη, ευρύτερη, στενότερη, επικαλυπτόμενη, ανεπίλυτη}. Η αξιολόγηση σε επίπεδο οικογένειας είναι διπλή: συγκρίνει το επαγόμενο λεξιλόγιο με το \textbf{επαληθευμένο οικογενειακό υποσύνολο της καταλογογραφικής υποδομής της Ιεράς Μονής Αγίου Αντωνίου} και, χωριστά, με τις οικογένειες που τεκμηριώνονται από δημοσιευμένες αναλύσεις συγκεκριμένων εμφανίσεων. Η αξιολόγηση αναφέρει \textbf{κάλυψη οικογενειών, ισχύ αντιστοίχισης, μη απόδοση οικογενειακής ετικέτας σε περιπτώσεις εκτός του γνωστού λεξιλογίου και δυνατότητα ανάκτησης του λειτουργικού ρόλου}. Η εναλλακτική συνθήκη κρίνεται τόσο ως προς τα όρια όσο και ως προς την ικανότητά της να ανασυγκροτεί την οικογενειακή οργάνωση της παράδοσης.

Η ισοτιμία ως προς την πληροφορία διατηρείται με αξιολόγηση της επαγόμενης τμηματοποίησης στα ίδια τέσσερα επίπεδα γνωρισμάτων
\[
F=\{N,\ N{+}T,\ N{+}T{+}M,\ N{+}T{+}M{+}L\}.
\]
Η πιθανή υπεροχή της δεν συγχέεται με πρόσθετη πρόσβαση σε τροπικά, χρονικά ή γλωσσικά γνωρίσματα. Στη μουσικολογική, την τροπική και την εκτελεστική δοκιμή χρησιμοποιούνται οι οριστικοποιημένες οριοθετήσεις που παράγονται στο αντίστοιχο πρωτόκολλο εκπαίδευσης και τελικής αξιολόγησης. Στην πληροφοριακή δοκιμή, όπου η πρόβλεψη χρησιμοποιεί αποκλειστικά προγενέστερο συμφραζόμενο, εφαρμόζεται αυστηρότερος κανόνας: το συμφραζόμενο της χρονικής στιγμής \(t\) μπορεί να χρησιμοποιεί μόνο τμηματοποίηση που έχει εξαχθεί από το πρόθεμα \(e_{\leq t-1}\). Όταν η επιλεγμένη μη παραμετρική μέθοδος απαιτεί πρόσβαση σε ολόκληρη την ακολουθία για την εξαγωγή της τμηματοποίησης, η πλήρης τμηματοποίησή της χρησιμοποιείται ως \emph{προκαθορισμένη δομική προεπεξεργασία πλήρους ακολουθίας} σε χωριστή διαγνωστική ανάλυση. Η κύρια ανάλυση της πληροφοριακής δοκιμής παραμένει αυστηρά περιορισμένη σε προηγούμενο συμφραζόμενο.

Η σύγκριση αυτή δημιουργεί ένα από τα κρισιμότερα πειράματα της ερευνητικής πρότασης: \textbf{αν μια μονάδα που βελτιστοποιείται από τα δεδομένα υπερέχει προβλεπτικά, διατηρεί επίσης τη μουσικολογική ταυτότητα, την τροπική λειτουργία και το εκτελεστικό αποτύπωμα της θέσης;} Η απάντηση διαχωρίζει την υπολογιστική αποτελεσματικότητα από την εγκυρότητα της μουσικής μονάδας και επιτρέπει άμεση σύγκριση της βυζαντινής περίπτωσης με το σύγχρονο πρόβλημα της μη επιβλεπόμενης τμηματοποίησης στο γρηγοριανό μέλος.

\subsection{Κοινός στόχος και χρονικό συμφραζόμενο}\label{app:repr-common-target}

Επειδή οι υποψήφιες αναπαραστάσεις παράγουν \emph{διαφορετικό αριθμό} συμβολικών μονάδων, η αρνητική λογαριθμική πιθανοφάνεια ανά μονάδα \textbf{δεν είναι συγκρίσιμη}: θα ευνοούσε αυτόματα τη συμπτυγμένη κωδικοποίηση. Η κοινή μονάδα μέτρησης συνδυάζεται με κοινό στόχο πρόβλεψης. Οι συνθήκες γνωρισμάτων \(N\), \(N{+}T\), \(N{+}T{+}M\), \(N{+}T{+}M{+}L\) προβλέπουν, στη φυσική τους μορφή, \emph{διαφορετικά εμπλουτισμένα σύμβολα}: μια πλουσιότερη συνθήκη κωδικοποιεί περισσότερη πληροφορία και επομένως εμφανίζει μεγαλύτερο μήκος κώδικα \emph{εξ ορισμού}. Σύγκριση της μορφής \(-\log P(z^R)/N_{\text{γεγονότων}}\) θα τιμωρούσε συστηματικά τις πλουσιότερες αναπαραστάσεις για λόγο άσχετο με την προβλεπτική τους αξία.

\textbf{Επιβάλλεται επομένως κοινός στόχος πρόβλεψης.} Ορίζεται ένας προκαθορισμένος και οριστικοποιημένος στόχος \(Y_t\) πάνω στο κοινό πλέγμα ατομικών γεγονότων --- η ελάχιστη μελωδική περιγραφή του επόμενου γεγονότος, κοινή για όλες τις συνθήκες --- και η αναπαράσταση εισέρχεται \emph{αποκλειστικά ως συμφραζόμενο}. Η κοινή μετρική είναι η \textbf{διασταυρούμενη εντροπία ανά ατομικό γεγονός}:
\[
\mathrm{CE}^{(\mathrm{event})}_R=\frac{\sum_t -\log_2 P\bigl(Y_t \mid C_R(t)\bigr)}{N_{\text{ατομικών γεγονότων}}}
\]
όπου \(C_R(t)\) το συμφραζόμενο που παρέχει η συνθήκη \(R\) στη χρονική στιγμή \(t\). Όλες οι συνθήκες προβλέπουν \emph{το ίδιο πράγμα} και διαφέρουν μόνο ως προς την πληροφορία που διαθέτουν για να το προβλέψουν.

Από την ίδια οριστικοποιημένη προβλεπτική κατανομή εξάγονται δύο διακριτές πληροφοριακές ποσότητες:
\[
h_R(t)=-\log_2 P\bigl(Y_t\mid C_R(t)\bigr),
\]
\[
\mathcal{H}_R(t)=-\sum_y P\bigl(y\mid C_R(t)\bigr)\log_2 P\bigl(y\mid C_R(t)\bigr).
\]
Το \(h_R(t)\) είναι το \emph{πληροφοριακό περιεχόμενο} του γεγονότος που πράγματι συνέβη —δηλαδή πόσο μη αναμενόμενο ήταν το παρατηρημένο γεγονός για το μοντέλο—, ενώ το \(\mathcal{H}_R(t)\) είναι η \emph{εντροπία της προβλεπτικής κατανομής}, δηλαδή η αβεβαιότητα του μοντέλου \emph{πριν} από την παρατήρηση του επόμενου γεγονότος. Η Υ3 εστιάζει πρωτίστως στην τροχιά της \(\mathcal{H}_R\), ενώ το \(h_R\) χρησιμοποιείται τόσο ως συμπληρωματικό πληροφοριακό προφίλ όσο και για τον ρητό έλεγχο αιχμών ορίου.

\subsection{Χρήση μόνο προηγούμενου συμφραζομένου}\label{app:repr-past-only}

Το \(C_R(t)\) ορίζεται με δύο κοινούς κανόνες πρόσβασης στην πληροφορία για όλες τις συγκρινόμενες αναπαραστάσεις.

\textbf{Χρήση μόνο προηγούμενης πληροφορίας.} Στην πρόβλεψη του \(Y_t\) χρησιμοποιείται αποκλειστικά πληροφορία από γεγονότα προγενέστερα του \(t\):
\[
C_R(t)\cap\{e_t,e_{t+1},\dots\}=\varnothing
\]
Καμία σύνθετη μονάδα της οποίας ο ορισμός ολοκληρώνεται στο \(t\) ή αργότερα, και κανένα γνώρισμα που απαιτεί μελλοντικό γεγονός για να υπολογιστεί, δεν εισέρχεται στο συμφραζόμενο.

\textbf{Κοινός ορίζοντας σε ατομικά γεγονότα.} Το διαθέσιμο ιστορικό ορίζεται σε κοινή έκταση \(H\) του πλέγματος ατομικών γεγονότων και \emph{όχι} σε σταθερό αριθμό μονάδων κάθε αναπαράστασης:
\[
C_R^{(H)}(t)=R\bigl(e_{t-H},\dots,e_{t-1}\bigr)
\]
Ο περιορισμός αυτός είναι ουσιώδης. Αν το ιστορικό μετρούνταν σε μονάδες, μια συμπτυγμένη αναπαράσταση θα έβλεπε πολύ μεγαλύτερο μουσικό παρελθόν με τον ίδιο αριθμό μονάδων και θα εμφανιζόταν ανώτερη για λόγο άσχετο με την ποιότητά της --- απλώς επειδή συμπιέζει. Με κοινό \(H\) όλες οι συνθήκες βλέπουν \emph{την ίδια μουσική ιστορία} και διαφέρουν μόνο ως προς το πώς την οργανώνουν.

\textbf{Ο περιορισμός δεσμεύει και την ενεργό μνήμη του εκτιμητή}, όχι μόνο το διαθέσιμο ιστορικό. Σε μοντέλο \(n\)-γραμμάτων σταθερής τάξης, οι τελευταίες \(n-1\) μονάδες μιας συμπτυγμένης αναπαράστασης καλύπτουν μεγαλύτερο ατομικό εύρος από τις αντίστοιχες μιας λεπτομερούς --- και το πλεονέκτημα της συμπίεσης θα επανεισαγόταν πλαγίως ως συγχυτικός παράγοντας. Το συμφραζόμενο ορίζεται επομένως ως η μέγιστη επιτρεπτή ακολουθία \emph{ολοκληρωμένων} μονάδων της \(R\) των οποίων η συνολική κάλυψη βρίσκεται αποκλειστικά εντός του διαστήματος \(e_{t-H},\dots,e_{t-1}\). Καμία οικογένεια μοντέλων δεν επιτρέπεται να χρησιμοποιεί πληροφορία εκτός του κοινού ατομικού ορίζοντα.

Το \(H\) οριστικοποιείται στο πρώτο πακέτο εργασίας μαζί με το \(Y_t\) και συνοδεύεται από προκαθορισμένη ανάλυση ευαισθησίας.

Ο ορισμός του \(Y_t\) οριστικοποιείται στο πρώτο πακέτο εργασίας μαζί με τα υπόλοιπα πεδία και δεν εξαρτάται από καμία υποψήφια αναπαράσταση. Χωρίς αυτή τη δέσμευση το πληροφοριακό πείραμα δεν θα ήταν έγκυρο: θα μετρούσε πόσα κωδικοποιεί κάθε αναπαράσταση, όχι πόσο καλά προβλέπει.

Για το \emph{προφίλ εντός του τμήματος}, όλες οι ποσότητες προβάλλονται στα ίδια σταθερά παράθυρα της κοινής συντεταγμένης ατομικών γεγονότων. Ορίζονται
\[
I_R(W)=\frac{1}{|W|}\sum_{t\in W} h_R(t),
\qquad
\mathcal{E}_R(W)=\frac{1}{|W|}\sum_{t\in W}\mathcal{H}_R(t).
\]
Κάθε τιμή αντιστοιχεί εξ ορισμού σε ατομικό γεγονός και επομένως οι διαφορετικές αναπαραστάσεις συγκρίνονται χωρίς αυθαίρετη κατανομή πληροφορίας από σύνθετες υπολογιστικές μονάδες προς τα γεγονότα που καλύπτουν. Το \(I_R\) αποτυπώνει το μέσο πληροφοριακό περιεχόμενο των γεγονότων που πράγματι παρατηρήθηκαν και το \(\mathcal{E}_R\) τη μέση αβεβαιότητα της προβλεπτικής κατανομής πριν από την παρατήρηση του επόμενου γεγονότος.

\textbf{Το πλάτος του παραθύρου οριστικοποιείται στο ΠΕ0} με βάση τη μετρημένη ακρίβεια ευθυγράμμισης και την εμπειρική έκταση των τεκμηριωμένων θέσεων. Προδηλώνεται εύρος ευαισθησίας γύρω από την πρωτεύουσα τιμή, ώστε το συμπέρασμα να ελέγχεται σε περισσότερες χωρικές αναλύσεις χωρίς εκ των υστέρων επιλογή παραθύρου.

\subsection{Ορισμός σύνθετης νευματικής μονάδας}\label{app:repr-compound}

Η μονάδα \(u_3\) ορίζεται ως \(u_3=f_{\mathrm{group}}(N)\): \textbf{ένας ντετερμινιστικός, εκδοτικά καταγεγραμμένος και αντιστρέψιμος κανόνας ομαδοποίησης}, που χρησιμοποιεί μόνο τα σημειογραφικά πεδία του \(N\) και τις αναγκαίες σχέσεις προέλευσής τους και είναι κοινός σε όλα τα επίπεδα του \(F\). Τα \(T\), \(M\) και \(L\) προστίθενται ως γνωρίσματα πάνω στην ίδια ομαδοποίηση, ώστε η μετάβαση μεταξύ επιπέδων του \(F\) να μεταβάλλει μόνο τη διαθέσιμη πληροφορία και όχι τις ίδιες τις μονάδες. Ο κανόνας είναι ανεξάρτητος από ετικέτες θέσεων και οικογενειών. Διατηρείται ελέγξιμη αντιστοίχιση κάθε μονάδας με τα ατομικά γεγονότα που καλύπτει· η προβολή στην κοινή ατομική κλίμακα επιτρέπει υποψήφιο όριο στο εσωτερικό της μονάδας και ελέγχεται ότι δεν επανεισάγει πληροφορία που έχει αφαιρεθεί από τη συνθήκη. Αν ο έλεγχος αντιστρεψιμότητας αποτύχει, η μονάδα αφαιρείται από το πλέγμα \emph{πριν} από κάθε επιβεβαιωτική ανάλυση.

\subsection{Έλεγχος γεωμετρίας αξιολόγησης}\label{app:repr-geometry}

Η κοινή γεωμετρία αξιολόγησης είναι το επίπεδο των ατομικών γεγονότων και επομένως συμπίπτει με την υποψήφια μονάδα \(u_2\). Η σύμπτωση αυτή μετατρέπεται σε \textbf{ελέγξιμη πρόβλεψη} και όχι σε άρρητη σχεδιαστική παραδοχή: πραγματικό πλεονέκτημα της \(u_2\) πρέπει να διατηρείται όταν η αξιολόγηση προβάλλεται και σε δύο εξωτερικές συντεταγμένες, τη συντεταγμένη πηγαίου σημείου και τη συντεταγμένη συλλαβής.

Εκτελείται επομένως \textbf{τριπλός έλεγχος ευαισθησίας προβολής}. Σταθερή υπεροχή της \(u_2\) και στις τρεις γεωμετρίες αποτελεί ισχυρή ένδειξη ότι το ατομικό γεγονός είναι καταλληλότερη υπολογιστική κλίμακα για τις προκαθορισμένες εργασίες και ότι το πλεονέκτημά του δεν οφείλεται απλώς στην κοινή συντεταγμένη αξιολόγησης. Μεταβολή της κατάταξης χαρτογραφεί, αντίστοιχα, ποια πτυχή της μουσικής δομής εξαρτάται από τη συγκεκριμένη προβολή και προσθέτει ερμηνευτική πληροφορία για τη σχέση σημείου, γεγονότος και συλλαβής.



\section{Απογραφή σωμάτων και κριτήρια ένταξης}\label{app:corpora}

\subsection{Καταλογογραφική βάση και τεκμηρίωση πηγών της Ιεράς Μονής Αγίου Αντωνίου}\label{sec:sta}\label{app:corpora-sta}

Η παπαδική συλλογή της Ιεράς Μονής Αγίου Αντωνίου συγκροτεί την \textbf{κύρια καταλογογραφική βάση και υποδομή τεκμηρίωσης των πηγών} της ερευνητικής πρότασης, σε συνδυασμό με τα πλήρη έργα, τις δημοσιευμένες αναλύσεις και την ανεξάρτητη επισημείωση. Η δομή της αποτυπώνεται ως
\[
\mathcal S_{\mathrm{StA}}=
\bigl(C_{\mathrm{StA}},O_{\mathrm{StA}},P_{\mathrm{StA}}\bigr),
\]
όπου $C_{\mathrm{StA}}$ είναι το καταλογογραφικό επίπεδο \textbf{υποψηφίων τύπων, κατηγοριών και ρητά επισημασμένων παραλλαγών}, $O_{\mathrm{StA}}$ το επίπεδο των επιμέρους καταγεγραμμένων εμφανίσεων/φορμουλών και $P_{\mathrm{StA}}$ το επίπεδο προέλευσης που συνδέει κάθε εγγραφή με τίτλο και καταλογογραφική απόδοση προσώπου και, μέσω της βιβλιογραφικής τεκμηρίωσης της συλλογής, με την αντίστοιχη έκδοση ή χειρόγραφη πηγή. Το $P_{\mathrm{StA}}$ κανονικοποιεί χωριστά \emph{αποδιδόμενο μελοποιό}, \emph{εξηγητή/μεταγραφέα}, \emph{προσαρμογέα}, \emph{έργο} και \emph{πηγή/μάρτυρα}. Η τριάδα αυτή παρέχει ρητή διαδρομή προέλευσης και επιτρέπει τη μετάβαση
\[
\begin{aligned}
\text{καταγεγραμμένη φόρμουλα}
&\rightarrow \text{υποψήφιος τύπος/παραλλαγή}\\
&\rightarrow \text{έργο και ρόλος προσώπου}
\rightarrow \text{πηγή/μάρτυρας}.
\end{aligned}
\]
Η ιδιότητα κοινής οικογένειας αποτελεί χωριστή μουσικολογική απόφαση. Η καταλογογραφική κατηγορία μετατρέπεται σε οικογενειακή ετικέτα μόνο μετά την επαλήθευση του ΠΕ0 και τον ανεξάρτητο έλεγχο της §\ref{app:corpora-supervision}. Η κάλυψη του καταλόγου αναφέρεται υποχρεωτικά ανά λειτουργικό ρόλο, με χωριστή αποτίμηση καταληκτικών και μη καταληκτικών θέσεων.

Η υποδομή τροφοδοτεί τέσσερις λειτουργίες του ίδιου ερευνητικού σχεδιασμού. Πρώτον, το επιλέξιμο τμήμα της συλλογής παρέχει μεγάλο συμβολικό σώμα για προεκπαίδευση. Δεύτερον, η οικογενειακή ένταξη παράγεται με προκαθορισμένη και εκδοχικά ελεγχόμενη διαδικασία. Για κάθε εμφάνιση $x$, υποψήφιο τύπο $T$ και κριτή $r$ καταγράφεται
\[
J(x,T;r)\in\{\mathrm{member},\mathrm{non\_member},\mathrm{unresolved}\}.
\]
Η κρίση \texttt{member} βασίζεται σε προδηλωμένη, επαρκή θετική τεκμηρίωση σύμφωνα με την παγωμένη υπογραφή. Σε οικογένειες με περισσότερες μουσικολογικά επικυρωμένες επιτρεπτές διαγνωστικές διαμορφώσεις, αρκεί να πληρούται πλήρως μία ενεργή διαμόρφωση, μαζί με τους κανόνες περιβάλλοντος που της αντιστοιχούν. Τα \texttt{required} είναι αναγκαία \emph{μέσα στην ενεργή διαμόρφωση και το δηλωμένο πεδίο της}, όχι καθολικά αναλλοίωτα όλων των εμφανίσεων. Τα \texttt{supportive} ενισχύουν την απόφαση χωρίς να αποτελούν μόνα τους επαρκή τεκμηρίωση. Η κρίση \texttt{non\_member} απαιτεί ενεργοποίηση ρητού κανόνα \texttt{exclusion} ή παρατηρήσιμη ασυμβατότητα προς \emph{όλες τις επιλέξιμες} διαγνωστικές διαμορφώσεις σύμφωνα με προκαθορισμένο κανόνα αποκλεισμού. Η απουσία υποστηρικτικού γνωρίσματος, η χαμηλή ομοιότητα και η απλή έλλειψη θετικής απόδειξης δεν αρκούν για αρνητική κρίση. Η κρίση \texttt{unresolved} χρησιμοποιείται όταν το τεκμήριο, η κλίμακα, τα όρια, η τροπική αβεβαιότητα ή η δομή της κατηγορίας δεν επιτρέπουν ασφαλή απόφαση. Εάν μια προτεινόμενη πολυθετική κατηγορία δεν επιτρέπει ακόμη αναπαραγώγιμο θετικό και αρνητικό κανόνα, παραμένει \texttt{candidate\_only} και δεν τροφοδοτεί το κύριο ground truth.

Κάθε τύπος περνά διαδοχικά από καταστάσεις \texttt{candidate}, \texttt{piloted}, \texttt{frozen} και, εφόσον περάσει τον ανεξάρτητο έλεγχο, \texttt{reliability\_passed}. Τύπος που δεν περνά παραμένει \texttt{candidate\_only} ή \texttt{singleton}. Κάθε αλλαγή κανόνα δημιουργεί νέα \texttt{type\_signature\_version}. Η κάρτα δεν θεωρεί εκ των προτέρων ότι όλες οι οικογένειες έχουν κλασική κατηγορική δομή με ένα κοινό αναγκαίο γνώρισμα. Αν ο ανεξάρτητος έλεγχος δείξει ότι αξιόπιστα \texttt{member} παραδείγματα παραβιάζουν γνώρισμα που είχε χαρακτηριστεί \texttt{required}, το αποτέλεσμα καταγράφεται ως εύρημα για τη δομή της κατηγορίας. Η υπογραφή επιστρέφει σε νέα έκδοση ανάπτυξης και δεν διορθώνεται πάνω στο reliability sample. Η αναθεωρημένη υπογραφή ελέγχεται σε νέο ανεξάρτητο reliability sample. Έτσι μπορεί να προκύψει περισσότερο πολυθετική δομή ή δομή οργανωμένη γύρω από χαρακτηριστικά παραδείγματα, χωρίς να προαποφασίζεται από το πρωτόκολλο. Μόνο σχέσεις κοινής οικογένειας από επιλέξιμες παγωμένες υπογραφές συγκροτούν το σύνολο
\[
E_{\mathrm{fam}}^{\mathrm{train}}
=\{(i,j): i,j\ \text{έχουν ανεξάρτητα επαληθευμένη κοινή οικογενειακή ταυτότητα}\}.
\]
Το σύνολο αυτό τροφοδοτεί την B2 και αποτελεί κοινή βάση επιλεξιμότητας για την EISE. Οι καταλογογραφικές επικεφαλίδες και οι λειτουργικές κατηγορίες δεν μετατρέπονται απευθείας σε θετικά ζεύγη. Τρίτον, η προέλευση ανά φόρμουλα οδηγεί στις πλήρεις συνθέσεις και στις ιστορικές πηγές. Οι πραγματικές ιστορικές σχέσεις αποθηκεύονται ως αυτοτελείς κατευθυνόμενες ακμές στο μητρώο $\mathcal H_{\mathrm{ex}}^{\mathrm{train}}$ με αναγνωριστικό, πηγή και κατεύθυνση. Όταν δύο κωδικοποιημένες αποδόσεις της Νέας Μεθόδου συνδέονται μέσω επαληθευμένων ιστορικών ακμών με το ίδιο τεκμηριωμένο χωρίο, παράγεται για εκπαιδευτικούς σκοπούς το ζεύγος
\[
E_{\mathrm{ex}}^{\mathrm{train}}\subset O_{\mathrm{StA}}\times O_{\mathrm{StA}},
\]
με ρητή καταγραφή των πεδίων \texttt{pair\_\allowbreak derivation}, \texttt{common\_\allowbreak locus\_\allowbreak id}, \texttt{historical\_\allowbreak edge\_\allowbreak ids} και \texttt{exemplar\_\allowbreak level}. Το $E_{\mathrm{ex}}$ είναι επομένως παράγωγο εκπαιδευτικό ζεύγος και δεν ταυτίζεται με το ίδιο το ιστορικό βέλος. Η παλαιότερη γραφή και η τεκμηρίωση από τις ιστορικές πηγές πιστοποιούν τις ιστορικές σχέσεις ως μεταδεδομένα προέλευσης και δεν εισέρχονται ως χωριστή υπολογιστική είσοδος. Για την κύρια M2 διακρίνονται επιπλέον: $E_{\mathrm{hom}}$, η επαληθευμένη αντιστοιχία του συγκεκριμένου ιστορικού/μουσικού χωρίου· $E_{\mathrm{coex}}$, η συνεξήγηση ή άλλη εξάρτηση που τεκμηριώνεται αποκλειστικά από απόδοση, επικεφαλίδα, γενεαλογία ή εξάρτηση μαρτύρων και \textbf{ποτέ} από μελωδική ομοιότητα· και $E_{\mathrm{copy}}$, η τεκμηριωμένη άμεση αντιγραφική ή παράγωγη εξάρτηση. Η ύπαρξη $E_{\mathrm{ex}}$ δεν υποκαθιστά ούτε την επαλήθευση $E_{\mathrm{hom}}$ ούτε την ανεξάρτητη οικογενειακή κρίση $E_{\mathrm{fam}}$. Τέταρτον, οι ίδιες οικογένειες αναζητούνται στα έργα του ΔΑΜΑΣΚΗΝΟΣ και στο σώμα έργων του Νικολάου, δημιουργώντας τον χάρτη επικάλυψης που συνδέει το μεγάλο ιστορικό σώμα με την εκτελεστική και ιστορική και αναλυτική εφαρμογή.

Η σύγκριση B2--EISE αποκτά έτσι ιδιαίτερα καθαρή πειραματική μορφή: \textbf{ίδιο σώμα, ίδιος κωδικοποιητής, ίδιος υπολογιστικός προϋπολογισμός και διαφορετικό είδος σχεσιακής εποπτείας}. Το B2 χρησιμοποιεί ζεύγη που τεκμηριώνουν κοινή οικογένεια, ενώ ο EISE ζεύγη που συνδέονται με συγκεκριμένη, εξηγητική σχέση τεκμηριωμένη από τις ιστορικές πηγές. Οι διαχωρισμοί εκπαίδευσης, ανάπτυξης και τελικού ελέγχου εφαρμόζονται στις φυσικές οντότητες και στις γραμμές προέλευσης πριν από τη δημιουργία των ζευγών. Το $C_{\mathrm{exp}}$, η τροπική δοκιμή, η εκτελεστική δοκιμή και το $P^{\star}$ παραμένουν ανεξάρτητοι εξωτερικοί έλεγχοι.

\paragraph{Σύνθεση της παπαδικής συλλογής και διαδρομή προς τις πηγές}
Η παπαδική συλλογή, οργανωμένη κατά ήχο, έχει εξαχθεί από 758 συνθέσεις: 248 ελληνικά Χερουβικά, 101 μη ελληνικά Χερουβικά, 351 Κοινωνικά και 58 άλλα έργα. Τα αρχεία Word παρέχουν ανά φόρμουλα καταλογογραφικό πεδίο πηγής με τίτλο και απόδοση προσώπου, ενώ η βιβλιογραφία της συλλογής απαριθμεί εκδόσεις και χειρόγραφα, μεταξύ τους εξηγήσεις Χουρμουζίου και τα ΕΒΕ--ΜΠΤ 704--705. Το ΠΕ0 μετατρέπει αυτά τα πεδία σε ρητό μητρώο ρόλων ---αποδιδόμενος μελοποιός, εξηγητής/μεταγραφέας, προσαρμογέας, έργο, έκδοση και μάρτυρας--- πριν χρησιμοποιηθούν για ιστορική ανεξαρτησία ή χρονολόγηση. Η προέλευση αυτή παρέχει διαδρομή ελέγχου προς το έργο και τον ιστορικό μάρτυρα. Το ΠΕ0 επαληθεύει τη σύνδεση κάθε εγγραφής με συγκεκριμένη πηγή και μετρά χωριστά εμφανίσεις, οικογένειες και ατομικά γεγονότα. Για ορισμένες ομάδες πηγών ο ενιαίος κατάλογος δηλώνει επιλεκτική παράλειψη ιδιαίτερα μεγάλων φόρμουλων, μεταξύ άλλων άνω των τεσσάρων γραμμών στην αντίστοιχη καταλογογραφική ενότητα (\emph{Papadic Formulae}, τυπ. σ.~1005 και ανάλογες ενότητες). Το όριο αυτό αφορά την παρουσίαση του καταλόγου και δεν μετατρέπεται σε μουσικό κατώφλι. Το ΠΕ0 καταγράφει την κάλυψη ανά πηγή και, όπου η πληρότητα είναι κρίσιμη για μήκος ή τεκμηριωμένο εύρος, επιστρέφει στο πλήρες έργο. Οι 10.000+ δηλωμένες φόρμουλες ανήκουν στη χωριστή ειρμολογική/στιχηραρική συλλογή, της οποίας το ενιαίο PDF εξαιρεί την παπαδική.

\subsection{Εγκυρότητα οικογενειακής και ιστορικής εποπτείας}\label{app:corpora-supervision}

Το μητρώο δεδομένων διαχωρίζει ρητά τα επίπεδα ταυτότητας, πραγμάτωσης, εμφάνισης και ιστορικής σχέσης:
{\small
\begin{xltabular}{\textwidth}{@{}p{0.27\textwidth}Y@{}}
\toprule
\textbf{Πεδίο} & \textbf{Ρόλος}\\\midrule
\endfirsthead
\toprule
\textbf{Πεδίο} & \textbf{Ρόλος}\\\midrule
\endhead
\midrule
\multicolumn{2}{r}{\footnotesize Συνέχεια στην επόμενη σελίδα}\tabularnewline
\endfoot
\bottomrule
\endlastfoot

\texttt{family\_id} & Αναλυτική ομαδοποίηση κοινής ταυτότητας θέσης. Αποδίδεται μόνο σε τύπο που έχει περάσει την προβλεπόμενη διαδικασία παγώματος και αξιοπιστίας\\
\texttt{type\_signature\_version} & Έκδοση της παγωμένης υπογραφής του αναλυτικού τύπου/οικογένειας\\
\texttt{family\_status} & Κατάσταση τύπου: \texttt{candidate}, \texttt{piloted}, \texttt{frozen}, \texttt{reliability\_passed}, \texttt{candidate\_only} ή \texttt{singleton}\\
\texttt{analysis\_scale} & Κλίμακα κρίσης: \texttt{elementary\_position}, \texttt{composite\_position}, \texttt{position\_sequence}, \texttt{component\_of\_larger\_unit} ή \texttt{unresolved}\\
\texttt{historical\_designation} & Ιστορική ονομασία με συγκεκριμένη πηγή και βαθμό βεβαιότητας, όπου υπάρχει\\
\texttt{reference\_exemplar} & Αντιπροσωπευτική τεκμηριωμένη εμφάνιση που βοηθά στην περιγραφή της υπογραφής κατά την ανάπτυξη. Δεν αποτελεί μόνη της κριτήριο ένταξης και δεν χρησιμοποιείται ως ανεξάρτητο reliability validation\\
\texttt{diagnostic\_relations} & Διαγνωστικές μουσικές σχέσεις της συγκεκριμένης έκδοσης υπογραφής. Κάθε σχέση φέρει ρόλο \texttt{required}, \texttt{supportive} ή \texttt{exclusion}\\
\texttt{diagnostic\_scope} & Πεδίο στο οποίο ενεργοποιείται η αντίστοιχη διαγνωστική σχέση, π.χ. ήχος/βαθμίδα, γένος, ιδίωμα, κειμενικός ή λειτουργικός τόπος ή σημειογραφικό περιβάλλον\\
\texttt{diagnostic\_profiles} & Προαιρετικό, μουσικολογικά τεκμηριωμένο σύνολο εναλλακτικών επαρκών διαγνωστικών διαμορφώσεων στην ίδια αναλυτική κλίμακα· παγώνει μαζί με την υπογραφή και ελέγχεται σε ανεξάρτητο υλικό\\
\texttt{realization\_rules} & Προδηλωμένοι κανόνες επιτρεπτής διαφορετικής πραγμάτωσης. Κάθε κανόνας συνδέεται με συγκεκριμένο τεκμήριο\\
\texttt{modal\_degree\_rule} & Ρόλος της τροπικής βαθμίδας. Δηλώνεται μόνο όταν είναι διαγνωστικός ή επιτρέπει τεκμηριωμένο μετασχηματισμό\\
\texttt{genus\_rule} & Διατονικό, χρωματικό ή εναρμόνιο γένος, μόνο όπου επηρεάζει τεκμηριωμένα την ταυτότητα ή την πραγμάτωση\\
\texttt{global\_echos / local\_modal\_state} & Γενικός ήχος έργου και τοπική μουσικοτροπική κρίση καταγράφονται χωριστά, με πηγή, σημειογραφική μαρτυρία, ενδεχόμενη διάκριση \emph{ἀπὸ παραλλαγῆς/ἀπὸ μέλους} και αβεβαιότητα· δεν αποτελούν εξ ορισμού οικογενειακή ετικέτα\\
\texttt{idiom\_scope} & Ειρμολογικό, στιχηραρικό ή παπαδικό πεδίο στο οποίο τεκμηριώνεται ο κανόνας\\
\texttt{section\_scope} & Θέση μέσα στη μορφή ή στην ενότητα του έργου, όταν αυτή είναι μουσικολογικά σχετική\\
\texttt{text\_locus\_rule} & Κειμενικός τόπος ή συλλαβική συνθήκη ως \texttt{required}, \texttt{permitted}, \texttt{typical} ή \texttt{unknown}\\
\texttt{tempo\_rule} & Αγωγική ή χρονική συνθήκη μόνο όταν υπάρχει ανεξάρτητη τεκμηρίωση\\
\texttt{notational\_context\_rule} & Ιστορική σημειογραφική συνθήκη. Δεν συμπληρώνεται τεχνητά όταν δεν παρατηρείται στη Νέα Μέθοδο\\
\texttt{functional\_role\_status} & Ρόλος λειτουργίας: \texttt{identity\_required}, \texttt{realization\_condition}, \texttt{non\_diagnostic} ή \texttt{unknown}\\
\texttt{context\_compatibility} & Περιγραφική κρίση χρήσης: \texttt{attested\_typical}, \texttt{attested\_exception}, \texttt{unattested\_or\_unknown} ή \texttt{not\_applicable}\\
\texttt{positive\_examples} & Πραγματικά τεκμηριωμένα θετικά παραδείγματα που χρησιμοποιήθηκαν για τη συγκρότηση ή τον έλεγχο της υπογραφής\\
\texttt{hard\_negatives} & Κοντινά αντιπαραδείγματα με ρητή καταγραφή του διαγνωστικού γνωρίσματος που παραβιάζεται\\
\texttt{unresolved\_rules} & Προκαθορισμένες συνθήκες στις οποίες η κρίση πρέπει να παραμένει ανεπίλυτη\\
\texttt{source\_provenance} & Πηγή, έργο, μάρτυρας, φύλλο ή σελίδα, ρόλος προσώπου και κατάσταση άμεσης ή έμμεσης τεκμηρίωσης\\
\texttt{membership\_judgement} & Κρίση $J(x,T;r)$ ανά κριτή, με τιμή \texttt{member}, \texttt{non\_member} ή \texttt{unresolved}, αιτιολογία, τεκμήρια και έκδοση κανόνα\\
\texttt{diagnostic\_check} & Έλεγχος κάθε ενεργού διαγνωστικού γνωρίσματος ως \texttt{satisfied}, \texttt{violated} ή \texttt{not\_assessable}, μαζί με τον ρόλο και το πεδίο εφαρμογής του. Κρίση \texttt{non\_member} απαιτεί τεκμηριωμένο \texttt{exclusion} ή ασυμβατότητα με όλες τις ενεργές επιλέξιμες διαμορφώσεις βάσει προδηλωμένου κανόνα, όχι μόνο απουσία υποστήριξης\\
\texttt{unresolved\_reason} & Αιτία ανεπίλυτης κρίσης: \texttt{vague}, \texttt{evidence}, \texttt{boundary} ή άλλη προκαθορισμένη κατηγορία\\
\texttt{realization\_id} & Συγκεκριμένη τροπικά/λειτουργικά εξαρτημένη πραγμάτωση\\
\texttt{occurrence\_id} & Συγκεκριμένη εμφάνιση σε έργο, πηγή ή εκτέλεση\\
\texttt{text\_locus\_id} & Χωριστό αναγνωριστικό κειμενικού τόπου. Δεν αποτελεί οικογενειακή ετικέτα\\
\texttt{boundary\_status} & Κατάσταση οριοθέτησης: ακριβής, σύνολο αποδεκτών ζευγών ή ανεπίλυτη\\
\texttt{boundary\_sensitivity} & \texttt{boundary\_robust} ή \texttt{boundary\_dependent}, μόνο για ανάλυση ευαισθησίας και όχι ως νέα ετικέτα ground truth\\
\texttt{admissible\_boundary\_pairs} & Μουσικολογικά αποδεκτά ζεύγη αρχής--τέλους στην κοινή ατομική κλίμακα\\
\texttt{coexistence\_relation} & Σχέση δύο εμφανίσεων: συνυπάρχουσες, εναλλακτικές ή ανεπίλυτη\\
\texttt{historical\_homology} & Επαληθευμένη αντιστοιχία του συγκεκριμένου ιστορικού/μουσικού χωρίου ($E_{\mathrm{hom}}$)\\
\texttt{exegesis\_observability} & Κατάσταση εξηγητικού τεκμηρίου: \texttt{clear}, \texttt{equivocal}, \texttt{omitted}, \texttt{unavailable} ή \texttt{unresolved}\\
\texttt{coexegesis\_status} & Τεκμηριωμένη συνεξήγηση ή εξάρτηση μαρτύρων ($E_{\mathrm{coex}}$): \texttt{confirmed}, \texttt{excluded} ή \texttt{uncertain}; ορίζεται μόνο από στοιχεία προέλευσης των πηγών\\
\texttt{copy\_dependence} & Τεκμηριωμένη άμεση αντιγραφική ή παράγωγη εξάρτηση ($E_{\mathrm{copy}}$)\\
\texttt{cross\_repertoire} & Σημαία ότι η περίπτωση προέρχεται από διαφορετικό ρεπερτοριακό στρώμα και χρησιμοποιείται μόνο με τον δηλωμένο ελεγκτικό ρόλο\\
\texttt{overlap\_status} & Καμία, επιβεβαιωμένη, εξαρτώμενη από την ερμηνεία ή ανεπίλυτη επικάλυψη\\
\texttt{functional\_role} & Λειτουργικός ρόλος αναγνωρισμένης θέσης, π.χ. \texttt{opening}, \texttt{connector}, \texttt{internal} ή \texttt{cadential}, όπου τεκμηριώνεται\\
\texttt{region\_status} & \texttt{documented\_occurrence}, \texttt{unassigned\_connective\_material}, \texttt{documented\_background}, \texttt{unresolved} ή \texttt{unchecked}\\
\texttt{historical\_relation\_id} & Αυτοτελής ιστορική σχέση μεταξύ συγκεκριμένων εμφανίσεων ή πραγματώσεων\\
\texttt{pair\_derivation} & Κανόνας με τον οποίο παράγεται εκπαιδευτικό ζεύγος από μία ή περισσότερες ιστορικές σχέσεις\\
\texttt{common\_locus\_id} & Αναγνωριστικό του κοινού τεκμηριωμένου μουσικού/ιστορικού χωρίου που στηρίζει την αντιστοίχιση\\
\texttt{alignment\_cardinality} & Σχέση ορίων παλαιότερης και αναλυτικότερης απόδοσης: \texttt{one\_to\_one}, \texttt{one\_to\_many}, \texttt{many\_to\_one}, \texttt{partial} ή \texttt{unresolved}\\
\texttt{comparison\_scope} & Επίπεδο ιστορικής σύγκρισης: \texttt{whole\_position}, \texttt{subregion}, \texttt{composite\_unit} ή \texttt{unresolved}\\
\texttt{historical\_edge\_ids} & Αναγνωριστικά των πραγματικών ιστορικών ακμών από τις οποίες προκύπτει το παράγωγο ζεύγος\\
\texttt{exemplar\_level} & Επίπεδο του τεκμηρίου από το οποίο παράγεται η αντιστοίχιση (π.χ. παλαιότερη καταγραφή, εξήγηση, εμφάνιση Νέας Μεθόδου)\\
\texttt{source\_relation\_term} & Όρος της ίδιας της πηγής, π.χ. \texttt{exegesis}, \texttt{hermeneia}, \texttt{analepsis}, \texttt{other} ή \texttt{unspecified}\\
\texttt{historical\_relation\_type} & Επιχειρησιακή κατηγορία της έρευνας: \texttt{explicit\_old\_to\_new\_rendering}, \texttt{co\_explanation}, \texttt{copy\_dependence}, \texttt{other} ή \texttt{unresolved}\\
\texttt{transformation\_profile} & Περιγραφικό προφίλ ανάπτυξης, σύντμησης, διαποίκιλσης, προσαρμογής, ανασυνδυασμού ή άλλης τεκμηριωμένης διαφοροποίησης\\
\texttt{evidence\_role} & Ρόλος τεκμηρίου: \texttt{definition}, \texttt{development}, \texttt{reliability\_validation} ή \texttt{historical\_relation\_only}\\
\texttt{dataset\_role} & Ρόλος της περίπτωσης στο συγκεκριμένο σύνολο δεδομένων ή πείραμα\\
\end{xltabular}
}\par

Η λειτουργία \texttt{connector} είναι ιδιότητα αναγνωρισμένης θέσης. Δεν χρησιμοποιείται ως υπολειμματική ετικέτα για κάθε υλικό ανάμεσα σε δύο θέσεις. Το \texttt{unassigned\_connective\_material} δηλώνει μόνο περιοχή που δεν έχει τεκμηριωθεί ως εμφάνιση τύπου. Η διάκριση αυτή ακολουθεί τις συνδετικές formulae που περιγράφει ο Gheorghiță (2010, Πίν.~8--9).

Το \texttt{evidence\_role} προστατεύει τον ανεξάρτητο έλεγχο. Παράδειγμα που χρησιμοποιήθηκε για να οριστεί ή να αναθεωρηθεί μια υπογραφή δεν μπορεί αργότερα να μετρήσει ως ανεξάρτητη επικύρωση της ίδιας έκδοσης. Οι τυπολογικές περιπτώσεις του Gheorghiță μπορούν επομένως να λειτουργήσουν ως \texttt{definition} ή \texttt{development} material. Δεν θεωρούνται από μόνες τους \texttt{reliability\_validation}.

Η διαθεσιμότητα εξηγητικού υλικού καταγράφεται ανά locus και συνοψίζεται ανά οικογένεια. Ως διαγνωστική κάλυψη αναφέρεται
\[
C_{\mathrm{ex}}(F)=\frac{\#\{\text{επιλέξιμες εμφανίσεις του }F\text{ με χρησιμοποιήσιμο εξηγητικό τεκμήριο}\}}{\#\{\text{επιλέξιμες επικυρωμένες εμφανίσεις του }F\}}.
\]
Το $C_{\mathrm{ex}}(F)$ δεν είναι νέο επιβεβαιωτικό εκτιμώμενο μέγεθος. Χρησιμοποιείται για να δηλώνει το πεδίο στο οποίο μπορεί να γενικευθεί η M2. Αναφέρονται χωριστά οι καταστάσεις \texttt{clear}, \texttt{equivocal}, \texttt{omitted}, \texttt{unavailable} και \texttt{unresolved}.

Η ιστορική ομολογία δεν συνεπάγεται ταυτόσημη τμηματοποίηση. Μία παλαιότερη θέση μπορεί να αντιστοιχεί σε περισσότερες αναλυτικές μονάδες της Νέας Μεθόδου ή το αντίστροφο. Τα πεδία \texttt{alignment\_cardinality} και \texttt{comparison\_scope} καταγράφουν αυτή τη σχέση πριν παραχθεί εκπαιδευτικό ζεύγος. Η περίπτωση της Formula V στον Gheorghiță (2010, Πίν.~5) χρησιμοποιείται ως μουσικολογικό παράδειγμα του γιατί δεν επιβάλλεται αντιστοίχιση 1:1.

Για κάθε εμφάνιση $x$ αποθηκεύεται το σύνολο $\mathcal B_x=\{(a_k,b_k)\}$ των μουσικολογικά αποδεκτών ζευγών αρχής--τέλους στην κοινή ατομική κλίμακα. Η ακριβής οριοθέτηση αντιστοιχεί σε ένα μόνο ζεύγος· ανεξάρτητα εύρη αρχής και τέλους χρησιμοποιούνται μόνο όταν όλοι οι συνδυασμοί τους είναι τεκμηριωμένα επιτρεπτοί. Αν η οικογενειακή απόδοση συνδέεται με συγκεκριμένη οριοθέτηση, διατηρούνται οι επιτρεπτές τριάδες $(a,b,F)$ χωρίς υπόθεση ότι όλοι οι συνδυασμοί είναι ισοδύναμοι. Δεν αποδίδονται βάρη πιθανότητας στις εναλλακτικές οριοθετήσεις χωρίς ανεξάρτητη εμπειρική βάση. Οι εναλλακτικές οριοθετήσεις της ίδιας εμφάνισης δεν αυξάνουν το πλήθος εμφανίσεων. Η σημαία \texttt{boundary\_robust}/\texttt{boundary\_dependent} χρησιμοποιείται μόνο για να δηλώσει αν ένα συμπέρασμα παραμένει σταθερό σε όλες τις αποδεκτές οριοθετήσεις.

Οι τύποι σχέσης δεν εξισώνονται μεταξύ τους. Η \texttt{explicit\_exegesis} δηλώνει συγκεκριμένη τεκμηριωμένη εξηγητική σύνδεση και είναι ο μόνος τύπος που μπορεί να τροφοδοτήσει τα θετικά ζεύγη του EISE, εφόσον η κοινή οικογενειακή ταυτότητα έχει επίσης επαληθευθεί. Οι \texttt{teacher\_student}, \texttt{succession} και \texttt{explicit\_mimesis} καταγράφονται ως μεταδεδομένα μετάδοσης· η \texttt{inferred\_similarity} ως αναλυτικό τεκμήριο και η \texttt{same\_family} ως χωριστή οικογενειακή σχέση. Τεκμηριωμένη σύντμηση ή άλλη ιστορική σχέση διατηρεί τον δικό της τύπο και δεν μετατρέπεται σε EISE-positive χωρίς ανεξάρτητη τεκμηρίωση εξηγητικής σύνδεσης.

Ο όρος που χρησιμοποιεί η ίδια η πηγή φυλάσσεται χωριστά από την επιχειρησιακή κατηγορία της έρευνας. Έτσι οι όροι \texttt{exegesis}, \texttt{hermeneia} ή \texttt{analepsis} δεν εξισώνονται αυτόματα με μία υπολογιστική σχέση. Επίσης, η παράθεση αποδόσεων Γρηγορίου και Χουρμουζίου στην ίδια μελέτη δεν δημιουργεί από μόνη της \texttt{co\_explanation}. Κάθε τέτοια σχέση απαιτεί επαλήθευση στο συγκεκριμένο \texttt{common\_locus\_id} και στις πραγματικές πηγές.

Η καταλογογραφική ένδειξη τύπου, κατηγορίας ή παραλλαγής αποθηκεύεται αυτούσια μαζί με τη θέση της στην πηγή. Το ΠΕ0 εξετάζει σε προκαθορισμένο, στρωματοποιημένο δείγμα εκπαίδευσης και ανάπτυξης τι κωδικοποιεί η οργάνωση του καταλόγου: κοινή μελωδική ταυτότητα, ευρύτερη κατηγορία, παραλλαγή ή λειτουργική συγγένεια. Η απογραφή στρωματοποιείται επιπλέον ως προς λειτουργικό ρόλο και αναφέρει χωριστά καταληκτική και μη καταληκτική κάλυψη. Ελέγχει επίσης πραγματικά καταλογογραφημένες μη λεξικές εμφανίσεις, την αντίστοιχη κατηγορία και την ερμηνεία της συλλαβικής κωδικοποίησης στο συγκεκριμένο λεκτικό καθεστώς. Ανεπίλυτες περιπτώσεις διατηρούνται ως υποψήφιες σχέσεις. Η μετατροπή μιας ένδειξης σε θετικό ζεύγος του B2 ακολουθεί γραπτό κριτήριο τεκμηριωμένης κοινής οικογένειας και ανεξάρτητο μουσικολογικό έλεγχο. Οι ευρύτερες και ανεπίλυτες σχέσεις διατηρούν χωριστή καταγραφή. Η επιλογή δείγματος καλύπτει ήχους, λειτουργικούς ρόλους, υποψήφιες οικογένειες, πηγές και βαθμούς επιφανειακής απόστασης και προηγείται της εκτίμησης οποιασδήποτε επίδοσης μοντέλου.

Στο μητρώο προέλευσης διακρίνονται η εμφάνιση, η υποψήφια μουσική πραγμάτωση και η γραπτή απόδοσή της. Ισοδύναμες αποδόσεις που τεκμηριώνονται από τις πηγές συνδέονται χωρίς να μετρώνται ως ανεξάρτητες μουσικές παραλλαγές· αβέβαιες ταυτίσεις διατηρούνται ανεπίλυτες. Η διάκριση αφορά την ερμηνεία και τις μονάδες ανεξαρτησίας· οι επαληθευμένες ακμές $E_{\mathrm{an}}$ παραμένουν στην υπάρχουσα εποπτεία και στις προκαθορισμένες αναλύσεις συμβολικής γεωμετρίας. Ιστορικές ονομασίες τύπων ή υποστάσεων φυλάσσονται αυτούσιες με πηγή και βαθμό βεβαιότητας, ανεξάρτητα από τα σύγχρονα αναγνωριστικά οικογένειας. Ο ρόλος προσώπου ανήκει στη συγκεκριμένη σχέση ή μαρτυρία, όχι σε καθολική ιδιότητα του προσώπου. Οι τελικές προτάσεις διάσπασης/συγχώνευσης διατηρούν τα αρχικά αναγνωριστικά και χωριστό μητρώο αντιστοιχίσεων, χωρίς επαναβαθμολόγηση του επιβεβαιωτικού ελέγχου.

Δεύτερος ανεξάρτητος γνώστης ελέγχει το ίδιο δείγμα χωρίς πρόσβαση στις προβλέψεις. Αναφέρονται παρατηρούμενη συμφωνία, κατάλληλος συντελεστής αξιοπιστίας, διαφωνίες και ο κανόνας επίλυσής τους. Όλα τα κύρια θετικά ζεύγη απαιτούν ήδη την ελάχιστη προκαθορισμένη επαλήθευση κοινής οικογένειας. Ένα αυστηρότερα διπλοελεγμένο υποσύνολο υψηλής βεβαιότητας χρησιμοποιείται επιπλέον σε προκαθορισμένο έλεγχο ευαισθησίας της ίδιας σύγκρισης B2--EISE. Η πρόσθετη αξία της εξηγητικής σχέσης ερμηνεύεται μαζί με την ποιότητα και την αναλυτική κλίμακα της οικογενειακής εποπτείας. Η κοινή οικογενειακή εποπτεία και τα πρότυπα αναφοράς των B1, B2 και EISE αντλούνται από το ίδιο οριστικοποιημένο λεξιλόγιο του συνόλου εκπαίδευσης. Το τελικό $C_{\mathrm{exp}}$ διατηρείται αποκλειστικά για ανεξάρτητη αξιολόγηση.

Για κάθε υποψήφια ιστορική αντιστοίχιση καταγράφονται τρεις \textbf{διακριτές} κρίσεις: (1) η εγκυρότητα και η έκταση της ιστορικής σχέσης, (2) η οικογενειακή ένταξη των εμπλεκόμενων εμφανίσεων βάσει της οριστικοποιημένης υπογραφής τύπου και (3) ο τύπος διαφοροποίησης $E_{\mathrm{an}}/E_{\mathrm{tr}}$, όπου μπορεί να διαγνωστεί. Καμία από τις τρεις κρίσεις δεν υποκαθιστά τις άλλες. Ανεξάρτητος δεύτερος έλεγχος σε προεπιλεγμένο δείγμα τεκμηριώνει την αξιοπιστία κάθε άξονα. Επαληθευμένη ιστορική σχέση με ανεπίλυτο τύπο διαφοροποίησης μπορεί να παραμένει στο μητρώο, αλλά τροφοδοτεί την εκπαιδευτική απώλεια ταυτότητας μόνο όταν πληρούνται επιπλέον οι ανεξάρτητοι κανόνες $E_{\mathrm{hom}}$ και $E_{\mathrm{fam}}$. Αμφίβολη ιστορική σχέση παραμένει στο μητρώο υποψηφίων έως την επίλυσή της από τα διαθέσιμα ιστορικά τεκμήρια.

Για την προκαθορισμένη δευτερεύουσα Υ2\textsubscript{E}, η επιφανειακή απόσταση μεταξύ τεκμηριωμένων εξηγήσεων υπολογίζεται \textbf{σε συντεταγμένη ουδέτερη ως προς την αναπαράσταση που αξιολογείται και χωρίς χρήση ετικετών οικογένειας}. Με προκαθορισμένο κανόνα τα ζεύγη ή οι ομάδες χωρίζονται σε χαμηλή, μέση και υψηλή επιφανειακή απόκλιση. Η στρωμάτωση αυτή ορίζει τη δυσκολία του ελέγχου ανεξάρτητα από το μοντέλο: η ισχυρότερη ένδειξη αναλλοιότητας είναι η διατήρηση της οικογενειακής ταυτότητας στο στρώμα μεγάλης επιφανειακής απόκλισης.

Στο υποσύνολο όπου η τεκμηρίωση από τις ιστορικές πηγές επιτρέπει ασφαλή διάκριση, οι επαληθευμένες εξηγητικές ακμές ταξινομούνται ως
\[
E_{\mathrm{ex}}^{\mathrm{class}}=E_{\mathrm{an}}\cup E_{\mathrm{tr}}.
\]
Το $E_{\mathrm{an}}$ δηλώνει τεκμηριωμένη αναλυτική διαφοροποίηση. Το $E_{\mathrm{tr}}$ δηλώνει μουσικά σημαντική διαφοροποίηση μέσα σε ρητά τεκμηριωμένη εξηγητική σχέση. Η ταξινόμηση δεν βασίζεται στον αριθμό σημείων. Απαιτεί συγκρίσιμα όρια, ελεγχόμενη αντιστοίχιση και επαρκή μουσική τεκμηρίωση. Όταν αυτά δεν επαρκούν, ο τύπος διαφοροποίησης παραμένει \texttt{unresolved}. Η ιστορική σχέση και η κοινή οικογενειακή ταυτότητα εξακολουθούν να κρίνονται χωριστά.

\subsection{Συμβολικά σώματα και εκτελεστική τεκμηρίωση}\label{app:corpora-symbolic}

{\small
\begin{xltabular}{\textwidth}{@{}p{0.15\textwidth}p{0.17\textwidth}Y@{}}
\toprule
& \textbf{Ενδεικτικό μέγεθος} & \textbf{Ρόλος}\tabularnewline
\midrule
\endfirsthead
\toprule
& \textbf{Ενδεικτικό μέγεθος} & \textbf{Ρόλος}\tabularnewline
\midrule
\endhead
\midrule
\multicolumn{3}{r}{\footnotesize Συνέχεια στην επόμενη σελίδα}\tabularnewline
\endfoot
\bottomrule
\endlastfoot

Ανάπτυξης & 40--60 έργα & Δοκιμή αναλυτών, ρύθμιση, βαθμονόμηση· ποτέ για επιβεβαιωτικό ισχυρισμό\tabularnewline
Ελέγχου & 60--80 έργα, \(\ge 3\) ομάδες συγγενών πηγών & Διατηρείται χωριστό και χρησιμοποιείται μόνο μετά την οριστικοποίηση του πρωτοκόλλου\tabularnewline
Εφαρμογής & Νικόλαος 15--25· κλασικός ορίζοντας 30--40 & Μουσικολογική εφαρμογή\tabularnewline
\end{xltabular}
}\par

Ένα παπαδικό μέλος αντιστοιχεί ενδεικτικά σε \(300\)--\(800\) ατομικά γεγονότα, με \(8\)--\(15\) θέσεις ανά έργο. Οι αριθμοί δηλώνονται ως \emph{σχεδιαστική παραδοχή} και αντικαθίστανται από τη μετρημένη απογραφή.

\paragraph{Εκτελεστικό σώμα ΔΑΜΑΣΚΗΝΟΣ}
Το σώμα που καθιστά εφικτή την εκτελεστική αξιολόγηση δημιουργήθηκε από τους \textbf{Δελβινιώτη} και \textbf{Κουρουπέτρογλου} στο Εργαστήριο Ομιλίας του ΕΚΠΑ, με ηχογραφήσεις από τον Ιούλιο 2005 έως τον Ιούλιο 2006.

Τα σχεδιαστικά του χαρακτηριστικά είναι ιδιαίτερα ευνοϊκά:
\begin{itemize}
\item \textbf{Είκοσι επαγγελματίες ψάλτες} από διαφορετικές περιοχές, επιλεγμένοι με ρητά κριτήρια που περιλάμβαναν την ηλικία, την καλλιτεχνική επάρκεια, τη δισκογραφική δραστηριότητα και \emph{τους δασκάλους} του κάθε ψάλτη.
\item \textbf{Κοινό προκαθορισμένο ρεπερτόριο} επιλεγμένων τροπαρίων, σε παπαδικό, στιχηραρικό και ειρμολογικό μέλος, και στους οκτώ ήχους· η πραγματική κάλυψη ανά εκτελεστή ελέγχεται στην απογραφή.
\item \textbf{Τυποποιημένες συνθήκες ηχογράφησης}: ίδιος εξοπλισμός, σταθερά επίπεδα προενισχυτών, ακουστικό στούντιο και δειγματοληψία 44,1 kHz.
\item \textbf{Συγχρονισμένο ηλεκτρογλωττογράφημα}, το οποίο καταγράφει την επιφάνεια επαφής των φωνητικών χορδών και επιτρέπει ανθεκτική εκτίμηση θεμελιώδους συχνότητας ανεξάρτητα από τον φωνητικό σωλήνα.
\item \textbf{Επισημείωση σε πέντε επίπεδα} με το πρόγραμμα PRAAT: όρια συλλαβών τοποθετημένα χειροκίνητα, χαρακτήρες της γραπτής σημειογραφίας, αναμενόμενο τονικό ύψος υπολογισμένο από τις βυζαντινές κλίμακες, χαρακτήρες βασισμένοι στην εκτέλεση και αντιληπτικό ύψος.
\item Τα μουσικά κείμενα προέρχονται από έντυπα βιβλία πατριαρχικών δασκάλων \textbf{μετά τη μεταρρύθμιση της σημειογραφίας}, δηλαδή είναι γραμμένα στη \emph{Νέα Μέθοδο της βυζαντινής μουσικής}.
\end{itemize}

Η πρωτογενής περιγραφή του σώματος τεκμηριώνει επίσης \textbf{χωριστές αεροδυναμικές μετρήσεις φωνητικών ασκήσεων} με ροή και πίεση. Η δημοσιευμένη μεθοδολογία επισημείωσης περιγράφει ημιαυτόματη τμηματοποίηση στο PRAAT και ενσωμάτωση μεταδεδομένων στίχου, σημειογραφίας και μουσικών διαστημάτων στα τμηματοποιημένα αρχεία. Το ΠΕ0 δεν θεωρεί αυτονόητη τη σημερινή διαθεσιμότητα αυτών των παραγώγων: επαληθεύει χωριστά αν παραδίδονται τα πρωτογενή αρχεία επισημείωσης PRAAT ή ισοδύναμη εξαγωγή τους, το σχήμα κωδικοποίησης και η κάλυψη ανά ηχογράφηση. Διακρίνει επίσης τις αεροδυναμικές μετρήσεις από τα συγχρονισμένα δεδομένα της ψαλμώδησης και καταγράφει την πραγματική πρόσβαση και αντιστοίχιση κάθε καναλιού. Κινηματικό ή βιντεοληπτικό υλικό συνυπολογίζεται μόνο εφόσον επιβεβαιωθεί αυτοτελώς από τον φορέα.

Η εκτελεστική αξιολόγηση στηρίζεται στα \textbf{πρωτογενή σήματα}: ακατέργαστο ήχο και ηλεκτρογλωττογράφημα. Τα επίπεδα επισημείωσης που βασίζονται στην εκτέλεση της αρχικής επισημείωσης διατηρούνται ως ερμηνευτικό υλικό και όχι ως μετρητικό πρότυπο. Η φωνητική κατηγορία και τα διαθέσιμα μεταδεδομένα των ψαλτών ενσωματώνονται στις προκαθορισμένες αναλύσεις ώστε η εκτιμώμενη μουσική δομή να μην συγχέεται με συστηματικές διαφορές φωνής ή προέλευσης.

\subsection{Σύνολο εκπαίδευσης και ιστορική προέλευση}\label{app:corpora-training}

Το βασικό σώμα προεκπαίδευσης $\mathcal D_{\mathrm{sym}}^{\mathrm{train}}$ συγκροτείται από το επιλέξιμο παπαδικό εκπαίδευσης τμήμα της $\mathcal S_{\mathrm{StA}}$ και από τα λοιπά παπαδικά συμβολικά σώματα που ήδη προβλέπει η πρόταση, εφόσον περάσουν τον ίδιο έλεγχο προέλευσης, κωδικοποίησης και διαχωρισμού. Περιλαμβάνει πραγματικά μετατρεμμένες ακολουθίες: αποσπάσματα καταλόγου και, όπου εξασφαλίζεται η πρόσβαση και η μετατροπή, πλήρη έργα των επιλέξιμων πηγών. Κάθε εγγραφή δηλώνει συλλογή, γένος, έργο, μάρτυρα, ομάδα συγγενών πηγών, τύπο αποσπάσματος, ποιότητα κωδικοποίησης και εκπαιδευτική χρήση. Η βιβλιογραφική παραπομπή λειτουργεί ως διαδρομή αναζήτησης· στον μετρημένο όγκο εισέρχονται τα γεγονότα που έχουν πράγματι κωδικοποιηθεί.

Η επέκταση προεκπαίδευσης προς την ειρμολογική/στιχηραρική συλλογή αποτελεί προκαθορισμένη δυνατότητα μόνο εφόσον ο πιλοτικός έλεγχος τεκμηριώσει περιορισμένο επιλέξιμο παπαδικό όγκο. Η απόφαση αξιολογεί, με κοινό συμβολικό στόχο ανάπτυξης και κοινό υπολογιστικό προϋπολογισμό, την ανάκτηση διαστηματικών και χρονικών ιδιοτήτων στο παπαδικό υλικό. Καταγράφονται οι διαφορές γένους, βαθμού μελισματικής ανάπτυξης και χρονικής οργάνωσης. Το τελικό σύνολο εκπαίδευσης και οι αποκλεισμοί του οριστικοποιούνται από κοινού για B1, B2 και EISE πριν από τη συγκριτική εκπαίδευση. Η επιλογή αυτή αφορά αποκλειστικά συμβολική προεκπαίδευση· η σχεσιακή εποπτεία και ο κύριος έλεγχος διατηρούν το παπαδικό πεδίο τους.

Το μητρώο διακρίνει \textbf{χρονολογία μελοποίησης, χρονολογία εξήγησης/καταγραφής και χρονολογία έκδοσης}, μαζί με τον ρόλο κάθε προσώπου και τη σχέση της έκδοσης με παλαιότερο μάρτυρα. Η επιλεκτική ένταξη μεταγενέστερων του 1900 συνθέσεων που δηλώνει η Μονή καταγράφεται ανά έργο και ήχο. Τα κριτήρια ένταξης εξειδικεύονται ανά χρήση:
\begin{itemize}
\item Ο κλασικός ορίζοντας του Νικολάου συγκροτείται από τεκμηριωμένα προγενέστερες μελοποιήσεις των δηλωμένων ιστορικών στρωμάτων. Οι νεότερες εκδόσεις τους ελέγχονται ως φορείς παλαιότερων τεκμηρίων.
\item Κάθε ακμή $G_{\mathrm{ex}}$ εισέρχεται στην ιστορική ανάλυση βάσει επαληθευμένης σχέσης και του προκαθορισμένου χρονολογικού πεδίου. Η μουσική ομοιότητα υποδεικνύει υποψήφια σχέση, της οποίας η τεκμηρίωση καταγράφεται χωριστά.
\item Η συμβολική προεκπαίδευση ακολουθεί τους κανόνες του $\mathcal D_{\mathrm{sym}}^{\mathrm{train}}$. Η σύγχρονη ημερομηνία έκδοσης ή η νεότερη μελοποίηση αξιολογείται σύμφωνα με αυτόν τον εκπαιδευτικό ρόλο.
\end{itemize}
Ανατυπώσεις, εξαρτημένες εκδόσεις και επαναλαμβανόμενες μεταγραφές συνδέονται με την ίδια γραμμή προέλευσης πριν από τους διαχωρισμούς. Τα έργα του Νικολάου που προορίζονται για την ιστορική εφαρμογή εξαιρούνται από τη μάθηση του κλασικού προτύπου και της μετρικής της συγκεκριμένης σύγκρισης.

\subsection{Απογραφή, επάρκεια και οριστικοποίηση κλίμακας}\label{app:corpora-gate}

Το ΠΕ0 παράγει μητρώο \textbf{ανά υποψήφια/επικυρωμένη οικογένεια και κατά ήχο}, με χωριστές μετρήσεις πριν και μετά την αποδιπλοποίηση και τους αποκλεισμούς. Πριν από την ενεργοποίηση των δευτερευόντων ισοκρατηματικών διαγνωστικών, το ίδιο μητρώο καταγράφει ποια έργα διαθέτουν τεκμηριωμένη ισοκρατηματική πληροφορία, την προέλευση και χρονολογία κάθε λύσης, καθώς και τη δυνατότητα ανεξάρτητης φραστικής οριοθέτησης. Οι εκδόσεις Μισαηλίδη (1902) και Στ. Χουρμουζίου (1924) αντιμετωπίζονται ως υποψήφια δεύτερα εκδοτικά/ρεπερτοριακά σώματα και απογράφονται ως προς αριθμό Χερουβικών, ήχους, πυκνότητα σημάνσεων, προέλευση και πιθανή εξάρτηση από την κύρια γραμμή. Η πλήρης αξιοποίησή τους για επανάληψη των I1/I2 ενεργοποιείται μόνο αν το πρώτο σώμα παρέχει επαρκή πληροφορία και η δεύτερη γραμμή είναι πραγματικά συγκρίσιμη. Καταγράφονται έργα, εμφανίσεις, υποψήφιες και επικυρωμένες οικογένειες/παραλλαγές, κωδικοποιήσιμα ατομικά γεγονότα, ανεξάρτητοι μάρτυρες και ομάδες συγγενών πηγών, επαληθευμένες ακμές, βαθμοί γράφου και εφικτές αντιστοιχισμένες ζευγοποιήσεις B2. Ανά χωρίο προστίθενται λειτουργικός ρόλος, λεκτικό καθεστώς, κατάσταση των πεδίων $L$ ως εφαρμοστέων ή μη, αναγνωριστικό σταθερού λειτουργικού κειμένου όπου υπάρχει, κειμενικό χωρίο, βαθμός μελισματικής ανάπτυξης, τύπος αποσπάσματος και πληρότητα κάλυψης της πηγής. Τα πρόσωπα καταγράφονται με ρόλο στη συγκεκριμένη μαρτυρία, ώστε μελοποιός, εξηγητής/μεταγραφέας και προσαρμογέας να μη συγχέονται ως μονάδες ανεξαρτησίας. Για την τομή με το ΔΑΜΑΣΚΗΝΟΣ προστίθενται συγκεκριμένο έργο και χωρίο, αριθμός εκτελεστών, διαθέσιμα κανάλια μέτρησης και ποιότητα αντιστοίχισης. Η αναφορά κατά ήχο και λειτουργικό ρόλο δείχνει ποια υποσύνολα στηρίζουν την τροπική δοκιμή, τη μη καταληκτική γενίκευση, την ανάλυση κοινού κειμενικού χωρίου και την εφαρμογή στον Νικόλαο. Το κύριο στατιστικό μέγεθος της μεταφοράς δίνει ίσο βάρος σε κάθε επιλέξιμη οικογένεια· η στρωμάτωση κάλυψης προσδιορίζει το πεδίο γενίκευσης.

Ο συνδυασμός των δύο υποδομών επιτρέπει στο ΠΕ0 να διαπιστώσει αν υπάρχει επαρκές κοινό υλικό για το ιστορικό--εκτελεστικό πείραμα. Η παπαδική συλλογή παρέχει το ιστορικό σημείο εκκίνησης και τη διαδρομή προς έργα και πηγές, ενώ το ΔΑΜΑΣΚΗΝΟΣ παρέχει καταγραφές του ίδιου ρεπερτορίου από πολλαπλούς εκτελεστές, σε τυποποιημένες συνθήκες και με συγχρονισμένα σήματα. Με βάση την ακρίβεια και την κάλυψη, το ΠΕ0 καθορίζει το λεπτομερέστερο επίπεδο ανάλυσης που υποστηρίζουν τα δεδομένα: εσωτερική συντεταγμένη, μεγαλύτερες μουσικά ομοιογενείς περιοχές ή, όταν τα επιβεβαιωτικά κριτήρια δεν ικανοποιούνται, το προκαθορισμένο διερευνητικό επίπεδο. Η επιλογή παραμένει ανεξάρτητη από την τελική κατάταξη των μοντέλων.

\textbf{Πριν από τον πιλοτικό έλεγχο} καθορίζονται οι ορισμοί των συγκριτικών μεγεθών, τα κριτήρια εγκυρότητας οικογένειας και ακμής, οι εναλλακτικές χρονικές κλίμακες, η μέθοδος προσομοίωσης και τα κατώφλια απόφασης. \textbf{Ο πιλοτικός έλεγχος} μετρά την πραγματική δυνατότητα μετατροπής, την ερμηνεία των καταλογογραφικών ενδείξεων, την ανεξαρτησία των μαρτύρων, την επικάλυψη και την αξιοπιστία των μετρήσεων. Η αποδοχή συγκεκριμένων εγγραφών και η επιλογή επιπέδου αποτελούν εφαρμογή των προκαθορισμένων κανόνων στα ευρήματα αυτής της απογραφής.

Για κάθε κύρια αντίθεση $j\in\{\mathrm{hist},\mathrm{learn}\}$ ορίζονται επιστημονικά ενδιαφέρουσα διαφορά $\delta_j>0$ στην κλίμακα $S$ και αποδεκτό ημιεύρος $\varepsilon_j>0$, με $\varepsilon_j\leq\delta_j/2$. Οι αριθμητικές τιμές και η αιτιολόγησή τους καταγράφονται πριν από τον πιλοτικό έλεγχο, με αναφορά στο όφελος έναντι του συγκεκριμένου συγκριτή και στη μουσικολογική σημασία της εκτελεστικής αντιστοιχίας. Παραμένουν σχεδιαστικές παράμετροι προς προκαθορισμό και δεν παρουσιάζονται ως ήδη μετρημένες επιδόσεις. Η προσομοίωση εξετάζει την ακρίβεια των δύο αντιθέσεων σε προκαθορισμένο εύρος παραδοχών για εξαρτήσεις και μετρητικό θόρυβο. Το κριτήριο επάρκειας είναι
\[
\Pr_{\mathrm{sim}}\!\left\{
\tfrac12\operatorname{width}\!\left[\mathrm{CI}_{0.95}^{\mathrm{simult}}(\Delta_j)\right]
\leq\varepsilon_j\ \text{για κάθε }j
\right\}\geq 0.90.
\]
Τα διαστήματα είναι ταυτόχρονα για τις δύο αντιθέσεις και η προσομοίωση ακολουθεί τη δομή ομαδοποίησης του τελικού ελέγχου. Ο σχεδιασμός περιλαμβάνει αβεβαιότητα ιστορικών γραμμών, έργων, οικογενειών και εκτελεστών, ολόκληρα προφίλ $u$, καθώς και προκαθορισμένα σενάρια διακύμανσης μεταξύ ρυθμίσεων τυχαιότητας και επιλέξιμων ζευγοποιήσεων B2. Η μέτρηση κάλυψης ανά ήχο συνοδεύει κάθε υποψήφιο επίπεδο, χωρίς να εισάγει πρόσθετες κύριες υποθέσεις ανά ήχο.

\begin{enumerate}
\item \textbf{Επίπεδο 1 --- λεπτομερής επιβεβαιωτική μεταφορά.} Αποτελεί την πρωτεύουσα σχεδιαστική στόχευση: η αντιστοίχιση σε λεπτομερή εσωτερική κλίμακα περνά το προκαθορισμένο κριτήριο αξιοπιστίας και οι δύο κύριες αντιθέσεις περνούν το κριτήριο ακρίβειας στο οριστικοποιημένο επιλέξιμο σύνολο. Επιτρέπει την ισχυρότερη ανά εσωτερική θέση ερμηνεία της εσωτερικής γεωμετρίας σταθερότητας και ελευθερίας.
\item \textbf{Επίπεδο 2 --- επιβεβαιωτική μεταφορά σε αδρότερη μουσική κλίμακα.} Αποτελεί πλήρως επιβεβαιωτική προκαθορισμένη κλίμακα όταν οι μεγαλύτερες μουσικά αντιστοιχισμένες περιοχές προσφέρουν ανθεκτικότερη μέτρηση από τη λεπτομερή συντεταγμένη. Οι ίδιες δύο αντιθέσεις, οι ίδιες απαιτήσεις ακρίβειας και η ίδια πολιτική πολλαπλότητας διατηρούνται. Η επιστημονική υπόθεση παραμένει η ίδια· μεταβάλλεται μόνο η χωρική ανάλυση στην οποία τεκμηριώνεται.
\item \textbf{Επίπεδο 3 --- προκαθορισμένη διερευνητική εκτελεστική επέκταση.} Ενεργοποιείται μόνο αν η πραγματικό κοινό πεδίο αντιστοίχισης ιστορικών πραγματώσεων και σύγχρονων εκτελέσεων ή η αξιοπιστία των μετρήσεων δεν ικανοποιούν τα προκαθορισμένα κριτήρια των δύο επιβεβαιωτικών επιπέδων. Το προφίλ εκτελεστικής μεταβλητότητας τότε αναφέρεται διερευνητικά, ενώ οι ιστορικές και τροπικές επιβεβαιωτικές εργασίες προχωρούν στα δικά τους επιλέξιμα υποσύνολα με τις ήδη δηλωμένες υποθέσεις.
\end{enumerate}

Η επιλογή χρησιμοποιεί αποκλειστικά απογραφή, κάλυψη, ανεξαρτησία πηγών, ποιότητα κωδικοποίησης/αντιστοίχισης, αξιοπιστία μέτρησης σε χωριστό μη επιβεβαιωτικό πιλοτικό έλεγχο και τις προκαθορισμένες προσομοιώσεις. Τα παρατηρούμενα $\rho_F^{(m)}$, $S_m$, $\Delta_{\mathrm{hist}}$ και $\Delta_{\mathrm{learn}}$, η κατάταξη μοντέλων και η ευνοϊκή ή δυσμενής έκβαση της μεταφοράς παραμένουν εκτός της απόφασης. Η έκδοση του μητρώου, οι παράμετροι προσομοίωσης, το επιλεγμένο επίπεδο, οι αποκλεισμοί και οι αιτιολογίες χρονοσημαίνονται πριν από το άνοιγμα των επιβεβαιωτικών αποτελεσμάτων. Ο εκτελεστικός πιλοτικός έλεγχος χρησιμοποιείται για βαθμονόμηση της διαδικασίας μέτρησης από χωριστή διαδικασία· τα δεδομένα της τελικής εκτελεστικής και της τροπικής δοκιμής, του $C_{\mathrm{exp}}$ και του $P^\star$ διατηρούνται αποκλειστικά για την τελική αξιολόγηση και καμία εκτελεστική παρατήρηση δεν εισέρχεται στη μάθηση του EISE.

\paragraph{Πύλη $G_{\mathrm{perf}}$ για το ανεξάρτητο εκτελεστικό κριτήριο (GREEN/AMBER/RED).}
Οι χρωματικές ετικέτες αποτελούν συντομογραφία των τριών ήδη ορισμένων επιπέδων επάρκειας και δεν εισάγουν νέο στατιστικό κριτήριο. Οι επιστημονικά ενδιαφέρουσες διαφορές $\delta_j$ και τα αποδεκτά ημιεύρη $\varepsilon_j$ ορίζονται πριν από τον πιλοτικό έλεγχο. Το ΠΕ0 εφαρμόζει αυτές τις απαιτήσεις στη μετρημένη κάλυψη, στην ανεξαρτησία πηγών και στην αξιοπιστία των μετρήσεων και οριστικοποιεί το επιλέξιμο δείγμα και την κλίμακα πριν από την τελική αξιολόγηση.
{\small
\begin{xltabular}{0.98\linewidth}{@{}p{0.13\linewidth}p{0.28\linewidth}p{0.27\linewidth}Y@{}}
\toprule
\textbf{Κατάσταση} & \textbf{Κάλυψη/αξιοπιστία} & \textbf{Επιτρεπόμενη ανάλυση} & \textbf{Συμπερασματικό καθεστώς}\tabularnewline
\midrule
\endfirsthead
\toprule
\textbf{Κατάσταση} & \textbf{Κάλυψη/αξιοπιστία} & \textbf{Επιτρεπόμενη ανάλυση} & \textbf{Συμπερασματικό καθεστώς}\tabularnewline
\midrule
\endhead
\midrule
\multicolumn{4}{r}{\footnotesize Συνέχεια στην επόμενη σελίδα}\tabularnewline
\endfoot
\bottomrule
\endlastfoot

GREEN & επαρκής λεπτομερής αντιστοίχιση και ακρίβεια & λεπτομερής εσωτερική κλίμακα & επιβεβαιωτικό\tabularnewline
AMBER & επαρκής κάλυψη μόνο σε αδρότερη μουσικά ομοιογενή κλίμακα & σύγκριση ιστορικού και ανεξάρτητου εκτελεστικού προφίλ στην προκαθορισμένη αδρότερη κλίμακα & επιβεβαιωτικό στην προκαθορισμένη αδρότερη κλίμακα\tabularnewline
RED & ανεπαρκές κοινό πεδίο αντιστοίχισης ή ανεπαρκής αξιοπιστία των μετρήσεων για τα δύο προηγούμενα επίπεδα & περιγραφική/διερευνητική εκτελεστική επέκταση & διερευνητικό\tabularnewline
\end{xltabular}
}\par

\subsection{Διαχρονικός έλεγχος ιστορικών ηχογραφήσεων με κριτήριο ετοιμότητας}\label{app:historic-audio}
Το αν μια ιστορική ηχογράφηση μπορεί να ενταχθεί στην ανάλυση καθορίζεται στο ΠΕ0 από την ασφαλή ταυτοποίηση έργου και εκτελεστή, την επικάλυψη με τις τεκμηριωμένες οικογένειες, την ποιότητα της διαθέσιμης μεταφοράς και την αξιοπιστία της μουσικής ευθυγράμμισης. Όπου ο έλεγχος απαιτεί σύγκριση μεταξύ εκτελεστών, προϋποτίθενται περισσότεροι του ενός ανεξάρτητοι ιστορικοί μάρτυρες· ένας μεμονωμένος μάρτυρας μπορεί να παρέχει χρονικό προφίλ, όχι εκτίμηση της διασποράς μεταξύ των εκτελεστών. Οι πρώιμοι μηχανικοί δίσκοι ελέγχονται ειδικά για αβεβαιότητα ταχύτητας, τοπικές μεταβολές ταχύτητας και λοιπές χρονικές παραμορφώσεις. Η αβεβαιότητα αυτή αποκλείει καταρχήν επιβεβαιωτική χρήση του απόλυτου τονικού ύψους και ενσωματώνεται στην αξιολόγηση της αξιοπιστίας των χρονικών μεγεθών. Οι ιστορικές ηχογραφήσεις δεν χρησιμοποιούνται για εκπαίδευση, ρύθμιση ή επιλογή μοντέλου και δεν αναπληρώνουν ανεπαρκή κύρια εκτελεστική κάλυψη· αποτελούν μόνο πρόσθετο διαχρονικό έλεγχο ανθεκτικότητας, εφόσον πληρούν τα αντίστοιχα κριτήρια ένταξης.

\subsection{Πρόσβαση, δεοντολογία και διάθεση}\label{app:corpora-rights}

Το ΠΕ0 καταγράφει ανά πηγή τους όρους πρόσβασης, επεξεργασίας και δημοσιοποίησης και την κατάσταση εξασφάλισής τους. Για τις συλλογές της Μονής και τα υπόλοιπα συμβολικά σώματα διακρίνονται τα βιβλιογραφικά μεταδεδομένα, οι νέες ερευνητικές επισημειώσεις, τα μουσικά αποσπάσματα και τα πρωτογενή αρχεία. Για τυχόν ιστορικές ηχογραφήσεις ελέγχονται χωριστά η κατάσταση του πρωτογενούς τεκμηρίου, οι όροι της συγκεκριμένης ψηφιακής μεταφοράς ή επανέκδοσης και το επιτρεπόμενο καθεστώς ερευνητικής χρήσης. Το μητρώο προσδιορίζει ποια στοιχεία μπορούν να διατεθούν ανοικτά και ποια απαιτούν ελεγχόμενη πρόσβαση ή πρόσθετη άδεια. Η δημόσια πρόσβαση σε μια πηγή καταγράφεται χωριστά από το δικαίωμα αναδιανομής της.

Η χρήση των ηχογραφήσεων του ΔΑΜΑΣΚΗΝΟΣ οργανώνεται με τον αρμόδιο φορέα, με έλεγχο των αρχικών όρων συλλογής και των επιτρεπόμενων ερευνητικών χρήσεων. Πριν από την επεξεργασία, ο αρμόδιος θεσμικός έλεγχος τεκμηριώνει τη νόμιμη βάση, την ελαχιστοποίηση δεδομένων, τους δικαιούχους πρόσβασης, τον χρόνο διατήρησης και τις εγγυήσεις επιστημονικής έρευνας σύμφωνα με τα άρθρα 5, 6 και 89 του ΓΚΠΔ.\footnote{Κανονισμός (ΕΕ) 2016/679, επίσημο κείμενο: \url{https://eur-lex.europa.eu/eli/reg/2016/679/oj}.} Τα αναλυτικά αρχεία χρησιμοποιούν ψευδωνυμοποιημένα αναγνωριστικά εκτελεστών και το αρχείο αντιστοίχισης τηρείται χωριστά με περιορισμένη πρόσβαση. Η πρωτογενής φωνή και το ηλεκτρογλωττογράφημα διατίθενται μόνο σύμφωνα με τους τεκμηριωμένους όρους τους. Η αρχική ενημέρωση και συγκατάθεση των ψαλτών ελέγχεται ρητά ως προς το αν καλύπτει τις σχεδιαζόμενες αναλύσεις και, όπου απαιτείται, εξασφαλίζεται η πρόσθετη συγκατάθεση ή θεσμική έγκριση πριν από τη χρήση των δεδομένων.

Ο πιλοτικός έλεγχος των οδηγιών της §~\ref{app:corpora-annotation-pilot} πραγματοποιείται μετά από ενημέρωση και συγκατάθεση των συμμετεχόντων. Τηρούνται μόνο τα στοιχεία που απαιτούνται για την αναθεώρηση του πρωτοκόλλου και οι πιλοτικές επισημειώσεις δεν εντάσσονται στα ερευνητικά σύνολα δεδομένων ή στις επιβεβαιωτικές αναλύσεις.

Οι ανεξάρτητοι ειδικοί ενημερώνονται για τον σκοπό, τη διαδικασία, τη χρήση των επισημειώσεων και τους όρους συμμετοχής τους, και η σχετική συναίνεση και θεσμική διαδικασία ολοκληρώνονται πριν από τη μελέτη. Το πακέτο αναπαραγωγιμότητας περιλαμβάνει κώδικα, σχήματα δεδομένων, πρωτόκολλο και τα επιτρεπόμενα μεταδεδομένα/παράγωγα. Για υλικό με περιορισμούς παρέχει σταθερές παραπομπές και διαδικασία νόμιμης πρόσβασης. Οι όροι αυτοί εντάσσονται στο μητρώο σώματος μαζί με την προέλευση και το καθεστώς ένταξης στο σώμα.

\subsection{Εξωτερικά μουσικολογικά πρότυπα}\label{app:corpora-extrefs}

Η μουσικολογική οριοθέτηση και η οικογενειακή αναφορά οργανώνονται σε \textbf{τρία συμπληρωματικά στρώματα}, καθένα με διακριτό ερευνητικό ρόλο και πλήρη καταγραφή προέλευσης:

\begin{enumerate}
\item \(C_{\mathrm{StA}}\): \textbf{καταλογογραφικό στρώμα υποψηφίων τύπων, παραλλαγών και προέλευσης} που εξάγεται από την ευρύτερη υποδομή $\mathcal S_{\mathrm{StA}}$. Αποτυπώνει την οργάνωση του καταλόγου κατά ήχο, λειτουργική/καταληκτική κατηγορία και ρητά επισημασμένες παραλλαγές, με χωριστή καταγραφή της κάλυψης ανά λειτουργικό ρόλο. Ως σημείο αναφοράς σε επίπεδο οικογένειας χρησιμοποιείται μόνο το διατηρημένο εκτός εκπαίδευσης υποσύνολο που έχει περάσει το κριτήριο επαληθευμένης κοινής μουσικής ταυτότητας. Το υπόλοιπο $C_{\mathrm{StA}}$ παραμένει στρώμα υποψηφίων σχέσεων και διαδρομή προς τα πλήρη έργα και τις πηγές.
\item \(C_{\mathrm{occ}}\): \textbf{στρώμα τεκμηριωμένων εμφανίσεων και δημοσιευμένων τυπολογικών προτάσεων} από αναλύσεις συγκεκριμένων έργων. Κύριες γραμμές είναι οι ένδεκα τυπολογικές περιπτώσεις του Gheorghiță σε σαράντα δύο Κοινωνικά και οι χαρτογραφήσεις εμφανίσεων του Χαραλαμπίδη σε τριάντα Κοινωνικά. Οι ένδεκα περιπτώσεις του Gheorghiță δεν θεωρούνται αυτομάτως ένδεκα επικυρωμένες οικογένειες. Παρέχουν πραγματικά μουσικά παραδείγματα, υποψήφιες οριοθετήσεις, λειτουργικές παρατηρήσεις και υλικό για συγκρότηση ή έλεγχο υπογραφών.
\item \(C_{\mathrm{exp}}\): \textbf{ανεξάρτητο στρώμα ειδικών} από την ειδικά σχεδιασμένη τυφλή μελέτη 3--5 αξιολογητών. Παρέχει νέα επισημείωση, μετρημένη αξιοπιστία μεταξύ αξιολογητών και εμπειρική κλίμακα αναφοράς για το επίπεδο συμφωνίας που επιτυγχάνεται από ανεξάρτητους γνώστες της παράδοσης.
\end{enumerate}

Τα τρία στρώματα διατηρούν αυτοτελή προέλευση και συνδέονται μέσω \textbf{ρητού πίνακα αντιστοίχισης υποψηφίων τύπων, επικυρωμένων οικογενειών και εμφανίσεων με διατήρηση της προέλευσης}. Για κάθε πηγή αναφέρονται χωριστά η κάλυψη, η δυνατότητα μεταφοράς στο κοινό επίπεδο αναφοράς και η κάλυψη ανά λειτουργικό ρόλο. Το διατηρημένο εκτός εκπαίδευσης \(C_{\mathrm{StA}}\) παρέχει καταλογογραφικές υποψήφιες σχέσεις και, στο επαληθευμένο υποσύνολο, οικογενειακό σημείο αναφοράς· το \(C_{\mathrm{occ}}\) εντοπίζει τεκμηριωμένες πραγματώσεις μέσα σε έργα και το \(C_{\mathrm{exp}}\) μετρά την ανεξάρτητη ανθρώπινη αναγνωρισιμότητα της ίδιας δομής. Η συνδυασμένη αξιολόγηση εξετάζει ποιες υποψήφιες ομαδοποιήσεις αποκτούν στήριξη από ανεξάρτητες πηγές και ειδικούς, με χωριστή αναφορά καταληκτικών και μη καταληκτικών θέσεων. Η χρήση της $\mathcal S_{\mathrm{StA}}$ ως υποδομής εκπαίδευσης και η χρήση του επαληθευμένου $C_{\mathrm{StA}}^{\mathrm{held}}$ ως μουσικολογικού σημείου αναφοράς καταγράφονται χωριστά στο μητρώο προέλευσης, ώστε κάθε αποτέλεσμα να έχει σαφή πηγή και ρόλο.

\subsection{Συμφωνία αναφορών}\label{app:corpora-agreement}

Τα τρία στρώματα αναφοράς και οι επιμέρους πηγές τους προέρχονται από \textbf{διακριτές αναλυτικές διαδικασίες}: καταλογογράφηση οικογενειών, δημοσιευμένη ανάλυση συγκεκριμένων έργων και ανεξάρτητη επισημείωση ειδικών. Η προέλευση, η κάλυψη και οι μεταξύ τους σχέσεις καταγράφονται ρητά. Η πολλαπλότητα των αναφορών γίνεται έτσι μετρήσιμος πόρος για την εκτίμηση της σταθερότητας της μουσικολογικής αναγνώρισης σε επίπεδο ορίου, οικογένειας και λειτουργικού ρόλου. Η επιλογή να αναφέρεται ρητά η συμφωνία και να διατηρούνται οι ανεπίλυτες περιπτώσεις ακολουθεί τη γενικότερη μεθοδολογική αρχή ότι η διαφωνία των επισημειωτών μπορεί να αποτελεί πληροφορία για το ίδιο το αντικείμενο και όχι απλώς θόρυβο προς εξάλειψη (Artstein \& Poesio, 2008· Uma κ.ά., 2021· Plank, 2022· Davani κ.ά., 2022).

Στο υλικό όπου δύο ή περισσότερες αναφορές επικαλύπτονται, μετριέται η \textbf{μεταξύ τους συμφωνία} ως προς τα όρια των εμφανίσεων, την οικογενειακή απόδοση και, όπου παρέχεται, τον λειτουργικό ρόλο. Χρησιμοποιούνται τα ίδια μέτρα που εφαρμόζονται στην αξιολόγηση του μοντέλου: F1 ορίων με ανοχή γεγονότος, δείκτης τομής προς ένωση, και συντελεστής συμφωνίας κατάλληλος για ανισομερείς κατανομές κατηγοριών. Όπου η αναφορά δεν προσδιορίζει μοναδικά όρια, αποθηκεύεται το σύνολο των αποδεκτών ζευγών αρχής--τέλους $\mathcal B_x$ και η σύγκριση διατηρεί τους περιορισμούς που συνδέουν τις επιτρεπτές οριοθετήσεις μεταξύ τους.

Επιβεβαιωμένη επικάλυψη καταγράφεται όταν δύο διακριτές, συνυπάρχουσες εμφανίσεις στο ίδιο επίπεδο ανάλυσης απαιτούν κοινά γεγονότα σε όλες τις αποδεκτές κοινές ερμηνείες ή όταν υπάρχει ρητή τεκμηρίωση κοινής συμμετοχής. Όταν η κοινή συμμετοχή εμφανίζεται μόνο σε ορισμένες αποδεκτές ερμηνείες, η επικάλυψη καταγράφεται ως εξαρτώμενη από την ερμηνεία. Η γεωμετρική τομή δύο αβέβαιων περιοχών δεν αρκεί από μόνη της για να τεκμηριώσει πραγματική επικάλυψη.

Η μέτρηση αυτή έχει τρεις συνέπειες.

\textbf{Πρώτον, δίνει εμπειρική κλίμακα αναφοράς στη μουσικολογική δοκιμή.} Η επίδοση του μοντέλου ερμηνεύεται σε σχέση με την παρατηρούμενη συμφωνία των ίδιων των μουσικολογικών αναφορών· έτσι μια τιμή, για παράδειγμα \(0{,}70\), αποκτά νόημα μέσα στην πραγματική κλίμακα συμφωνίας μεταξύ ανεξάρτητων αξιολογητών στο συγκεκριμένο υλικό.

Η κλίμακα αυτή λειτουργεί ως \textbf{συγκριτικό σημείο αναφοράς}. Η σχέση
\[
\mathrm{Agr}(M,R_1)>\mathrm{Agr}(R_1,R_2)
\]
ερμηνεύεται ως ισχυρή προσαρμογή του μοντέλου στην αναλυτική σύμβαση της \(R_1\) και αξιολογείται ταυτόχρονα έναντι της \(R_2\) και των υπόλοιπων στρωμάτων. Με αυτόν τον τρόπο η αξιολόγηση διακρίνει την αναπαραγωγή μιας συγκεκριμένης αναλυτικής σύμβασης από τη γενίκευση της οικογενειακής ταυτότητας σε περισσότερα μουσικολογικά συστήματα αναφοράς.

\textbf{Δεύτερον, αποτελεί αυτοτελές μουσικολογικό εύρημα.} Η συμφωνία ποσοτικοποιεί πόσο συνεπής είναι η αναγνώριση θέσεων μεταξύ ανεξάρτητων αξιολογητών στο συγκεκριμένο δείγμα και πρωτόκολλο. Υψηλή συμφωνία δείχνει σταθερή αναγνώριση υπό αυτές τις συνθήκες, ενώ συστηματικές αποκλίσεις χαρτογραφούν τα σημεία όπου η οικογενειακή ή λειτουργική απόδοση διαφοροποιείται μεταξύ αναλυτικών πλαισίων.

\textbf{Τρίτον, παρέχει δομημένο στόχο για την αβεβαιότητα και την αποχή του μοντέλου.} Για την οικογενειακή απόδοση, η εντροπία της προβλεπτικής κατανομής (predictive entropy) και η διαφορά πιθανότητας των δύο επικρατέστερων κλάσεων συγκρίνονται πρωτίστως με τη διαφωνία ανά περίπτωση στο ανεξάρτητο $C_{\mathrm{exp}}$ και, χωριστά, με τις αποκλίσεις μεταξύ των στρωμάτων αναφοράς που διατηρούν διακριτή την προέλευσή τους. Για τα όρια, οι οριακές πιθανότητες του τμηματικού μοντέλου συγκρίνονται με τη διασπορά των ανεξάρτητων ανθρώπινων ορίων. Η επιλεκτική αξιοπιστία αξιολογείται με καμπύλες κινδύνου--κάλυψης και το αντίστοιχο εμβαδόν (AURC), ώστε να ελεγχθεί αν η αποχή απομακρύνει πράγματι τις μουσικολογικά δυσδιάκριτες περιπτώσεις. Η αντιστοιχία προβλεπτικής αβεβαιότητας και μουσικολογικής διαφωνίας αναφέρεται ως αυτοτελές αποτέλεσμα και διαχωρίζεται από την πιθανοτική βαθμονόμηση, η οποία αξιολογείται χωριστά με κατάλληλο ορθό κανόνα βαθμολόγησης (proper scoring rule), όπου υπάρχει σαφώς ορισμένος πιθανοτικός στόχος. Η ιστορική παραλλαγή διακρίνεται από την αβεβαιότητα: μεγάλη ενδοοικογενειακή διαφοροποίηση μπορεί να συνυπάρχει με υψηλή βεβαιότητα οικογενειακής ταυτότητας.

Η σύγκριση οικογενειακής απόδοσης μεταξύ διαφορετικών αναφορών χρησιμοποιεί ρητό \textbf{κανόνα αντιστοίχισης με διατήρηση προέλευσης}. Οι οικογένειες παραμένουν συνδεδεμένες με το δικό τους αναλυτικό σύστημα και, για κάθε ζεύγος οικογενειών \(F_i^{(A)}\) και \(F_j^{(B)}\) καταγράφεται σχέση από κλειστό σύνολο: ισοδύναμη, ευρύτερη, στενότερη, επικαλυπτόμενη, ανεπίλυτη. Οι σχέσεις παράγονται στο δεύτερο πακέτο εργασίας, τεκμηριώνονται και δημοσιεύονται· η κατηγορία \emph{ανεπίλυτη} παραμένει θεμιτή. Οι αντίστοιχες περιπτώσεις εξαιρούνται μόνο από τις μετρικές που απαιτούν οριστική οικογενειακή απόδοση, αλλά \textbf{παραμένουν πάντοτε στον παρονομαστή της αναφοράς κάλυψης} και αναφέρονται χωριστά ως μουσικολογική αβεβαιότητα, ώστε η συμφωνία να μην βελτιώνεται τεχνητά με αφαίρεση των δύσκολων περιπτώσεων.

Η μέτρηση προϋποθέτει επικάλυψη μεταξύ των αναφορών, η οποία υπολογίζεται στο πρώτο πακέτο εργασίας μαζί με τα υπόλοιπα κριτήρια κάλυψης και βαθμονόμησης. Όπου η επικάλυψη είναι μικρή, η κλίμακα εκτιμάται στο πραγματικό κοινό υποσύνολο.

\subsection{Εγχειρίδιο επισημείωσης και πλήρως ελεγμένο υποσύνολο ανάπτυξης}\label{app:corpora-annotation-pilot}

Το εγχειρίδιο διαχωρίζει τέσσερα επίπεδα. Το μητρώο τύπων ορίζει την υπογραφή πριν από την κρίση εμφανίσεων. Το μητρώο εμφανίσεων καταγράφει όρια, διαγνωστικά γνωρίσματα και $J(x,T;r)$. Ο γράφος ιστορικών σχέσεων καταγράφει εξήγηση, ομολογία χωρίου και εξάρτηση πηγών. Τα παράγωγα ζεύγη B2/EISE δημιουργούνται μόνο μετά τους προηγούμενους ελέγχους.

Για τις πρώτες πιλοτικές κάρτες χρησιμοποιούνται τρεις υποψήφιες περιπτώσεις: παρακλητική, κύλισμα και πίασμα. Στόχος δεν είναι να επιβεβαιωθούν εκ των προτέρων ως οικογένειες. Στόχος είναι να ελεγχθεί αν μπορούν να αποκτήσουν σαφή και εφαρμόσιμη υπογραφή. Κάθε κάρτα περιλαμβάνει πραγματικά θετικά παραδείγματα, δύσκολα αρνητικά και ρητές συνθήκες αποχής. Οι αριθμοί παραδειγμάτων ορίζονται από την απογραφή και τη χρονική δοκιμή του φόρτου. Δεν λειτουργούν ως ήδη επικυρωμένο στατιστικό κατώφλι.


Στο ΠΕ0, πριν από οποιαδήποτε εκπαίδευση ή σύγκριση μοντέλων της M1, συντάσσεται γραπτό εγχειρίδιο επισημείωσης που ορίζει τι θεωρείται θέση, πώς καταγράφονται τα αποδεκτά ζεύγη ορίων, πώς διακρίνεται η πραγματική επικάλυψη από εναλλακτικές ερμηνείες και πώς καταγράφονται το συνδετικό υλικό, το τεκμηριωμένο υπόβαθρο και οι ανεπίλυτες περιοχές. Με βάση το εγχειρίδιο συγκροτείται \textbf{πλήρως ελεγμένο υποσύνολο ανάπτυξης} από αντιπροσωπευτικά έργα ή συνεχόμενα αποσπάσματα. Κάθε περιοχή εξετάζεται και καταγράφεται ως τεκμηριωμένη εμφάνιση, χωριστό συνδετικό υλικό, τεκμηριωμένο υπόβαθρο ή ανεπίλυτη περιοχή, με χωριστή διατήρηση τυχόν κοινών γεγονότων δύο εμφανίσεων. Ο πλήρης έλεγχος σημαίνει ότι έχει καταγραφεί η κατάσταση τεκμηρίωσης κάθε περιοχής και όχι ότι κάθε μουσικολογική αμφιβολία έχει υποχρεωτικά λυθεί. Η έκταση του υποσυνόλου καθορίζεται από την απογραφή του ΠΕ0 ώστε να καλύπτει διαφορετικές οικογένειες, ήχους, καταληκτικές και μη καταληκτικές θέσεις, ποικιλία μήκους και επαρκές τεκμηριωμένο υπόβαθρο, χωρίς να προκαθορίζεται αυθαίρετος αριθμός έργων. Η κύρια επισημείωση πραγματοποιείται από τον υποψήφιο χωρίς πρόσβαση σε αποτελέσματα μοντέλων. Προκαθορισμένο μέρος του ίδιου υποσυνόλου επισημειώνεται ανεξάρτητα από δεύτερο ειδικό, κατά προτίμηση διαφορετικής εκπαιδευτικής ή ψαλτικής προέλευσης και διαφορετικό από την επιβεβαιωτική ομάδα του $C_{\mathrm{exp}}$. Η συμφωνία αναφέρεται χωριστά για τα ζεύγη ορίων, την οικογενειακή απόδοση και την κατάσταση της περιοχής· οι διαφωνίες επιλύονται με καταγεγραμμένο κανόνα ή διατηρούνται ως ανεπίλυτες όταν η τεκμηρίωση δεν επιτρέπει μοναδική απόφαση. Η εμπειρική διαφωνία χρησιμοποιείται για την αιτιολόγηση του περιθωρίου πρακτικής ισοδυναμίας στον μεταγενέστερο έλεγχο της τμηματικής διαδικασίας. Το υποσύνολο και το εγχειρίδιο οριστικοποιούνται πριν από τη M1 και χρησιμοποιούνται για ακρίβεια, ανάκληση και F1 ορίων και τμημάτων, για την εκτίμηση της εμπειρικής κάλυψης και για τη βαθμονόμηση του τεκμηριωμένου υποβάθρου. Η απογραφή της επικάλυψης αναφέρεται συνολικά και ανά οικογένεια, λειτουργικό ρόλο, έργο, ιστορικό στρώμα και κατηγορία μήκους, μαζί με την έκταση των ανεπίλυτων περιπτώσεων. Ασθενής κοινός περιορισμός κάλυψης ενεργοποιείται μόνο αν ο πιλοτικός έλεγχος δείξει συστηματική υπερκάλυψη· διαφορετικά δεν αποτελεί μέρος της βασικής αντικειμενικής συνάρτησης. Προαιρετικά, το εγχειρίδιο επιτρέπει να καταγράφεται το σημειογραφικό ή μελικό γνώρισμα που στήριξε την αναγνώριση μιας θέσης, χωριστά από την πλήρη έκταση της εμφάνισης, αποκλειστικά για έλεγχο της επισημείωσης και ερμηνεία σφαλμάτων.

Χωριστά, πριν από τη συλλογή του $C_{\mathrm{exp}}$, η αρχική εκδοχή των οδηγιών της ανεξάρτητης μελέτης ειδικών δοκιμάζεται σε \textbf{2--4 ενεργούς ψάλτες} που δεν συμμετέχουν στην επιβεβαιωτική ομάδα. Χρησιμοποιείται μικρό σύνολο χωρίων αποκλειστικά από το υλικό ανάπτυξης. Οι συμμετέχοντες εκτελούν πρώτα την ίδια εργασία που προβλέπεται για τους ανεξάρτητους ειδικούς ---εντοπισμό ορίων, οικογενειακή απόδοση και, όπου τεκμηριώνεται, λειτουργικό ρόλο, με δυνατότητα αβέβαιης απάντησης--- χωρίς προηγούμενη συζήτηση των θεωρητικών κριτηρίων. \textbf{Μόνο μετά την ολοκλήρωση της εργασίας} ακολουθεί σύντομη αποσαφήνιση για όρους που παρερμηνεύθηκαν, σημεία στα οποία οι οδηγίες δεν επαρκούσαν και διατυπώσεις που δεν ήταν μουσικολογικά φυσικές. Οι παρατηρήσεις χρησιμοποιούνται αποκλειστικά για γλωσσική και διαδικαστική αναθεώρηση του πρωτοκόλλου. Οι πιλοτικές επισημειώσεις απορρίπτονται από την επιβεβαιωτική αξιολόγηση, δεν αποτελούν ετικέτες αναφοράς και δεν χρησιμοποιούνται στην επιλογή αναπαράστασης ή στην εκπαίδευση του EISE. Η τελική έκδοση των οδηγιών οριστικοποιείται πριν από το άνοιγμα του $C_{\mathrm{exp}}$.

\subsection{Τυφλή μελέτη ειδικών}\label{app:corpora-experts}

Οι δημοσιευμένες αναλύσεις συμπληρώνονται από \textbf{ειδικά σχεδιασμένη τυφλή μελέτη 3--5 ανεξάρτητων ειδικών} πάνω σε προκαθορισμένο υποσύνολο που δεν χρησιμοποιείται για ανάπτυξη ή ρύθμιση των μοντέλων. Οι οδηγίες της ανεξάρτητης μελέτης ειδικών οριστικοποιούνται μετά τον χωριστό πιλοτικό έλεγχο της §~\ref{app:corpora-annotation-pilot}. Οι 2--4 συμμετέχοντες του πιλοτικού ελέγχου δεν συμμετέχουν στο $C_{\mathrm{exp}}$ και οι πιλοτικές τους επισημειώσεις δεν μετατρέπονται σε ετικέτες ή πρότυπα αναφοράς. Οι ειδικοί επισημειώνουν όρια εμφανίσεων, οικογενειακή ταυτότητα και, όπου τεκμηριώνεται επαρκώς, λειτουργικό ρόλο. Εργάζονται ανεξάρτητα, χωρίς πρόσβαση στις προβλέψεις των μοντέλων και χωρίς αμοιβαία γνώση των επισημειώσεων. Η κατηγορία \emph{ανεπίλυτη} παραμένει διαθέσιμη ώστε η πραγματική μουσικολογική αβεβαιότητα να καταγράφεται αντί να εξαναγκάζεται σε τεχνητή συμφωνία. Το ποσοστό ανεπίλυτων περιπτώσεων αναφέρεται ως αυτοτελές μέτρο κάλυψης ακόμη όταν οι συγκεκριμένες περιπτώσεις δεν εισέρχονται σε μετρική συμφωνίας που απαιτεί οριστική ετικέτα.

Η μελέτη παράγει \textbf{μετρημένη αξιοπιστία μεταξύ αξιολογητών} και εμπειρική κλίμακα αναφοράς για τη μουσικολογική οριοθέτηση και την οικογενειακή αναφορά. Η σύνθεση της ομάδας επιδιώκει εύρος εκπαιδευτικών και ερευνητικών καταβολών· τυχόν συστηματικές διαφορές μεταξύ σχολών εξετάζονται ως δευτερεύον μουσικολογικό εύρημα. Ο υποψήφιος σχεδιάζει και τεκμηριώνει το πρωτόκολλο, ενώ το επιβεβαιωτικό σύνολο επισημειώνεται αποκλειστικά από την ανεξάρτητη ομάδα ειδικών.

\subsection{Θέση του ερευνητή}\label{app:corpora-researcher}

Η μακρόχρονη πρακτική ενασχόληση του υποψηφίου με τη βυζαντινή ψαλτική αξιοποιείται ως \textbf{εξειδικευμένη γνώση πεδίου} για τη διατύπωση της εννοιολογικής κατασκευής, τον σχεδιασμό του σχήματος δεδομένων, την επιλογή μουσικολογικά ουσιαστικών συγκρίσεων και την ερμηνεία των σφαλμάτων. Η επιβεβαιωτική μέτρηση οργανώνεται σε ανεξάρτητα στρώματα αναφοράς, τυφλή επισημείωση ειδικών και προκαθορισμένα πρωτόκολλα. Ο συνδυασμός βιωματικής και εξειδικευμένης γνώσης της παράδοσης στη διαμόρφωση της υπόθεσης και ανεξάρτητης μέτρησης στην αξιολόγηση καθιστά τη θέση του ερευνητή \textbf{ρητό και ελέγξιμο στοιχείο του μεθοδολογικού σχεδιασμού}.



\subsection{Αναλυτική βιβλιογραφική θεμελίωση}\label{app:literature-full}

Η βιβλιογραφική θεμελίωση ακολουθεί δύο συγκλίνουσες γραμμές. Η πρώτη προσδιορίζει το μουσικό αντικείμενο και τις ιστορικά τεκμηριωμένες σχέσεις που γίνονται μεταβλητές της υπολογιστικής μελέτης: θέση, εκδοχή, οικογένεια, ακολουθία και εξηγητική σχέση. Η δεύτερη εξετάζει τις διαθέσιμες μεθόδους αναπαράστασης, τμηματοποίησης, ανάκτησης, κατηγοριοποίησης και πρόβλεψης με τις οποίες αυτές οι μεταβλητές μπορούν να ελεγχθούν. Το ερευνητικό κενό προκύπτει ακριβώς από το σημείο όπου οι δύο γραμμές συναντώνται.

Η σειρά της ανασκόπησης είναι επομένως λειτουργική: πρώτα ορίζεται με ιστορική και αναλυτική ακρίβεια η μονάδα και οι μορφές συγγένειάς της, κατόπιν εξετάζεται τι ανακτούν ήδη οι υπολογιστικές μέθοδοι και, τέλος, ποια ανεξάρτητα ανθρώπινα και μουσικά δεδομένα μπορούν να κρίνουν την εγκυρότητα της αναπαράστασης. Η μουσικολογική θεμελίωση στη βυζαντινή μουσική δεν αποτελεί προοίμιο ξένο προς το MIR· ορίζει τις κατηγορίες, τις σχέσεις και τα εξωτερικά τεκμήρια πάνω στα οποία στηρίζονται οι αποφασιστικές συγκρίσεις της ερευνητικής πρότασης.

\subsubsection{Η μουσική θέση ως ιστορική και αναλυτική μονάδα}

Η \emph{θέσις} απαντά ως ιστορική μουσικολογική έννοια σε διαφορετικά σημειογραφικά και διδακτικά περιβάλλοντα. Το κεντρικό χωρίο του Μανουήλ Χρυσάφη, στην έκδοση και μετάφραση του Dimitri E.~Conomos (1985b, σσ.~40--41), συνδέει την ένωση σημείων με τη συγκρότηση μέλους. Η πηγή τεκμηριώνει οργανωμένο μελικό σχηματισμό, ενώ ο σημερινός επιχειρησιακός ορισμός της οικογένειας απαιτεί δικά του ανεξάρτητα κριτήρια.

Ο \textbf{Γαβριήλ ιερομόναχος} προσφέρει μια δεύτερη, ιδιαίτερα χρήσιμη θεωρητική διατύπωση. Στην κριτική έκδοση των Christian Hannick και Gerda Wolfram (1985, σ.~72), οι θέσεις παρουσιάζονται ως σχηματισμοί που προκύπτουν από συνδυασμούς φωνητικών και άφωνων σημείων\footnote{Στην ιστορική ορολογία, οι φωνητικοί/ποσοτικοί χαρακτήρες δηλώνουν τη βασική μελωδική κίνηση, ενώ άφωνα ή ποιοτικά σημεία μεταφέρουν άλλη σημειογραφική/εκτελεστική πληροφορία. Οι κατηγορίες μεταβάλλονται ιστορικά και δεν εξισώνονται μηχανικά μεταξύ παλαιάς γραφής και Νέας Μεθόδου.} και παρομοιάζονται με τις λέξεις της γραμματικής. Η αναλογία αυτή έχει άμεση σημασία για την αναπαράσταση: η ίδια η ψαλτική θεωρία περιγράφει μία ενδιάμεση μουσική μονάδα ανάμεσα στο μεμονωμένο σημείο και στη σύνθεση.

Στον \textbf{Χρύσανθο τον εκ Μαδύτων}, το \emph{Θεωρητικόν Μέγα} (§400) συνδέει τη διδασκαλική μετάδοση με τον τρόπο μελοποιίας. Στην §408, αναφερόμενος στην παλαιά γραφή, περιγράφει ότι οι ίδιες υποστάσεις μπορούσαν να μεταβάλλουν τη δύναμη ή την ενέργειά τους ανάλογα με τη βαθμίδα εφαρμογής τους. Η μαρτυρία αυτή τεκμηριώνει ιστορικά εξάρτηση της εκφοράς από το μουσικό περιβάλλον· η μεταφορά της από την υπόσταση στη σύγχρονη ερευνητική έννοια της οικογένειας θέσεων αποτελεί αναλογική υπόθεση που ελέγχεται και δεν θεωρείται δεδομένη.

Η μεσοβυζαντινή σημειογραφία και η τροπική της γραμματική παρέχουν το ευρύτερο πλαίσιο μέσα στο οποίο λειτουργούν αυτές οι μονάδες. Ο \textbf{Jørgen Raasted} (1966), ο \textbf{Christian Troelsgård} (2011) και οι \textbf{Maria Alexandru} και Christian Troelsgård (2020) τεκμηριώνουν τη λειτουργία των σημείων, των συνδυασμών τους και της τροπικής οργάνωσης, χωρίς να επιβάλλουν έναν ενιαίο διαχρονικό ορισμό της θέσεως. Η λειτουργική ταξινόμηση των εμφανίσεων ---όπως εναρκτήρια, συνδετική ή καταληκτική χρήση--- καταγράφεται ως χωριστό αναλυτικό επίπεδο και συνδέεται με συγκεκριμένα τεκμήρια ανά περίπτωση. Οι λεπτότερες κατηγορίες δεν αποδίδονται αναδρομικά στην ιστορική θεωρία χωρίς άμεση επαλήθευση της αντίστοιχης πηγής.

Νεότερη αναλυτική προσέγγιση της Νέας Μεθόδου (Δεσπότης, 2007a) περιγράφει τη μελωδικο-ρυθμική λειτουργία των σημαδοφώνων ως σχεσιακή: εξαρτάται από τον συνδυασμό τους με γειτονικά στοιχεία, την τοποθέτησή τους μέσα στη μουσική φράση και την ένταξή τους σε ευρύτερες μορφολογικές ενότητες. Η διάκριση ανάμεσα σε τοπικό συνδυασμό, μουσική φράση και ευρύτερη μορφολογική ενότητα παρέχει ειδικό μουσικολογικό κίνητρο για την πειραματική σύγκριση των κλιμάκων αναπαράστασης της M1, χωρίς να προδικάζει ποια κλίμακα θα αποδειχθεί υπολογιστικά καταλληλότερη.

Για το εμπειρικό σώμα της ερευνητικής πρότασης, η \textbf{μουσική θέση} εξετάζεται ως ιστορικά προσδιορισμένος, οργανωμένος μελικός σχηματισμός μεταβλητής έκτασης, που μπορεί να αναλαμβάνει διαφορετικές φραστικές και τροπικές λειτουργίες αναλόγως τύπου και εμφάνισης. Η \textbf{αναλυτική οικογένεια} συγκροτείται χωριστά ως ελέγξιμος τύπος συγγενών πραγματώσεων, με προδηλωμένη υπογραφή και ανεξάρτητη επικύρωση. Ο ορισμός συνδέει την ιστορική θεωρητική έννοια με τη νεότερη αναλυτική πρακτική σε ένα πεδίο όπου τα συγκεκριμένα μουσικά χωρία μπορούν να οριοθετηθούν, να συγκριθούν και να συνδεθούν με πηγή, έργο, ήχο και λειτουργικό ρόλο.

Η μεγάλης κλίμακας έρευνα του \textbf{Κωνσταντίνου Χαριλ. Καραγκούνη} (2003) στα Χερουβικά προσφέρει το κεντρικό νεότερο μουσικολογικό υπόβαθρο. Η διατριβή του καταγράφει εκτεταμένο σώμα μελοποιήσεων, μελουργών και χειρογράφων, εντοπίζει διαχρονικές επανεμφανίσεις ονομασμένων θέσεων και καταγράφει παραλλαγές, διαφορετικές τροπικές πραγματώσεις και ιστορικές σχέσεις μετάδοσης. Στη σ.~722 ο ίδιος ο συγγραφέας οριοθετεί τη λεπτομερή διεξαγωγή αυτού του είδους συγκρίσεων ως αντικείμενο αυτοτελούς μορφολογικής έρευνας. Η παρούσα ερευνητική πρόταση τοποθετείται σε αυτή τη συνέχεια και μετατρέπει τη σχέση θέσης--πραγμάτωσης--παραλλαγής σε ρητό, επαναλήψιμο σύστημα δεδομένων. Με αυτό το σύστημα ελέγχεται αν οι μουσικολογικά τεκμηριωμένες διακρίσεις ανακτώνται υπολογιστικά και εφαρμόζονται με συνέπεια από ανεξάρτητους ειδικούς.

Η μουσική θέση ως ενδιάμεσο επίπεδο μεταξύ γραφικού σημείου και ευρύτερης μουσικής οργάνωσης αποτελεί ιστορικά ενημερωμένη επιλογή μονάδας· η υπεροχή της έναντι άλλων υπολογιστικών κλιμάκων εξετάζεται στη M1 και δεν συνάγεται από τη γραμματική αναλογία.

\paragraph{Μελικά πρότυπα και τροπική θεωρία στον Αλυγιζάκη.}
Ο Αντώνιος Ε.~Αλυγιζάκης, στο \emph{Η οκταηχία στην ελληνική λειτουργική υμνογραφία} (έντυπη έκδοση 1985, σσ.~81--82), περιγράφει τις θέσεις ως μελικά πρότυπα που συγκροτούν φράσεις και μεγαλύτερες μελοποιητικές ενότητες. Στο ίδιο έργο (ιδίως σσ.~149--158, 198--201 και 235) εξετάζονται ιστορικές θεωρίες τροπικής οργάνωσης και οι διακρίσεις ανάμεσα σε παλαιότερη σημειογραφία, εξήγηση και νεότερες καταγραφές. Η σ.~235 περιλαμβάνει, κατά τη διαθέσιμη αποδελτίωση, εκτενέστερη εκδοχή της \emph{Ἑρμηνείας τῆς Παραλλαγῆς} του Πλουσιαδηνού από τις καταχωρισμένες μεταγραφές του Μποτωνάκη (2014, σ.~352). Η ύπαρξη της ίδιας φράσης σε διαφορετικές δημοσιευμένες αποδόσεις οδηγεί σε αίτημα παραβολής των κωδίκων Διονυσίου~570 και ΕΒΕ~968 και όχι σε εκ των προτέρων συμπέρασμα για διαφορετική χειρογραφική παράδοση. Για τον παρόντα σχεδιασμό η πηγή θεμελιώνει τη σχέση μελικού τύπου, τροπικής λειτουργίας και σύνθεσης, ενώ οι περιορισμοί της οικογενειακής ένταξης ορίζονται από τα δεδομένα της έρευνας.

\subsubsection{Μορφολογία, επανάληψη και οργανωμένη παραλλαγή}

Οι μελέτες του Κ.~Χ. Καραγκούνη (2003, 2005) παρέχουν συγκεκριμένα παραδείγματα άμεσα συναφή με το πρόβλημα της οικογένειας θέσεων. Στη σ.~788 της διατριβής του, περιπτώσεις με κοινή ή πολύ συγγενή μετροφωνική\footnote{Μετροφωνία: η βασική ανάγνωση της μελωδικής κίνησης που προκύπτει από τα φωνητικά/ποσοτικά σημεία. Στην παλαιότερη γραφή η μετροφωνική συγγένεια δεν αρκεί από μόνη της για να ορίσει ταυτότητα θέσης. Βλ. Καραγκούνης 2003, σ.~788.} επιφάνεια διακρίνονται ως διαφορετικές θέσεις όταν συνδέονται με διαφορετικές μεγάλες υποστάσεις\footnote{Μεγάλες υποστάσεις: ιστορικά σημειογραφικά στοιχεία της παλαιότερης γραφής των οποίων η λειτουργία δεν εξαντλείται στην ποσοτική μετροφωνική κίνηση. Στη βιβλιογραφία συνδέονται με ποιοτική/εκτελεστική και μελική πληροφορία· η ακριβής λειτουργία τους εξετάζεται κατά περίπτωση. Βλ. Στάθης 1978· Καραγκούνης 2003.}. Στις σσ.~810--812 παρουσιάζονται ονομασμένες θέσεις, εξηγήσεις και τροπικά συμφραζόμενα· η μαρτυρία της σ.~811 περιγράφει ρητά σταθερότητα της μεταγραφής της παρακλητικής στις αντίστοιχες περιπτώσεις. Η διαφορά τύπων ή πραγματώσεων τεκμηριώνεται ανά συγκεκριμένο παράδειγμα και όχι από γενική υπόθεση μεταβλητότητας. Η μουσικολογική συνέπεια αποδίδεται συνοπτικά ως
\[
\boxed{\text{ισχυρή επιφανειακή ομοιότητα}\not\Rightarrow\text{ίδια θέση}}
\]
\[
\boxed{\text{διαφορετική αναλυτική επιφάνεια}\not\Rightarrow\text{διαφορετική θέση}}.
\]
Το πρώτο κριτήριο εμποδίζει την ταύτιση της θέσης με απλή μετροφωνική ή μελωδική ομοιότητα· το δεύτερο εμποδίζει την ταύτιση της οικογενειακής ταυτότητας με μία μοναδική αναλυτική ακολουθία. Μαζί ορίζουν το κεντρικό πρόβλημα μάθησης αναπαραστάσεων της ερευνητικής πρότασης: η κατάλληλη αναπαράσταση πρέπει να διαχωρίζει μουσικολογικά διαφορετικές θέσεις ακόμη και όταν η επιφανειακή τους μορφή είναι συγγενής, και να αναγνωρίζει κοινή μουσική ταυτότητα ακόμη και όταν οι συγκεκριμένες πραγματώσεις διαφοροποιούνται αισθητά.

Η σχέση θέσης και συλλαβής τεκμηριώνεται επίσης ως μεταβλητή και προς τις δύο κατευθύνσεις. Σε μορφολογικές αναλύσεις του ίδιου σώματος, μία εναρκτήρια θέση καλύπτει τις τρεις πρώτες συλλαβές «Οἱ τὰ Χε-», ενώ σε άλλα σημεία μία συλλαβή αναπτύσσεται σε περισσότερες διαδοχικές θέσεις (Καραγκούνης, 2003, σσ.~207, 209). Η παρατήρηση αυτή στηρίζει τη μεταβλητού μήκους αναπαράσταση της θέσης και τη χωριστή διατήρηση της συλλαβικής πληροφορίας στο κοινό συμβολικό επίπεδο αναφοράς.

Στην ίδια μορφολογική έρευνα καταγράφονται βραχέα μουσικά σχήματα με συνδετική λειτουργία ανάμεσα σε εκτενέστερες ακολουθίες θέσεων (Καραγκούνης, 2003, σ.~192). Το τεκμήριο παρέχει μουσικολογική βάση για τη διατήρηση συνδετικών περιοχών στον χάρτη πλήρους έργου· η κατάσταση κάθε συγκεκριμένου χωρίου καταγράφεται στη συνέχεια με τις χωριστές κατηγορίες του προκαθορισμένου πρωτοκόλλου επισημείωσης.

Στα \emph{Παραλειπόμενα περὶ τοῦ Χερουβικοῦ Ὕμνου} καταγράφονται, μεταξύ άλλων, ονομασμένες θέσεις ---Λυγίσματος, Ψηφιστοῦ καὶ Ἀντικενώματος και Κυλίσματος---\footnote{Λύγισμα, Κύλισμα και συγγενείς ονομασίες συνδέονται ιστορικά με ποιοτικούς χαρακτήρες/σημειογραφικές κατηγορίες, αρκετές από τις οποίες ανήκουν στην παλαιότερη γραφή. Η «θέση Λυγίσματος» ή «θέση Κυλίσματος» δηλώνει μελικό σχηματισμό που συνδέεται ιστορικά με την αντίστοιχη ονομασία και δεν ταυτίζεται με απλή παρουσία ενός συμβόλου. Το Ψηφιστόν είναι ποιοτικός χαρακτήρας που διατηρείται και στη Νέα Μέθοδο. Βλ. Καραγκούνης 2003/2005· Στάθης 1978.} και εκτενείς αλυσίδες θέσεων που επανέρχονται σε διαφορετικά σημεία του μέλους (Καραγκούνης, 2005, σσ.~157--163). Οι επαναλήψεις μπορούν να περιλαμβάνουν περιορισμένες προσθήκες ή απουσίες μεμονωμένων θέσεων, καθώς και μεταβολές στον συνδυασμό, στην κατάληξη, στην τροπική πορεία ή στην έκταση, χωρίς να χάνεται η αναγνωρίσιμη δομική σχέση. Το εύρημα δεν ταυτίζει την αλυσίδα με την οικογένεια θέσεων· δείχνει όμως γιατί η ταυτότητα μιας θέσης, η επανάληψη μιας ακολουθίας θέσεων και ο ανασυνδυασμός γνωστού μελικού υλικού πρέπει να παραμένουν αναλυτικά διακριτά.

Η ίδια μελέτη συνδέει την επιλογή και τον συνδυασμό των θέσεων με τη μορφή του έργου, τον συλλαβισμό και την καταληκτική οργάνωση, ενώ μικρότερα τμήματα ακολουθιών επανέρχονται σε διαφορετικό κειμενικό περιβάλλον (Καραγκούνης, 2005, σσ.~158--163). Η επανάληψη δεν αντιμετωπίζεται ως μηχανική αντιγραφή: σε διαφορετικά σημεία μπορεί να αλλάζει ο συνδυασμός, η κατάληξη, η τροπική πορεία ή η έκταση, διατηρώντας παράλληλα αναγνωρίσιμη δομική σχέση.

Η ανάλυση του Χερουβικού του Μανουήλ Αγαλλιανού ενισχύει αυτή τη διάκριση από διαφορετική κατεύθυνση: η μελοποιητική πρακτική του αντιπαραβάλλεται με εκείνη των Ιωάννη Γλυκύ και Ιωάννη Κουκουζέλη, με έμφαση στην αποφυγή εκτενών περιοδικών αλληλουχιών και στην επανεμφάνιση επιμέρους θέσεων σε διαφορετικούς συνδυασμούς (Καραγκούνης, 2005, σσ.~175--176). Για την παρούσα έρευνα, το παράδειγμα αυτό διαχωρίζει τρία επίπεδα που πρέπει να παραμείνουν αναλυτικά αυτόνομα: την \emph{ταυτότητα μιας θέσης}, την \emph{επανάληψη μιας ακολουθίας θέσεων} και τον \emph{ανασυνδυασμό γνωστού λεξιλογίου σε νέο συνθετικό περιβάλλον}.

Στις ίδιες αναλύσεις η παραλλαγή εμφανίζεται λειτουργικά οργανωμένη. Στον Αγαλλιανό, για παράδειγμα, η εναρκτήρια θέση διαφοροποιείται ανάλογα με τη σημειωμένη μαρτυρία\footnote{Μαρτυρία: σημειογραφική ένδειξη που δηλώνει ή επιβεβαιώνει τροπική και φθογγική κατάσταση· διακρίνεται από τη «μαρτυρία» με τη συνήθη ιστοριογραφική έννοια. Βλ. Χρύσανθος 1832.} ώστε να διατηρείται σταθερό το τονικό σημείο εκκίνησης του μέλους (Καραγκούνης, 2005, σ.~175). Αυτό προσφέρει μουσικολογικό προηγούμενο για την ιδέα ότι διαφορετικές επιφανειακές πραγματώσεις μπορούν να υπηρετούν κοινή λειτουργία. Η παρούσα έρευνα μεταφέρει αυτή την αρχή σε συστηματική μελέτη οικογενειών: το ζητούμενο δεν είναι απλώς να εντοπιστούν επαναλαμβανόμενα διαστήματα ή ακολουθίες, αλλά να προσδιοριστεί ποια χαρακτηριστικά στηρίζουν την κοινή μουσική ταυτότητα και σε ποιες περιοχές συγκεντρώνεται η τεκμηριωμένη παραλλαγή. Για την υπολογιστική μελέτη, οι παρατηρήσεις αυτές επιβάλλουν να διακριθούν η ταυτότητα μιας θέσης, η εσωτερική διαφοροποίηση των πραγματοποιήσεών της και η διαδοχή των θέσεων μέσα στο έργο.

\subsubsection{Από την εκδοχή στην οικογένεια θέσεων}

Ο όρος \textbf{οικογένεια θέσεων} ορίζεται στην παρούσα ερευνητική πρόταση ως αναλυτική κατηγορία για σύνολο συγκεκριμένων πραγματοποιήσεων των οποίων η κοινή μουσική ταυτότητα προτείνεται από αναλυτικά ή ιστορικά τεκμήρια και ελέγχεται με δηλωμένα κριτήρια. Η θεωρητική της βάση προκύπτει πρώτα από το ίδιο το φαινόμενο της βυζαντινής μουσικής: ονομασμένες θέσεις, επανεμφανίσεις, διαφοροποιήσεις, διαφορετικές σημειογραφικές αποδόσεις, λειτουργικά σταθερές μεταμορφώσεις και ανασυνδυασμοί έχουν ήδη περιγραφεί σε αναλύσεις έργων. Η οικογένεια καθιστά ρητή τη σχέση ανάμεσα στις διαφορετικές πραγματώσεις και επιτρέπει να μελετηθούν τα όρια της κοινής ταυτότητας σε κλίμακα σώματος.

Συγκριτικό προηγούμενο για αυτή την αναλυτική κίνηση προσφέρει η έρευνα των \emph{tune families}. Η \textbf{Anja Volk} και ο \textbf{Peter van Kranenburg}, στο ``Melodic Similarity among Folk Songs: An Annotation Study on Similarity-Based Categorization in Music'' (2012), μελέτησαν 360 μελωδίες κατανεμημένες σε 26 οικογένειες. Το πρωτόκολλό τους χρησιμοποίησε ρητά κριτήρια, μελωδία αναφοράς, επαναληπτική βελτίωση και μέτρηση συμφωνίας των ειδικών. Η συνολική συμφωνία Fleiss' $\kappa$ ήταν 0,74, ενώ η διαγνωστική συμβολή των επιμέρους χαρακτηριστικών διέφερε αισθητά. Οι ίδιοι καταγράφουν ότι δεν υπάρχει ένα χαρακτηριστικό ή ένα σταθερό σύνολο χαρακτηριστικών που να μοιράζονται όλες οι μελωδίες μιας οικογένειας με τη μελωδία αναφοράς. Για την παρούσα πρόταση αυτό λειτουργεί ταυτόχρονα ως πρότυπο πρωτοκόλλου και ως εμπειρική προειδοποίηση: η υπογραφή κάθε οικογένειας πρέπει να ελέγχεται ως type-specific και μπορεί να αποδειχθεί πολυθετική. Η εργασία των Finkensiep κ.ά. (2020) προσθέτει ότι ακόμη και υποψήφιες περιπτώσεις με το βασικό σχηματικό πρότυπο μπορεί να παραμένουν αμφίσημες ως προς την κατηγορική ένταξη. Η τριμερής κρίση $J(x,T;r)$ και τα δύσκολα αρνητικά παραδείγματα χρησιμοποιούνται για να καταγράψουν ακριβώς αυτή την περιοχή αβεβαιότητας.

Στο επίπεδο της μάθησης αναπαραστάσεων, οι \textbf{Folgert Karsdorp}, \textbf{Peter van Kranenburg} και \textbf{Enrique Manjavacas}, στο ``Learning Similarity Metrics for Melody Retrieval'' (ISMIR 2019), χρησιμοποίησαν αναδρομικά νευρωνικά δίκτυα και συναρτήσεις απώλειας για ζεύγη ή τριάδες για να μάθουν συνεχή χώρο συμβολικών μελωδιών στον οποίο η συγγένεια υποστηρίζει ανάκτηση. Η εργασία αυτή είναι βασική προηγούμενη υπολογιστική εργασία για την παρούσα ερευνητική πρόταση, επειδή δείχνει πώς η μελωδική συγγένεια μπορεί να μετατραπεί σε μετρική αναπαράσταση. Το ανοικτό πρόβλημα στη βυζαντινή περίπτωση είναι διαφορετικό: η οικογενειακή ένταξη πρέπει να τεκμηριωθεί μουσικολογικά και ιστορικά, να διαχωριστεί από την απλή επιφανειακή ομοιότητα και να αξιολογηθεί σε άγνωστα έργα και πηγές.

Συγγενείς εργασίες ανακάλυψης προτύπων σε παραδόσεις με ισχυρό φορμουλαϊκό λεξιλόγιο επιβεβαιώνουν ότι οι πολιτισμικά ορισμένες μελωδικές μονάδες απαιτούν αξιολόγηση με βάση τη γνώση της μουσικής παράδοσης. Οι \textbf{Thomas Nuttall}, \textbf{Miguel García Casado}, \textbf{Andres Ferraro}, \textbf{Darrell Conklin} και \textbf{Rafael Caro Repetto}, στο ``A Computational Exploration of Melodic Patterns in Arab-Andalusian Music'' (2021), αξιολογούν υπολογιστικές μεθόδους εντοπισμού χαρακτηριστικών προτύπων σε σχέση με θεωρία \emph{centos} που προέρχεται από ειδικό της ίδιας της παράδοσης. Οι Thomas Nuttall, Génis Plaja-Roglans, Lara Pearson και Xavier Serra, στο ``In Search of Sañcāras'' (ISMIR 2022), ενσωματώνουν επίσης γνώση της καρνατικής παράδοσης στην αναγνώριση επαναλαμβανόμενων μελωδικών προτύπων. Οι εργασίες αυτές χρησιμοποιούνται ως συγκριτικά προηγούμενα: ενισχύουν την αρχή ότι η υπολογιστική μονάδα πρέπει να αντιστοιχεί σε πολιτισμικά και μουσικολογικά καθορισμένο αντικείμενο, χωρίς να υποκαθιστούν τη βυζαντινή θεωρητική τεκμηρίωση.

Η κοινή ταυτότητα στηρίζεται σε διακριτά στρώματα τεκμηρίων: καταγεγραμμένες εμφανίσεις και αναλύσεις, καταλογογραφικές σχέσεις, τεκμηριωμένες ιστορικές συνδέσεις και ανεξάρτητες κρίσεις ειδικών. Η υπολογιστική ομοιότητα χρησιμοποιείται για ανάκτηση, γενίκευση και ποσοτικοποίηση, ενώ η συγκρότηση των οικογενειών και η κατοπινή επικύρωσή τους παραμένουν οργανωτικά χωριστές. Έτσι, η οικογένεια αποτελεί ελέγξιμη αναλυτική υπόθεση και όχι αποτέλεσμα της ίδιας της γεωμετρίας που πρόκειται να αξιολογηθεί.

\subsubsection{Εξήγηση, σημειογραφική μορφή και ιστορική μετάδοση}

Η ιστοριογραφία και η παλαιογραφία εξετάζουν την εξήγηση ως αυτοτελές ερευνητικό πρόβλημα. Ο Γ.~Θ. Στάθης, στο \emph{Η εξήγησις της Παλαιάς Βυζαντινής Σημειογραφίας} (1978), συστηματοποίησε τη μελέτη των σχέσεων παλαιότερων σημειογραφιών με τις αναλυτικές αποδόσεις. Στο πλαίσιο αυτής της προσέγγισης, η μετάβαση προς τη Νέα Μέθοδο ερμηνεύεται ως αναλυτικότερη καταγραφή μελικού περιεχομένου που στην παλαιότερη γραφή μπορούσε να αποδίδεται συνοπτικότερα. Η προσέγγιση αυτή υιοθετείται και εξειδικεύεται σε αναλύσεις ονομασμένων θέσεων και των αποδόσεών τους στη Νέα Μέθοδο (Καραγκούνης, 2003, σσ.~721, 810--812). Η παρούσα ερευνητική πρόταση δεν χρησιμοποιεί την ερμηνεία αυτή ως δεδομένο ιστορικής ακουστικής ταυτότητας· αξιοποιεί τις τεκμηριωμένες εξηγητικές σχέσεις ως ιστορικά εντοπισμένες σχέσεις μεταξύ συγκεκριμένων αποδόσεων, ενώ η μουσική διατήρηση ή μεταβολή εξετάζεται χωριστά.

Στην ίδια ιστορική συζήτηση, ο \textbf{Απόστολος Κώνστας ο Χίος} προσφέρει άμεση μαρτυρία θεωρητικού και εξηγητή της εποχής· η συνολική θεωρητική και κωδικογραφική συμβολή του έχει μελετηθεί ειδικά από τον Αποστολόπουλο (2002). Το κείμενό του του 1820 σώζεται στον κώδικα ΕΒΕ 1867 και έχει ενταχθεί στη γ΄ έκδοση του Στάθη (1993). Η βιβλιογραφική ταυτότητα του χωρίου της σ.~90 είναι ασφαλής· η ακριβής λεκτική αντιπαραβολή του περιεχομένου του παραμένει προϋπόθεση πριν αποδοθούν σε αυτό ειδικότεροι ισχυρισμοί για τον σκοπό της εξήγησης. Η μαρτυρία συνεξετάζεται με τις επιμέρους διαφωνίες που καταγράφει ο Στάθης για άλλες κρίσεις του Κώνστα (1993, σσ.~36--38). Στην παρούσα έρευνα, η εξηγητική σχέση τεκμηριώνει τη σύνδεση και την προέλευση συγκεκριμένων αποδόσεων, ενώ η κοινή μουσική ταυτότητα και η διατήρηση ή διαφοροποίηση των χαρακτηριστικών τους ελέγχονται χωριστά.

Ο \textbf{Ιωάννης Αρβανίτης}, στη διδακτορική του \emph{Ο ρυθμός των εκκλησιαστικών μελών μέσα από την παλαιογραφική έρευνα και την εξήγηση της παλαιάς σημειογραφίας} (2010) και στο άρθρο ``The Transcriptions of the Three Teachers: An Old or a New Sound?'' (2022, σσ.~143--149), επανεξετάζει τη σχέση των αποδόσεων των Τριών Διδασκάλων με το παλαιότερο άκουσμα. Η διαφορετική αυτή προβληματική κρατά ανοικτό το ιστορικό ακουστικό ερώτημα. Για τον λόγο αυτό, στην παρούσα έρευνα η τεκμηριωμένη εξηγητική σύνδεση χρησιμοποιείται ως δεδομένο προέλευσης ανάμεσα σε συγκεκριμένες αποδόσεις, ενώ η μουσική διατήρηση ή μεταβολή μετράται χωριστά.

\textbf{Ο σχετικός έλεγχος της M1 εξετάζει αν οι αναλυτικές αποδόσεις της Νέας Μεθόδου διατηρούν διακρίσεις που τεκμηριώνονται στην παλαιότερη θεωρία της θέσης· δεν εξετάζει αν οι εξηγήσεις αναπαράγουν το ιστορικό άκουσμα της παλαιάς σημειογραφίας.}

Η ίδια αρχή εμφανίζεται σε ειδικές αναλύσεις του ρεπερτορίου. Ο \textbf{Nicolae Gheorghiță}, στο ``Observations on the Technique of Transcription (Exegesis) into the New Method of Analytical Music Notation of the Sunday Koinonikon of the 18th Century'' (2010), εξετάζει τεχνικές εξήγησης σε συγκεκριμένο σώμα Κυριακάτικων Κοινωνικών. Ο \textbf{Alexander Lingas}, στο ``Performance Practice and the Politics of Transcribing Byzantine Chant'' (2003), αναδεικνύει τις ιστοριογραφικές προϋποθέσεις που συνδέουν γραπτή αναπαράσταση και εκτελεστική πραγμάτωση. Οι πηγές αυτές ορίζουν το πεδίο μέσα στο οποίο η εξήγηση λειτουργεί ως ιστορικά εντοπισμένη σχέση και όχι ως αφηρημένο χαρακτηριστικό ομοιότητας.

Οι εξηγήσεις του \textbf{Γρηγορίου Πρωτοψάλτου} και του \textbf{Χουρμουζίου Χαρτοφύλακος} εντάσσονται στο ίδιο ιστορικό περιβάλλον διαμόρφωσης και διάδοσης της Νέας Μεθόδου. Στη M2 αξιοποιούνται ως διακριτές εξηγητικές μαρτυρίες, ενώ η μεταξύ τους συμφωνία ερμηνεύεται λαμβάνοντας υπόψη αυτή την ιστορική συνάφεια και την προέλευση κάθε πηγής.

Η σχέση ανάμεσα στη σημειογραφική μορφή και στη μελική πραγμάτωση εξετάζεται επομένως μέσω τεκμηριωμένων εξηγητικών σχέσεων και συγκριτικής ανάλυσης πολλαπλών αποδόσεων. Οι εξηγήσεις παρέχουν ιστορικά εντοπισμένες μαρτυρίες για τον τρόπο με τον οποίο συγκεκριμένες παλαιότερες γραφές αποδίδονται στη Νέα Μέθοδο και επιτρέπουν να εξεταστεί σε επίπεδο θέσης τι διατηρείται, τι μεταβάλλεται και πώς οργανώνεται η μορφολογική διαφοροποίηση. Για τη M2, η συνέπεια είναι άμεση: η οικογενειακή συγγένεια και η τεκμηριωμένη ιστορική σχέση αποτελούν διαφορετικές πηγές πληροφορίας και μπορούν να συγκριθούν υπό αντιστοιχισμένους όρους.

\subsubsection{Χερουβικά και Κοινωνικά ως ιστορικό σώμα}

Η επιλογή Χερουβικών και Κοινωνικών στηρίζεται σε καθιερωμένη ιστορική έρευνα του ίδιου του ρεπερτορίου. Ο \textbf{Dimitri E. Conomos}, στο \emph{Byzantine Trisagia and Cheroubika of the Fourteenth and Fifteenth Centuries: A Study of Late Byzantine Liturgical Chant} (1974), μελετά συστηματικά τα Τρισάγια και Χερουβικά του ύστερου βυζαντινού ρεπερτορίου και παρέχει ιστορικό πλαίσιο για τις πηγές, τις συνθέσεις και τη μετάδοσή τους. Η μονογραφία δεν χρησιμοποιείται ως κατάλογος οικογενειών της παρούσας έρευνας· θεμελιώνει το ιστορικό πεδίο μέσα στο οποίο οι επιμέρους θέσεις και οι παραλλαγές τους πρέπει να ερμηνευθούν.

Για τα Κοινωνικά, ο ίδιος ερευνητής στο \emph{The Late Byzantine and Slavonic Communion Cycle: Liturgy and Music} (1985a) εξετάζει το σωζόμενο ελληνικό και σλαβικό ρεπερτόριο Κοινωνικών από τον 12ο έως τον 16ο αιώνα και οργανώνει τη μελέτη του σε σχέση με τη λειτουργία και τις χειρόγραφες παραδόσεις. Η εργασία αυτή είναι κρίσιμη για την ιστορική οριοθέτηση του δεύτερου βασικού σώματος της ερευνητικής πρότασης και για τη διάκριση λειτουργικού κύκλου, πηγής και μουσικής παράδοσης.

Για τα Χερουβικά, οι μελέτες του 2003 και 2005 προσθέτουν στο ιστορικό αυτό υπόβαθρο λεπτομερή ανάλυση της εσωτερικής οργάνωσης του ρεπερτορίου (Καραγκούνης, 2003, 2005). Στο επίπεδο των Κοινωνικών, ο Nicolae Gheorghiță (2010) προσφέρει συγκεκριμένη εξηγητική μελέτη, ενώ η δημοσιευμένη εργασία του \textbf{Χρυσοβαλάντη Ιωαννίδη}, ``Η σύνθεση του Νικολάου Πρωτοψάλτου Σμύρνης στο Κοινωνικό `Εἰς πᾶσαν τὴν γῆν' σε ήχο πλ. Δ΄ και ο βαθμός εξάρτησής της από την προγενέστερη παράδοση της εκκλησιαστικής μελοποιίας'' (2022), παρέχει οριοθετημένη ιστορική και δομική αντιπαραβολή που χρησιμοποιείται αργότερα ως ανεξάρτητη περίπτωση εφαρμογής. Οι εργασίες του \textbf{Αρσενίου Χαραλαμπίδη} για Κοινωνικά σε πρώτο ήχο (2024, 2026) χρησιμοποιούνται ως συμπληρωματικές σύγχρονες μορφολογικές χαρτογραφήσεις με σαφώς διαφορετικό τεκμηριωτικό βάρος από τις δημοσιευμένες με κρίση ομοτίμων ή μονογραφικές πηγές.

Η διαδικτυακή υποδομή \emph{Byzantine Music Formulae} και η χωριστή \emph{Papadic Formulae} της Ιεράς Μονής Αγίου Αντωνίου παρέχουν οργανωμένο κατάλογο υποψηφίων τύπων και διαδρομή προς έργα και πηγές. Η παπαδική συλλογή αφορά 758 συνθέσεις και χρησιμοποιείται ως υποδομή υποψηφίων σχέσεων, αναζήτησης και ασθενούς εποπτείας. Δεν αντιμετωπίζεται ως χρυσό πρότυπο: κάθε σχέση που χρησιμοποιείται ως οικογενειακή αναφορά περνά χωριστή τεκμηρίωση των πηγών και μουσικολογική τεκμηρίωση. Έτσι το ιστορικό σώμα και το καταλογογραφικό σώμα έχουν διακριτό ρόλο. Για τον υπολογιστικό σχεδιασμό, η πηγή, το έργο και η ιστορική γραμμή αποτελούν μεταδεδομένα που οργανώνουν τους διαχωρισμούς εκπαίδευσης και αξιολόγησης και εμποδίζουν τη διαρροή συγγενικού υλικού μεταξύ συνόλων.

\subsubsection{Υπολογιστική αναπαράσταση, τμηματοποίηση και μελωδική συγγένεια}

Η υπολογιστική μελέτη της ελληνικής εκκλησιαστικής μουσικής έχει ουσιαστικό προηγούμενο και πρέπει να διακρίνεται από τη γενική MIR βιβλιογραφία. Ο \textbf{Panayotis Mavromatis}, στο ``A Hidden Markov Model of Melody Production in Greek Church Chant'' (2005a), διατυπώνει πιθανοτικό μοντέλο της μελωδικής διαδικασίας της σύγχρονης ελληνικής εκκλησιαστικής ψαλτικής με στόχο την επαγωγή λανθανόντων κανόνων από σώμα μελών και τη σύνδεσή τους με τη νοητική αναπαράσταση της μελωδίας. Στη διδακτορική του \emph{The Echoi of Modern Greek Church Chant in Written and Oral Transmission: A Computational Model and Its Cognitive Implications} (2005b) συνδέει γραπτή και προφορική μετάδοση με υπολογιστική μοντελοποίηση και γνωσιακές προεκτάσεις. Στο ``Minimum Description Length Modelling of Musical Structure'' (2009a) και στη συναφή εργασία για topology-sensitive HMM model selection (2009b) αναπτύσσει περαιτέρω την επαγωγή μουσικής δομής από δεδομένα. Η γραμμή αυτή αποτελεί άμεσο προηγούμενο για την υπολογιστική μελέτη της βυζαντινής μελωδικής γνώσης. Η παρούσα έρευνα εστιάζει ειδικότερα σε οικογένειες πραγματώσεων μιας θέσης μεταβλητού μήκους, με χωριστή καταγραφή οικογενειακής ταυτότητας και τεκμηριωμένης εξηγητικής σχέσης και με ανεξάρτητους μουσικολογικούς, ανθρώπινους και εκτελεστικούς ελέγχους.

Η διδακτορική διατριβή του \textbf{Πολύκαρπου Γ. Πολυκαρπίδη}, \emph{Υπολογιστική προσέγγιση της Βυζαντινής Μουσικής: Η περίπτωση του Ειρμολογίου του Θεοφάνη Καρύκη} (Εθνικό και Καποδιστριακό Πανεπιστήμιο Αθηνών, 2023), αναπτύσσει υπολογιστική αναπαράσταση της Μεσοβυζαντινής σημειογραφίας με ιεραρχική οργάνωση σε επίπεδα έργου, φράσης, συλλαβής και Φωνητικής Μονάδας και με πληροφορίες για τη μελωδική κίνηση, τα ποιοτικά σημάδια, τον ήχο και το κείμενο. Περιγράφει επίσης τη γλώσσα BNL και την υλοποίηση του αντίστοιχου αναλυτή, καθώς και εφαρμογή εξόρυξης επαναλαμβανόμενων μοτίβων σε σώμα ειρμών. Η φρασεοποίηση και οι συγκριτικές αναλύσεις μελωδικών τόξων, σημειογραφικής υφής και ακολουθιών γνωρισμάτων αποτελούν επίσης μέρος της εφαρμοσμένης μεθοδολογίας της διατριβής. Σε συναφή εργασία, οι Polykarpidis, Kalofonos, Balageorgos και Anagnostopoulou (2022) συγκρίνουν τρία σώματα ειρμολογικού μέλους σε Μέση Βυζαντινή σημειογραφία. Η παρούσα έρευνα θέτει διαφορετικό αλλά συναφές ερώτημα: εξετάζει τη μουσική θέση ως αναγνωριζόμενη μονάδα μεταβλητής έκτασης και την οικογενειακή συγγένεια των πραγματώσεών της υπό τεκμηριωμένες εξηγητικές σχέσεις, με αντιστοιχισμένη οικογενειακή εποπτεία και ανεξάρτητη μουσικολογική, ανθρώπινη και εκτελεστική αξιολόγηση.

Οι εργασίες των E.~Cambouropoulos, \emph{Musical Parallelism and Melodic Segmentation: A Computational Approach} (2006), και D.~Conklin--C.~Anagnostopoulou, \emph{Segmental Pattern Discovery in Music} (2006), θεμελιώνουν δύο συναφή υπολογιστικά προβλήματα. Η πρώτη εξετάζει πώς οι επαναλαμβανόμενες και παράλληλες μελωδικές δομές μπορούν να συνεισφέρουν στον εντοπισμό πιθανών ορίων, ενώ η δεύτερη αναζητεί επαναλαμβανόμενα πρότυπα όταν η μελωδία έχει ήδη αναπαρασταθεί ως ακολουθία τμημάτων. Νεότερες εργασίες σε μονοφωνικές παραδόσεις εξετάζουν επίσης φυσικές μονάδες και εναλλακτικές τμηματοποιήσεις σε μεγάλα σώματα ψαλμωδίας (Cornelissen κ.ά., 2020a, 2020b), ενώ οι Lanz και Hajič (2023, 2025) δείχνουν ότι η επιλογή ορίων και η μη παραμετρική πιθανοτική τμηματοποίηση αποτελούν αυτοτελές ερευνητικό πρόβλημα και δεν ταυτίζονται αναγκαστικά με μια προκαθορισμένη μουσικολογική μονάδα. Η παρούσα έρευνα χρησιμοποιεί αυτές τις δυνατότητες ως συγκριτικό υπόβαθρο της M1 και εξετάζει επιπλέον αν μια μουσικολογικά τεκμηριωμένη μονάδα μεταβλητής έκτασης διατηρεί αναγνωρίσιμη οικογενειακή ταυτότητα μέσα σε διαφορετικές πραγματώσεις. Για τον λόγο αυτό οι αλγοριθμικά επαγόμενες μονάδες παραμένουν συγκριτικό μοντέλο της M1 και δεν υποκαθιστούν την οικογένεια θέσεων.

Η έρευνα εξετάζει τη συμβολική κωδικοποίηση ως μεταβλητή στο MIR. Οι Fradet και συνεργάτες (2023) εξετάζουν την προσαρμογή Byte Pair Encoding στη συμβολική μουσική, ενώ οι Kaliakatsos-Papakostas και συνεργάτες (2025), στο \emph{HarmonyTok}, συγκρίνουν διαφορετικές στρατηγικές κωδικοποίησης της αρμονίας σε κοινές εργασίες Transformer και δείχνουν ότι ο τρόπος διάσπασης της μουσικής πληροφορίας μεταβάλλει τόσο την προβλεπτική επίδοση όσο και τις δομικές ιδιότητες που διατηρούνται. Για την M1, οι εργασίες αυτές αποτελούν συγκριτικό προηγούμενο για την ελεγχόμενη αξιολόγηση εναλλακτικών μονάδων αναπαράστασης υπό κοινούς στόχους, χωρίς να προδικάζουν ποια μονάδα είναι κατάλληλη για τη βυζαντινή θέση.

Η διεθνής βιβλιογραφία της μελωδικής συγγένειας και των schemata παρέχει συμπληρωματικά εργαλεία. Η εργασία των Karsdorp, van Kranenburg και Manjavacas (2019) μελετά τη μάθηση μετρικής για την ανάκτηση συγγενών μελωδιών· η εργασία των Volk και van Kranenburg (2012) συνδέει την ομοιότητα με ανθρώπινες κατηγορίες tune families· οι σχετικές computational μελέτες schemata και tradition-informed ανακάλυψη προτύπων δείχνουν πώς πολιτισμικά ορισμένοι τύποι μπορούν να αναζητηθούν σε μεγάλα σώματα χωρίς να ανάγονται σε απλά $n$-grams. Το ακριβές όριο της παρούσας ερευνητικής πρότασης βρίσκεται στη σύνθεση αυτών των δυνατοτήτων με ταυτότητα ορισμένη με κριτήρια της βυζαντινής μουσικολογίας και με συγκεκριμένες ιστορικές σχέσεις μεταξύ εκδοχών.

\subsubsection{Ανθρώπινη κατηγοριοποίηση και μουσική πρόβλεψη}

Η M3b στηρίζεται στη διάκριση ανάμεσα στην αντιληπτή ομοιότητα και στην κατηγορική ταυτότητα. Η Anja Volk και ο Peter van Kranenburg (2012) έδειξαν εμπειρικά ότι η ανθρώπινη κατηγοριοποίηση σε tune families σχετίζεται με πολλαπλά μουσικά χαρακτηριστικά και ότι η βαρύτητά τους μεταβάλλεται εντός και μεταξύ οικογενειών. Η εργασία αυτή είναι σημαντική επειδή καθιστά εύλογο ένα ερώτημα που δεν μπορεί να λυθεί μόνο από απόσταση συμβολικών ακολουθιών: αν δύο εκδοχές μπορούν να αναγνωριστούν ως συγγενείς παρά μεγάλη επιφανειακή απόσταση και αν το αντίστροφο, τοπική ομοιότητα, αρκεί για κοινή κατηγορική ένταξη. Η M3b μεταφέρει αυτό το γενικό πρόβλημα σε οικογένειες θέσεων που έχουν προηγουμένως οριστεί και οριστικοποιηθεί με μουσικολογικά κριτήρια. Η υπογραφή της οικογένειας λειτουργεί ως αναλυτικός κανόνας της έρευνας· δεν θεωρείται εκ των προτέρων μοντέλο της γνωσιακής κατηγοριοποίησης των ψαλτών. Η σχέση των δύο εξετάζεται εμπειρικά στην ίδια τη M3b.

Η νευροεπιστημονική μελέτη της μελωδικής ταυτότητας παρέχει ένα στενότερο συγκριτικό σημείο. Οι Andreas Schindler, Marcus Herdener και Andreas Bartels (2013) έδειξαν ότι η ίδια μελωδική Gestalt μπορεί να αναπαρίσταται με σχετική αναλλοιότητα στον ακουστικό φλοιό όταν αλλάζει η τονικότητα ή το όργανο, δηλαδή όταν διατηρείται το σχετικό προφίλ των φθόγγων. Η οικογένεια θέσεων θέτει διαφορετικό και αυστηρότερο κατηγορικό ερώτημα: αν διατηρείται κοινή ταυτότητα όταν μπορούν να μεταβάλλονται οι ίδιες οι ακολουθίες φθόγγων και τα μέλη της κατηγορίας έχουν οριστεί με ανεξάρτητη ιστορική και μουσικολογική τεκμηρίωση. Αυτό είναι το ειδικό αντικείμενο της M3b.

Η μουσική πρόβλεψη έχει διαφορετική βιβλιογραφική αφετηρία. Οι \textbf{Marcus T. Pearce}, \textbf{Daniel Müllensiefen} και \textbf{Geraint A. Wiggins}, στο ``The Role of Expectation and Probabilistic Learning in Auditory Boundary Perception'' (2010), χρησιμοποιούν πιθανοτικές προσδοκίες για να εξηγήσουν ανθρώπινη αντίληψη μουσικών ορίων. Οι Benjamin Gingras, Marcus T. Pearce, Meghan Goodchild, Roger T. Dean, Geraint A. Wiggins και Stephen McAdams (2016) συνδέουν μελωδική προσδοκία με εκφραστικό χρονισμό και αντιληπτή ένταση, ενώ οι Niels C. Hansen, Haley E. Kragness, Peter Vuust, Laurel J. Trainor και Marcus T. Pearce (2021) εξετάζουν την προβλεπτική αβεβαιότητα στην ακουστική αντίληψη ορίων. Οι Koelsch κ.ά. (2013) και οι Steffen A. Herff και συνεργάτες (2021) δείχνουν ότι ιεραρχική μουσική δομή σχετίζεται με ανθρώπινη επεξεργασία και αποκρίσεις ολοκλήρωσης ακολουθιών. Οι Giovanni M. Di Liberto και συνεργάτες (2020) συνδέουν υπολογιστικά μέτρα μελωδικής προσδοκίας με χρονικά εξελισσόμενη νευρική δραστηριότητα. Οι Kern, Heilbron, de Lange και Spaak (2022) έδειξαν επιπλέον ότι, σε φυσική ακρόαση δυτικής μουσικής, η νευρική απόκριση στο πληροφοριακό περιεχόμενο αποδόθηκε καλύτερα από μοντέλα μακροχρόνιας στατιστικής μάθησης, ενώ το άμεσο πλαίσιο που εξηγούσε καλύτερα τη νευρική απόκριση ήταν σχετικά σύντομο. Η διάκριση αυτή αιτιολογεί στη M3c και στη M3ν τον χωριστό έλεγχο της μακροχρόνιας μάθησης του μοντέλου και του μήκους του άμεσου συμφραζομένου. Οι εργασίες αυτές θεμελιώνουν τη μέτρηση της μουσικής προσδοκίας και του πληροφοριακού περιεχομένου, αλλά δεν εξετάζουν τη βυζαντινή οικογένεια θέσεων ως συγκεκριμένο επίπεδο πρόβλεψης.

Η μουσικολογική ανάλυση της βυζαντινής μουσικής συνδέει συγκεκριμένες παρατηρήσεις με το γνωσιακό ερώτημα. Σε μορφολογική ανάλυση του Χερουβικού του ψευδο-Ιωάννου Δαμασκηνού περιγράφονται ενδιάμεσες καταλήξεις που αφήνουν τη μελωδία μετέωρη και μεταθέτουν την οριστική ανάπαυση στη συνολική κατάληξη, καθώς και επαναλαμβανόμενη αλυσίδα θέσεων που καλύπτει σχεδόν όλο το πρώτο μέρος πριν από τη Μεγάλη Είσοδο· στο επόμενο τμήμα η επαναληπτικότητα μειώνεται και η μελωδική οργάνωση διαφοροποιείται (Καραγκούνης, 2005, σσ.~158--159). Η παρατήρηση αυτή δεν ισοδυναμεί με πείραμα ακροατών ούτε με καθολικό κανόνα των Χερουβικών· παρέχει όμως σαφές αντικείμενο ελέγχου σε κλίμακα σώματος και στην ανθρώπινη ακρόαση της σχέσης ανάμεσα σε οργάνωση της μουσικής διαδοχής, λειτουργικό τόπο και μουσική προσδοκία.

Η M3c διατηρεί χωριστά δύο επίπεδα αυτής της ερώτησης. Στο ανθρώπινο πείραμα, η πιθανότητα της επόμενης οικογένειας από το προηγούμενο μουσικό συμφραζόμενο συγκρίνεται με το πληροφοριακό περιεχόμενο των επιμέρους μελωδικών γεγονότων. Σε δευτερεύουσα ανάλυση του συμβολικού σώματος, με αφετηρία την παραπάνω μορφολογική παρατήρηση, εξετάζεται αν η επανάληψη και η προβλεψιμότητα ακολουθιών οικογενειών διαφοροποιούνται συστηματικά μεταξύ τεκμηριωμένων λειτουργικών τμημάτων πλήρων Χερουβικών. Η δεύτερη ανάλυση δεν υποκαθιστά την πειραματική ένδειξη λατρευτικού συμφραζομένου $R$· εξετάζει την οργανωτική δομή του ίδιου του ρεπερτορίου, με έλεγχο μήκους, αριθμού θέσεων, κειμενικού τόπου, καταληκτικής πυκνότητας, ήχου, μελοποιού\footnote{Μελοποιός: ο δημιουργός/συνθέτης ενός μέλους στην ορολογία της ψαλτικής· μελοποίηση είναι η πράξη και το αποτέλεσμα της μουσικής σύνθεσης ενός υμνογραφικού ή λειτουργικού κειμένου.} και ιστορικού στρώματος.

\section{Διαχωρισμοί δεδομένων και αποτροπή διαρροής}\label{app:splits}

Οι διαχωρισμοί δεδομένων προστατεύουν την ανεξαρτησία της τελικής αξιολόγησης. Η μονάδα διαχωρισμού επιλέγεται σύμφωνα με τις πραγματικές εξαρτήσεις του μουσικού υλικού: διαφορετικές θέσεις από το ίδιο έργο, διαφορετικές εκδόσεις της ίδιας ομάδας συγγενών πηγών ή τεκμηριωμένα συγγενείς πραγματώσεις μπορούν να μοιράζονται πληροφορία. Για τον λόγο αυτό η τυχαία κατανομή μεμονωμένων γεγονότων σε εκπαίδευση και αξιολόγηση δεν παρέχει επαρκή προστασία. Οι διαχωρισμοί κατά ολόκληρο έργο (work-held-out), ομάδα συγγενών πηγών (source-line-held-out) και μελοποιό εφαρμόζονται όπου αντιστοιχούν στο εκτιμώμενο μέγεθος της συγκεκριμένης ανάλυσης.

\subsection{Κόμβοι, ομάδες συγγενών πηγών και διαχωρισμός συνόλων}\label{app:splits-nodes}

Ο έλεγχος διαρροής οργανώνεται ως \textbf{μήτρα πεδίου ανά εργασία}: η ίδια μεταβλητή μπορεί να είναι νόμιμη είσοδος σε μία εργασία και απαγορευμένο υποκατάστατο του στόχου σε άλλη.

{\small
\begin{xltabular}{\textwidth}{@{}p{0.235\textwidth}p{0.135\textwidth}p{0.115\textwidth}p{0.145\textwidth}Y@{}}
\toprule
\textbf{Πεδίο εισόδου} & \textbf{Όρια} & \textbf{Οικογένεια} & \textbf{Λειτ. ρόλος} & \textbf{Εκτέλεση}\tabularnewline
\midrule
\endfirsthead
\toprule
\textbf{Πεδίο εισόδου} & \textbf{Όρια} & \textbf{Οικογένεια} & \textbf{Λειτ. ρόλος} & \textbf{Εκτέλεση}\tabularnewline
\midrule
\endhead
\midrule
\multicolumn{5}{r}{\footnotesize Συνέχεια στην επόμενη σελίδα}\tabularnewline
\endfoot
\bottomrule
\endlastfoot

Οικογένεια θέσης & Μη επιτρεπτό & Μη επιτρεπτό & Μη επιτρεπτό & Μη επιτρεπτό\tabularnewline
Όρια εμφάνισης & Μη επιτρεπτό & επιτρεπτά & επιτρεπτά & επιτρεπτά\tabularnewline
\textbf{Φραστικό όριο} & \textbf{Μη επιτρεπτό} & με προϋποθέσεις & \textbf{Μη επιτρεπτό} & επιτρεπτό\tabularnewline
Τύπος κατάληξης & Μη επιτρεπτό & επιτρεπτός & Μη επιτρεπτό & επιτρεπτός\tabularnewline
Μελισματική πυκνότητα & έλεγχος & έλεγχος & Μη επιτρεπτό & επιτρεπτή\tabularnewline
Ήχος, μαρτυρία, φθορά & επιτρεπτά & επιτρεπτά & επιτρεπτά & επιτρεπτά\tabularnewline
Χρονικά σημάδια & επιτρεπτά & επιτρεπτά & επιτρεπτά & επιτρεπτά\tabularnewline
Επισημείωση εκτέλεσης & Μη επιτρεπτό & Μη επιτρεπτό & Μη επιτρεπτό & Μη επιτρεπτό\tabularnewline
\end{xltabular}
}\par

\textbf{Το φραστικό όριο αποτελεί τον κρισιμότερο εναλλακτικό εξηγητικό παράγοντα} και ελέγχεται ρητά. Τα όρια των μουσικών θέσεων συσχετίζονται ισχυρά με τα φραστικά όρια, χωρίς να ταυτίζονται. Επιβάλλεται επομένως \textbf{κανόνας προέλευσης}: το πεδίο γίνεται δεκτό ως είσοδος \emph{μόνο} όταν παράγεται ντετερμινιστικά από το κείμενο --- στίξη, συντακτική δομή, όρια υμνογραφικής ενότητας --- και \emph{ποτέ} όταν προέρχεται από μουσικολογική τμηματοποίηση ή από τον ίδιο αναλυτικό μηχανισμό που παρήγαγε τα όρια των θέσεων. Σε μη λεξικά χωρία δεν κατασκευάζονται τεχνητές λέξεις ή φράσεις: όταν ο προκαθορισμένος κανόνας δεν ορίζει έγκυρο $S_{\mathrm{phrase}}$, η τιμή δηλώνεται μη εφαρμοστέα και η πλήρης αντίθεση Υ8 περιορίζεται στο κοινό επιλέξιμο σύνολο όπου ορίζονται όλες οι απαιτούμενες εναλλακτικές τμηματοποιήσεις. Η προέλευση κάθε τιμής καταγράφεται και ελέγχεται.

Διαφορετικά λεκτικά καθεστώτα του ίδιου έργου παραμένουν στην ίδια ομάδα διαχωρισμού· το ότι ένα χωρίο δεν είναι κατάλληλο για συγκεκριμένο έλεγχο που βασίζεται στο κείμενο δεν διαγράφει το χωρίο από την κανονική ακολουθία ούτε συγκολλά τεχνητά τα γειτονικά τμήματα.

Παρόμοιος κανόνας ισχύει για τη μελισματική πυκνότητα και τους δείκτες ανά συλλαβή: αποτελούν έγκυρα διαγνωστικά, αλλά όταν ο στόχος είναι ο λειτουργικός ρόλος λειτουργούν ως υποκατάστατα και αποκρύπτονται.

Οι ελεγχόμενες συμβολικές παρεμβάσεις της προαιρετικής διερευνητικής υπομελέτης συγκροτούν \textbf{χωριστό σύνολο μετά την οριστικοποίηση του κύριου μοντέλου EISE}. Τα ίδια ερεθίσματα αξιολογούνται από EISE, B2 και B0 και από ανεξάρτητες τυφλές κρίσεις ειδικών, χωρίς ενημέρωση της κύριας μάθησης από την έκβασή τους. Η επιλογή των τελικών ερεθισμάτων γίνεται ανεξάρτητα από τη βαθμολογία του EISE, ώστε η ανθρώπινη αναγνώριση να λειτουργεί ως ανεξάρτητο κριτήριο αξιολόγησης πρόβλεψης.

{\small
\begin{xltabular}{\textwidth}{@{}p{0.25\textwidth}Y@{}}
\toprule
\textbf{Πηγή κινδύνου στα δεδομένα} & \textbf{Μέτρο}\tabularnewline
\midrule
\endfirsthead
\toprule
\textbf{Πηγή κινδύνου στα δεδομένα} & \textbf{Μέτρο}\tabularnewline
\midrule
\endhead
\midrule
\multicolumn{2}{r}{\footnotesize Συνέχεια στην επόμενη σελίδα}\tabularnewline
\endfoot
\bottomrule
\endlastfoot

Συγγενείς μάρτυρες σε εκπαίδευση και έλεγχο & Εξαίρεση ολόκληρης ομάδας συγγενών πηγών\tabularnewline
Ίδιος εκτελεστής σε εκπαίδευση και έλεγχο & Εξαίρεση εκτελεστή\tabularnewline
Ίδια ηχογράφηση σε δύο ρόλους & Εξαίρεση συνεδρίας\tabularnewline
Επικάλυψη αναφοράς με έργα αξιολόγησης & Ρητή καταγραφή και χωριστή αναφορά\tabularnewline
Προεκπαιδευμένες αναπαραστάσεις & Έλεγχος επικάλυψης με τις πηγές του συνόλου τελικής αξιολόγησης\tabularnewline
\end{xltabular}
}\par

\paragraph{Αποκλεισμός της ιστορικής εφαρμογής από ανάπτυξη και εκπαίδευση.}
Τα έργα του Νικολάου αποτελούν σύνολο εφαρμογής και εξαιρούνται, μαζί με τις ισοδύναμες γραπτές αποδόσεις και τους εξαρτημένους μάρτυρές τους, από προεκπαίδευση, επαγωγή μονάδων, επιλογή των $R^\star$, $\mathcal R^\star$ και $R_{\mathrm{sens}}$, οικογενειακή/εξηγητική εκπαίδευση και βαθμονόμηση. Για την περίπτωση Ιωαννίδη εφαρμόζεται ο ίδιος αποκλεισμός στις συγκεκριμένες συγκρινόμενες συνθέσεις και στις ομάδες συγγενών πηγών τους, πριν δημιουργηθούν ζεύγη ή λεξιλόγια. Άλλα επιλέξιμα έργα των ίδιων συνθετών μπορούν να εκπροσωπούν τα ιστορικά στρώματα. Μετά την οριστικοποίηση, οι συνθέσεις αναφοράς που έχουν διατηρηθεί εκτός ανάπτυξης μπορούν να εισέλθουν μόνο σε χωριστό κατάλογο ανάκτησης για τη συγκεκριμένη δοκιμή, χωρίς ενημέρωση προτύπων, μετρικής ή παραμέτρων. Η κάλυψη που χάνεται από αυτούς τους αποκλεισμούς προσμετράται στο ΠΕ0. Αν προηγούμενη χρήση δεν μπορεί να αποκλειστεί, η περίπτωση χαρακτηρίζεται αναδρομική εφαρμογή και όχι ανεξάρτητη επικύρωση.

\subsection{Ανεξαρτησία της πληροφοριακής δοκιμής}\label{app:splits-info}

Η ακολουθία είναι αυστηρή και εξαλείφει την κυκλικότητα. Τα μοντέλα ακολουθίας εκπαιδεύονται \textbf{αποκλειστικά στη συμβολική ακολουθία}, χωρίς καμία πρόσβαση σε ετικέτες, όρια ή οικογένειες θέσεων. Παράγουν τα προφίλ πληροφοριακού περιεχομένου και εντροπίας της προβλεπτικής κατανομής σε υλικό που έχει διατηρηθεί εκτός εκπαίδευσης και τα προφίλ \textbf{οριστικοποιούνται πριν από οποιαδήποτε σύγκριση με τις μουσικολογικές ετικέτες}. Μόνο μετά συγκρίνονται με τις εξωτερικές μουσικολογικές οριοθετήσεις.

\subsection{Κυκλικότητα και επιτρεπτές υποκατάστατες μεταβλητές στην τροπική δοκιμή}\label{app:splits-trop}

Τρεις κανόνες διασφαλίζουν ότι το κριτήριο παραμένει ανεξάρτητο.

\textbf{Πρώτον}, οι ετικέτες θέσεων \emph{απαγορεύονται} ως είσοδος σε αυτή την εργασία, όπως ακριβώς η οικογένεια απαγορεύεται στη μουσικολογική δοκιμή.

\textbf{Δεύτερον}, τα ισοκρατηματικά στρώματα πρέπει να έχουν παραχθεί \emph{χωρίς} χρήση ορίων θέσεων ως κανόνα αναγωγής. Όπου η προέλευση δεν το εγγυάται, η πηγή αποκλείεται από την επιβεβαιωτική ανάλυση και αναφέρεται χωριστά.

\textbf{Τρίτον}, όπου η πηγή δεν προσδιορίζει ακριβές σημείο μετακίνησης αλλά περιοχή, το όριο αποθηκεύεται ως \emph{διαστηματικά προσδιορισμένο} και η αξιολόγηση χρησιμοποιεί κριτήρια ένταξης στο διάστημα αντί να επιβάλλει τεχνητή σημειακή ετικέτα.

\subsection{Επιλογή εκτελεστικών δεδομένων και ανεξαρτησία}\label{app:splits-perf}

Από την επισημείωση του σώματος χρησιμοποιούνται \textbf{μόνο τα επίπεδα που προέρχονται από την παρτιτούρα}: όρια συλλαβών, χαρακτήρες γραπτής σημειογραφίας, αναμενόμενο τονικό ύψος.

Τα επίπεδα που αφορούν την εκτέλεση \textbf{δεν χρησιμοποιούνται ως πρότυπο αναφοράς}, επειδή αποτελούν ερμηνευτική ανθρώπινη μεταγραφή. Η εκτελεστική πλευρά υπολογίζεται εκ νέου από τον \textbf{ακατέργαστο ήχο} και το \textbf{ηλεκτρογλωττογράφημα}. Έτσι το κριτήριο είναι πράγματι ανεξάρτητο, και το ηλεκτρογλωττογράφημα δίνει ανθεκτική εκτίμηση θεμελιώδους συχνότητας διαχωρίζοντας τη γλωττιδική πηγή από τον φωνητικό σωλήνα.

\subsection{Ανεξαρτησία του τροπικού πυρήνα από τον ήχο}\label{app:splits-trop-audio}

Η τροπική δοκιμή είναι \textbf{συμβολική}. Στηρίζεται σε πηγές που σημειώνουν ισοκράτημα και \emph{δεν} εξαρτάται από την πρόσβαση στα πρωτογενή ηχητικά δεδομένα. Μαζί με τη μουσικολογική και την πληροφοριακή δοκιμή, διατηρεί \textbf{τρεις πλήρεις και διακριτούς μηχανισμούς εγκυρότητας ακόμη και χωρίς ηχητικό υλικό}: μουσικολογικό, πληροφοριακό και τροπικό.


\subsection{Ανθρώπινα δεδομένα, διαχωρισμός ερεθισμάτων και αντιστάθμιση}\label{app:splits-human}

Η M3c και η M3b χρησιμοποιούν κοινό σύνολο διαθέσιμων συμμετεχόντων αλλά διαφορετικές συνεδρίες. Η M3c προηγείται σταθερά της M3b. Κάθε μουσική συνέχεια της M3c παρουσιάζεται μία φορά ανά συμμετέχοντα και αντιστοιχίζεται σε μία μόνο τιμή του $R$. Η λειτουργική και η ουδέτερη συνθήκη κατανέμονται με σχήμα λατινικού τετραγώνου και αντισταθμισμένες λίστες, ώστε κάθε μουσική συνέχεια να παρουσιάζεται και στις δύο συνθήκες σε διαφορετικούς συμμετέχοντες.

Τα ερεθίσματα που χρησιμοποιούνται για γενίκευση εκτός εκπαίδευσης δηλώνονται πριν από την κύρια συλλογή. Η εκτίμηση του ανθρώπινου χώρου των αναπαραστάσεων διαχωρίζει ρητά τα ερεθίσματα που χρησιμοποιούνται για προσαρμογή από εκείνα που χρησιμοποιούνται για CCGP ή ευθυγράμμιση εκτός δείγματος. Η τράπεζα ερεθισμάτων διαθέτει σταθερά αναγνωριστικά και συνδέεται με έργο, πηγή, οικογένεια, εκτελεστή, χαρακτηριστικά επιφανειακής απόστασης, $I_F$, $I_{\mathrm{event}}$ και πειραματική λίστα. Οι δοκιμές που αποτυγχάνουν στον τυφλό έλεγχο μουσικής φυσικότητας αποκλείονται πριν από την κύρια συλλογή.

\section{Στατιστική ανάλυση και πολλαπλότητα}\label{app:stats}

\paragraph{Κοινή σύμβαση πρωτοκόλλου.} Εκτός αν αναφέρεται διαφορετικά, τα κατώφλια, οι μονάδες αναδειγματοληψίας, οι κανόνες αντιστοίχισης και οι επιλογές μοντέλου που εξαρτώνται από την απογραφή καθορίζονται στο ΠΕ0 με τους κανόνες που έχουν οριστεί εκ των προτέρων και δεν μεταβάλλονται μετά την πρόσβαση στα αντίστοιχα αποτελέσματα. Οι επιμέρους ενότητες επαναλαμβάνουν τον χρονισμό αυτό μόνο όταν αποτελεί ειδικό μέρος του ελέγχου διαρροής ή του αν τα διαθέσιμα δεδομένα επιτρέπουν αξιόπιστη εκτίμηση.

Τα στατιστικά μοντέλα λαμβάνουν υπόψη ότι οι παρατηρήσεις έχουν οργανωμένη εξάρτηση. Οι επαναλαμβανόμενες απαντήσεις του ίδιου ακροατή δεν είναι ανεξάρτητες, επειδή μοιράζονται τη μεταβλητότητα που συνδέεται με τον συγκεκριμένο συμμετέχοντα· αντίστοιχα, πολλές εμφανίσεις του ίδιου έργου μοιράζονται ιδιότητες του έργου. Οι \textbf{σταθερές επιδράσεις} εκφράζουν τις προκαθορισμένες σχέσεις ενδιαφέροντος, ενώ οι \textbf{τυχαίες επιδράσεις} αποτυπώνουν τη συστηματική μεταβλητότητα που συνδέεται με επαναλαμβανόμενους συμμετέχοντες, ερεθίσματα, έργα ή οικογένειες. Έτσι η στατιστική ανάλυση δεν αντιμετωπίζει επαναλαμβανόμενες μετρήσεις από την ίδια πηγή σαν ισάριθμα ανεξάρτητα δείγματα.

Όπου χρησιμοποιείται αναδειγματοληψία bootstrap, η μονάδα αναδειγματοληψίας διατηρεί τη δομή εξάρτησης του ερευνητικού υλικού. Όταν πολλές θέσεις προέρχονται από το ίδιο έργο, αναδειγματοληπτείται το έργο ως ομάδα αντί να θεωρούνται οι θέσεις ανεξάρτητες. Στις συγκρίσεις μεταξύ μελοποιών εφαρμόζεται αντίστοιχα μονάδα μελοποιού όπου το εκτιμώμενο μέγεθος το απαιτεί.

\subsection{Επιλογή αναπαράστασης}\label{app:stats-selection}
Η M1 εφαρμόζει τον οριστικό κανόνα J στο πλέγμα $U\times\Phi$ με προκαθορισμένη κοινή αρχιτεκτονική A (joint semi-Markov CRF). Οι B/C αξιολογούνται αυτοτελώς σε όλες τις δώδεκα αναπαραστάσεις. Η βασική μετρική συνδυάζει την ορθή αναγνώριση πλήρους τμήματος και την οικογενειακή ταυτότητα.

\paragraph{Αποδεκτά ζεύγη, αντιστοίχιση και περιοχές αξιολόγησης.}
Κάθε εμφάνιση $x$ έχει οικογένεια $f_x$ και προκαθορισμένο σύνολο πλήρων αποδεκτών ζευγών ορίων $\mathcal B_x$. Η κύρια μετρική $F_{1,\mathrm{seg+fam}}$ χρησιμοποιεί πρόσθετη ανοχή $\tau=0$, μέγιστη αντιστοίχιση ένα προς ένα και ταυτόχρονο έλεγχο $(\hat s_y,\hat e_y)\in\mathcal B_x$ και $\hat f_y=f_x$. Η μεγιστοποίηση γίνεται επί των νόμιμων κοινών αναλύσεων $\mathcal C_{\mathrm{legal}}$, όταν υπάρχουν περιορισμοί μεταξύ γειτονικών εμφανίσεων.
\[
TP=m^\star,\quad FP=P-m^\star,\quad FN=X-m^\star,\qquad
F_{1,\mathrm{seg+fam}}=\frac{2m^\star}{P+X}.
\]
Η κύρια αξιολόγηση εκτελείται σε πλήρως επισημειωμένους μη επικαλυπτόμενους πυρήνες $C$, με επισημειωμένα συμφραζόμενα που καλύπτουν το μέγιστο μήκος τμήματος $L_{\max}$. Κάθε πρόβλεψη και εμφάνιση αναφοράς προσμετράται στον πυρήνα όπου αρχίζει. Πρόβλεψη που αρχίζει έξω από τον πυρήνα εξαιρείται από τον βαθμό του, ενώ πρόβλεψη μέσα στον πυρήνα που υπερβαίνει το $L_{\max}$ καταγράφεται ως μη έγκυρη και ψευδώς θετική. Με μεταβλητό πλήθος εμφανίσεων αναφοράς ανά νόμιμη ανάλυση $c$ ο στόχος γίνεται $\max_{c\in\mathcal C_{\mathrm{legal}}}2m(c)/(P+X(c))$. Η εκπαίδευση σε μερικώς επισημειωμένο υλικό ακολουθεί χωριστό στόχο μερικής εποπτείας.

\paragraph{Οικογενειακή επάρκεια σε γνωστά όρια.}
Το $G_A(R)$ είναι η macro ανάκληση γνωστών οικογενειών, με υποχρεωτική επιλογή ετικέτας και χωρίς αποχή. Από τις εναλλακτικές του $\mathcal B_x$ επιλέγεται μία μουσικολογικά προκαθορισμένη αναφορά $b_x^0$ ανά εμφάνιση, η οποία μαζί με τις υπόλοιπες συνθέτει νόμιμη ανάλυση. Κάθε εμφάνιση μετρά μία φορά· οι λοιπές οριοθετήσεις εξετάζονται σε χωριστή ανάλυση ευαισθησίας. Η συμφωνία επισημειωτών και η κοινή κλίμακα ατομικών γεγονότων ελέγχονται στο ΠΕ0. Τα $F_{1,\mathrm{bnd}}$ και F1 τμημάτων με ανοχή $\tau_b$ αποτελούν διαγνωστικές μετρικές.

\paragraph{Προκαθορισμένο κατώφλι και λειτουργική επιλογή.}
Δύο ειδικοί δηλώνουν ανεξάρτητα το $g_{\mathrm{use}}$ και το $\delta_{\mathrm{use}}>0$ πριν από οποιαδήποτε επίδοση μοντέλου· για κάθε ποσότητα λαμβάνεται η μεγαλύτερη δήλωση. Ένας πιλοτικός medoid αξιολογείται σε χωριστές ομάδες πηγών, με καταγραφή της αβεβαιότητας. Πριν από την κατάταξη ορίζεται
\[
g_{\min}=\max\{g_{\mathrm{use}},\widehat G_0^{\mathrm{pilot}}+\delta_{\mathrm{use}}\}.
\]
Για τουλάχιστον πέντε αντιστοιχισμένες αρχικοποιήσεις ορίζονται
\[
\widehat J(R)=\frac1{n_{\mathrm{rep}}}\sum_{r=1}^{n_{\mathrm{rep}}}
F_{1,\mathrm{seg+fam}}^{(r)}(R,A),
\]
\[
\widehat{\mathcal A}_{\mathrm{dec}}=
\{R\in U\times\Phi:\widehat G_A(R)\ge g_{\min}\},\qquad
R^\star=\arg\max_{R\in\widehat{\mathcal A}_{\mathrm{dec}}}\widehat J(R).
\]
Το κενό επιλέξιμο σύνολο δηλώνει ανεπαρκή ετοιμότητα. Σε ακριβείς ισοπαλίες προτιμώνται λιγότερα πρόσθετα στρώματα και η σταθερή σειρά $U$. Η εμφωλευμένη διασταυρούμενη επικύρωση αξιολογεί τη διαδικασία επιλογής σε εξωτερικές ομάδες. Η μοναδική $R^\star$ οριστικοποιείται στο πλήρες υλικό ανάπτυξης πριν από την Α2/Μ2.

\paragraph{Κοινή στατιστική περιοχή.}
Η ταυτόχρονη περιοχή επιπέδου $1-\alpha$ αφορά τις 66 διαφορές $J(R_i)-J(R_j)$ και τις 12 τιμές $G_A(R)$. Με $L_{iR}$ κάτω όριο του $J(R_i)-J(R)$ και $L_G,U_G$ όρια της οικογενειακής επίδοσης ορίζουμε
\[
\widehat{\mathcal A}^{-}=\{R:L_G(R)\ge g_{\min}\},\quad
\widehat{\mathcal A}^{+}=\{R:U_G(R)\ge g_{\min}\},
\]
\[
\mathcal R^\star=\left\{R\in\widehat{\mathcal A}^{+}:
L_{iR}\le0\ \text{για κάθε }i\in\widehat{\mathcal A}^{-}\setminus\{R\}\right\}.
\]
Όταν $\widehat{\mathcal A}^{-}=\varnothing$, ισχύει $\mathcal R^\star=\widehat{\mathcal A}^{+}$. Ο υπολογισμός της κοινής περιοχής διατηρεί τις συνεκτικές ομάδες έργων--πηγών και κοινές αρχικοποιήσεις για όλες τις συνθήκες. Σε αναδειγματοληψίες με $n_f^{(b)}=0$, η ανάκληση είναι μη ορισμένη: η συχνότητα και η επάρκεια της μεθόδου καταγράφονται χωρίς μεταβολή του πληθυσμιακού στόχου. Η κάλυψη βαθμονομείται στο ΠΕ1. Όταν δεν υπάρχει αξιόπιστη κάλυψη, αναφέρεται συντηρητικά $\mathcal R^\star=U\times\Phi$, χωρίς αλλαγή της λειτουργικής επιλογής $R^\star$ και χωρίς απόδοση ονομαστικής κάλυψης στα επιμέρους διαστήματα.

\paragraph{Οικογενειακή απώλεια και μετάβαση στη Μ2.}
\[
\Delta_G^{\mathrm{loss}}=G_A(R^\star)-\max_{R\in U\times\Phi}G_A(R)\le0,
\]
με συντηρητικό διάστημα προβολής
\[
\Bigl[L_G(R^\star)-\max_RU_G(R),\
\min\{0,U_G(R^\star)-\max_RL_G(R)\}\Bigr].
\]
Το πολύ μία περιγραφική $R_{\mathrm{sens}}$ επιλέγεται από το
$(\mathcal R^\star\cap\widehat{\mathcal A}_{\mathrm{dec}})\setminus\{R^\star\}$ με προτίμηση σε διαφορετικό $U$ και έπειτα στο μέγιστο $\widehat J$. Η κύρια $\Delta_{\mathrm{hist}}$ υπολογίζεται αποκλειστικά στην $R^\star$. Οι A/B/C συγκρίνονται σε όλο το πλέγμα, με κύριο μέγεθος τη μέση διαφορά A--C και συμπληρωματικές τις A--B και boundary A--C. Η B0 ορίζεται ως $A(R^\star)$ και η A διατηρείται ως κοινή υποδομή των Α2/Μ2.


\subsection{Πολλαπλότητα}\label{app:stats-multiplicity}

Η ιεραρχία διαχωρίζει τον προκαθορισμένο υπολογιστικό πυρήνα από την οικογένεια επιβεβαιωτικών ανθρώπινων υποθέσεων και από τα κύρια δευτερεύοντα αποτελέσματα. Το $\Delta_{\mathrm{hist}}$ της M2 διατηρεί το δικό του προκαθορισμένο καθεστώς σύμφωνα με το πρωτόκολλο B2--EISE. Στην ανθρώπινη ενότητα υπάρχουν ακριβώς δύο επιβεβαιωτικές υποθέσεις:
\[
H_{\mathrm{cat}}:\beta_F^{\mathrm{cat}}>0,
\qquad
H_{\mathrm{pred}}:\beta_F^{\mathrm{pred}}>0.
\]
Η Holm διόρθωση εφαρμόζεται στα δύο αντίστοιχα $p$-values. Το $\Delta_{\mathrm{ctx}}$, το $\mathrm{CCGP}_F$, η υπολειμματική ενδοοικογενειακή ευθυγράμμιση EISE--B2 και οι αλληλεπιδράσεις εμπειρίας δηλώνονται ως κύρια δευτερεύοντα και αναφέρονται με μεγέθη επίδρασης, διαστήματα αβεβαιότητας και το προκαθορισμένο καθεστώς ανάλυσής τους. Η βιωματική απόκριση αναλύεται διερευνητικά.

Οι προϋπάρχουσες οικογένειες μουσικολογικών, πληροφοριακών, τροπικών και εκτελεστικών ελέγχων διατηρούν τις δικές τους διορθώσεις εκεί όπου παραμένουν ενεργές. Η ανθρώπινη οικογένεια ελέγχων έχει αυτοτελή διόρθωση πολλαπλότητας, όπως και οι μουσικολογικές και εκτελεστικές οικογένειες σύμφωνα με το δικό τους πρωτόκολλο. Κάθε πίνακας και σύνοψη χρησιμοποιεί την ίδια ονομασία και το ίδιο καθεστώς ανάλυσης για $H_{\mathrm{cat}}$, $H_{\mathrm{pred}}$ και $\Delta_{\mathrm{ctx}}$.

\subsection{Μουσικολογική οριοθέτηση και οικογενειακή αναφορά}\label{app:stats-axisA}

Ορίζονται \textbf{τρεις διακριτοί στόχοι}, με ρητή πηγή ετικετών και ρητό επιλέξιμο υποσύνολο για τον καθένα:

{\small
\begin{xltabular}{\textwidth}{@{}p{0.185\textwidth}Yp{0.215\textwidth}@{}}
\toprule
\textbf{Στόχος} & \textbf{Πηγή ετικετών} & \textbf{Επιλέξιμο υποσύνολο}\tabularnewline
\midrule
\endfirsthead
\toprule
\textbf{Στόχος} & \textbf{Πηγή ετικετών} & \textbf{Επιλέξιμο υποσύνολο}\tabularnewline
\midrule
\endhead
\midrule
\multicolumn{3}{r}{\footnotesize Συνέχεια στην επόμενη σελίδα}\tabularnewline
\endfoot
\bottomrule
\endlastfoot

Όρια εμφάνισης & Τα όρια των τεκμηριωμένων εμφανίσεων στις δημοσιευμένες αναλύσεις & Έργα που καλύπτονται από τουλάχιστον μία αναφορά\tabularnewline
\midrule
Οικογένεια & Το αναγνωριστικό τύπου της αναφοράς & Οικογένειες με επαρκείς εμφανίσεις σε περισσότερα του ενός έργα\tabularnewline
\midrule
Λειτουργικός ρόλος & Ο ρητός λειτουργικός χαρακτηρισμός, όπου η αναφορά τον παρέχει & \emph{Μόνο} οι αναφορές που χαρακτηρίζουν λειτουργία· κλειστό λεξιλόγιο με ρητή τιμή για απούσα ετικέτα\tabularnewline
\end{xltabular}
}\par

Ο λειτουργικός ρόλος αποτελεί \emph{κανονικό τρίτο στόχο} και όχι περιγραφικό χαρακτηριστικό. Επειδή όμως δεν τον παρέχουν όλες οι αναφορές, το επιλέξιμο υποσύνολό του είναι μικρότερο και δηλώνεται χωριστά· η απουσία ετικέτας δεν αντιμετωπίζεται ως αρνητικό παράδειγμα.

Η μορφή της κύριας εξόδου οριστικοποιείται στο ΠΕ0 πριν από τη M1. Όταν η απογραφή στηρίζει τη μη επικαλυπτόμενη διαδρομή, η κύρια εργασία είναι ανίχνευση τμημάτων μεταβλητού μήκους χωρίς κοινά γεγονότα, με δυνατότητα ακάλυπτων περιοχών. Αν ενεργοποιηθεί η προκαθορισμένη διαδρομή επικαλυπτόμενων τμημάτων, ο ίδιος στόχος αξιολογείται πάνω στον κοινό χάρτη εμφανίσεων με την αντίστοιχη μορφή εξόδου. Η αντιστοίχιση οικογένειας γίνεται με χειρισμό \emph{ανοικτού κόσμου}: το σύστημα δεν επιτρέπεται να υποθέτει ότι όλες οι πιθανές οικογένειες θέσεων είναι ήδη γνωστές ή παρούσες στο λεξιλόγιο εκπαίδευσης. Πρέπει επομένως να μπορεί να δηλώσει ότι η σωστή οικογένεια απουσιάζει από το σύνολο υποψηφίων, και η αξιολόγηση αναφέρει πάντοτε κάλυψη και αποχή μαζί με την ακρίβεια.

Η κύρια αντιστοίχιση $F_{1,\mathrm{seg+fam}}$ είναι ένα προς ένα με έλεγχο ολόκληρου αποδεκτού ζεύγους ορίων και σωστής οικογένειας. Η γεωμετρική αντιστοίχιση χωρίς οικογενειακή πρόβλεψη διατηρείται χωριστά για το διαγνωστικό F1 ορίων. Για κάθε αντιστοιχισμένη εμφάνιση επιλέγεται ένα αποδεκτό ζεύγος του $\mathcal B_x$, με τήρηση των κοινών περιορισμών ώστε οι οριοθετήσεις διαφορετικών εμφανίσεων να μπορούν να συνυπάρχουν στην ίδια αποδεκτή ανάλυση. Περισσότερες προβλέψεις γύρω από την ίδια εμφάνιση δεν δημιουργούν περισσότερες σωστές ανιχνεύσεις. Αναφέρονται χωριστά το F1 ορίων, στο οποίο τα δύο άκρα αξιολογούνται διακριτά, και το F1 τμημάτων, στο οποίο απαιτείται επιτυχία και των δύο ορίων ή του προκαθορισμένου κριτηρίου τομής προς ένωση. Η επίδοση στο κοινό πεδίο που μπορεί να αναπαραστήσει η επιλεγμένη έξοδος συνοδεύεται από ανάκληση και F1 στο σύνολο των επαρκώς τεκμηριωμένων εμφανίσεων και από την απώλεια κάλυψης λόγω μήκους, επικάλυψης ή ανεπίλυτης αναφοράς. Οι τεκμηριωμένες επικαλυπτόμενες εμφανίσεις παραμένουν στον παρονομαστή της πλήρους αξιολόγησης.

Για την επιλογή της M1, το κύριο $J(R)$ και το οικογενειακό $G_A(R)$ ορίζονται στην §~\ref{app:stats-selection}. Το F1 πλήρους τμήματος με σωστή οικογένεια καθορίζει τον κανόνα J· τα F1 ορίων, F1 με ανοχή και IoU παραμένουν διαγνωστικά.

\subsection{Πληροφοριακό προφίλ και μέγεθος επίδρασης}\label{app:stats-axisB}

Τα προφίλ \(I_R(u)\) και \(\mathcal{E}_R(u)\) εκφράζονται ως συνάρτηση της κανονικοποιημένης θέσης \(u\in[0,1]\). Η κύρια ανάλυση χρησιμοποιεί ολόκληρη την ενδοτμηματική τροχιά και όχι μόνο τον συνολικό μέσο όρο. Στα ίδια προκαθορισμένα παράθυρα \(W_1,\ldots,W_J\) της κοινής συντεταγμένης ατομικών γεγονότων ορίζεται, για κάθε τμήμα \(s\), ο δείκτης μεταβολής της προβλεπτικής εντροπίας
\[
PG_{R,\mathcal H}(s)
=
\frac{1}{J-1}
\sum_{j=2}^{J}
\left[
\mathcal E_R(W_1)-\mathcal E_R(W_j)
\right].
\]
Θετικό \(PG_{R,\mathcal H}\) σημαίνει ότι η αβεβαιότητα της προβλεπτικής κατανομής είναι κατά μέσο όρο μικρότερη μετά την αρχική περιοχή του τμήματος, περίπου μηδενική τιμή ότι δεν εμφανίζεται αντίστοιχη συστηματική μεταβολή και αρνητική τιμή ότι η αβεβαιότητα αυξάνεται κατά μέσο όρο. Επειδή στον υπολογισμό συμμετέχουν όλα τα μεταγενέστερα παράθυρα, ενδεχόμενη αύξηση της εντροπίας προς το τέλος του τμήματος διατηρείται στην εκτίμηση και δεν αφαιρείται με επιλογή μιας μόνο εσωτερικής περιοχής.

Ο δείκτης συνοψίζεται χωριστά για τις τεκμηριωμένες θέσεις και για καθεμία από τις προκαθορισμένες οικογένειες τμημάτων ελέγχου, με την ίδια ομαδοποίηση και αναδειγματοληψία σε επίπεδο έργου. Το πρωτεύον μέγεθος επίδρασης της Υ3 είναι
\[
\Delta_R^{(\mathcal H)}=
PG_{R,\mathcal H}^{(\text{θέση})}-PG_{R,\mathcal H}^{(\text{ελέγχου})}.
\]
Θετικό \(\Delta_R^{(\mathcal H)}\) δηλώνει ισχυρότερη εσωτερική μείωση της προβλεπτικής αβεβαιότητας στις τεκμηριωμένες θέσεις από ό,τι στα αντιστοιχισμένα τμήματα ελέγχου. Ο ίδιος ορισμός εφαρμόζεται στο πληροφοριακό περιεχόμενο,
\[
PG_{R,h}(s)
=
\frac{1}{J-1}
\sum_{j=2}^{J}
\left[
I_R(W_1)-I_R(W_j)
\right],
\]
και δίνει το συμπληρωματικό \(\Delta_R^{(h)}\). Η προκαθορισμένη ανάλυση με εξαίρεση των πρώτων \(k\) γεγονότων επανυπολογίζει τους ίδιους δείκτες στην περικομμένη τροχιά. Έτσι ελέγχεται αν η διαφοροποίηση παραμένει μετά την αφαίρεση της άμεσης περιοχής του ορίου και δεν οφείλεται αποκλειστικά σε αρχική πληροφοριακή κορυφή.

\subsection{Εκτελεστική μεταβλητότητα και εξωτερική πρόβλεψη}\label{app:stats-performance}

Η εκτελεστική δοκιμή ελέγχει \textbf{ποια αναπαράσταση προβλέπει τη δομή αυτής της εκτελεστικής μεταβλητότητας} σε έργα και εκτελεστές που έχουν διατηρηθεί εκτός εκπαίδευσης.

Η Υ5 εξετάζει ειδικά αν η αναπαράσταση που έχει οριστικοποιηθεί από τη μουσικολογική δοκιμή προσφέρει, χωρίς επαναπροσαρμογή με βάση την εκτελεστική δοκιμή, πρόσθετη εκτός δείγματος πληροφορία για τις εκτελεστικές αποκλίσεις και τη διασπορά μεταξύ των εκτελεστών \textbf{σε σύγκριση με μοντέλο αναφοράς που χρησιμοποιεί αποκλειστικά το αναμενόμενο τονικό ύψος}. Η Υ5 εφαρμόζεται στην $R^\star$· η αβεβαιότητα της επιλογής αναφέρεται μέσω του $\mathcal R^\star$ (§~\ref{app:stats-selection}) και η εκτελεστική δοκιμή δεν χρησιμοποιείται για εκ των υστέρων επιλογή ή αντικατάσταση της $R^\star$. Τυχόν ισχυρισμός στατιστικής ισοδυναμίας απαιτεί προκαθορισμένο αμφίπλευρο περιθώριο, ενώ έλεγχος μη κατωτερότητας δηλώνεται χωριστά με δικό του μονόπλευρο περιθώριο.

Για κάθε επίπεδο \(R\) και με \emph{κοινό} στατιστικό μοντέλο εξετάζεται πόσο καλά οι μονάδες και τα πεδία του \(R\) εξηγούν τη μέση απόκλιση και, κυρίως, τη \emph{διασπορά} της:
\[
\mathbb{E}[d_{pe}\mid R],\qquad \operatorname{Var}(d_{pe}\mid R)
\]
με αξιολόγηση σε \textbf{έργα και εκτελεστές που έχουν διατηρηθεί εκτός εκπαίδευσης}.

Επειδή το ερώτημα αφορά \emph{και} τη διασπορά, ένα συμβατικό μοντέλο μικτών επιδράσεων δεν επαρκεί: μοντελοποιεί τη μέση τιμή και αντιμετωπίζει τη διασπορά ως θόρυβο. Χρησιμοποιείται επομένως \textbf{ιεραρχικό μοντέλο κεντρικής τάσης και κλίμακας (location--scale)}, στο οποίο και οι δύο παράμετροι εξαρτώνται από την αναπαράσταση:
\[
\begin{aligned}
d_{pe}&\sim\mathcal{N}(\mu_{pe},\sigma^2_{pe}),\\
\mu_{pe}&=f(R_e)+u_p+v_w+\beta_C C_p+\boldsymbol\beta^{\top}\mathbf z_e,\\
\log\sigma_{pe}&=g(R_e)+u^{(\sigma)}_p+v^{(\sigma)}_w+\boldsymbol\gamma^{\top}\mathbf z_e.
\end{aligned}
\]
όπου \(R_e\) τα γνωρίσματα της αναπαράστασης στο γεγονός \(e\), \(u_p,v_w\) τυχαίες επιδράσεις εκτελεστή και έργου, \(C_p\) η κλάση φωνής και \(\mathbf z_e\) το προκαθορισμένο διάνυσμα συγχυτικών παραγόντων: εσωτερική θέση, εγγύτητα σε κατάληξη, ήχος, μελισματική πυκνότητα και πυκνότητα μελωδικών γεγονότων, τοπική πολυπλοκότητα, θέση στη φράση και η αντίστοιχη άμεση συμβολική εκτίμηση. Η συνάρτηση \(g\) είναι το ουσιώδες: επιτρέπει να δηλωθεί ότι μια αναπαράσταση προβλέπει τη \emph{δομή της διασποράς} και όχι απλώς τη μέση απόκλιση. Η κανονική οικογένεια αποτελεί το αρχικό μοντέλο εργασίας· πριν από την αξιολόγηση στο οριστικοποιημένο σύνολο ελέγχου προκαθορίζεται έλεγχος καταλοίπων και, εφόσον απαιτείται, ανθεκτική εναλλακτική ή κατανομή με βαριές ουρές, ώστε η επιλογή κατανομής να μην θεωρηθεί αυτονόητη. Η ακριβής οικογένεια μοντέλων οριστικοποιείται πριν από την τελική αξιολόγηση.

Ο \textbf{έλεγχος της κλάσης φωνής είναι υποχρεωτικός}: επειδή η δειγματοληψία των είκοσι ψαλτών ήταν στρωματοποιημένη κατά κλάση φωνής, οποιοδήποτε φαινόμενο μεταβλητότητας τονικού ύψους θα μπορούσε να αντανακλά φωνητική περιοχή αντί για μουσική δομή.

Πρόσθετος \textbf{διαγνωστικός έλεγχος}: τα επιλεγμένα ηχοχρωματικά χαρακτηριστικά χρησιμοποιούνται ως συμπληρωματικοί δείκτες φωνητικής ταυτότητας. Τυχόν σύνδεσή τους με τη μουσική δομή συνυπολογίζεται στην ερμηνεία της ειδικότητας των αποτελεσμάτων, ώστε η σχέση προς τις μεταβλητές ύψους και χρονισμού να αξιολογείται απέναντι και σε χαρακτηριστικά που εξαρτώνται ισχυρά από τον εκτελεστή.

\subsection{Τροπικός εκτιμητής και κύρια εργασία}\label{app:stats-trop}

Για κάθε συνθήκη του πλέγματος, η εργασία είναι \textbf{πρόβλεψη των σημείων μετακίνησης της ισοκρατηματικής βάσης} σε έργα που έχουν διατηρηθεί εκτός εκπαίδευσης, με το ίδιο κοινό πρωτόκολλο, την ίδια γεωμετρία εξόδου και τους ίδιους διαχωρισμούς που χρησιμοποιούνται σε όλες τις σχετικές δοκιμές. Η ισοκρατηματική τροχιά εκφράζεται στο κοινό πλέγμα ως διαδοχή τμημάτων σταθερής βάσης, με τα όρια στον σημειογραφικό χρόνο.

Το κύριο μοντέλο αναφοράς είναι \textbf{πλήρες τοπικό τροπικό μοντέλο αναφοράς} που διαθέτει όλα τα επιτρεπτά γνωρίσματα τα οποία μπορούν να εξηγήσουν μετακίνηση της ισοκρατηματικής βάσης χωρίς να χρησιμοποιούν δομή θέσης: τοπική διαστηματική τροχιά, ενεργή τροπική κατάσταση που είναι ήδη διαθέσιμη στη γραφή, εγγύτητα σε κατάληξη, θέση στη φράση και βασικά τοπικά γνωρίσματα πυκνότητας. Η οικογενειακή ταυτότητα και τα όρια θέσεων αποκλείονται.

Το επιβεβαιωτικό ερώτημα είναι η \textbf{πρόσθετη εκτός δείγματος συνεισφορά της αναπαράστασης που προκύπτει από την εκπαίδευση της θέσης σε σύγκριση με αυτό το μοντέλο αναφοράς}, στο πλήρες επιλέξιμο σύνολο. Η ίδια οριστικοποιημένη σύγκριση επαναλαμβάνεται ως προκαθορισμένος έλεγχος ευαισθησίας στο μη καταληκτικό υποσύνολο, με αναφορά μεγέθους επίδρασης και διαστήματος εμπιστοσύνης. Η ερμηνεία της τροπικής λειτουργίας σε μη καταληκτικά σημεία στηρίζεται στο αποτέλεσμα αυτού του υποσυνόλου. Τα γνωρίσματα φράσης και κατάληξης παράγονται από τις επιτρεπτές πηγές της γραφής και του κειμένου με τον κανόνα προέλευσης της ενότητας διαρροής· οι ετικέτες ισοκρατήματος διατηρούνται αποκλειστικά ως στόχος.

\subsection{Πολυμεθοδική σύγκλιση}\label{app:stats-convergence}

Η πολυμεθοδική σύγκλιση θέτει το αυστηρότερο ερώτημα της αρχιτεκτονικής: αν \textbf{η ίδια} αναπαράσταση αναδεικνύεται από \textbf{τέσσερις διαφορετικούς μηχανισμούς τεκμηρίωσης}, με διαφορετικές πηγές παρατήρησης και σφάλματος. Η σύγκλιση ανεξάρτητων μεθόδων αποτελεί έτσι άμεσο αντικείμενο μέτρησης και ενοποιεί τα επιμέρους αποτελέσματα σε κοινό ισχυρισμό για τη μουσική μονάδα.

Εξετάζονται έξι προκαθορισμένες κατά ζεύγη σχέσεις: η μουσικολογική οριοθέτηση με την πληροφοριακή, την τροπική και την εκτελεστική δοκιμή, η πληροφοριακή δοκιμή με την τροπική και την εκτελεστική, και η τροπική με την εκτελεστική δοκιμή.
δηλαδή: συμπίπτουν τα όρια που αναδεικνύει η μείωση της προβλεπτικής αβεβαιότητας με τα τεκμηριωμένα μουσικολογικά όρια; συμπίπτουν τα τεκμηριωμένα όρια με τις μετακινήσεις της τροπικής βάσης και με τις μεταβολές της εκτελεστικής διασποράς; και συμπίπτουν μεταξύ τους οι περιοχές που αναδεικνύουν τα κριτήρια που δεν χρησιμοποιούν καθόλου δημοσιευμένα όρια;

\textbf{Η σύγκλιση μουσικολογικής οριοθέτησης και τροπικής δοκιμής έχει ξεχωριστή μουσικολογική βαρύτητα.} Η σύμπτωση, πάνω από το τυχαίο, των τεκμηριωμένων ορίων θέσεων με τα σημεία όπου η ψαλτική ή εκδοτική πράξη μετακινεί την τροπική βάση θα συνδέσει δύο διακριτές παραδόσεις μουσικής κρίσης --- την αναλυτική και την ισοκρατηματική --- στην ίδια δομική διαίρεση του μέλους. Ο βαθμός αυτής της σύμπτωσης ποσοτικοποιεί άμεσα τη σχέση ανάμεσα στη θέση και στην τοπική τροπική οργάνωση. Επειδή οι καταλήξεις αποτελούν φυσικά σημεία τόσο δομικής όσο και τροπικής άρθρωσης, η ίδια σχέση αναφέρεται και σε \textbf{μη καταληκτικά όρια θέσεων}: η διατήρηση της σύγκλισης σε αυτό το υποσύνολο υποστηρίζει την ερμηνεία της θέσης ως τροπικά λειτουργικής δομικής μονάδας και όχι απλώς ως επαναλαμβανόμενου καταληκτικού σχήματος.

Η σύγκλιση πληροφοριακής και εκτελεστικής δοκιμής εξετάζεται σε δύο συμπληρωματικές διατυπώσεις. Η \textbf{πρωτεύουσα, μη επισημειωμένη διατύπωση} υπολογίζει πάνω στα ίδια προκαθορισμένα παράθυρα το πληροφοριακό προφίλ \(I_R(W)\), το προφίλ εντροπίας της προβλεπτικής κατανομής \(\mathcal E_R(W)\) και το εκτελεστικό προφίλ \(V(W)\), χωρίς χρήση μουσικολογικών ετικετών. Η αντιστοιχία τους μετράται ως προς τιμές, τοπικές μεταβολές και θέσεις ακροτάτων με στρωματοποιημένη κατανομή αναφοράς της μηδενικής υπόθεσης που ελέγχει θέση στο έργο, εγγύτητα σε κατάληξη, πυκνότητα και ήχο. Η σύμπτωση πληροφοριακής δομής και εκτελεστικής συμπεριφοράς στα ίδια σημεία αποτελεί τον αυστηρότερο έλεγχο που δεν εξαρτάται από μουσικολογικές επισημειώσεις του σχεδιασμού.

Η \textbf{συμπληρωματική επισημειωμένη διατύπωση} χρησιμοποιεί το μέγεθος \(\Delta_R\) της πληροφοριακής δοκιμής και ποσοτικοποιεί πόσο η πληροφοριακή διαφοροποίηση γύρω από τεκμηριωμένες θέσεις συνδέεται με την εκτελεστική μεταβλητότητα. Οι δύο διατυπώσεις απαντούν σε διαφορετικά αλλά συμπληρωματικά ερωτήματα: η πρώτη ελέγχει ανεξάρτητη τοπική σύγκλιση και η δεύτερη συνδέει τη σύγκλιση με τη μουσικολογικά τεκμηριωμένη θέση.

\subsection{Υ8: διακρίνουσα εγκυρότητα}\label{app:stats-disc}

Η επικύρωση με ανεξάρτητες μεθόδους ολοκληρώνεται με \textbf{τεκμηρίωση διακρίνουσας εγκυρότητας}. Οι κύριες τοπικές αναλύσεις της πολυμεθοδικής σύγκλισης επαναλαμβάνονται σε τρεις εύλογες εναλλακτικές τμηματοποιήσεις: \(S_{\mathrm{fixed}}\), ισομήκη τμήματα σταθερού μήκους· \(S_{\mathrm{cadence}}\), όρια που καθορίζονται αποκλειστικά από καταλήξεις· και \(S_{\mathrm{phrase}}\), ευρύτερα φραστικά όρια. Η πλήρης τριπλή αντίθεση εκτιμάται στο κοινό επιλέξιμο σύνολο όπου και οι τρεις εναλλακτικές τμηματοποιήσεις ορίζονται με τους προκαθορισμένους κανόνες· η κάλυψη αναφέρεται ρητά και δεν συμπληρώνεται με μουσικολογικά ή προβλεπόμενα όρια. Όπου οι μετρικές είναι συγκρίσιμες, τα μεγέθη σύγκλισης τυποποιούνται στην ίδια κλίμακα.

Για να συγκριθούν οι διαφορετικές μορφές σύγκλισης σε κοινή κλίμακα, κάθε προκαθορισμένη κατά ζεύγη σχέση \(r\) μετατρέπεται σε τυποποιημένο μέτρο υπερβάλλουσας σύγκλισης ως προς τη δική της στρωματοποιημένη κατανομή αναφοράς:
\[
C_r^{\mathrm{std}}(S)
=
\frac{C_r(S)-\mu_{0,r}(S)}{\sigma_{0,r}(S)},
\]
όπου \(C_r(S)\) είναι το προκαθορισμένο μέτρο σύγκλισης για την τμηματοποίηση \(S\), ενώ \(\mu_{0,r}(S)\) και \(\sigma_{0,r}(S)\) είναι αντίστοιχα ο μέσος και η τυπική απόκλιση της κατανομής αναφοράς της μηδενικής υπόθεσης. Όλες οι επιμέρους μετρικές προσανατολίζονται πριν από την τελική αξιολόγηση έτσι ώστε μεγαλύτερη τιμή να σημαίνει ισχυρότερη σύγκλιση.

Για τις έξι προκαθορισμένες σχέσεις
\[
\mathcal R_{\Sigma}=
\{r_{\mathrm{μουσ.-πληρ.}},\,r_{\mathrm{μουσ.-τροπ.}},\,r_{\mathrm{μουσ.-εκτ.}},\,
r_{\mathrm{πληρ.-τροπ.}},\,r_{\mathrm{πληρ.-εκτ.}},\,r_{\mathrm{τροπ.-εκτ.}}\},
\]
ο συνοπτικός δείκτης ορίζεται χωρίς προσαρμοζόμενα βάρη ως
\[
\mathrm{Conv}(S)
=
\frac{1}{6}
\sum_{r\in\mathcal R_{\Sigma}}C_r^{\mathrm{std}}(S).
\]
Ο δείκτης εκφράζει τη μέση σύγκλιση των ανεξάρτητων μεθόδων της συγκεκριμένης τμηματοποίησης σε σχέση με τις κατάλληλες μηδενικές κατανομές. Μαζί με τον συνοπτικό δείκτη αναφέρονται υποχρεωτικά και οι έξι τυποποιημένες συνιστώσες του, ώστε αντίθετη ή ασθενής συμπεριφορά μιας συγκεκριμένης σχέσης να μην αποκρύπτεται από τον μέσο όρο.

Η διακρίνουσα σύγκριση πραγματοποιείται στο κοινό επιλέξιμο σύνολο και ορίζει
\[
\Delta_{\mathrm{Conv},c}
=
\mathrm{Conv}(S_{\mathrm{θέσης}})-\mathrm{Conv}(S_c),
\qquad
S_c\in\{S_{\mathrm{fixed}},S_{\mathrm{cadence}},S_{\mathrm{phrase}}\}.
\]
Η ισχυρή πρόβλεψη της Υ8 είναι ότι και οι τρεις αντιθέσεις έχουν θετική κατεύθυνση με την προκαθορισμένη ταυτόχρονη αποτίμηση αβεβαιότητας:
\[
\Delta_{\mathrm{Conv,fixed}}>0,\qquad
\Delta_{\mathrm{Conv,cadence}}>0,\qquad
\Delta_{\mathrm{Conv,phrase}}>0.
\]
Η σύνοψη αυτή δεν αντικαθιστά την ερμηνεία των έξι επιμέρους σχέσεων. Η τελευταία παραμένει αναγκαία για να προσδιορίζεται ποια ανεξάρτητα είδη τεκμηρίωσης συγκλίνουν και ποια όχι.

\subsection{\texorpdfstring{Πρόσθετη εξωτερική εγκυρότητα και \(\Delta_{\mathrm{ext}}\)}{Πρόσθετη εξωτερική εγκυρότητα και Delta ext}}\label{app:stats-ext}

Η γενικότερη διάκριση μεταξύ επαναλαμβανόμενου προτύπου και πολιτισμικά θεμελιωμένης δομικής μονάδας αποκτά ενιαίο προκαθορισμένο κριτήριο. Για κάθε υποψήφια μονάδα $Q$ επιλέγεται από το αλγοριθμικά επαγόμενο λεξιλόγιο μία αντιστοιχισμένη συνθήκη επαναλαμβανόμενου προτύπου αναφοράς $P^{\star}$ με αντίστοιχο μήκος, συχνότητα εμφάνισης, προβλεπτική χρησιμότητα, ήχο και σχετική θέση στο έργο. Η αντιστοίχιση πραγματοποιείται στο υλικό εκπαίδευσης και ρύθμισης και εφαρμόζεται αυτούσια στο οριστικοποιημένο σύνολο αξιολόγησης.

Ορίζεται το \textbf{διάνυσμα πρόσθετης εξωτερικής εγκυρότητας}
\[
\boldsymbol{\Delta}_{\mathrm{ext}}(Q,P^{\star})=
\bigl(\Delta_A,\Delta_{\Delta},\Delta_{\Gamma}\bigr),
\]
όπου $\Delta_A$ είναι η τυποποιημένη αύξηση στην οικογενειακή ταυτότητα/κάλυψη της μουσικολογικής δοκιμής, $\Delta_{\Delta}$ η αύξηση στην πρόβλεψη ισοκρατηματικής μετακίνησης και $\Delta_{\Gamma}$ η αύξηση στην πρόβλεψη του εκτελεστικού αποτυπώματος. Οι τρεις αυξήσεις υπολογίζονται στο ίδιο επιλέξιμο υλικό και συνοδεύονται από ταυτόχρονα διαστήματα εμπιστοσύνης από αναδειγματοληψία bootstrap κατά συστάδες.

Η \textbf{κύρια επιβεβαιωτική διάκριση ``επαναλαμβανόμενο πρότυπο $\rightarrow$ μουσική μονάδα''} τεκμηριώνεται όταν
\[
\min_{j\in\{A,\Delta,\Gamma\}}
\inf\,\mathrm{CI}^{\mathrm{simult}}_{1-\alpha}(\Delta_j)>0.
\]
Το πλήρες $\boldsymbol{\Delta}_{\mathrm{ext}}$ αναφέρεται σε κάθε περίπτωση ως προφίλ της μουσικής μονάδας. Έτσι η ερευνητική πρόταση ποσοτικοποιεί \textbf{ποια εξωτερική μουσική πληροφορία προσθέτει η θέση σε σύγκριση με όσα εξηγούν η επαναληπτικότητα, το μήκος και η προβλεπτική χρησιμότητα ενός προτύπου}.
Εφόσον ικανοποιηθεί το προκαθορισμένο κριτήριο σύγκλισης, το εύρημα θα υποστηρίζει τη μουσική θέση ως δομική μονάδα που διατηρεί περισσότερη εξωτερικά τεκμηριωμένη μουσική πληροφορία από τις εναλλακτικές τμηματοποιήσεις.

\subsection{\texorpdfstring{Διπλή διατήρηση και \(\Delta V\)}{Διπλή διατήρηση και Delta V}}\label{app:stats-double}

Στις οικογένειες θέσεων με ιστορική κάλυψη και κάλυψη από πολλαπλές εκτελέσεις εξετάζεται αν η ίδια μονάδα διατηρεί οικογενειακή ταυτότητα μαζί με την εσωτερική οργάνωση της παραλλαγής. Η Υ2\textsubscript{E} ελέγχει το πρώτο επίπεδο. Το δεύτερο ελέγχεται κυρίως με το συγκριτικό πείραμα EISE της \S\ref{sec:eise-transfer}, το οποίο μετατρέπει τις οικογενειακές συσχετίσεις σε δύο προκαθορισμένες αντιθέσεις πρόσθετης αξίας. Η παρούσα ενότητα ορίζει τα μουσικολογικά προφίλ και τις προκαθορισμένες εξειδικεύσεις τους. Όλες χρησιμοποιούν την κοινή μουσικά τεκμηριωμένη συντεταγμένη $u$ της \S\ref{sec:latent-metric}.

Για οικογένεια $F$ ορίζεται το συνολικό ιστορικό προφίλ
\[
\tau_{\mathrm{hist}}^{F}(u)=\text{μεταβλητότητα μεταξύ ιστορικά συνδεδεμένων εξηγητικών μαρτύρων},
\]
και το ανεξάρτητο εκτελεστικό προφίλ
\[
\tau_{\mathrm{perf}}^{F}(u)=\text{μεταβλητότητα μεταξύ σύγχρονων εκτελεστών}.
\]
Στο υποσύνολο όπου η προέλευση των εξηγητικών σχέσεων έχει ταξινομηθεί με ασφάλεια, το ιστορικό προφίλ αναλύεται περαιτέρω σε
\[
\tau_{\mathrm{an}}^{F}(u)\quad\text{και}\quad\tau_{\mathrm{tr}}^{F}(u),
\]
που αντιστοιχούν αντίστοιχα σε αναλυτική απόκλιση $E_{\mathrm{an}}$ και μελωδική παραλλαγή μέσα σε τεκμηριωμένη ιστορική σχέση $E_{\mathrm{tr}}$.

Από το σημείο αυτό, όταν το $F$ εμφανίζεται ως δείκτης σε ποσότητες όπως $V_F$, $\rho_F$ ή $\overline h_F$, δηλώνει οικογένεια θέσεων. Η χρήση αυτή διακρίνεται από τον παράγοντα γνωρισμάτων $F$ του πλέγματος $U\times F$ της M1.

Η περιγραφική οικογενειακή σχέση γράφεται
\[
\rho_F=\operatorname{Spearman}_{u}\!\left(\tau_{\mathrm{hist}}^{F}(u),\tau_{\mathrm{perf}}^{F}(u)\right),
\]
με έλεγχο ήχου, εγγύτητας σε κατάληξη, μελισματικής πυκνότητας και σχετικής θέσης στο έργο και με στατιστική διαδικασία που λαμβάνει υπόψη την αυτοσυσχέτιση κατά μήκος του $u$. Ο κύριος επιβεβαιωτικός ισχυρισμός στηρίζεται στα συγκριτικά μεγέθη του ΠΕ3β· οι μεμονωμένες οικογενειακές συσχετίσεις αναφέρονται ως εκτιμήσεις προφίλ. Η ταξινόμηση προέλευσης επιτρέπει δύο επιπλέον, μουσικολογικά ερμηνεύσιμες εκτιμήσεις:
\[
\rho_F^{\mathrm{an}}=\operatorname{Spearman}_{u}\!\left(\tau_{\mathrm{an}}^{F}(u),\tau_{\mathrm{perf}}^{F}(u)\right),
\qquad
\rho_F^{\mathrm{tr}}=\operatorname{Spearman}_{u}\!\left(\tau_{\mathrm{tr}}^{F}(u),\tau_{\mathrm{perf}}^{F}(u)\right).
\]
Θετική $\rho_F^{\mathrm{an}}$ σημαίνει ότι οι περιοχές στις οποίες οι τεκμηριωμένες αναλυτικές γραφές αποκλίνουν συστηματικά είναι επίσης περιοχές μεγαλύτερης μετρημένης μεταβλητότητας μεταξύ σύγχρονων εκτελεστών. Θετική $\rho_F^{\mathrm{tr}}$ σημαίνει ότι η σχετική κατανομή της ιστορικά τεκμηριωμένης μεταβλητότητας παρουσιάζει αντίστοιχη κατάταξη στο προφίλ μεταβλητότητας των σύγχρονων εκτελέσεων. Η ταυτόχρονη παρουσία των δύο σχέσεων παρέχει προκαθορισμένη μουσικολογική εξειδίκευση της ένδειξης ότι η πολιτισμική ταυτότητα της θέσης περιλαμβάνει \textbf{δομημένο χώρο πραγματώσεων που υπερβαίνει μία μοναδική επιφανειακή συμβολική ανάλυση}.

Η ερμηνεία παραμένει άμεσα συνδεδεμένη με τα παρατηρούμενα δεδομένα: η αναλυτική και η εκτελεστική διαφοροποίηση συσχετίζονται ως ανεξάρτητες εκδηλώσεις της ίδιας μονάδας, ενώ η ιστορική βιβλιογραφία για τη σχέση γραπτής και προφορικής μετάδοσης παρέχει το μουσικολογικό πλαίσιο αυτής της σύγκλισης. Τα προφίλ εντροπίας της προβλεπτικής κατανομής, η ισοκρατηματική γεωμετρία και τα αντιστοιχισμένα $P^{\star}$ χρησιμοποιούνται ως δευτερεύουσες ερμηνευτικές και διακρίνουσες δοκιμές.
Στο υλικό όπου και οι τέσσερις μηχανισμοί είναι ταυτόχρονα διαθέσιμοι, εξετάζεται η τετραπλή τομή \(\text{Α}\cap\text{Β}\cap\text{Δ}\cap\text{Γ}\). Το πληρέστερο δυνατό αποτέλεσμα της έρευνας είναι η ανάδειξη της ίδιας μονάδας και της ίδιας δομικής κλίμακας από \textbf{μουσικολογική ανάλυση, πληροφοριακή οργάνωση, τροπική πρακτική και πραγματική εκτέλεση}.

Οι ζευγαρωτές συγκρίσεις που περιλαμβάνουν τη M4 εκτελούνται \textbf{μόνο} στο υλικό όπου τα αντίστοιχα κριτήρια επικαλύπτονται· η σχέση Α\(\leftrightarrow\)Β είναι διαθέσιμη σε ολόκληρο το συμβολικό σώμα. Η διαδρομή επικάλυψης είναι συγκεκριμένη και ελέγξιμη: το ηχητικό σώμα περιλαμβάνει \textbf{Χερουβικά} ---εκτενή μελισματικά μέλη του ύμνου των Χερουβείμ στη Θεία Λειτουργία---, τα οποία ανήκουν στο παπαδικό μέλος, και οι κατάλογοι παπαδικών θέσεων περιλαμβάνουν εκτενές τμήμα Χερουβικών. Η πραγματική τομή υπολογίζεται στο πρώτο πακέτο εργασίας.

Ως συμπληρωματική διερευνητική αντίθεση ορίζεται η διαφορά
\[
\Delta V=V_{\text{όριο}}-V_{\text{εσωτερικό}},
\]
όπου \(V\) είναι το μέτρο μεταβλητότητας μεταξύ εκτελεστών. Η μηδενική υπόθεση ελέγχεται με αμφίπλευρο στρωματοποιημένο έλεγχο μεταθέσεων (permutation test), χρησιμοποιώντας το \(|\Delta V|\) ως στατιστικό. Η κατεύθυνση της επίδρασης αναφέρεται ως αποτέλεσμα.

\textbf{Η κατανομή αναφοράς της μηδενικής υπόθεσης ελέγχει ρητά τη φραστική και καταληκτική γεωμετρία.} Συνεκτιμά τον αριθμό ορίων, το μήκος του έργου και την ακόλουθη εναλλακτική διαδρομή:
\[
\text{όριο θέσης}\ \leftrightarrow\ \text{κατάληξη ή φραστικό τέλος}\ \leftrightarrow\ \text{εκτελεστική μεταβλητότητα}
\]
Η κατανομή αναφοράς της μηδενικής υπόθεσης στρωματοποιείται επομένως ως προς τη θέση μέσα στο έργο, τη συλλαβική και μελισματική πυκνότητα, την εγγύτητα σε κατάληξη ή φραστικό όριο, τον ήχο και τη σημειογραφική διάρκεια. Χωριστή ανάλυση των μη καταληκτικών ορίων θέσεων απομονώνει περαιτέρω τη συνεισφορά της θέσης από τη γενική καταληκτική συμπεριφορά.

\subsection{Μεταφορά EISE: πλήρης στατιστικός έλεγχος}\label{app:stats-eise-transfer}

Για κάθε επιλέξιμη οικογένεια και κάθε συνθήκη $m$ εκτιμάται
\[
\rho_F^{(m)}=\operatorname{Spearman}_{u}\!\left(
\widehat\tau_{m,\mathrm{hist}}^F(u),
\tau_{\mathrm{perf}}^F(u)
\right),
\]
στο κοινό επιλέξιμο σύνολο οικογενειών $\mathcal F_{\mathrm{eval}}$, με κοινό πεδίο αντιστοίχισης για όλες τις συνθήκες. Το σύνολο οριστικοποιείται βάσει τεκμηρίωσης, κάλυψης και ποιότητας μέτρησης. Η κύρια σύγκριση B0/B1/B2/EISE χρησιμοποιεί κοινό, εξωτερικά τεκμηριωμένο μητρώο εμφανίσεων με σταθερό αναγνωριστικό εμφάνισης, σύνθεση, μάρτυρα, πηγή, κοινό χειρισμό των αποδεκτών ζευγών ορίων, οικογενειακή απόδοση αναφοράς, διαθεσιμότητα της αντιστοίχισης $u$ και βάρος στη σύνοψη. Η αυτόματη ανίχνευση και η κάλυψή της αξιολογούνται χωριστά και δεν επιτρέπεται να ορίζουν διαφορετικό κύριο υποσύνολο για κάθε μοντέλο. Οι αποκλεισμοί λόγω ανεπαρκούς τεκμηρίωσης ή αντιστοίχισης εφαρμόζονται κοινά πριν από την εξέταση των αποτελεσμάτων και αναφέρονται ως περιορισμοί του πεδίου της αντίθεσης. Για κάθε ομάδα επανάληψης $k$ και ρύθμιση τυχαιότητας $\ell$ ορίζεται
\[
S_{m,k,\ell}=\frac{1}{|\mathcal F_{\mathrm{eval}}|}
\sum_{F\in\mathcal F_{\mathrm{eval}}}\rho_F^{(m,k,\ell)}.
\]
Κάθε ομάδα $k$ αντιστοιχεί σε μία από $K$ επιλέξιμες ζευγοποιήσεις του B2. Οι B1 και EISE εκτελούνται επίσης σε $K\times L$ επαναλήψεις, με κοινή αρχικοποίηση και κοινή τυχαιότητα εκπαίδευσης ανά $(k,\ell)$ μεταξύ των συνθηκών, ώστε να αντιστοιχίζεται και ο συνολικός υπολογιστικός προϋπολογισμός. Ο EISE διατηρεί τις ίδιες επαληθευμένες ιστορικές ακμές σε όλες τις ομάδες. Για κάθε συνθήκη που προκύπτει από εκπαίδευση ορίζεται
\[
S_m=\frac{1}{KL}\sum_{k=1}^{K}\sum_{\ell=1}^{L}S_{m,k,\ell},
\qquad m\in\{B1,B2,\mathrm{EISE}\}.
\]
Το $S_{B0}$ είναι ο ίδιος ισοβαρής μέσος οικογενειακών συσχετίσεων της άμεσης συμβολικής εκτίμησης. Τα $K$, $L$, οι κανόνες δειγματοληψίας ζευγοποιήσεων και ο συνολικός υπολογιστικός προϋπολογισμός και προϋπολογισμός ρύθμισης υπερπαραμέτρων οριστικοποιούνται πριν από τον έλεγχο. Οι πολλαπλές εκπαιδεύσεις του B2 αποτελούν επαναλήψεις του συγκριτή και εισέρχονται στην αβεβαιότητα· το B2 εκπροσωπείται από τον παραπάνω μέσο και από την πλήρη κατανομή των επαναλήψεών του. Τα δύο αποφασιστικά μεγέθη είναι
\[
\Delta_{\mathrm{hist}}=S_{\mathrm{EISE}}-S_{B2},
\qquad
\Delta_{\mathrm{learn}}=S_{\mathrm{EISE}}-S_{B0}.
\]
Θετικό $\Delta_{\mathrm{hist}}$, υποστηριζόμενο από τον προκαθορισμένο έλεγχο αβεβαιότητας, τεκμηριώνει υψηλότερη επίδοση της εποπτείας μέσω τεκμηριωμένων εξηγητικών σχέσεων στο προκαθορισμένο ανεξάρτητο εκτελεστικό κριτήριο αξιολόγησης έναντι της αντιστοιχισμένης οικογενειακής εποπτείας. Θετικό $\Delta_{\mathrm{learn}}$ με την αντίστοιχη στατιστική τεκμηρίωση δείχνει ότι η αναπαράσταση που προκύπτει από την εκπαίδευση αποτυπώνει περισσότερη δομή που συμφωνεί με το ανεξάρτητο εκτελεστικό κριτήριο από ό,τι η πλήρης άμεση συμβολική περιγραφή.

Η \textbf{προκαθορισμένη εκτελεστική μεταβλητή αξιολόγησης} για την κεντρική αξιολόγηση του EISE είναι η διασπορά μεταξύ των εκτελεστών του κανονικοποιημένου λογαριθμικού χρονισμού, εφόσον το ΠΕ0 επιβεβαιώσει επαρκή αξιοπιστία μέτρησης:
Ο χρονικός εκτιμητής αφαιρεί από την εκτελεστική διάρκεια δύο διαφορετικές πηγές κλίμακας. Για τον εκτελεστή $p$, το έργο ή τη φράση $w$ και το γεγονός $e$, το $d_{pwe}$ είναι η μετρημένη διάρκεια σε χρόνο, ενώ το $q_{we}$ είναι η σχετική σημειογραφική χρονική αξία σε κοινή μονάδα ανάγνωσης. Ο λόγος $d_{pwe}/q_{we}$ εκφράζει επομένως εκτελεστικό χρόνο ανά μονάδα σημειογραφικής αξίας. Ο λογάριθμος μετατρέπει τις αναλογικές μεταβολές σε προσθετική κλίμακα και το $\alpha_{pw}$ αφαιρεί τον συνολικό όρο αγωγής της συγκεκριμένης εκτέλεσης ή φράσης. Το $y_{pwe}$ εκφράζει έτσι τοπική σχετική χρονική απόκλιση αφού ληφθούν υπόψη η σημειογραφική αξία και η συνολική αγωγή. Το $q_{we}$ είναι συμβολική σχετική αξία και δεν αποτελεί πρόβλεψη διάρκειας σε δευτερόλεπτα.

\[
y_{pwe}=\log\!\frac{d_{pwe}}{q_{we}}-\alpha_{pw},
\]
όπου $d_{pwe}$ είναι η πραγματοποιημένη διάρκεια, $q_{we}$ η σημειογραφική χρονική αξία στην κοινή μονάδα ανάγνωσης και $\alpha_{pw}$ ο συνολικός όρος αγωγής της ηχογράφησης, εκτιμημένος χωρίς χρήση ορίων θέσεων. Το τονικό ύψος παραμένει προκαθορισμένη δευτερεύουσα μεταβλητή αξιολόγησης. Ως ανάλυση ευαισθησίας, ο κανονικοποιημένος χρονισμός επανεκτιμάται με τοπικό όρο αγωγής σε επίπεδο φράσης, υπολογισμένο χωρίς χρήση της εστιαζόμενης εσωτερικής περιοχής όπου αυτό είναι εφικτό. Η σύγκριση με τον κύριο όρο $\alpha_{pw}$ ελέγχει αν μέρος της διασποράς που αποδίδεται σε μια θέση εξηγείται από τοπική επιτάχυνση ή επιβράδυνση της φράσης· δεν αντικαθιστά την κύρια προκαθορισμένη μεταβλητή αξιολόγησης. Χωριστή προκαθορισμένη ανάλυση ευαισθησίας μπορεί να χρησιμοποιεί όρο αγωγής ανά τύπο τμήματος· αν αυτή μεταβάλλει την κατεύθυνση των συγκριτικών διαφορών ή την κατάταξη των $S_m$, η κύρια και η εναλλακτική εκτίμηση αναφέρονται ισότιμα. Αν ο πιλοτικός έλεγχος δείξει ανώτερη αξιοπιστία άλλης ήδη προβλεπόμενης μεταβλητής, η αλλαγή οριστικοποιείται στο ΠΕ0 πριν ενεργοποιηθεί ο επιβεβαιωτικός εκτελεστικός έλεγχος.

Ο προκαθορισμένος έλεγχος συγχυτικών παραγόντων στηρίζει την πλήρη ερμηνεία της μεταφοράς. Το ίδιο οριστικοποιημένο ιστορικό προφίλ εισέρχεται, αποκλειστικά για εξωτερική στατιστική αξιολόγηση, σε κοινό ιεραρχικό μοντέλο κεντρικής τάσης και κλίμακας (location--scale) που ελέγχει εσωτερική θέση, εγγύτητα σε κατάληξη, ήχο, μελισματική πυκνότητα και πυκνότητα μελωδικών γεγονότων, τοπική πολυπλοκότητα, θέση στη φράση και την άμεση συμβολική περιγραφή B0. Οι τυχαίες επιδράσεις και οι ομαδοποιήσεις αντανακλούν έργο, οικογένεια, ιστορική γραμμή και εκτελεστή.

Η αναδειγματοληψία διατηρεί ολόκληρα προφίλ $u$ και τις ομαδοποιήσεις έργων, πηγών και εκτελεστών, χρησιμοποιώντας τις ίδιες μονάδες αναδειγματοληψίας σε όλες τις συνθήκες. Ενσωματώνει τις κοινές ρυθμίσεις τυχαιότητας και τις επιλέξιμες ζευγοποιήσεις του B2. Η ακριβής διαδικασία πολυεπίπεδης αναδειγματοληψίας οριστικοποιείται στο ΠΕ0 με βάση την απογεγραμμένη δομή εξαρτήσεων. Για τις δύο αντιθέσεις αναφέρονται ταυτόχρονα διαστήματα εμπιστοσύνης 95\% από μέγιστο στατιστικό και διορθωμένες τιμές $p$ με Holm. Ο έλεγχος σφάλματος αφορά χωριστή οικογένεια στατιστικών ελέγχων για τις δύο κύριες αντιθέσεις EISE. Η σύγκριση B2--B1 παραμένει διαγνωστική και το κριτήριο επιτυχίας του κύριου πειράματος καθορίζεται από τα $\Delta_{\mathrm{hist}}$ και $\Delta_{\mathrm{learn}}$.

Η πλήρης διαχρονική--συγχρονική ερμηνεία απαιτεί ταυτόχρονα: θετική συμφωνία με το ανεξάρτητο εκτελεστικό προφίλ, θετικά $\Delta_{\mathrm{hist}}$ και $\Delta_{\mathrm{learn}}$, και διατήρηση του αποτελέσματος στον προκαθορισμένο έλεγχο με προσαρμογή για συγχυτικούς παράγοντες. Οι $\rho_F^{\mathrm{an}}/\rho_F^{\mathrm{tr}}$ παραμένουν προκαθορισμένες εξειδικεύσεις όπου η απογραφή των πηγών παρέχει επαρκή ανεξάρτητη κάλυψη. Αναφέρονται με μεγέθη επίδρασης και διαστήματα· εφόσον χρησιμοποιηθούν τυπικοί στατιστικοί έλεγχοι, όλες οι προκαθορισμένες συγκρίσεις κατά τύπο σχέσης συγκροτούν χωριστή δευτερεύουσα οικογένεια με FDR. Οι συσχετίσεις που υπολογίζονται χωριστά ανά οικογένεια περιγράφουν τα προφίλ, ενώ ο κύριος ισχυρισμός συνάγεται από τις συγκριτικές αντιθέσεις. Ο επιβεβαιωτικός χαρακτήρας ενεργοποιείται στις καταστάσεις GREEN και AMBER της \S\ref{sec:pilot-decision}· στη RED η αντίστοιχη σύγκριση παραμένει διερευνητική.

Η επιτυχής πλήρης μεταφορά τεκμηριώνει ότι \textbf{η εποπτεία μέσω τεκμηριωμένων σχέσεων ιστορικής συμβολικής μετάδοσης παράγει προφίλ που συμφωνεί περισσότερο με την οργάνωση των ανεξάρτητων σύγχρονων εκτελέσεων, όπως αυτή αποτυπώνεται στο προκαθορισμένο ανεξάρτητο κριτήριο αξιολόγησης}. Το εύρημα τεκμηριώνει δομική αντιστοιχία των δύο προφίλ· οποιαδήποτε ευρύτερη ιστορική ερμηνεία απαιτεί συνεκτίμηση της τεκμηρίωσης από τις ιστορικές πηγές και δεν συνάγεται μόνο από το $\Delta_{\mathrm{hist}}$. Ο εμπειρικός ισχυρισμός αφορά τη συγκριτική προβλεπτική σχέση· ο αισθησιοκινητικός μηχανισμός μέσω του οποίου οι ψάλτες αποκτούν και πραγματώνουν αυτή τη γνώση αποτελεί διακριτό ερευνητικό αντικείμενο.



\subsection{M3b: τριαδικό μοντέλο επιλογής, CCGP και ενδοοικογενειακή ευθυγράμμιση}\label{app:stats-m3b}

Η υπόθεση $H_{\mathrm{cat}}$ ελέγχεται με τις μεικτές τριάδες και μοντέλο επιλογής συμβατό με τη δοκιμασία επιλογής του διαφορετικού ερεθίσματος. Η πιθανότητα επιλογής κάθε ερεθίσματος παραμετροποιείται από οικογενειακές σχέσεις, συμβολικές και ακουστικές αποστάσεις και προκαθορισμένες συμμεταβλητές ερεθισμάτων. Η εξάρτηση λόγω επαναλαμβανόμενων ακροατών και ερεθισμάτων αντιμετωπίζεται με κατάλληλη διασταυρούμενη ή πολλαπλής συμμετοχής (multi-membership) δομή. Η ακριβής συνάρτηση σύνδεσης (link function) επιλέγεται στον πιλοτικό έλεγχο με κριτήριο την ποιότητα προσαρμογής και τη σαφήνεια της ερμηνείας και οριστικοποιείται πριν από την επιβεβαιωτική συλλογή.

Η οικογενειακή μεταβλητή κωδικοποιείται ώστε θετικό $\beta_F^{\mathrm{cat}}$ να δηλώνει αυξημένη πιθανότητα επιλογής της εκδοχής διαφορετικής οικογένειας ως διαφορετικής, με συνεκτίμηση των επιφανειακών αποστάσεων. Οι ενδοοικογενειακές τριάδες εκτιμούν τη γεωμετρία παραλλαγής και την κύρια δευτερεύουσα ενδοοικογενειακή ευθυγράμμιση. Οι προκαθορισμένες συμβολικές και ακουστικές RDM αποτυπώνουν την επιφανειακή απόσταση. Τυχόν αναπαραστάσεις από προεκπαιδευμένα θεμελιώδη μοντέλα (foundation models) εντάσσονται ως εναλλακτικές υποθέσεις αναπαράστασης, με χωριστή σύγκριση προς τον ανθρώπινο χώρο.

Ο χώρος των ανθρώπινων αναπαραστάσεων εκτιμάται από τις τριαδικές επιλογές, με τη διάσταση και την κανονικοποίηση να επιλέγονται μέσω πρόβλεψης τριάδων που δεν χρησιμοποιήθηκαν στην προσαρμογή του μοντέλου. Οι ετικέτες οικογένειας εισάγονται στο στάδιο αξιολόγησης της κατηγορικής γενίκευσης. Το $\mathrm{CCGP}_F$ εκπαιδεύει γραμμικό αποκωδικοποιητή σε προκαθορισμένες συνθήκες και τον αξιολογεί σε χωριστές συνθήκες. Η ενδοοικογενειακή ανάλυση υπολογίζει RDM ανά οικογένεια και εκτιμά την υπολειμματική συμφωνία μεταξύ ανθρώπινου και υπολογιστικού χώρου με προεγγεγραμμένη μερική συσχέτιση ή συμμετρική αφαίρεση της συμβολικής και ακουστικής συνιστώσας. Η συμπερασματική συγκέντρωση γίνεται στο επίπεδο της οικογένειας.

\subsection{M3c: ιεραρχικό μοντέλο μουσικής προβλεψιμότητας και αντίθεση πλαισίου}\label{app:stats-m3c}

Η κύρια εξαρτημένη μεταβλητή είναι η αξιολόγηση του πόσο απρόσμενη φαίνεται η συνέχεια στον ακροατή. Το βασικό μοντέλο μικτών επιδράσεων έχει τη μορφή
\[
\mathrm{Surprise}_{it}=\beta_F^{\mathrm{pred}} z(I_F)_t+\beta_E^{\mathrm{pred}} z(I_{\mathrm{event}})_t+\mathbf X_{it}\gamma+u_i+v_t+\epsilon_{it},
\]
όπου $\mathbf X_{it}$ περιλαμβάνει προκαθορισμένα ακουστικά και επιφανειακά χαρακτηριστικά, διάρκεια/$K_t$, εξοικείωση και δείκτες συνθήκης. Η δομή τυχαίων επιδράσεων για συμμετέχοντες και ερεθίσματα οριστικοποιείται με βάση τη δομή των δεδομένων και προσομοίωση στατιστικής ισχύος πριν από την κύρια συλλογή.

Για την κύρια δευτερεύουσα ανάλυση της παροχής λατρευτικού συμφραζομένου, το μοντέλο περιλαμβάνει τον κύριο όρο $R_{it}$ και ταυτόχρονα τους δύο όρους αλληλεπίδρασης $R_{it}\times z(I_F)_t$ και $R_{it}\times z(I_{\mathrm{event}})_t$, μαζί με τους αντίστοιχους κύριους προβλεπτικούς όρους. Η κωδικοποίηση ορίζεται ως $R=0$ για την ουδέτερη και $R=1$ για την ένδειξη λατρευτικού πλαισίου. Η τυποποίηση των δύο πληροφοριακών δεικτών υπολογίζεται μία φορά στην κοινή οριστικοποιημένη τράπεζα ερεθισμάτων και εφαρμόζεται σταθερά σε όλες τις λίστες και συνθήκες.

Το κύριο δευτερεύον contrast είναι
\[
\Delta_{\mathrm{ctx}}=\beta_{R\times z(I_F)}-\beta_{R\times z(I_{\mathrm{event}})}.
\]
Η σύγκριση των συντελεστών κλίσης χρησιμοποιεί τυποποιημένες μεταβλητές πρόβλεψης. Το μέσο πληροφοριακό περιεχόμενο γεγονότος είναι η κύρια σύνοψη, ενώ το αθροιστικό και το μέγιστο πληροφοριακό περιεχόμενο αποτελούν αναλύσεις ευαισθησίας. Η αλληλεπίδραση με $E_{\mathrm{Byz}}$ χρησιμοποιεί συνεχή δείκτη εμπειρίας και διατηρεί χωριστά γενική μουσική εκπαίδευση και λειτουργική έκθεση όταν η διακύμανση των δύο μεταβλητών επιτρέπει την εκτίμησή τους.

\section{Αντιστοίχιση και αναδειγματοληψία}\label{app:matching}

\subsection{Αντιστοίχιση B2--EISE}\label{app:match-b2}

Η αντιστοίχιση επιδιώκει να κάνει τις B2 και EISE συνθήκες συγκρίσιμες ως προς παράγοντες που θα μπορούσαν να προσφέρουν εναλλακτική εξήγηση της διαφοράς τους. Για την κύρια σύγκριση χρησιμοποιούνται μόνο εμφανίσεις των οποίων η οικογενειακή ετικέτα προέρχεται από παγωμένη και επιλέξιμη υπογραφή. Η παραγωγή ζευγών ακολουθεί χρονικά τον έλεγχο αξιοπιστίας και δεν χρησιμοποιείται για να διορθώσει εκ των υστέρων την ίδια την υπογραφή. Ένας \textbf{συγχυτικός παράγοντας (confounder)} είναι χαρακτηριστικό που συνδέεται συστηματικά τόσο με τη συνθήκη όσο και με τον υπό εξέταση στόχο. \textbf{Κοινό πεδίο τιμών (common support)} υπάρχει στην περιοχή όπου οι δύο συνθήκες διαθέτουν πραγματικά συγκρίσιμες περιπτώσεις ως προς τα προκαθορισμένα χαρακτηριστικά. Η αντιστοίχιση περιορίζεται σε παράγοντες που πρέπει να εξισωθούν για τη συγκεκριμένη αιτιολόγηση της σύγκρισης. Στη χρονική διάσταση, η πλήρης εξίσωση του εσωτερικού χρονικού προφίλ χρησιμοποιείται μόνο ως ανάλυση ευαισθησίας, επειδή ένα τέτοιο κριτήριο θα μπορούσε να αφαιρέσει από την κύρια σύγκριση μέρος της ίδιας οργανωμένης μεταβολής που εξετάζει το πείραμα.

Οι τέσσερις συνθήκες B0/B1/B2/EISE ορίζονται στην §~\ref{sec:eise} και δεν επαναλαμβάνονται εδώ. Για την αντιστοίχιση B2--EISE, οι B1, B2 και EISE χρησιμοποιούν την ίδια οριστικοποιημένη μορφή εξόδου και το ίδιο πρωτόκολλο τμηματικής εποπτείας, κοινή εποπτεία οικογενειακής ένταξης μέσω προτύπων αναφοράς, τα ίδια αρνητικά παραδείγματα, αντιστοιχισμένες αρχικοποιήσεις και ρυθμίσεις τυχαιότητας, τους ίδιους διαχωρισμούς δεδομένων και αντιστοιχισμένη έκθεση στις εμφανίσεις, τον ίδιο βοηθητικό έλεγχο τοπικής πληροφορίας, καθώς και κοινό υπολογιστικό προϋπολογισμό εκπαίδευσης και ρύθμισης. Κάθε συνθήκη εκπαιδεύεται χωριστά. Η κύρια αντίθεση B2--EISE μεταβάλλει την προέλευση των θετικών σχεσιακών ζευγών σύμφωνα με τον προκαθορισμένο σχεδιασμό.

Οι δεξαμενές θετικών σχέσεων ορίζονται πριν από την εκπαίδευση. Το $E_{\mathrm{coex}}^{C}$ περιέχει τις επιβεβαιωμένες σχέσεις συνεξήγησης ή άλλης εξάρτησης μαρτύρων και το $E_{\mathrm{coex}}^{U}$ τις περιπτώσεις με κατάσταση \texttt{uncertain}. Και τα δύο σύνολα ορίζονται αποκλειστικά από στοιχεία προέλευσης και εξάρτησης των πηγών, ποτέ από μελωδική ομοιότητα. Η κύρια, συντηρητική προδιαγραφή είναι
\[
E_{\mathrm{EISE}}^+=E_{\mathrm{elig}}\cap E_{\mathrm{ex}}\cap E_{\mathrm{hom}}\cap E_{\mathrm{fam}},
\qquad
E_{\mathrm{B2,main}}^+=(E_{\mathrm{elig}}\cap E_{\mathrm{fam}})\setminus(E_{\mathrm{ex}}\cup E_{\mathrm{coex}}^{C}\cup E_{\mathrm{coex}}^{U}\cup E_{\mathrm{copy}}).
\]
Το $E_{\mathrm{fam}}$ για ένα EISE-positive ζεύγος πρέπει να έχει προκύψει χωρίς χρήση της συγκεκριμένης εξηγητικής ακμής ως κριτηρίου ένταξης. Όπου η ανεξάρτητη οικογενειακή κρίση διαφωνεί με την ιστορική σύνδεση, η περίπτωση διατηρείται ως μουσικολογικό εύρημα και δεν διορθώνεται η ετικέτα για να δημιουργηθεί θετικό ζεύγος.

Σε μία προκαθορισμένη χαλαρότερη ανάλυση ευαισθησίας χρησιμοποιείται
\[
E_{\mathrm{B2,relaxed}}^+=(E_{\mathrm{elig}}\cap E_{\mathrm{fam}})\setminus(E_{\mathrm{ex}}\cup E_{\mathrm{coex}}^{C}\cup E_{\mathrm{copy}}),
\]
ώστε να ελεγχθεί η επίδραση των αμφίβολων σχέσεων εξάρτησης. Η κύρια $\Delta_{\mathrm{hist}}$ χρησιμοποιεί το $E_{\mathrm{B2,main}}^+$· η εκτίμηση με $E_{\mathrm{B2,relaxed}}^+$ είναι μόνο sensitivity analysis. Το εύρος μεταξύ των δύο εκτιμήσεων αναφέρεται ως εύρος ευαισθησίας στην ταξινόμηση των εξαρτήσεων των πηγών, δεν αποτελεί διάστημα εμπιστοσύνης και δεν δημιουργεί δεύτερη επιβεβαιωτική αντίθεση.

Μετά τους αποκλεισμούς σχηματίζεται το κοινό πεδίο
\[
\begin{aligned}
\mathcal S_{\mathrm{common}}=\{s:\;& s\ \text{ πληροί τους προκαθορισμένους κανόνες κάλυψης}\\
&\text{και αντιστοίχισης και στις δύο συνθήκες}\}.
\end{aligned}
\]
Το ΠΕ0 αξιολογεί την κάλυψη του $\mathcal S_{\mathrm{common}}$ χωρίς πρόσβαση στις τιμές της $\Delta_{\mathrm{hist}}$. Η πύλη $G_{\mathrm{pair}}$ χαρακτηρίζεται GREEN όταν το προκαθορισμένο κύριο πεδίο υποστηρίζεται, AMBER όταν η αξιόπιστη αντιστοίχιση περιορίζεται σε προκαθορισμένο υποσύνολο οικογενειών ή στρωμάτων και RED όταν δεν υπάρχει επαρκές κοινό πεδίο για αξιόπιστη εκτίμηση της κύριας αντίθεσης. Οι κανόνες ένταξης δεν χαλαρώνονται μετά την εξέταση αποτελεσμάτων.

Η απογραφή αναφέρει επιπλέον το $C_{\mathrm{ex}}(F)$ και την κατανομή του \texttt{exegesis\_observability} ανά οικογένεια. Η κύρια ερμηνεία της $\Delta_{\mathrm{hist}}$ αφορά μόνο το προκαθορισμένο κοινό πεδίο οικογενειών και εμφανίσεων με επαρκές χρησιμοποιήσιμο εξηγητικό υλικό. Χαμηλή ή επιλεκτική ιστορική κάλυψη περιορίζει το πεδίο γενίκευσης. Δεν ερμηνεύεται ως ιδιότητα όλων των θέσεων του παπαδικού μέλους.


Η νέα ζευγοποίηση του B2 αντιστοιχίζεται με τον EISE ως προς \textbf{αριθμό ζευγών, έκθεση κάθε εμφάνισης, σταθμισμένο βαθμό γράφου, οικογένεια, συνδυασμό πηγών, επιφανειακή απόσταση και δυσκολία}. Για την κύρια αντιστοίχιση, η \textbf{επιφανειακή απόσταση ορίζεται μόνο στο στρώμα $N$}: σχετική μελωδική κίνηση και σημειογραφικά και ποιοτικά στοιχεία, χωρίς τα λεπτομερή χρονικά γνωρίσματα του $T$. Παράλληλα αναφέρεται ισορροπία σε αδρά χρονικά χαρακτηριστικά που δεν αποτελούν το εσωτερικό χρονικό προφίλ της θέσης: συνολική σημειογραφική ανάπτυξη, αριθμός γεγονότων, μελισματική πυκνότητα και, όπου εφαρμόζεται, συλλαβική διάρθρωση, μαζί με τις ήδη προβλεπόμενες τροπικές, λειτουργικές και πηγικές συνθήκες. Έτσι η κύρια σύγκριση δεν επιτρέπει σε προφανείς διαφορές κλίμακας να καθορίσουν τη δυσκολία, αλλά δεν αφαιρεί εκ των προτέρων τη λεπτομερή χρονική οργάνωση που μπορεί να αποτελεί μέρος της χρήσιμης πληροφορίας των εξηγητικών σχέσεων. Σε δευτερεύουσα προκαθορισμένη ανάλυση ευαισθησίας, και μόνο όπου υπάρχει επαρκές κοινό πεδίο τιμών, η αντιστοίχιση αυστηροποιείται με το χρονικό προφίλ που υπολογίζεται μετά την οριοθέτηση. Αν η $\Delta_{\mathrm{hist}}$ εξασθενεί σε αυτή την ανάλυση, το αποτέλεσμα ερμηνεύεται ως ένδειξη ότι μέρος της πρόσθετης πληροφορίας της εποπτείας που βασίζεται σε τεκμηριωμένες ιστορικές σχέσεις είναι χρονικό, όχι ως εκ των υστέρων αποτυχία του κύριου πρωτοκόλλου. Δημιουργούνται πολλαπλές επιλέξιμες ζευγοποιήσεις για το B2 και το $\Delta_{\mathrm{hist}}$ αναφέρεται ως προς την κατανομή αυτών των επαναλήψεων του συγκριτή. Η κύρια αντίθεση χρησιμοποιεί οικογένειες με επαρκείς επιλέξιμες εναλλακτικές ζευγοποιήσεις εντός της ίδιας οικογένειας. Οι περιπτώσεις χωρίς τέτοιες εναλλακτικές δεν χρησιμοποιούνται για αριθμητική εκτίμηση της ειδικής αντίθεσης και παραμένουν ορατές στην αναφορά κάλυψης.

Η αντιστοίχιση χρησιμοποιεί επιπλέον, όπου είναι διαθέσιμα και συγκρίσιμα, την ακριβή ταυτότητα κειμένου ή κοινού χωρίου, τον λειτουργικό τόπο, το λεκτικό καθεστώς και την προέλευση από κοινή σύνθεση. Ελέγχεται επίσης η ισορροπία ως προς την αβεβαιότητα των ορίων, την κατάσταση επικάλυψης και την επαναχρησιμοποίηση κοινών γεγονότων, χωρίς οι εναλλακτικές οριοθετήσεις της ίδιας εμφάνισης να αντιμετωπίζονται ως ανεξάρτητες εμφανίσεις. Η σύνθεση, ο μάρτυρας και η πηγή διατηρούν χωριστά αναγνωριστικά. Για κάθε πρόσθετο περιορισμό αναφέρεται πόσες πραγματικές εναλλακτικές B2 ζευγοποιήσεις παραμένουν και ποιο είναι το αντίστοιχο κοινό πεδίο σύγκρισης. Δεν χαλαρώνονται οι περιορισμοί μετά την εξέταση των αποτελεσμάτων· αν η αυστηρότερη αντιστοίχιση περιορίζει τον πληθυσμό αναφοράς, η μεταβολή καταγράφεται στο χάρτη εκτιμώμενων μεγεθών. Οι θεολογικές κατηγορίες της H.7 δεν χρησιμοποιούνται για την κατασκευή των ιστορικών ή οικογενειακών γράφων.

\paragraph{Δευτερεύων συγκριτής αυτόματης αντιστοίχισης βάσει μελωδικής ομοιότητας $B2_{\mathrm{sim}}$.}
Εφόσον η απογραφή του ΠΕ0 δείξει ότι η κύρια B2--EISE αντίθεση διατηρεί επαρκή κάλυψη, κατασκευάζεται μία πρόσθετη δευτερεύουσα συνθήκη χωρίς χρήση ιστορικών ακμών. Πρώτα εκπαιδεύεται και οριστικοποιείται ο B1 μόνο στα δεδομένα εκπαίδευσης. Η οριστικοποιημένη αναπαράσταση B1 χρησιμοποιείται έπειτα για να υπολογιστεί ενδοοικογενειακή μελωδική ομοιότητα και να επιλεγούν, με τους ίδιους περιορισμούς αριθμού ζευγών, έκθεσης εμφανίσεων, πηγής, επιφανειακής απόστασης, μουσικής δυσκολίας και προϋπολογισμού εκπαιδευτικών ζευγών, τα διαθέσιμα μη ιστορικά ζεύγη με τη μεγαλύτερη μελωδική ομοιότητα. Ο γράφος αυτός οριστικοποιείται πριν από την εκπαίδευση της συγκρινόμενης συνθήκης και δεν χρησιμοποιεί $E_{\mathrm{ex}}$, $C_{\mathrm{exp}}$, εκτελεστικά δεδομένα ή αποτελέσματα του συνόλου ελέγχου. Η σύγκριση EISE--$B2_{\mathrm{sim}}$ απαντά στο αυστηρότερο ερώτημα αν η ιστορική σχέση περιέχει πρόσθετη πληροφορία σε σύγκριση με τη βέλτιστη αυτόματα επαγόμενη μελωδική ομοιότητα. Παραμένει δευτερεύουσα και δεν μεταβάλλει την κύρια επιβεβαιωτική οικογένεια.

\paragraph{Δευτερεύων συγκριτής ενδοοικογενειακής εποπτείας με αντιστοιχισμένο περιβάλλον.}
Μόνο αν η απογραφή δείξει ότι το κύριο B2 δεν εξισώνει ήδη επαρκώς το κειμενικό και λειτουργικό περιβάλλον και ότι υπάρχουν αρκετές πραγματικές εναλλακτικές ακμές χωρίς αύξηση του σώματος, μπορεί να προστεθεί δευτερεύων συγκριτής $B2_{\mathrm{context}}$. Ο συγκριτής διατηρεί την ενδοοικογενειακή εποπτεία και αντιστοιχίζει στενότερα το περιβάλλον χωρίς να χρησιμοποιεί τεκμηριωμένες εξηγητικές ακμές. Παρόμοια επίδοση με τον EISE περιορίζει την τεκμηρίωση πρόσθετης αξίας της ιστορικής οργάνωσης έναντι αυτού του αυστηρότερου συγκριτή· δεν αποτελεί από μόνη της απόδειξη ότι η διαφορά εξηγείται από ένα συγκεκριμένο γνώρισμα, όπως το κείμενο.

\paragraph{Υποχρεωτική ευαισθησία ως προς το κειμενικό στρώμα όταν $L\subset R^\star$.}
Αν η επιλεγμένη αναπαράσταση $R^\star$ περιλαμβάνει το στρώμα $L$, επαναλαμβάνεται η σχεσιακή σύγκριση με την ίδια επιλεγμένη κλίμακα $U$, τα ίδια ζεύγη και την ίδια αρχιτεκτονική, αλλά χωρίς το $L$ στην οδό οικογενειακής ταυτότητας. Η συνθήκη συμβολίζεται $R^\star_{\setminus L}$ και αποτελεί προκαθορισμένη ανάλυση ευαισθησίας. Δεν επανεκτελείται η επιλογή M1 και δεν δημιουργείται νέα υποψήφια αναπαράσταση. Η ανάλυση ελέγχει αν η πρόσθετη επίδοση μπορεί να αποδοθεί σε shortcut από το κειμενικό locus. Ο $B2_{\mathrm{context}}$ παραμένει χωριστός και ενεργοποιείται μόνο όταν υπάρχουν αρκετές πραγματικές εναλλακτικές ακμές για στενότερη αντιστοίχιση περιβάλλοντος.

Ο B0 χρησιμοποιεί την ίδια μουσικά τεκμηριωμένη αντιστοίχιση με τον EISE και αποτυπώνει την πλήρη άμεση συμβολική περιγραφή της ιστορικής παραλλαγής. Ένα Profile-HMM ή συγγενές δομημένο συμβολικό μοντέλο χρησιμοποιείται μόνο σε συμπληρωματική ανάλυση αφαίρεσης όταν το ΠΕ0 δείξει ουσιαστική αβεβαιότητα στην ευθυγράμμιση.

\subsection{Δύο οικογένειες του πληροφοριακού ελέγχου}\label{app:match-info-controls}

Η πρώτη οικογένεια αποτελείται από \textbf{μετατοπισμένα τμήματα ελέγχου} στο ίδιο τοπικό στρώμα, αντιστοιχισμένα ως προς μήκος, ήχο, λειτουργική περιοχή, σχετική θέση στο έργο, πυκνότητα γεγονότων, συχνότητα τοπικών υποακολουθιών και απόσταση από κατάληξη. Κάθε υποψήφιο τμήμα αποκλείεται αν επικαλύπτει τεκμηριωμένη θέση σε οποιαδήποτε διαθέσιμη εξωτερική αναφορά.

Η δεύτερη οικογένεια είναι ο αυστηρότερος \textbf{έλεγχος μηδενικής αναφοράς ευθυγραμμισμένος σε πληροφοριακή κορυφή}. Επιλέγονται τμήματα που αρχίζουν σε τοπικές κορυφές \(h_R\) ή \(\mathcal H_R\) αντίστοιχου μεγέθους με τις κορυφές που παρατηρούνται κοντά σε τεκμηριωμένα όρια, αλλά τα σημεία αυτά δεν αποτελούν δημοσιευμένα όρια θέσεων. Τα τμήματα αντιστοιχίζονται επιπλέον ως προς ήχο, μήκος, θέση στο έργο, πυκνότητα και εγγύτητα σε κατάληξη. Η σύγκριση απομονώνει την ειδική ενδο-τμηματική οργάνωση των πραγματικών θέσεων από τη γενική επίδραση μιας υψηλής πληροφοριακής κορυφής στο όριο.

Η Υ3 απαιτεί η χαρακτηριστική τροχιά των τεκμηριωμένων θέσεων να υπερβαίνει και τις δύο οικογένειες ελέγχου. Επιπλέον εκτελείται \textbf{προκαθορισμένη ανάλυση ευαισθησίας με εξαίρεση των πρώτων \(k\) γεγονότων}: τα πρώτα \(k\) ατομικά γεγονότα κάθε τμήματος αφαιρούνται και η ενδο-τμηματική τροχιά επανεκτιμάται. Η τιμή ή το μικρό σύνολο τιμών του \(k\) οριστικοποιείται στο ΠΕ0. Η επιμονή του αποτελέσματος μετά την αφαίρεση της άμεσης περιοχής του ορίου δείχνει ότι η διαφοροποίηση δεν περιορίζεται στην αρχική πληροφοριακή κορυφή και παραμένει μέσα στο τμήμα.

\subsection{Έλεγχος συχνότητας}\label{app:match-frequency}

Πρόκειται για τον σοβαρότερο συγχυτικό παράγοντα και αντιμετωπίζεται ρητά. Μια θέση που εμφανίζεται πολλές φορές στην εκπαίδευση θα προβλέπεται καλύτερα \emph{επειδή έχει ιδωθεί}, όχι επειδή αποτελεί ιδιαίτερα οργανωμένο μελωδικό τύπο.

Διακρίνονται \textbf{τρία επίπεδα έκθεσης}: συχνότητα μεμονωμένου γεγονότος, συχνότητα βραχέων τοπικών υποακολουθιών, και συχνότητα της οικογένειας θέσης. Οι δύο οικογένειες τμημάτων ελέγχου αντιστοιχίζονται ως προς τα δύο πρώτα, και τα \(\Delta_R^{(\mathcal H)}\) και \(\Delta_R^{(h)}\) αναφέρονται \textbf{στρωματοποιημένα ως προς το τρίτο}.

\textbf{Κρίσιμη λεπτομέρεια:} οι δύο πρώτες συχνότητες υπολογίζονται \emph{στο κοινό πλέγμα ατομικών γεγονότων} και όχι στις μονάδες κάθε αναπαράστασης. Διαφορετικά κάθε αναπαράσταση θα όριζε μόνη της τι σημαίνει «συχνό πρότυπο» και η αντιστοίχιση των τμημάτων ελέγχου θα έπαυε να είναι κοινή μεταξύ των συνθηκών, καθιστώντας τα \(\Delta_R\) μη συγκρίσιμα. Η συχνότητα οικογένειας είναι εξ ορισμού ανεξάρτητη από την αναπαράσταση. Η συχνότητα μονάδων ανά αναπαράσταση αναφέρεται μόνο ως διαγνωστικό, ποτέ ως μεταβλητή αντιστοίχισης. Έτσι φαίνεται αν η επίδραση υπάρχει μόνο στις πολύ συχνές οικογένειες ή αν παραμένει και στις μεσαίας και χαμηλής συχνότητας.

\subsection{Γεωμετρικός συγχυτικός παράγοντας}\label{app:match-geometry}

Η προβλεψιμότητα αυξάνεται γενικά καθώς μια φράση πλησιάζει σε κατάληξη. Για να μην αποδοθεί στη θέση ένα γενικό φαινόμενο, η κανονικοποιημένη θέση \(u\), η θέση μέσα στο έργο και η εγγύτητα σε κατάληξη εισάγονται ως συμμεταβλητές ελέγχου· το ενδιαφέρον εστιάζεται στον δείκτη τμήματος και στις αλληλεπιδράσεις του με το \(u\).



\subsection{Κοινή αναδειγματοληψία και κοινά προφίλ}\label{app:match-resampling}

Η αναδειγματοληψία για τις κύριες συγκρίσεις διατηρεί ολόκληρα προφίλ \(u\), τις ομαδοποιήσεις έργων, πηγών και εκτελεστών και χρησιμοποιεί τις ίδιες μονάδες αναδειγματοληψίας στις συγκρινόμενες συνθήκες. Για τη μεταφορά του EISE, η διαδικασία ενσωματώνει τις κοινές ρυθμίσεις τυχαιότητας και τις επιλέξιμες ζευγοποιήσεις του \(B2\), όπως ορίζεται στην §~\ref{app:stats-eise-transfer}. Για τις σχέσεις της πολυμεθοδικής σύγκλισης οι έλεγχοι μεταθέσεων και τα σχήματα αναδειγματοληψίας bootstrap εφαρμόζονται στις ίδιες προκαθορισμένες μονάδες και δεν αντιμετωπίζουν ως ανεξάρτητα διαφορετικά σημεία του ίδιου προφίλ. Η ακριβής διαδικασία οριστικοποιείται στο ΠΕ0 από την απογεγραμμένη δομή εξαρτήσεων.

\subsection{M3b: επιλογή ερεθισμάτων σε κοινή περιοχή αποστάσεων}\label{app:match-m3b}

Πριν από τη συλλογή δεδομένων από ακροατές υπολογίζονται οι προκαθορισμένες συμβολικές και βασικές ακουστικές αποστάσεις, χωρίς χρήση των οικογενειακών ετικετών. Τα επιβεβαιωτικά ερεθίσματα επιλέγονται από κοινή περιοχή συμβολικών και ακουστικών αποστάσεων, στην οποία εκπροσωπούνται τόσο σχέσεις μέσα στην ίδια οικογένεια όσο και σχέσεις μεταξύ διαφορετικών οικογενειών. Η επιλογή ακολουθεί τα προκαθορισμένα κριτήρια αντιστοίχισης και, όπου απαιτείται, τις προκαθορισμένες ομάδες αντιστοίχισης. Η περιοχή, τα ελάχιστα πλήθη και οι κανόνες αποκλεισμού οριστικοποιούνται πριν από την κύρια συλλογή. Τα δύσκολα θετικά και δύσκολα αρνητικά παραδείγματα αναφέρονται χωριστά ως ερμηνευτικός έλεγχος.

\subsection{M3c: αντιδιαγώνια δειγματοληψία και μουσική φυσικότητα}\label{app:match-m3c}

Η τράπεζα ερεθισμάτων κατασκευάζεται από το δισδιάστατο πεδίο $z(I_F)\times z(I_{\mathrm{event}})$. Η κύρια συλλογή απαιτεί επαρκή κάλυψη στους δύο αντιδιαγώνιους συνδυασμούς, όπου η πρόβλεψη σε επίπεδο οικογένειας και η πρόβλεψη σε επίπεδο γεγονότος δίνουν αντίθετη ένδειξη: υψηλό πληροφοριακό περιεχόμενο στο ένα επίπεδο και χαμηλό στο άλλο. Τα κατώφλια καθορίζονται από την εμπειρική κατανομή των υποψηφίων περιπτώσεων εκτός εκπαίδευσης και οριστικοποιούνται πριν από τον πιλοτικό έλεγχο. Κάθε κατασκευασμένη συνέχεια περνά τυφλό έλεγχο μουσικής φυσικότητας από ανεξάρτητους αξιολογητές και κάθε πραγματική συνέχεια συνοδεύεται από τεκμηριωμένη προέλευση προς το έργο και την πηγή της.

\subsection{Νικόλαος: αντιστοίχιση και κοινή μετρική}\label{app:match-nikolaos}

\paragraph{Περιβάλλον αναφοράς και κλείδωμα του μοντέλου.}
Για κάθε εφαρμογή ορίζεται πριν από την εξέταση του έργου-στόχου ένα $\Omega=(D_{\Omega},\rho,\Pi,\mathcal K_{\Omega})$. Το $D_{\Omega}$ περιέχει πλήρη έργα με ελεγμένη προέλευση και χάρτη εμφανίσεων από τον οποίο προκύπτουν οι τεκμηριωμένες συνεχείς διαδοχές. Κάθε έργο συνοδεύεται από υποχρεωτική αναφορά κάλυψης. Μη αναγνωρισμένες, συνδετικές ή ανεπίλυτες περιοχές παραμένουν ρητές διακοπές και δεν γεφυρώνονται κατά τον υπολογισμό μεταβάσεων. Μετά από διακοπή το ιστορικό οικογενειακού συμφραζομένου επανεκκινεί, ενώ η εμφάνιση που ακολουθεί διακρίνεται στα μεταδεδομένα από την πραγματική αρχή του έργου. Η ίδια προκαθορισμένη πολιτική συμμετοχής της στο $I_{\Omega}$ εφαρμόζεται στο έργο-στόχο και στο σώμα αναφοράς. Η $\rho$ περιλαμβάνει το οριστικοποιημένο $R^\star$, το επικυρωμένο λεξιλόγιο οικογενειών και τους κανόνες αποχής. Το $\Pi$ είναι το PPM της οικογενειακής ακολουθίας με προκαθορισμένο κανόνα επιστροφής σε συντομότερο συμφραζόμενο, κοινό με τη M3c ως προς την αρχιτεκτονική. Το $\mathcal K_{\Omega}$ ορίζει το χρονικό όριο ένταξης, τον αποκλεισμό έργων του εξεταζόμενου μελοποιού και των δικών του εξηγήσεων, καθώς και τους κανόνες για εξαρτημένους μάρτυρες. Κάθε μεταβολή του $D_{\Omega}$ αποτελεί χωριστή ανάλυση ευαισθησίας και δεν επαναρυθμίζει τη $\rho$ πάνω στα αποτελέσματα του Νικολάου.

\paragraph{Τέσσερις διαστάσεις χαρακτηρισμού και κατανομή αναφοράς μελοποιών.}
Οι τέσσερις διαστάσεις $D_{\mathrm{fam}}$, $D_{\mathrm{seq}}$, $D_{\mathrm{func}}$ και $D_{\mathrm{within}}$ ορίζονται στην §~\ref{sec:nikolaos}. Εδώ εξειδικεύονται μόνο η αβεβαιότητα και η κατανομή αναφοράς σε επίπεδο μελοποιού. Για σύνολα έργων η κύρια μονάδα αναδειγματοληψίας είναι το έργο. Για μεμονωμένο έργο-στόχο αναδειγματοληπτούνται τα έργα ή οι μελοποιοί του σώματος αναφοράς και μεταφέρεται χωριστά η αβεβαιότητα οικογενειακής απόδοσης και ορίων. Κάθε επιλέξιμο έργο ή κάθε επιλέξιμος μελοποιός αντιμετωπίζεται διαδοχικά ως στόχος και συγκρίνεται με το υπόλοιπο σώμα με τους ίδιους κανόνες. Εμπειρική θέση ή εκατοστημόριο αναφέρεται μόνο αν $K\ge K_{\min}$ ανεξάρτητοι επιλέξιμοι μελοποιοί παραμένουν μετά τους αποκλεισμούς. Το $K_{\min}$ καθορίζεται στο ΠΕ0 με προσομοίωση πριν από την M5. Με $K<K_{\min}$ αναφέρονται οι τιμές των μεγεθών, bootstrap διαστήματα σε επίπεδο έργου και η πραγματική κάλυψη, χωρίς κατάταξη του Νικολάου σε εμπειρική κατανομή μελοποιών. Στη κατανομή αναφοράς ο εκάστοτε στόχος αφαιρείται από το $D_{\Omega}$. Αν το επιτρέπουν τα δεδομένα και ο υπολογιστικός προϋπολογισμός, η διαδικασία επαναλαμβάνεται ως αυστηρότερη ανάλυση ευαισθησίας με επανεκπαίδευση της αναπαράστασης μετά τον αποκλεισμό του μελοποιού-στόχου (composer-held-out).

\paragraph{Χρονικό διαγνωστικό ενδοοικογενειακής πραγμάτωσης.}
Το $D_{\mathrm{within}}$ παραμένει το ενιαίο κύριο μέγεθος ενδοοικογενειακής πραγμάτωσης. Για ερμηνεία της χρονικής συνιστώσας του παρατηρούμενου φαινομένου, και μόνο σε οικογένειες με αξιόπιστες κοινές εσωτερικές περιοχές, υπολογίζεται χωριστό περιγραφικό διαγνωστικό. Οι περιοχές $z_1,\ldots,z_J$ ορίζονται χωρίς χρονική πληροφορία, κατά προτεραιότητα από τις ίδιες ζώνες των $A(z)$/$\Gamma(z)$ που έχουν οριστεί εκ των προτέρων· αν απαιτείται λεπτότερη ευθυγράμμιση, χρησιμοποιείται αποκλειστικά το στρώμα $N$. Το $J$ παραμένει ίδιο για όλες τις συγκρινόμενες πραγματώσεις της ίδιας οικογένειας. Αβεβαιότητα ορίων ή αντιστοίχισης μεταφέρεται ως $U_{\mathrm{seg}}$ και περιοχές χωρίς αξιόπιστη αντιστοίχιση εξαιρούνται με ρητή αναφορά κάλυψης.

Για πραγμάτωση $x$, έστω $q_{x,r}$ το άθροισμα των σημειογραφικών χρονικών αξιών στην κοινή περιοχή $r$. Ορίζονται
\[
Q_x=\sum_{r=1}^{J}q_{x,r},
\qquad
p_{x,r}=\frac{q_{x,r}}{Q_x},
\]
όπου το $Q_x$ περιγράφει τη συνολική γραπτή χρονική ανάπτυξη και το $p_x=(p_{x,1},\ldots,p_{x,J})$ την εσωτερική κατανομή του χρόνου. Η συνολική ανάπτυξη αναλύεται επιπλέον ως
\[
Q_x=n_x\,\bar v_x,
\]
με $n_x$ τον αριθμό ατομικών μελωδικών γεγονότων της πραγμάτωσης και $\bar v_x$ τη μέση σημειογραφική χρονική αξία ανά γεγονός. Η διάσπαση αυτή διακρίνει αύξηση της ανάπτυξης μέσω περισσότερων γεγονότων από αύξηση μέσω μεγαλύτερων χρονικών αξιών.

Η σύγκριση του $Q_x$ γίνεται μόνο μέσα σε προκαθορισμένο συγκρίσιμο κελί $c$ που απαιτεί το ίδιο είδος/τύπο μέλους και το ίδιο στρώμα εξήγησης, μαζί με τις λοιπές αναγκαίες συνθήκες της κύριας M5. Η αναφορά είναι
\[
\delta_Q(x;\Omega,c)=\ln\!\frac{Q_x}{\widetilde Q_{\Omega,c}},
\]
όπου $\widetilde Q_{\Omega,c}$ είναι η διάμεσος της συνολικής ανάπτυξης στο επιλέξιμο σώμα αναφοράς, με ίσο βάρος ανά μελοποιό.

Χρησιμοποιείται ο φυσικός λογάριθμος για πλήρη αναπαραγωγιμότητα του ορισμού· θετική τιμή δηλώνει μεγαλύτερη συνολική σημειογραφική χρονική ανάπτυξη από την τιμή αναφοράς, μηδενική τιμή ισότητα και αρνητική τιμή μικρότερη ανάπτυξη. Πρώτα συνοψίζονται οι επιλέξιμες πραγματώσεις μέσα σε κάθε μελοποιό και κατόπιν οι μελοποιοί λαμβάνουν ίσο βάρος. Αν το ίδιο είδος μέλους και στρώμα εξήγησης δεν παρέχουν επαρκή κάλυψη, το $Q_x$ αναφέρεται μόνο περιγραφικά, χωρίς σχετική κατάταξη ή ισχυρισμό αυξημένης/μειωμένης ανάπτυξης.

Το προφίλ αναφοράς $\bar p_{\Omega,r}$ υπολογίζεται με την ίδια αρχή ίσου βάρους ανά ανεξάρτητο μελοποιό. Ορίζεται
\[
d_{\mathrm{time}}(x;\Omega)=
\frac12\sum_{r=1}^{J}\left|p_{x,r}-\bar p_{\Omega,r}\right|.
\]
Το $d_{\mathrm{time}}$ είναι απόσταση ολικής μεταβολής μεταξύ σχετικών χρονικών προφίλ και αποτελεί βοηθητικό ερμηνεύσιμο διαγνωστικό, όχι ιστορικό ορισμό ρυθμού ούτε νέα κύρια διάσταση M5. Επειδή εξαρτάται από τον αριθμό περιοχών, ερμηνεύεται μέσα στην ίδια οικογένεια με σταθερό $J$· δεν συγκρίνεται ακατέργαστο μεταξύ οικογενειών διαφορετικής κατάτμησης. Ο ίδιος κανόνας $K_{\min}$ ισχύει για κάθε κατάταξη του $d_{\mathrm{time}}$ σε επίπεδο μελοποιού: με ανεπαρκείς ανεξάρτητους μελοποιούς αναφέρονται η τιμή, η κάλυψη και διαστήματα αβεβαιότητας από το σώμα αναφοράς, χωρίς εκατοστημόριο ή πληθυσμιακή κατάταξη. Το χρονικό διαγνωστικό συνοδεύει το $D_{\mathrm{within}}$ και δεν προστίθεται αριθμητικά σε αυτό.

\paragraph{Ιστορικά στρώματα και κριτήριο ενεργοποίησης.}
Το κύριο $D_{\Omega}$ συγκροτείται από συγκρίσιμη Νέα Μέθοδο. Το $C_{\mathrm{pre}}$ περιλαμβάνει προγενέστερους μελοποιούς μόνο μέσω τεκμηριωμένων εξηγήσεων που διατηρούνται ως χωριστό στρώμα, ενώ το $C_{\mathrm{contemp}}\setminus \mathrm{Nik}$ περιλαμβάνει σύγχρονους και ανεξάρτητους μελοποιούς της επιλέξιμης περιόδου. Πρόσθετο $C_{\mathrm{pre}}$ χρησιμοποιείται όταν υπάρχει επαρκής κοινή κάλυψη και οι εξηγήσεις του ίδιου του Νικολάου έχουν αποκλειστεί. Επειδή ορισμένοι εξηγητές του $C_{\mathrm{pre}}$ είναι ταυτόχρονα μελοποιοί του $C_{\mathrm{contemp}}$, η σύγκριση επαναλαμβάνεται με αφαίρεση των κοινών εξηγητών και, όπου είναι εφικτό, με leave-one-exegete-out. Η διαφορά $I_{\mathrm{pre}}-I_{\mathrm{contemp}\setminus \mathrm{Nik}}$ περιγράφεται ως μετατόπιση προβλεψιμότητας της ίδιας μουσικής επιλογής μεταξύ δύο σωμάτων αναφοράς. Ένα $C_{\mathrm{post}}$ ενεργοποιείται μόνο διερευνητικά. Η ανάλυση διατηρεί ίσο ή εξισορροπημένο αριθμό οικογενειακών tokens όπου συγκρίνονται corpora διαφορετικού μεγέθους και επαναλαμβάνει τον υπολογισμό σε εναλλακτικές επιλέξιμες πηγές για έλεγχο ευαισθησίας.

\paragraph{Πληθυσμός αναφοράς και επάρκεια δεδομένων.}
Για κάθε οικογένεια και ήχο δηλώνονται χωριστά συνθέτης, έργο, χρονολόγηση σύνθεσης και μάρτυρα, ιστορικό στρώμα, λειτουργικός ρόλος και αριθμοί εμφανίσεων, έργων και ανεξάρτητων ομάδων συγγενών πηγών. Η χρονολόγηση ορίζει αν ένα στρώμα είναι προγενέστερο ή απλώς συγκριτικό. Οι 15--25 συνθέσεις Νικολάου και οι 30--40 του ορίζοντα αναφοράς παραμένουν στόχοι απογραφής, όχι επιβεβαιωμένο δείγμα. Το ΠΕ0 προσδιορίζει ποιες τομές οικογένειας--ήχου--στρώματος διαθέτουν επαρκή κάλυψη για αξιόπιστη εκτίμηση, σύμφωνα με τις υπάρχουσες προσομοιώσεις. Αραιά στρώματα αναφέρονται περιγραφικά, χωρίς υποχρεωτική σύγκριση και χωρίς συγχώνευση συνθετών σε τεχνητό μέσο πρότυπο. Στρώματα ή πηγές με γνωστή επιλεκτική παράλειψη μεγάλων φόρμουλων δεν χρησιμοποιούνται για βαθμονόμηση «τεκμηριωμένου εύρους» πριν ελεγχθεί η αντίστοιχη κάλυψη στα πλήρη έργα· διαφορετικά το συμπέρασμα περιορίζεται ρητά στο καταλογογραφικά παρατηρούμενο εύρος. Οι εκτιμήσεις αφορούν τα εξετασμένα στρώματα και δεν γενικεύονται σε πληθυσμό συνθετών.

Οι συγκρίσεις αντιστοιχίζονται ως προς λειτουργία, μήκος, συχνότητα οικογένειας και τεκμηρίωση πηγών. Κοινοί μάρτυρες, ανατυπώσεις και εξαρτημένες αποδόσεις δεν αυξάνουν το ανεξάρτητο δείγμα. Η κοινή αναδειγματοληψία της §~\ref{app:match-resampling} διατηρεί έργα, ομάδες συγγενών πηγών και ολόκληρα εσωτερικά προφίλ· οι ίδιες οικογένειες που συμμετέχουν σε περισσότερες συγκρίσεις παραμένουν συνδεδεμένες. Η άνιση έκθεση και το πλήθος εμφανίσεων αντιμετωπίζονται με την προκαθορισμένη αντιστοίχιση και κοινή υποδειγματοληψία, ώστε διαφορές διασποράς να μη συγχέονται με τη διαφορετική πεπερασμένη κάλυψη. Αναφέρονται μέγεθος επίδρασης, περιγραφικό διάστημα και κάλυψη, χωρίς χωριστή σειρά ελέγχων σημαντικότητας για κάθε παράγωγο μέγεθος.

\paragraph{Αναλυτική επανεγγραφή και μουσικολογικός χαρακτηρισμός.}
Η πιστότητα εισαγωγής και η κανονική απόδοση της §~\ref{app:repr} ελέγχονται πριν από την ταξινόμηση διαφορών. Η κοινή μουσική πραγμάτωση τεκμηριώνεται από πηγές και μουσική αντιπαραβολή· η μικρή απόσταση σε ανηγμένα γνωρίσματα αποτελεί ένδειξη προς έλεγχο, όχι αυτόματο κριτήριο ταύτισης. Ο χαρακτηρισμός σημειογραφικής επανεγγραφής παραμένει χωριστός από την ιστορική κατηγορία ακμής $E_{\mathrm{an}}/E_{\mathrm{tr}}$. Μικτές ή ανεπίλυτες περιπτώσεις διατηρούν την αβεβαιότητά τους και αναφέρονται χωριστά σε έλεγχο ευαισθησίας. Η σταθερότητα της ερμηνείας ελέγχεται με τα ήδη διαθέσιμα άμεσα/ανηγμένα μελωδικά γνωρίσματα του B0, χωρίς επιλογή νέων γνωρισμάτων από τα αποτελέσματα του Νικολάου.

Η κωδικοποίηση είναι ιεραρχική και επιτρέπει συμβατές πολλαπλές περιγραφές: πρώτα σημειογραφική, μουσική, μικτή ή ανεπίλυτη διαφορά· στη συνέχεια σύντμηση, ανάπτυξη, προσαρμογή και ανασυνδυασμός. Η διατήρηση δηλώνει απουσία εντοπισμένου μουσικού μετασχηματισμού στο συγκεκριμένο επίπεδο. Η διαποίκιλση σημειώνεται ως εμπλουτισμός διατηρούμενου κορμού και δεν συνεπάγεται αυτομάτως δομική ανάπτυξη. Οι υποτύποι προσαρμογής περιορίζονται σε κατάληξη, τροπικό περιβάλλον και συλλαβική και ρυθμική κατανομή. Οι ορισμοί δοκιμάζονται μόνο σε υλικό ανάπτυξης και δεύτερος γνώστης ελέγχει προκαθορισμένο δείγμα χωρίς προβλέψεις μοντέλου. Ο τελικός μουσικολογικός χαρακτηρισμός δεν εκπαιδεύει πρόσθετο ταξινομητή και δεν επηρεάζει το $C_{\mathrm{exp}}$.

\paragraph{Μετατόπιση και διαφορά διασποράς.}
Στην κοινή μετρική της §~\ref{app:eise-metric}, για κάθε στρώμα $s$ το πρότυπο $\widetilde\pi_{F,s}(u)$ εκτιμάται αποκλειστικά από επιλέξιμο ιστορικό υλικό εκπαίδευσης. Στο κοινό αξιόπιστο πεδίο $u$ ορίζονται
\[
M_{\mathrm{Nik},F,s}(u)=\|\overline h_{\mathrm{Nik},F}(u)-\widetilde\pi_{F,s}(u)\|^2,
\qquad
\delta_{\mathrm{disp},\mathrm{Nik},F,s}(u)=d_h\,[V_{\mathrm{Nik},F}(u)-V_{F,s}(u)].
\]
Το $d_hV$ παραμένει συνολική διασπορά, ενώ το $\delta_{\mathrm{disp}}$ δηλώνει τη διαφορά της από το συγκεκριμένο στρώμα. Το δεύτερο δεν ταυτίζεται με το ήδη ορισμένο $\Delta V$ της Α\(\leftrightarrow\)Γ. Η εμπειρική αποσύνθεση διατηρείται ακριβώς:
\[
D_{\mathrm{Nik},F,s}(u)=d_hV_{\mathrm{Nik},F}(u)+M_{\mathrm{Nik},F,s}(u).
\]
Και οι δύο όροι εκτιμώνται με την υπάρχουσα σύμβαση δεύτερης ροπής και στο κοινό τμήμα όπου η ευθυγράμμιση είναι αξιόπιστη. Η διασπορά μετρά εξάπλωση γύρω από το κέντρο, όχι πλήρη όγκο ή σχήμα πολυδιάστατου χώρου πραγματώσεων. Τα κριτήρια πρακτικά μικρής μετατόπισης ή διαφοράς καθορίζονται σε ιστορικό υλικό ανάπτυξης· αν δεν βαθμονομούνται αξιόπιστα, αναφέρονται συνεχείς εκτιμήσεις αντί τεσσάρων διακριτών κατηγοριών. Αντίστοιχα, συγκρίσεις αρχικών, εσωτερικών και καταληκτικών περιοχών ή υποτύπων ανά ήχο/λειτουργία παραμένουν δευτερεύουσες περιγραφές, με το $u$ και τα όρια περιοχών οριστικοποιημένα χωρίς χρήση της περίπτωσης Ιωαννίδη.

\paragraph{Υπέρβαση τεκμηριωμένου εύρους και ανασυνδυασμός.}
Η απόσταση μιας πραγμάτωσης εξετάζεται ως προς τη βαθμονομημένη ιστορική κατανομή της ίδιας οικογένειας, ήχου και στρώματος, σε συγκρίσιμη κλίμακα και με δηλωμένες ανεξάρτητες μονάδες αναφοράς. Το όριο προκύπτει από τον προκαθορισμένο έλεγχο κάλυψης και σφάλματος σε ιστορικές γραμμές ανάπτυξης, χωρίς σταθερό εκ των προτέρων q90/q95 και χωρίς επιλογή από τις αποστάσεις του Νικολάου. Ο κανόνας, η αβεβαιότητα και η διαθέσιμη υποστήριξη οριστικοποιούνται πριν από την εφαρμογή. Χωρίς επαρκή βαθμονόμηση αναφέρονται απόσταση και σειρά κατάταξης, όχι χαρακτηρισμός υπέρβασης. Ο παρονομαστής κάθε ποσοστού είναι όλες οι επιλέξιμες εμφανίσεις της συγκεκριμένης τομής· η αποχή και η έλλειψη αναφοράς αναφέρονται χωριστά.

Για τον ανασυνδυασμό παρουσιάζεται ακολουθία πλησιέστερων συγγενειών από την υπάρχουσα ανάκτηση, με αβεβαιότητα για κάθε εμφάνιση. Η εναλλαγή πλησιέστερου στρώματος προκαλεί μουσική αντιπαραβολή και δεν ισοδυναμεί αυτομάτως με ανασυνδυασμό ή διαδοχή επιρροών. Διακρίνονται άτυπη πραγμάτωση, υπέρβαση του τεκμηριωμένου εύρους, νέος συνδυασμός ως προς το σώμα αναφοράς και πρώτη γνωστή εμφάνιση στην εξετασμένη χρονολογία. Η ιστορική καινοτομία θεμελιώνεται με θετικά τεκμήρια από τις πηγές και τη χρονολογική αντιπαραβολή των μουσικών σχέσεων. Η υπολογιστική απόκλιση, η μουσικολογική διαφορά και η ιστορική καινοτομία αναφέρονται ως διαφορετικά επίπεδα τεκμηρίωσης.

\paragraph{Εξωτερική περίπτωση Ιωαννίδη: ευρήματα και νέες προβλέψεις.}
Η πρωτογενής δημοσίευση είναι Ιωαννίδης (2022), \emph{Θεολογία} 93(4), 173--193. Ως δημοσιευμένα ευρήματα καταγράφονται ο εκτεταμένος παραλληλισμός με τον Πέτρο Λαμπαδάριο (σ.~182), η ιδιαίτερη σύγκλιση της περιόδου III (σ.~183), η τοπική συγγένεια με σύνθεση του Χουρμουζίου (σσ.~185--188) και η διαφοροποίηση του συγκεκριμένου Κοινωνικού από τη γενική εικόνα που προτείνει ο συγγραφέας για τον Νικόλαο (σ.~175). Η αρίθμηση IV στη σ.~186 διαφέρει από VI στη σ.~180 και από τα αναγνωριστικά 6β/6γ των πινάκων· απαιτείται ταυτοποίηση με μουσικό χωρίο και πηγή, όχι σιωπηρή επιλογή αριθμού.

Σε χωριστό μητρώο διατυπώνονται οι νέες ποσοτικές συνέπειες: κατάταξη των δημοσιευμένων συγγενών στην ανάκτηση, εντοπισμός των κοινών και διαφοροποιούμενων περιοχών και συμφωνία με τις δημοσιευμένες αντιστοιχίσεις στην κοινή μουσική κλίμακα. Η υπόθεση υψηλότερης οικογενειακής κάλυψης από τη διάμεσο των έργων του Νικολάου θα ήταν νέα διερευνητική πρόβλεψη και δεν εντάσσεται στον παρόντα κύριο έλεγχο. Δεν μεταφέρεται αυτούσια μία κατάταξη ολόκληρων περιόδων σε τοπικές αποστάσεις χωρίς προσδιορισμό του κοινού αντικειμένου μέτρησης.

Οι συνθέσεις, οι πηγές, οι ακριβείς αντιστοιχίσεις, οι επιλέξιμες περιπτώσεις σύγκρισης και το κριτήριο επιτυχίας ή ασυμφωνίας καταγράφονται πριν εξεταστούν οι υπολογιστικές έξοδοι. Το άρθρο παρέχει γνωστή ανεξάρτητη μουσικολογική ανάλυση για τον προκαθορισμένο έλεγχο ανάκτησης των συγκεκριμένων σχέσεων. Οι δημοσιευμένες κρίσεις φυλάσσονται σε χωριστό αρχείο αξιολόγησης· δεν τροφοδοτούν την οντολογία, τα γνωρίσματα, τα κατώφλια, την επιλογή αναπαράστασης ή την εκπαίδευση. Αν κάποιο αναλυτικό λεξιλόγιο είχε αντληθεί από το άρθρο, η επανάληψή του δεν προσμετράται ως επικύρωση. Αξιολογείται μόνο η ανεξάρτητη ανάκτηση συγκεκριμένων μουσικών σχέσεων, με πλήρη αναφορά ασυμφωνιών και μη καλυπτόμενων χωρίων. Η περίπτωση είναι περιγραφικός έλεγχος εφαρμογής της M5 και δεν προστίθεται στην επιβεβαιωτική πολλαπλότητα. Οι αποκλεισμοί δεδομένων ορίζονται στην §~\ref{app:splits-nodes}.


\section{Αναλυτική εκπαίδευση EISE}\label{app:eise}

\subsection{Βιβλιογραφική και αρχιτεκτονική θεμελίωση της επιλογής EISE}\label{app:eise-rationale}

Η αρχιτεκτονική του EISE επιλέγεται με βάση τις απαιτήσεις του ίδιου του μουσικολογικού αντικειμένου και όχι με κριτήριο τη γενική πολυπλοκότητα ενός μοντέλου. Η θέση είναι μονάδα μεταβλητής έκτασης, τα όριά της δεν είναι πάντοτε γνωστά, οι οικογενειακές ετικέτες είναι άνισα και μερικώς διαθέσιμες, και η έρευνα χρειάζεται ταυτόχρονα αναγνώριση κοινής ταυτότητας και διατήρηση της εσωτερικής παραλλαγής. Επιπλέον, η M2 απαιτεί η B2 και η EISE να μπορούν να χρησιμοποιούν την ίδια μαθησιακή υποδομή και να διαφέρουν κυρίως στην προέλευση των θετικών σχεσιακών ζευγών.

Η βιβλιογραφική σύγκριση καλύπτει πέντε συγγενείς οικογένειες μεθόδων. Τα γραμμικά CRF προσφέρουν ώριμη ακολουθιακή μοντελοποίηση και μπορούν να εκπαιδευθούν με ελλιπείς επισημειώσεις (Tsuboi κ.ά., 2008), αλλά η βασική τους μονάδα είναι το επιμέρους γεγονός και όχι το τμήμα. Τα ημι-Μαρκοβιανά CRF μεταφέρουν τη βαθμολόγηση στο επίπεδο ολόκληρου τμήματος και επιτρέπουν ακριβή δυναμικό προγραμματισμό με μεταβλητό μήκος (Sarawagi--Cohen, 2004). Τα νευρωνικά τμηματικά μοντέλα επεκτείνουν αυτή την ιδέα με μαθημένες αναπαραστάσεις τμημάτων και κοινή εκτίμηση ορίων και ετικετών (Kong κ.ά., 2016). Στη μουσική, η χρήση νευρωνικού ημι-Μαρκοβιανού CRF για άμεση πρόβλεψη ολοκληρωμένων μουσικών γεγονότων αντί ανεξάρτητων επιμέρους χρονικών πλαισίων παρέχει άμεσο συγκριτικό προηγούμενο για τη χρησιμότητα της τμηματικής βαθμολόγησης (Yan κ.ά., 2021). Η νεότερη παραλλαγή με προεπιλογή υποψήφιων τμημάτων (filtered semi-CRF) δείχνει ταυτόχρονα ότι η μεγαλύτερη εκφραστικότητα του τμηματικού χώρου δεν αποτελεί από μόνη της εγγύηση καλύτερης επίδοσης και ότι ο έλεγχος του χώρου υποψηφίων είναι ουσιώδης (Zaratiana κ.ά., 2023).

Η μερική επισημείωση εξετάζεται χωριστά από την επιλογή τμηματικής αρχιτεκτονικής. Οι Tsuboi κ.ά. (2008) δείχνουν ότι η εκπαίδευση CRF μπορεί να αθροίζει πάνω σε συμβατές λανθάνουσες ετικέτες όταν η επισημείωση είναι ελλιπής, ενώ οι Effland--Collins (2021) αναδεικνύουν το πρόσθετο πρόβλημα της συστηματικά ελλιπούς επισήμανσης. Οι εργασίες αυτές δεν μεταφέρονται ως μουσικά αποτελέσματα· θεμελιώνουν όμως την αρχή ότι το μη επισημειωμένο υλικό δεν πρέπει να εξισώνεται μηχανικά με αρνητική τάξη.

Στην Ανάκτηση Μουσικής Πληροφορίας, η ανακάλυψη τμηματικών προτύπων (Conklin--Anagnostopoulou, 2006) και η μάθηση μετρικής για την ανάκτηση συγγενών μελωδιών (Karsdorp--van Kranenburg--Manjavacas, 2019) προσφέρουν δύο συμπληρωματικές δυνατότητες: εντοπισμό επαναλαμβανόμενων τμημάτων και μάθηση χώρου συγγένειας. Καμία από τις δύο δεν λύνει μόνη της το κοινό πρόβλημα ορίων, οικογενειακής απόδοσης και ιστορικής σχεσιακής εποπτείας. Αντίστοιχα, η μη παραμετρική πιθανοτική τμηματοποίηση του γρηγοριανού μέλους από τους Lanz--Hajič (2025) είναι σημαντικό εναλλακτικό μοντέλο για την M1, αλλά το ίδιο το αποτέλεσμά τους δείχνει ότι μια στατιστικά ή μνημονικά αποτελεσματική τμηματοποίηση δεν ταυτίζεται αναγκαστικά με την ιστορικο-μουσικολογική έννοια της φόρμουλας. Η διάκριση αυτή είναι κεντρική και για την παρούσα έρευνα.

\begin{xltabular}{0.98\linewidth}{@{}p{0.20\linewidth}p{0.31\linewidth}Y@{}}
\toprule
\textbf{Οικογένεια μεθόδων} & \textbf{Ισχυρό σημείο για το πρόβλημα} & \textbf{Ρόλος στην παρούσα διατριβή}\\
\midrule
Γραμμικό CRF & Ακολουθιακή οριοθέτηση, ώριμη εκπαίδευση με μερικές επισημειώσεις & Απλούστερος συγκριτής για να ελεγχθεί αν αρκεί ανεξάρτητη οριοθέτηση.\\
\midrule
Ημι-Μαρκοβιανό CRF / νευρωνικό τμηματικό μοντέλο & Βαθμολόγηση ολόκληρων τμημάτων, μεταβλητό μήκος, κοινή εκτίμηση ορίων και ετικέτας & \textbf{Κύρια μη επικαλυπτόμενη διαδρομή}, εφόσον η απογραφή του ΠΕ0 τεκμηριώνει επαρκή κάλυψη του μουσικολογικού αντικειμένου.\\
\midrule
Μάθηση μετρικής / ανάκτηση & Οργάνωση της συγγένειας σε συνεχή χώρο αναπαράστασης & Συστατικό της οικογενειακής ταυτότητας και της σχεσιακής εποπτείας, όχι αυτοτελής λύση οριοθέτησης.\\
\midrule
Μη επιβλεπόμενη τμηματοποίηση & Ανακάλυψη μονάδων χωρίς έτοιμα όρια ή οικογενειακές ετικέτες & Συγκριτικό μοντέλο της M1· δεν υποκαθιστά την πολιτισμικά τεκμηριωμένη οικογένεια.\\
\midrule
Παραγωγικά πιθανοτικά μοντέλα/HMM και συγγενείς προσεγγίσεις & Διαφανής διαδοχική πιθανοτική δομή και ισχυρό ιστορικό προηγούμενο στη μελέτη της ψαλτικής & Συμπληρωματικό σημείο αναφοράς· δεν προσφέρει τόσο άμεσο χειρισμό της ελεγχόμενης B2--EISE σχεσιακής αντίθεσης.\\
\bottomrule
\end{xltabular}

Με βάση αυτή την κριτική σύγκριση, ο EISE χρησιμοποιείται ως τεκμηριωμένη βασική αρχιτεκτονική κωδικοποίησης: συνδυάζει αναπαράσταση τμημάτων μεταβλητού μήκους, μερική επισημείωση, οικογενειακή γεωμετρία, σχεσιακή εποπτεία και τοπικές καταστάσεις μέσα σε ενιαίο πλαίσιο. Η μορφή της δομημένης εξόδου οριστικοποιείται στο ΠΕ0 πριν από τη M1, και η μεταγενέστερη επιλογή της συγκεκριμένης διαδικασίας αποκωδικοποίησης γίνεται μόνο μέσα στη μορφή εξόδου που έχει ήδη επιλεγεί.

\subsection{Κωδικοποιητής, μερική πιθανοφάνεια και αναγνώριση ανοικτού συνόλου}\label{app:eise-encoder}

Έστω $X=(e_1,\ldots,e_T)$ η ακολουθία ατομικών γεγονότων στην οριστικοποιημένη αναπαράσταση. Ο κοινός κωδικοποιητής παράγει τοπικές καταστάσεις
\[
h_t=g_{\theta}\!\left(R^{\star}(X)\right)_t.
\]
Η υποψήφια τμηματοποίηση ορίζεται στο κοινό επίπεδο ατομικών γεγονότων:
\[
\mathcal S_X=\{(a,b):1\le a\le b\le T,\ b-a+1\le L_{\max}\}.
\]
Το $L_{\max}$ αποτελεί όριο υπολογιστικής κάλυψης, οριστικοποιείται στα δεδομένα ανάπτυξης και είναι κοινό για B1, B2 και EISE. Αναφέρεται χωριστά το ποσοστό των τεκμηριωμένων θέσεων που βρίσκεται μέσα στο σύνολο υποψηφίων. Όταν η $R^\star$ χρησιμοποιεί μονάδες μεγαλύτερες από το ατομικό γεγονός, οι καταστάσεις και τα επιτρεπτά όρια προβάλλονται ρητά στο κοινό ατομικό επίπεδο, ώστε οι συγκρίσεις να αναφέρονται στο ίδιο σύνολο δυνατών ορίων.

Για κάθε υποψήφιο τμήμα $s=(a,b)$ ορίζεται σύνοψη ανεξάρτητη από την οικογενειακή αντιστοίχιση,
\[
r_s=g_{\mathrm{seg}}(h_a,\ldots,h_b),
\qquad
z_{\mathrm{id}}(s)=\operatorname{Normalize}\!\left(q_{\mathrm{id}}(r_s)\right).
\]
Η απλή αρχική υλοποίηση του $g_{\mathrm{seg}}$ χρησιμοποιεί τον μέσο όρο των τοπικών καταστάσεων μαζί με τις καταστάσεις αρχής και τέλους,
\[
r_s=\left[\frac{1}{b-a+1}\sum_{t=a}^{b}h_t;\,h_a;\,h_b\right],
\]
ενώ οποιαδήποτε πλουσιότερη σύνοψη της εσωτερικής σειράς επιλέγεται μόνο στα δεδομένα ανάπτυξης και εφαρμόζεται κοινά στις συγκρινόμενες συνθήκες. Η εσωτερική μουσική συντεταγμένη $u$ δεν απαιτείται για τη βαθμολόγηση ενός αυθαίρετου υποψήφιου τμήματος.

Στη μη επικαλυπτόμενη διαδρομή η δομημένη έξοδος υλοποιείται αρχικά με ημι-Μαρκοβιανό CRF μηδενικής τάξης. Για κάθε υποψήφιο τμήμα $s=(a,b)$ ορίζεται το σύνολο ατομικών γεγονότων που καλύπτει,
\[
\mathcal I(s)=\{a,a+1,\ldots,b\}.
\]
Έστω $\mathcal F_{\mathrm{known}}$ το λεξιλόγιο γνωστών οικογενειών και
\[
\mathcal F^{+}=\mathcal F_{\mathrm{known}}\cup\{\varnothing_F\},
\]
όπου $\varnothing_F$ δηλώνει έγκυρο τμήμα χωρίς ασφαλή απόδοση σε γνωστή οικογένεια. Ο χώρος μη επικαλυπτόμενων εξόδων είναι
\[
\mathcal Y_{\mathrm{seg}}(X)=
\left\{
 y\subseteq\mathcal S_X\times\mathcal F^{+}:
 \mathcal I(s)\cap\mathcal I(s')=\varnothing
 \ \text{για κάθε δύο διαφορετικά }(s,f),(s',f')\in y
\right\},
\]
και περιλαμβάνει και την κενή έξοδο $y=\varnothing$. Για γνωστή οικογένεια $F$ ορίζεται
\[
\phi_{\theta}(s,F)=\psi_{\mathrm{seg}}(r_s)
-\frac{d\!\left(z_{\mathrm{id}}(s),p_F\right)}{\tau},
\qquad \tau>0,
\]
ενώ για έγκυρο τμήμα χωρίς ασφαλή οικογενειακή απόδοση χρησιμοποιείται
\[
\phi_{\theta}(s,\varnothing_F)=
\psi_{\mathrm{seg}}(r_s)+\psi_{\mathrm{nov}}\!\left(\nu_s\right),
\qquad
\nu_s=\min_F d\!\left(z_{\mathrm{id}}(s),p_F\right).
\]
Κάθε ατομικό γεγονός που δεν καλύπτεται από επιλεγμένη εμφάνιση έχει μία μοναδική δομική απόφαση παράλειψης με βαθμό $\chi_{\theta,t}$. Επομένως
\[
\operatorname{score}_{\theta}(y,X)
=
\sum_{(s,f)\in y}\phi_{\theta}(s,f)
+
\sum_{t\notin\bigcup_{(s,f)\in y}\mathcal I(s)}\chi_{\theta,t},
\]
και
\[
P_{\theta}(y\mid X)
=
\frac{\exp\!\left(\operatorname{score}_{\theta}(y,X)\right)}
{\displaystyle\sum_{y'\in\mathcal Y_{\mathrm{seg}}(X)}
 \exp\!\left(\operatorname{score}_{\theta}(y',X)\right)}.
\]
Η κατασκευή δίνει μία μοναδική δομική αναπαράσταση σε κάθε έξοδο: το ακάλυπτο συμπλήρωμα δεν τεμαχίζεται σε εναλλακτικές ισοδύναμες περιοχές. Η δομική παράλειψη $\chi_{\theta,t}$ δεν αποτελεί μουσικολογική απόδοση υποβάθρου. Το τεκμηριωμένο υπόβαθρο παραμένει χωριστή κατάσταση αναφοράς για εποπτεία και βαθμονόμηση, ενώ η αποχή παραμένει μεταγενέστερος κανόνας απόφασης. Η $\psi_{\mathrm{nov}}$ και οι αντίστοιχοι κανόνες βαθμονόμησης οριστικοποιούνται αποκλειστικά στα δεδομένα ανάπτυξης με οικογένειες που έχουν εξαιρεθεί από την οικογενειακή εκπαίδευση. Δεν χρησιμοποιείται πρόσθετος πίνακας μετάβασης μεταξύ διαδοχικών οικογενειών στην αρχική προδιαγραφή. Ένας κοινός ασθενής όρος μήκους εξετάζεται μόνο ως ανάλυση αφαίρεσης συνιστώσας· δεν χρησιμοποιείται εξ αρχής οικογενειακά ειδική κατανομή μήκους.

Η διαθέσιμη μουσικολογική επισημείωση είναι εκ φύσεως \textbf{μερική}. Η πλήρης μουσικολογική αναφορά $A_X^{\mathrm{mus}}$ διατηρεί όλες τις τεκμηριωμένες εμφανίσεις, τα αποδεκτά ζεύγη ορίων, τις σχέσεις συνύπαρξης ή εναλλακτικής ερμηνείας, τις ελεγμένες περιοχές και τις ανεπίλυτες περιπτώσεις. Η προκαθορισμένη προβολή
\[
\Pi_{\mathrm{seg}}(A_X^{\mathrm{mus}})
=\{(W_j,A_j)\}_{j=1}^{J_X}
\]
παράγει τις περιοχές $W_j$ στις οποίες μπορεί να εφαρμοστεί νόμιμα η μη επικαλυπτόμενη δομημένη απώλεια και τους αντίστοιχους περιορισμούς αναφοράς $A_j$. Ο κωδικοποιητής μπορεί να χρησιμοποιεί το διαθέσιμο μουσικό συμφραζόμενο έξω από το $W_j$, αλλά οι δομημένοι βαθμοί και ο κανονικοποιητής της συγκεκριμένης συνεισφοράς ορίζονται μόνο μέσα στο $W_j$. Τα άκρα των περιοχών εφαρμογής επιλέγονται ώστε να προστατεύουν τις πλήρεις τεκμηριωμένες εμφανίσεις και δεν θεωρούνται μουσικά όρια. Αν μία εμφάνιση τέμνεται αναπόφευκτα από περιοχή εξαίρεσης, ο χειρισμός της και η αντίστοιχη απώλεια κάλυψης καταγράφονται· η εμφάνιση δεν μετατρέπεται σε περικομμένο θετικό ή αρνητικό παράδειγμα και παραμένει στο πλήρες μητρώο.

Για κάθε $W_j$ έστω $\mathcal Y_j$ ο αντίστοιχος χώρος επιτρεπτών μη επικαλυπτόμενων εξόδων και $\mathcal C_j(A_j)\subseteq\mathcal Y_j$ οι έξοδοι που ικανοποιούν τους εφαρμοστέους περιορισμούς της αναφοράς. Οι περιοχές εφαρμογής είναι ξένες ανά δύο. Κάθε κοινός περιορισμός μεταξύ εμφανίσεων εξετάζεται μέσα στην ίδια περιοχή· περιπτώσεις των οποίων οι αποδεκτές αναλύσεις εξαρτώνται μεταξύ τους διατηρούνται σε κοινή περιοχή ή καταγράφονται ως κοινή εξαίρεση από τον συγκεκριμένο δομημένο όρο. Οι περιορισμοί διατηρούν τα επιτρεπτά ζεύγη $\mathcal B_x$, τις επιτρεπτές οικογενειακές αποδόσεις και τους κοινούς περιορισμούς μεταξύ εμφανίσεων, αποκλείουν επιλεγμένες εμφανίσεις μέσα σε τεκμηριωμένο υπόβαθρο και δεν επιβάλλουν αρνητική ετικέτα σε μη ελεγμένες ή ανεπίλυτες περιοχές. Ορίζονται
\[
Z_{\theta,j}(X)=
\sum_{y\in\mathcal Y_j}
\exp\!\left(\operatorname{score}_{\theta,j}(y,X)\right),
\qquad
Z_{\theta,j}^{A}(X)=
\sum_{y\in\mathcal C_j(A_j)}
\exp\!\left(\operatorname{score}_{\theta,j}(y,X)\right),
\]
όπου $\operatorname{score}_{\theta,j}$ είναι η προηγούμενη βαθμολόγηση περιορισμένη σε επιλεγμένα τμήματα που περιέχονται στο $W_j$ και σε παραλείψεις γεγονότων $t\in W_j$. Η αντίστοιχη τοπική πιθανότητα είναι
\[
P_{\theta,j}(y\mid X)=
\frac{\exp\!\left(\operatorname{score}_{\theta,j}(y,X)\right)}
{Z_{\theta,j}(X)},
\]
και
\[
\mathcal L_j
=-\log\frac{Z_{\theta,j}^{A}(X)}{Z_{\theta,j}(X)}.
\]
Ισοδύναμα, για κάθε περιοχή
\[
\mathcal L_{\mathrm{seg},j}^{\mathrm{partial}}
=-\log\!\sum_{y\in\mathcal C_j(A_j)}P_{\theta,j}(y\mid X).
\]
Πριν από τη χρήση μιας περιοχής απαιτείται $\mathcal C_j(A_j)\neq\varnothing$. Ως συνοπτικός έλεγχος συμβατότητας γράφεται
\[
\mathcal C(A_X^{\mathrm{seg}})
=\prod_{j=1}^{J_X}\mathcal C_j(A_j)\neq\varnothing.
\]
Αν ο έλεγχος αποτύχει, καταγράφεται αν η αιτία είναι ασυμβατότητα επισημείωσης, υπέρβαση του $L_{\max}$, μη επιτρεπτός συνδυασμός ορίων ή πραγματική επικάλυψη. Πλήρως μη περιορισμένη περιοχή έχει $\mathcal C_j(A_j)=\mathcal Y_j$, άρα $Z_{\theta,j}^{A}=Z_{\theta,j}$ και $\mathcal L_j=0$.

Οι επιμέρους συνεισφορές συγκεντρώνονται με κοινό κανόνα κανονικοποίησης ανά έργο, ώστε ο αριθμός των περιοχών που παράγει η $\Pi_{\mathrm{seg}}$ να μην αλλάζει αυθαίρετα το βάρος του έργου. Με
\[
N_X^{\mathrm{elig}}=\sum_{j=1}^{J_X}|W_j|
\]
ορίζεται
\[
\mathcal L_{X,\mathrm{seg}}^{\mathrm{partial}}
=
\frac{\displaystyle\sum_{j=1}^{J_X}\mathcal L_j}
{N_X^{\mathrm{elig}}},
\]
και ο ίδιος κανόνας εφαρμόζεται σε όλες τις συγκρινόμενες νευρωνικές συνθήκες. Αν $N_X^{\mathrm{elig}}=0$, το συγκεκριμένο έργο δεν συνεισφέρει στον δομημένο όρο της απώλειας. Η κανονικοποίηση αυτή αφορά τον συγκεκριμένο δομημένο όρο· οι υπόλοιποι όροι της αντικειμενικής συνάρτησης διατηρούν τους δικούς τους προκαθορισμένους κανόνες στάθμισης.

\subsection{Συναρτήσεις απώλειας, πρότυπα αναφοράς, αρνητικά παραδείγματα και τοπική πληροφορία}\label{app:eise-loss}

Ο γράφος εξηγήσεων $G_{\mathrm{ex}}$ λειτουργεί ως πηγή εποπτείας. Η $\mathcal S_{\mathrm{StA}}$ προσφέρει την οργανωμένη αφετηρία για τη συγκρότησή του: το $P_{\mathrm{StA}}$ οδηγεί από την καταγεγραμμένη φόρμουλα στην πλήρη σύνθεση και στην ιστορική πηγή, και η επαληθευμένη σχέση μεταξύ μαρτύρων καταγράφεται ως ακμή με προέλευση, τύπο σχέσης και βαθμό τεκμηρίωσης. Παράλληλα, οι σχέσεις κοινής οικογένειας ή τεκμηριωμένης παραλλαγής του συνόλου εκπαίδευσης συγκροτούν το $E_{\mathrm{fam}}^{\mathrm{train}}$ του B2.

Για κάθε συνθήκη $m$ ορίζεται το σύνολο σχεσιακών ακμών $E_m$: για την B2 είναι οι αντιστοιχισμένες ενδοοικογενειακές σχέσεις και για την EISE οι τεκμηριωμένες εξηγητικές σχέσεις. Η σχεσιακή απώλεια εφαρμόζεται μόνο στην προβολή ταυτότητας,
\[
\mathcal L_{\mathrm{pairs}}^{\mathrm{id}}
=\sum_{(i,j)\in E_m}\omega_{ij}
 d\!\left(z_{\mathrm{id}}^{i},z_{\mathrm{id}}^{j}\right),
\qquad
\sum_{(i,j)\in E_m}\omega_{ij}=1.
\]
Τα βάρη, η έκθεση κάθε εμφάνισης και η δειγματοληψία ακμών ακολουθούν κοινό κανόνα στις δύο συνθήκες. Το επιβεβαιωτικό $P^{\star}$ χρησιμοποιείται αποκλειστικά στον έλεγχο $\boldsymbol{\Delta}_{\mathrm{ext}}$ και ποτέ ως αρνητικό παράδειγμα κατά την εκπαίδευση.

Για εμφανίσεις με αξιόπιστη οικογενειακή ετικέτα ο βοηθητικός όρος προτύπων αναφοράς ορίζεται ως
\[
\mathcal L_{\mathrm{proto}}
=-\sum_{i\in\mathcal I_{\mathrm{lab}}}\omega_i
\log
\frac{\exp[-d(z_i,p_{F_i})/\tau]}
{\sum_{F'}\exp[-d(z_i,p_{F'})/\tau]},
\qquad \sum_i\omega_i=1.
\]
Κάθε εκπαιδευμένο μοντέλο εκτιμά τα δικά του πρότυπα από τα ίδια επιλέξιμα αναγνωριστικά εκπαίδευσης, με κοινό κανόνα ενημέρωσης. Στη μη επικαλυπτόμενη διαδρομή, επειδή η οικογενειακή συμβατότητα εισέρχεται ήδη στη βαθμολογία του τμηματικού μοντέλου, η χρησιμότητα του βοηθητικού $\mathcal L_{\mathrm{proto}}$ ελέγχεται με προκαθορισμένη ανάλυση αφαίρεσης συνιστώσας.

Η συνολική αντικειμενική συνάρτηση είναι
\[
\mathcal L_{\mathrm{EISE}}
=\mathcal L_{X,\mathrm{seg}}^{\mathrm{partial}}
+\lambda_{\mathrm{rel}}\mathcal L_{\mathrm{pairs}}^{\mathrm{id}}
+\lambda_{\mathrm{proto}}\mathcal L_{\mathrm{proto}}
+\lambda_{\mathrm{loc}}\mathcal L_{\mathrm{local}}.
\]
Οι όροι υπολογίζονται με δηλωμένη κανονικοποίηση ανά έργο, οικογένεια και αριθμό διαθέσιμων ακμών, ώστε τα βάρη $\lambda$ να έχουν συγκρίσιμη σημασία μεταξύ B2 και EISE. Η κύρια τιμή του $\lambda_{\mathrm{rel}}$ και μία περιορισμένη περιοχή ευαισθησίας οριστικοποιούνται στα συμβολικά δεδομένα ανάπτυξης. Η επιλογή δεν χρησιμοποιεί το εκτελεστικό κριτήριο $\Gamma$, τα $S_m$ ή τις τελικές ιστορικές αντιθέσεις.

Πριν από την ερμηνεία της $\Delta_{\mathrm{hist}}$ πραγματοποιείται \textbf{διαγνωστικός έλεγχος της σχεσιακής εποπτείας} σε εξηγητικές ακμές ανάπτυξης που δεν χρησιμοποιήθηκαν στη σχεσιακή απώλεια. Πρώτα ελέγχεται αν οι αποστάσεις των αντίστοιχων πραγματώσεων στον χώρο οικογενειακής ταυτότητας $z_{\mathrm{id}}$ είναι μικρότερες στον EISE από ό,τι στη B2. Σε δεύτερο επίπεδο ελέγχεται αν η διαφοροποίηση μεταφέρεται και στις ευθυγραμμισμένες τοπικές καταστάσεις $\widetilde h(u)$. Ο έλεγχος είναι διαγνωστικός και δεν αποτελεί νέο επιβεβαιωτικό κριτήριο. Αν η σχεσιακή εποπτεία δεν παράγει μετρήσιμη διαφοροποίηση στο $z_{\mathrm{id}}$, ένα μηδενικό $\Delta_{\mathrm{hist}}$ δεν ερμηνεύεται ως ισχυρή ένδειξη ότι η ιστορική προέλευση των σχέσεων δεν προσθέτει πληροφορία· αναφέρεται μαζί με το αποτέλεσμα του διαγνωστικού ελέγχου. Αν η διαφοροποίηση εμφανίζεται στο $z_{\mathrm{id}}$ αλλά όχι στις τοπικές καταστάσεις, το αποτέλεσμα ερμηνεύεται ως οργάνωση της οικογενειακής ταυτότητας χωρίς αντίστοιχη μεταβολή της γεωμετρίας της συγκεκριμένης πραγμάτωσης.

Το $\mathcal L_{\mathrm{local}}$ διατηρεί, με μικρό προκαθορισμένο βάρος, τον στόχο προεκπαίδευσης με απόκρυψη συμβολικών γεγονότων. Η επάρκειά του ελέγχεται ρητά: από τις τοπικές καταστάσεις πρέπει να ανακτώνται προκαθορισμένες διαστηματικές και χρονικές ιδιότητες της \emph{συγκεκριμένης εμφάνισης} καλύτερα από μοντέλο αναφοράς που γνωρίζει μόνο την οικογένεια και την εσωτερική θέση $u$. Η σταθερότητα των προφίλ ανά εσωτερική θέση αξιολογείται μεταξύ διαφορετικών ρυθμίσεων τυχαιότητας. Αν ο πιλοτικός έλεγχος δείξει κατάρρευση της ενδοοικογενειακής γεωμετρίας, αντικειμενικές συναρτήσεις διατήρησης της διασποράς τύπου VICReg/Barlow Twins εξετάζονται μόνο ως αναλύσεις αφαίρεσης και όχι ως μέρος του βασικού μοντέλου.

Η αυτοεπιβλεπόμενη προεκπαίδευση χρησιμοποιεί το κοινό οριστικοποιημένο $\mathcal D_{\mathrm{sym}}^{\mathrm{train}}$ της \S\ref{sec:training-pool}, μετά τους αποκλεισμούς έργων, μαρτύρων και ομάδων συγγενών πηγών. Ο ίδιος αυστηρός διαχωρισμός εφαρμόζεται σε λεξιλόγιο/BPE, κανονικοποίηση, πρότυπα αναφοράς, ευθυγραμμίσεις, συχνότητες, επιλογή υπερπαραμέτρων και κάθε μετασχηματισμό που εκτιμάται από τα δεδομένα. Οι φυσικές οντότητες του γράφου χωρίζονται πριν από τη δημιουργία ακμών και στενά εξαρτημένες εκδόσεις ομαδοποιούνται πριν από τον διαχωρισμό. Η βαθμονόμηση της διαδικασίας μέτρησης του εκτελεστικού πιλοτικού ελέγχου παραμένει έξω από τη μάθηση και τη ρύθμιση των αναπαραστάσεων.

\subsection{\texorpdfstring{Λανθάνουσα μετρική και μουσική αντιστοίχιση \(u\)}{Λανθάνουσα μετρική και μουσική αντιστοίχιση u}}\label{sec:latent-metric}\label{app:eise-metric}

Η εσωτερική συντεταγμένη $u$ χρησιμοποιείται σε οριοθετημένες πραγματώσεις για τις οποίες υπάρχει εξωτερικά τεκμηριωμένη ή επαρκώς ελεγμένη μουσική αντιστοίχιση. Μπορεί επομένως να εφαρμόζεται σε επισημειωμένες πραγματώσεις εκπαίδευσης, σε ιστορικές εμφανίσεις της ανεξάρτητης αξιολόγησης και σε αποδεκτές αυτόματες αποδόσεις. Δεν αποτελεί όμως προϋπόθεση για την αναπαράσταση και τη βαθμολόγηση ενός αυθαίρετου υποψήφιου τμήματος, οι οποίες προηγούνται της οικογενειακής αντιστοίχισης.

Ο βασικός μετασχηματισμός είναι \textbf{κανονικοποιημένη λεύκανση με πλήρη συνδιακύμανση}. Έστω $d_h$ η κοινή διάσταση των τοπικών λανθανουσών καταστάσεων $h$ για B1, B2 και EISE. Για κάθε εκπαιδευμένο μοντέλο και ρύθμιση τυχαιότητας εκτιμώνται η μέση κατάσταση $\mu_{\mathrm{train}}$ και η συνδιακύμανση $\Sigma_{\mathrm{train}}$ στις επιλέξιμες καταστάσεις του συνόλου εκπαίδευσης, με μία συμβολή ανά ατομικό γεγονός. Ορίζονται
\[
a_{\mathrm{train}}=\frac{1}{d_h}\operatorname{tr}\Sigma_{\mathrm{train}}>0,
\qquad
W_{\mathrm{train}}
=\bigl(\Sigma_{\mathrm{train}}+\lambda a_{\mathrm{train}}I_{d_h}\bigr)^{-1/2},
\]
\[
\widetilde h_x(u)=W_{\mathrm{train}}\bigl(h_x(u)-\mu_{\mathrm{train}}\bigr).
\]
Η συμμετρική αντίστροφη τετραγωνική ρίζα ορίζει τη μετρική και το $\lambda>0$ ελέγχει την αριθμητική ευστάθεια. Η ίδια τιμή $\lambda$ επιλέγεται για τις συνθήκες που προκύπτουν από εκπαίδευση με προκαθορισμένο κριτήριο ευστάθειας στα συμβολικά δεδομένα εκπαίδευσης και ανάπτυξης και οριστικοποιείται πριν από την εξωτερική αξιολόγηση. Οι B1, B2 και EISE έχουν κοινή διάσταση $d_h$ και κοινή διαδικασία εκτίμησης· κάθε μοντέλο χρησιμοποιεί τον δικό του πίνακα, ενιαία για όλες τις οικογένειες και τις τιμές $u$. Μηδενική ολική διασπορά στο σύνολο εκπαίδευσης καταγράφεται ως αποτυχία διατήρησης πληροφορίας. Η διαγώνια τυποποίηση αποτελεί διαφορετικό μετασχηματισμό και μπορεί να αναφερθεί μόνο ως χωριστή προκαθορισμένη ευαισθησία της μετρικής.

Η συνολική θετική αλλαγή κλίμακας αφήνει τη συσχέτιση Spearman αμετάβλητη. Η ανισοτροπία, αντίθετα, μεταβάλλει τη συνεισφορά των διαστάσεων στη διασπορά και μπορεί να αλλάξει την κατάταξη των περιοχών κατά μήκος του $u$. Η κοινή διαδικασία λεύκανσης ορίζει ελεγχόμενη γεωμετρία μέτρησης· η μουσική σημασία της ελέγχεται από τους διαγνωστικούς ελέγχους της τοπικής πληροφορίας και την ανεξάρτητη μεταφορά.

Για τις $n_F(u)$ διαθέσιμες, αποδιπλοποιημένες ιστορικές πραγματώσεις μιας οικογένειας στην κοινή θέση $u$, ορίζονται
\[
\overline h_F(u)=\frac{1}{n_F(u)}\sum_{x\in F(u)}\widetilde h_x(u),
\]
\[
\widehat\Sigma_F(u)=\frac{1}{n_F(u)}
\sum_{x\in F(u)}
\bigl(\widetilde h_x(u)-\overline h_F(u)\bigr)
\bigl(\widetilde h_x(u)-\overline h_F(u)\bigr)^\top,
\qquad
V_F(u)=\frac{1}{d_h}\operatorname{tr}\widehat\Sigma_F(u).
\]
Ο παρονομαστής $n_F(u)$ ορίζει την εμπειρική δεύτερη κεντρική ροπή που χρησιμοποιείται σε όλες τις περιγραφικές αποσυνθέσεις. Η αβεβαιότητα εκτιμάται με τη δηλωμένη ομαδοποίηση πηγών και έργων· το πλήθος εμφανίσεων καταγράφεται χωριστά από το πλήθος ανεξάρτητων μαρτύρων.

Το $p_F$ είναι πρότυπο αναφοράς στον χώρο ταυτότητας $z_{\mathrm{id}}$. Το τοπικό πρότυπο $\widetilde\pi_F(u)$ εκτιμάται ως μέση ευθυγραμμισμένη τοπική κατάσταση των επιλέξιμων ιστορικών πραγματώσεων εκπαίδευσης της οικογένειας και οριστικοποιείται χωριστά. \textbf{Στην εφαρμογή στον Νικόλαο, εκτιμάται χωριστό $\widetilde\pi_{F,s}(u)$ για κάθε επιλέξιμο ιστορικό στρώμα $s$, αποκλειστικά από το ιστορικό σύνολο εκπαίδευσης, χωρίς έργα του Νικολάου ή της οριστικοποιημένης μελέτης περίπτωσης στην εκτίμηση προτύπων και μετρικής.} Ορίζονται
\[
D_F(u)=\frac{1}{n_F(u)}\sum_{x\in F(u)}
\|\widetilde h_x(u)-\widetilde\pi_F(u)\|^2,
\qquad
M_F(u)=\|\overline h_F(u)-\widetilde\pi_F(u)\|^2.
\]
Η ακριβής εμπειρική αποσύνθεση είναι
\[
\boxed{D_F(u)=d_h\,V_F(u)+M_F(u).}
\]
Αν χρησιμοποιηθεί αμερόληπτη εκτίμηση $V_F^{\mathrm{unb}}(u)$ με παρονομαστή $n_F(u)-1$, η αντίστοιχη ταυτότητα γίνεται $D_F=d_h\,[n_F(u)-1]/n_F(u)\,V_F^{\mathrm{unb}}+M_F$. Η βασική αναφορά διατηρεί παντού την εμπειρική σύμβαση με $n_F(u)$.

Η ιστορική γεωμετρία παραλλαγής του μοντέλου είναι
\[
\widehat\tau_{\mathrm{EISE,hist}}^F(u)=\sqrt{V_F(u)}.
\]
Η συνολική απόσταση $D_F$ περιγράφει τη θέση των πραγματώσεων ως προς το οριστικοποιημένο πρότυπο, ενώ το $V_F$ μετρά την ενδοοικογενειακή μεταβλητότητα. Η εφαρμογή στον Νικόλαο χρησιμοποιεί το ίδιο σύστημα με δείκτες $\mathrm{Nik},F,s$ για την απόσταση από κάθε στρώμα και αναφέρει δίπλα-δίπλα \textbf{$d_hV_{\mathrm{Nik},F}$ και $M_{\mathrm{Nik},F,s}$}. Η διαφορά διασποράς από το αντίστοιχο στρώμα ορίζεται χωριστά στην §~\ref{app:match-nikolaos}, ώστε το γινόμενο $d_hV$ να διατηρεί παντού τη σημασία του. Η σταθερότητα των προφίλ ελέγχεται μεταξύ ρυθμίσεων τυχαιότητας και με τη μέση μη τετραγωνισμένη ζευγαρωτή απόσταση στην ίδια οριστικοποιημένη μετρική.

Η εσωτερική συντεταγμένη $u$ ορίζεται από κοινές συλλαβές και τεκμηριωμένες μελωδικές αντιστοιχίες. Μεταξύ διαδοχικών σημείων αναφοράς εφαρμόζεται περιορισμένη μονοτονική ευθυγράμμιση ακολουθιών με Dynamic Time Warping (DTW) και προκαθορισμένη συνάρτηση κόστους βασισμένη σε μουσικολογικά χαρακτηριστικά. Για το χωριστό χρονικό διαγνωστικό της M5, οι αντίστοιχες εσωτερικές περιοχές δεν επαναπροσδιορίζονται με χρονικά κόστη: χρησιμοποιούνται κατά προτεραιότητα οι ίδιες ζώνες που έχουν οριστεί εκ των προτέρων $z$ των $A(z)$/$\Gamma(z)$ και, όταν απαιτείται λεπτότερη αντιστοίχιση, αυτή υπολογίζεται αποκλειστικά από το στρώμα $N$, χωρίς χρήση των χρονικών αξιών $T$. Η αβεβαιότητα της αντιστοίχισης μεταφέρεται ως $U_{\mathrm{seg}}$ στο χρονικό διαγνωστικό. Αναπτύξεις και συντμήσεις επιτρέπουν αντιστοιχίες πολλά-προς-ένα ή ένα-προς-πολλά μέσα στην ίδια μουσικά τεκμηριωμένη περιοχή. Οι τοπικές καταστάσεις συγκεντρώνονται με βάρη που αθροίζουν στη μονάδα ανά περιοχή· η επανάληψη μιας αντιστοίχισης δεν δημιουργεί πρόσθετο ανεξάρτητο μάρτυρα. Περιοχές χωρίς αξιόπιστη αντιστοιχία φέρουν μάσκα διαθεσιμότητας και αναφέρονται στην κάλυψη. Η ίδια αντιστοίχιση, τα βάρη και το κοινό πεδίο $u$ χρησιμοποιούνται από B0, B1, B2 και EISE. Η αβεβαιότητα της ευθυγράμμισης μεταφέρεται στα διαστήματα εκτίμησης και η μετάβαση σε αδρότερες περιοχές ακολουθεί αποκλειστικά τον κανόνα του ΠΕ0. Ο συμβολικός κανόνας αντιστοίχισης οριστικοποιείται πριν από οποιοδήποτε επιβεβαιωτικό εκτελεστικό αποτέλεσμα.

\subsection{Μοντέλα αναφοράς, αναλύσεις αφαίρεσης και όριο πολυπλοκότητας}\label{app:eise-baselines}

\subsubsection{Περιορισμένος έλεγχος τμηματικής αρχιτεκτονικής}\label{app:eise-model-selection}
Η αρχιτεκτονική A οριστικοποιείται πριν από την κατάταξη των αναπαραστάσεων και χρησιμοποιείται στην επιλογή $R^\star$ και ως κοινή τμηματική υποδομή των Α2 και Μ2. Οι B/C αποτελούν αυτοτελείς συγκριτές, που αξιολογούνται σε όλο το πλέγμα $U\times\Phi$ με κύριο μέγεθος τη μέση αντίθεση $\Delta_{A-C}$ στο $F_{1,\mathrm{seg+fam}}$. Συμπληρωματικά αναφέρονται η A--B και η διαφορά A--C στο διαγνωστικό F1 ορίων, με κοινή αναδειγματοληψία συζευγμένων προφίλ ανά ομάδα έργων--πηγών. Οι επιδόσεις B/C στην $R^\star$ έχουν περιγραφικό χαρακτήρα. Η B0 ταυτίζεται με $A(R^\star)$· οι A0/AL αποτελούν προδηλωμένους διαγνωστικούς ελέγχους της ποινής μήκους. Αν το ΠΕ0 ενεργοποιήσει τη διαδρομή επικαλυπτόμενων τμημάτων, η προδιαγραφή εξόδου και οι κατάλληλοι συγκριτές οριστικοποιούνται πριν από τις κύριες εκπαιδεύσεις.



\subsubsection{Ισχυρός Transformer στο επίπεδο ατομικών γεγονότων}\label{app:event-transformer}

Ο $T_{\mathrm{evt}}$ είναι αυτοπαλίνδρομος Transformer που εκπαιδεύεται στο ενιαίο επίπεδο αναφοράς ατομικών μελωδικών γεγονότων. Χρησιμοποιεί την ίδια επιλέξιμη διαμόρφωση γνωρισμάτων με την κύρια σύγκριση, αποκλειστικά προγενέστερο μουσικό συμφραζόμενο και τους ίδιους διαχωρισμούς έργων και ομάδων συγγενών πηγών. Η επιλογή χωρητικότητας, μήκους συμφραζομένου, κανονικοποίησης και ρυθμίσεων βελτιστοποίησης γίνεται μόνο στο σύνολο ανάπτυξης, με τουλάχιστον τον ίδιο προϋπολογισμό ρύθμισης που παρέχεται στις κύριες νευρωνικές συνθήκες.

Η σύγκριση με αναπαραστάσεις διαφορετικού λεξιλογίου βασίζεται σε κοινούς στόχους πάνω στο ενιαίο επίπεδο αναφοράς και σε ίδιες οριστικοποιημένες μονάδες αξιολόγησης. Οι ακατέργαστες τιμές perplexity ή εντροπίας ανά συμβολική μονάδα (token) αναφέρονται μόνο ως διαγνωστικά, επειδή μεταβάλλονται με το μέγεθος και τη δομή του λεξιλογίου. Για εργασίες που απαιτούν αναπαράσταση τμήματος, οι καταστάσεις των γεγονότων μέσα στο ίδιο υποψήφιο τμήμα συγκεντρώνονται με οριστικοποιημένο κανόνα συγκέντρωσης (pooling) που επιλέγεται στα δεδομένα ανάπτυξης· ο ίδιος κανόνας εφαρμόζεται χωρίς επαναπροσαρμογή στο σύνολο ελέγχου.

Ο $T_{\mathrm{evt}}$ λειτουργεί ως αυστηρός έλεγχος της ανάγκης για ρητή μονάδα θέσης. Η θέση αποκτά πρόσθετη υπολογιστική τεκμηρίωση όταν οι αναπαραστάσεις σε επίπεδο θέσης διατηρούν πλεονέκτημα ή συμπληρωματική πληροφορία σε κοινούς στόχους, σε άγνωστα έργα και ομάδες συγγενών πηγών, απέναντι σε αυτό το μοντέλο σε επίπεδο γεγονότων. Αναφέρονται μαζί η επίδοση, ο αριθμός παραμέτρων, ο πραγματικός υπολογιστικός χρόνος και η ευαισθησία στο μήκος συμφραζομένου, ώστε η ερμηνεία να διαχωρίζει το όφελος της δομικής μονάδας από το όφελος της χωρητικότητας.


Οι $B0/B1/B2/\mathrm{EISE}$ αποτελούν το κεντρικό πείραμα. Συμπληρωματικά μοντέλα αναφοράς είναι n-gram/TF--IDF και medoid οικογενειών, HMM/MDL, BiLSTM--CRF, Transformer--CRF και η αλγοριθμικά επαγόμενη Pitman--Yor συνθήκη. Οι κύριες αναλύσεις αφαίρεσης εξετάζουν εναλλακτική επιβλεπόμενη αντιθετική συνάρτηση απώλειας με πολλαπλά θετικά παραδείγματα, βάρη τεκμηρίωσης και χωριστή αξιοποίηση των κατηγοριών σχέσεων $E_{\mathrm{an}}/E_{\mathrm{tr}}$. Πρόσθετοι μηχανισμοί ενεργοποιούνται μόνο όταν ο έλεγχος στο σύνολο ανάπτυξης εντοπίσει συγκεκριμένη αστοχία, όπως κατάρρευση της γεωμετρίας παραλλαγής ή ανεπαρκή βαθμονόμηση ανοικτού συνόλου. Ο υπολογιστικός προϋπολογισμός παραμένει συγκεντρωμένος στο πείραμα με αντιστοιχισμένες συνθήκες και στις αναλύσεις αφαίρεσης που μπορούν να αλλάξουν την ερμηνεία του.


\section{Ανθεκτικότητα και λοιπές διερευνητικές αναλύσεις}\label{app:robust}

\subsection{Διαγνωστικοί έλεγχοι αναπαραστάσεων}\label{app:robust-repr}

\textbf{Έλεγχος αφαίρεσης χρονικών γνωρισμάτων.} Σύγκριση της ίδιας συνθήκης με και χωρίς χρονικά πεδία. Ελέγχει αν η παράλειψη του χρόνου καθιστά την αναπαράσταση τεχνητά φτωχότερη από τα δυτικά μοντέλα αναφοράς, στα οποία η διάρκεια αποτελεί βασικό πεδίο.

\textbf{Αφαιρέσεις ομάδων γνωρισμάτων γύρω από την πλήρη περιγραφή.} Μετά την οριστικοποίηση της \(R^\star\), στο \(U\) της \(R^\star\) εκπαιδεύονται δύο προκαθορισμένες διαγνωστικές συνθήκες, \(N{+}M{+}L\) και \(N{+}T{+}L\), με τον ίδιο κοινό εκτιμητή, τις ίδιες \(n_{\mathrm{rep}}\) αντιστοιχισμένες επαναλήψεις και τους ίδιους δείκτες αναδειγματοληψίας. Μαζί με την \(N{+}T{+}M\), που ήδη ανήκει στο πλέγμα, σχηματίζουν αφαιρέσεις μίας ομάδας (\(T\), \(M\) ή \(L\)) από την \(N{+}T{+}M{+}L\) και δείχνουν αν η πλήρης περιγραφή εξαρτάται από κάθε ομάδα υπό την παρουσία των άλλων δύο. Οι συνθήκες αυτές δεν ανήκουν στο πλέγμα, δεν ανακατατάσσουν την \(R^\star\) και δεν εισέρχονται στο \(\mathcal R^\star\) και δεν μπορούν να επιλεγούν ως \(R_{\mathrm{sens}}\)· η ερμηνεία τους είναι υπό συνθήκη ως προς το επιλεγμένο \(U\). Αναφέρονται \(\widehat M\) και \(\widehat G\) με διαστήματα. Αν αποδοθούν στατιστικοί ισχυρισμοί στις τρεις αφαιρέσεις, χρησιμοποιείται χωριστή προκαθορισμένη οικογένεια συγκρίσεων. Το πρόσθετο υπολογιστικό κόστος είναι \(2\,n_{\mathrm{rep}}\) εκπαιδεύσεις.

\textbf{Μοντέλο γεωμετρίας.} Χρησιμοποιεί μόνο τη δομή της ακολουθίας χωρίς μουσική πληροφορία και λειτουργεί ως αρνητικός έλεγχος. Υψηλή επίδοση αυτού του μοντέλου θα έδειχνε ότι σημαντικό μέρος του σήματος εξηγείται από τη γεωμετρία του υλικού και θα περιόριζε την ερμηνεία των πλήρων συνθηκών.

\textbf{Έλεγχος ανακατασκευής.} Κάθε συνθήκη συνοδεύεται από μέτρηση του τι ανακτάται και τι χάνεται ως προς την αρχική σημειογραφική μορφή. Λειτουργεί ως δικλείδα ορθότητας, \emph{όχι} ως κριτήριο επιλογής: η πιο λεπτομερής αναπαράσταση κερδίζει εξ ορισμού.

\textbf{Λεξιλόγιο και αραιότητα.} Μέγεθος λεξιλογίου, ποσοστό σπάνιων μονάδων και μήκος ακολουθίας ανά συνθήκη, ως διαγνωστικά και όχι ως κριτήρια.

\subsection{Δεύτερες οικογένειες μοντέλων}\label{app:robust-models}

Χρησιμοποιούνται \textbf{μοντέλα σταθερής χωρητικότητας}: παρεμβαλλόμενο μοντέλο \(n\)-γραμμάτων με εξομάλυνση ως πλήρως ερμηνεύσιμος κύριος εκτιμητής, και μοντέλο μεταβλητής τάξης ή πολύ μικρό νευρωνικό μοντέλο ακολουθίας ως έλεγχος ανθεκτικότητας.

\textbf{Η πληροφοριακή επίδραση αποδίδεται στην αναπαράσταση μόνο όταν η κατεύθυνσή της αναπαράγεται σε δύο μοντέλα διαφορετικής επαγωγικής μεροληψίας.} Το κριτήριο εστιάζει έτσι στην αναπαράσταση ως κοινό παράγοντα των δύο εκτιμητών. Αναπαραγωγή του φαινομένου και σε καλά εξομαλυμένο μοντέλο \(n\)-γραμμάτων αποτελεί ισχυρή ένδειξη ότι το αποτέλεσμα είναι ανθεκτικό στην επιλογή μοντέλου και δεν εξαρτάται εμφανώς από τη μεγάλη χωρητικότητα του Transformer.

Εφαρμόζεται επιπλέον το μοντέλο γεωμετρίας ως αρνητικός έλεγχος. Η σύγκριση με τα τμήματα ευθυγραμμισμένα σε πληροφοριακές κορυφές και η ανάλυση με εξαίρεση των πρώτων \(k\) γεγονότων διαχωρίζουν ειδικότερα τη γενική δυναμική ενός ορίου από την ενδο-τμηματική οργάνωση που αποδίδεται στη θέση.

\subsection{Διερευνητική πληροφοριακή ανάλυση}\label{app:robust-info}

Διερευνητικά θα εξεταστεί αν τα προφίλ επιτρέπουν τον εντοπισμό χαρακτηριστικού σημείου ή περιοχής εντός της θέσης όπου η αβεβαιότητα της προβλεπτικής κατανομής μεταβάλλεται συστηματικά. Η ανάλυση αυτή παραμένει \textbf{διερευνητική επέκταση} και λειτουργεί συμπληρωματικά προς την προκαθορισμένη δευτερεύουσα Υ3.

\subsection{Διερευνητική μάθηση από ιστορικές αντιστοιχίες}\label{app:robust-reverse}

Η ίδια ιστορική υποδομή επιτρέπει μία διερευνητική επέκταση χωρίς να μετατρέπει την ιστορική σχέση σε αυτόματη οικογενειακή ισοδυναμία. Οι τεκμηριωμένες ιστορικές σχέσεις συγκροτούν γράφο αντιστοιχιών $G_{\mathrm{hist}}=(V,E_{\mathrm{hist}})$, του οποίου οι ακμές δηλώνουν συγκεκριμένη ιστορικά τεκμηριωμένη σύνδεση και όχι, από μόνες τους, κοινή οικογενειακή ταυτότητα. Η συνεκτικότητα του γράφου δεν αποδεικνύει σχέση ισοδυναμίας και δεν χρησιμοποιείται για να επιβάλει μεταβατικά family labels.

Για τη διερευνητική μάθηση αναπαράστασης ταυτότητας χρησιμοποιείται μόνο το υποσύνολο επιλέξιμων identity-positive ακμών
\[
E_{\mathrm{id}}^{+}=E_{\mathrm{hist}}\cap E_{\mathrm{hom}}\cap E_{\mathrm{fam}},
\]
όπου η αντιστοιχία του συγκεκριμένου χωρίου και η κοινή οικογενειακή ένταξη έχουν επαληθευθεί ανεξάρτητα. Αναζητείται αναπαράσταση $R$ που μειώνει τη διασπορά στα επιλέξιμα ιστορικά συνδεδεμένα ζεύγη και ταυτόχρονα διατηρεί ικανότητα διάκρισης μεταξύ διαφορετικών οικογενειών:
\[
\mathcal{L}_{\mathrm{inv}}(R)=
\sum_{(i,j)\in E_{\mathrm{id}}^{+}} d_R(x_i,x_j)
+\lambda\,\mathcal{L}_{\mathrm{sep}}(R).
\]
Οι συνεκτικές περιοχές του $G_{\mathrm{hist}}$ εξετάζονται μόνο διερευνητικά ως \textbf{δίκτυα ιστορικής αντιστοίχισης}. Ελέγχεται αν οι αναπαραστάσεις που μαθαίνονται από τις επιλέξιμες ακμές γενικεύουν σε μη χρησιμοποιημένες εμφανίσεις και αν η γεωμετρία τους συμφωνεί με τις ανεξάρτητα επικυρωμένες οικογένειες και με την $R^\star$. Η ανάλυση δεν παράγει νέες οικογενειακές κλάσεις, δεν ανατροφοδοτεί το ground truth και παραμένει εκτός της επιβεβαιωτικής ιεραρχίας.

\subsection{Εκτελεστικά επαγόμενες περιοχές}\label{app:robust-performance}

Σε χωριστό διερευνητικό πείραμα η κατεύθυνση της ανάλυσης αντιστρέφεται. Χωρίς μουσικολογικές ετικέτες, εφαρμόζεται ανίχνευση μεταβολών στα προφίλ μεταβλητότητας μεταξύ εκτελεστών \(V(W)\) ώστε να προκύψουν \textbf{εκτελεστικά επαγόμενες δομικές περιοχές}. Οι περιοχές αυτές συγκρίνονται εκ των υστέρων με τις τεκμηριωμένες θέσεις και με την αλγοριθμικά επαγόμενη συμβολική τμηματοποίηση.

Το πείραμα ελέγχει αν οι μεταβολές που προκύπτουν από πολλαπλές επαγγελματικές εκτελέσεις συμπίπτουν με μέρος της θεωρητικής δομής χωρίς να χρησιμοποιούνται μουσικολογικές ετικέτες στην ανίχνευσή τους. Παραμένει διερευνητικό και δεν συμμετέχει στην επιλογή της αναπαράστασης ή στις επιβεβαιωτικές διορθώσεις πολλαπλότητας.

\subsection{Ιστορικά τεκμηριωμένη μίμηση και μελλοντική ιστορική ανάκτηση}\label{app:robust-ανάκτηση}

Ο Αχιλλέας Χαλδαιάκης καταγράφει ως χαρακτηριστικό παράδειγμα της επίδρασης του Δανιήλ σε μεταγενέστερους τη μελοποίηση του δοξαστικού της Κασσιανής από τον Πέτρο Λαμπαδάριο «κατὰ μίμησιν Δανιὴλ πρωτοψάλτου» (Chaldaeakes, 2010, σ.~47)· στη σχετική παράδοση απαντά και η διατύπωση «κατὰ παράδοσιν». Η μαρτυρία τεκμηριώνει σχέση μίμησης ή παράδοσης σε επίπεδο μελοποιητικού τρόπου, χωρίς να ταυτοποιεί, με τα διαθέσιμα τεκμήρια, συγκεκριμένο πρότυπο έργο του Δανιήλ.

Η ρήτρα χρησιμοποιείται ως \textbf{διερευνητικός έλεγχος μεταφοράς} της επικυρωμένης αναπαράστασης σε επίπεδο μελοποιητικού τρόπου και ιστορικά τεκμηριωμένης σχέσης παράδοσης.

\textbf{Πρώτον, η διαθέσιμη μαρτυρία τεκμηριώνει σχέση με τον μελοποιητικό τρόπο και την παράδοση του Δανιήλ.} Αυτό το επίπεδο ιστορικής τεκμηρίωσης ορίζει τον στόχο της ανάλυσης και επιτρέπει να εξεταστεί αν η επικυρωμένη αναπαράσταση ανακτά υφολογική συγγένεια σε κλίμακα οικογενειών θέσεων. Η ταυτοποίηση συγκεκριμένου προτύπου έργου παραμένει αυτοτελές ερώτημα κριτικής των πηγών.

\textbf{Δεύτερον, η αλλαγή είδους παρέχει αυστηρό έλεγχο γενίκευσης σε διαφορετικό ρεπερτόριο.} Το κύριο σώμα της εφαρμογής είναι παπαδικό, ενώ το δοξαστικό\footnote{Δοξαστικό: υμνογραφικό και μελοποιητικό είδος που ψάλλεται σε συγκεκριμένες θέσεις της ακολουθίας και μπορεί να ανήκει στο στιχηραρικό ρεπερτόριο.} της Κασσιανής είναι στιχηραρικό. Επιτυχής μεταφορά θα δείξει ότι η αναπαράσταση αποτυπώνει υφολογική συγγένεια και σε διαφορετικό είδος από εκείνο στο οποίο θεμελιώθηκε· ο βαθμός μεταφοράς αναφέρεται ως μέτρο γενίκευσης μεταξύ διαφορετικών μορφολογικών περιβαλλόντων.

Ο έλεγχος αναφέρεται ως διερευνητικός, με ρητή αντιστοίχιση ως προς ήχο, λειτουργική περιοχή και μορφολογική κλίμακα, και δεν συμμετέχει σε καμία επιβεβαιωτική διόρθωση πολλαπλότητας.

Η ιστορική ανάλυση παραμένει πηγοκεντρική: το υπολογιστικό σύστημα επιταχύνει την ανάκτηση υποψηφίων και επεκτείνει την ανάλυση σε ευρεία κλίμακα ρεπερτορίου, ενώ η απόδοση ιστορικής σχέσης θεμελιώνεται σε χρονολογία, πατρότητα, εκδοτική ή εξηγητική σχέση και ανεξάρτητα τεκμήρια.

Μια φυσική επόμενη βαθμίδα είναι η \textbf{αντίστροφη ιστορική ανάκτηση}: από μια πραγμάτωση της Νέας Μεθόδου να ανακτώνται παλαιότερες μορφές συμβατές με τις ιστορικές πηγές ή οικογένειες θέσεων. Η αντίστροφη σχέση ανακτά σύνολα ιστορικά συμβατών υποψηφίων, αποτυπώνοντας τις πολλαπλές διαδρομές αναλυτικότερης απόδοσης που τεκμηριώνει η εξηγητική παράδοση. Πρώτο μετρήσιμο βήμα είναι η κάλυψη ενός σαφώς ορισμένου σώματος της Νέας Μεθόδου από τεκμηριωμένες οικογένειες θέσεων, με το ποσοστό να αναφέρεται πάντοτε μαζί με τον κανόνα ένταξης που ορίζει τον παρονομαστή. Μια επόμενη δοκιμή ανάκτησης με δηλωμένη συνθήκη πηγής μπορεί να συγκρίνει EISE, B2 και B0, κρατώντας εκτός ολόκληρες γραμμές προέλευσης, ώστε οι δεσμοί που τεκμηριώνονται από τις πηγές και ορίζουν το κριτήριο αναφοράς να μην έχουν χρησιμοποιηθεί στην εκπαίδευση. Η κατεύθυνση αυτή επαναχρησιμοποιεί την υπάρχουσα υποδομή και τη συγκριτική λογική του ερευνητικού προγράμματος χωρίς να μεταβάλλει το παρόν ερευνητικό πρόγραμμα.



\subsection{Δευτερεύουσα ανάλυση κοινού κειμενικού χωρίου στο παπαδικό μέλος}\label{sec:text-locus}

Η ανάλυση χρησιμοποιεί πλήρη έργα και κοινά κειμενικά χωρία των οποίων τα όρια καθορίζονται από το κείμενο πριν από την εξέταση των μουσικών αποτελεσμάτων. Η κύρια εφαρμογή αφορά Χερουβικά του ίδιου ήχου από ανεξάρτητες μελοποιήσεις και τεκμηριωμένες ομάδες συγγενών πηγών. Τα Κοινωνικά εξετάζονται μόνο σε ομάδες με ακριβώς κοινό στίχο. Κάθε μοναδικό χωρίο κωδικοποιείται μία φορά, με σταθερό αναγνωριστικό, όρια λέξεων και συλλαβών και θέση μέσα στο κείμενο, και η ίδια κωδικοποίηση χρησιμοποιείται σε όλες τις μελοποιήσεις. Δεν επανατεμαχίζεται ανά έργο για να αυξηθεί τεχνητά το δείγμα.

Η ποσοτική μελέτη της μελισματικότητας στα κεκραγάρια από τον Σπανουδάκη (2017, ιδίως το κεφάλαιο της στατιστικής ανάλυσης) προσφέρει μουσικολογικό προηγούμενο για την άνιση ανάπτυξη του μέλους σε επίπεδο συλλαβής, λέξης και μεγαλύτερου κειμενικού χωρίου. Η παρούσα δευτερεύουσα ανάλυση χρησιμοποιεί αυτή την προβληματική για να περιγράψει τις μουσικές επιλογές σε προκαθορισμένους κειμενικούς τόπους των Χερουβικών και, όταν υπάρχει κοινός στίχος, των Κοινωνικών. Οι συγκρίσεις διατηρούν τους υφιστάμενους ελέγχους για μήκος, μελισματική πυκνότητα, ήχο, έργο και λειτουργικό τμήμα.

Η κειμενική περιγραφή διακρίνει δύο άξονες. Ο πρώτος αφορά τη \textbf{λειτουργία του λόγου}, όπως αίνος ή δοξολογία, προτροπή, περιγραφική διατύπωση και λειτουργική επευφημία. Ο δεύτερος αφορά τη \textbf{θεολογική αναφορά}, όπως τριαδική, χριστολογική, αγγελική, ευχαριστιακή ή αναφορά στη σύναξη, και μπορεί να δέχεται περισσότερες από μία ετικέτες όταν αυτές συνυπάρχουν. Οι τελικές τιμές και οι κανόνες τους ορίζονται σε μικρό εγχειρίδιο κωδικοποίησης με βάση το ίδιο το λειτουργικό κείμενο και τις σχετικές πηγές. Για την ποσοτική σύνοψη επιλέγεται εκ των προτέρων μία μόνο κειμενική αντίθεση όταν η απογραφή δείχνει επαρκή κάλυψη· οι υπόλοιπες ετικέτες παραμένουν μεταδεδομένα.

Δύο ανεξάρτητοι αξιολογητές, ο ένας με λειτουργιολογική ή θεολογική και ο άλλος με μουσικολογική κατάρτιση, κωδικοποιούν τα μοναδικά χωρία χωρίς πρόσβαση στις παρτιτούρες, στις ηχογραφήσεις, στα ονόματα των μελοποιών, στις οικογενειακές αποδόσεις ή στα μουσικά αποτελέσματα. Βλέπουν το κείμενο, το αναγκαίο γλωσσικό συμφραζόμενο και κοινό πακέτο λειτουργικών πηγών και αιτιολογούν κάθε κρίση μόνο με κειμενικά ή λειτουργιολογικά τεκμήρια. Οι αρχικές κρίσεις αποθηκεύονται πριν από συζήτηση. Αναφέρονται η απλή συμφωνία, ο πίνακας διαφωνιών και ο συντελεστής Krippendorff~$\alpha$ για ονομαστικές κατηγορίες, όπου το πλήθος και η διακύμανση των μοναδικών χωρίων το επιτρέπουν. Ανεπίλυτες διαφωνίες εξετάζονται από τρίτο ειδικό που έχει οριστεί εκ των προτέρων με τους ίδιους περιορισμούς πρόσβασης ή παραμένουν αβέβαιες. Καμία κατηγορία δεν προσδιορίζεται από μήκος μελίσματος, χρονική ανάπτυξη, τροπική μεταβολή, κατάληξη, πλήθος γεγονότων ή υπολογιστικό δείκτη.

Ως \textbf{προφίλ μουσικής έμφασης} νοείται η πολυδιάστατη περιγραφή της μουσικής μεταχείρισης ενός προκαθορισμένου κειμενικού χωρίου. Δεν κατασκευάζεται συνολικός βαθμός έμφασης. Για κάθε χωρίο αναφέρονται χωριστά η σημειογραφική ανάπτυξη, ο αριθμός ατομικών γεγονότων, η μελισματική ανάπτυξη, η σχετική περιοχή και το εύρος φωνής, η ενεργή τροπική κατάσταση και οι μεταβολές της, ο μουσικός λειτουργικός ρόλος και ο τύπος κατάληξης. Ο μουσικός λειτουργικός ρόλος ---όπως εναρκτήρια, καταληκτική ή συνδετική χρήση--- διακρίνεται από το λειτουργικό τμήμα της λατρευτικής ακολουθίας, το οποίο καταγράφεται χωριστά. Η οικογενειακή ταυτότητα εξετάζεται μόνο για ανεξάρτητα επικυρωμένες οικογένειες και αποδίδεται από το εξωτερικό μουσικολογικό στρώμα ή από την οριστικοποιημένη μουσική συνθήκη $N{+}T{+}M$, ώστε οι νέες κειμενικές ετικέτες να μην εισέρχονται στην οικογενειακή αναγνώριση.

Οι υπάρχουσες ποσότητες $I_{\Omega}$, $D_{\mathrm{within}}$, $p_x$ και $d_{\mathrm{time}}$ χρησιμοποιούνται μόνο όπου πληρούνται οι προϋποθέσεις κάλυψης και αντιστοίχισης της M5. Το $I_{\Omega}$ υπολογίζεται με το οικογενειακό συμφραζόμενο της τεκμηριωμένης συνεχούς διαδοχής και δεν επανεκκινεί απλώς λόγω αλλαγής κειμενικού χωρίου. Στις διακοπές εφαρμόζεται ο κοινός προκαθορισμένος κανόνας της M5 για επανέναρξη του συμφραζομένου, συμμετοχή της πρώτης επόμενης εμφάνισης στο $I_{\Omega}$ και αναφορά της διαθέσιμης ακολουθιακής κάλυψης. Το $D_{\mathrm{within}}$ παραμένει ενδοοικογενειακή ποσότητα και δεν μετατρέπεται σε ενιαίο βαθμό αυθαίρετης φράσης. Τα $p_x$ και $d_{\mathrm{time}}$ αναφέρονται μόνο σε οικογένειες με αξιόπιστες κοινές εσωτερικές περιοχές. Αν μια τοπική αναπαράσταση έχει εκπαιδευτεί με πληροφορία του $L$, η χρήση της σε σχέση με κειμενική κατηγορία δηλώνεται ρητά και δεν παρουσιάζεται ως ανεξάρτητη από το κείμενο.

Για έργο $w$ και μοναδικό κειμενικό χωρίο $\ell$ ορίζεται η παράγωγη σημειογραφική ανάπτυξη
\[
Q^{\mathrm{txt}}_{w\ell}=\sum_{e\in(w,\ell)} v_e,
\]
όπου $v_e$ είναι η ήδη κωδικοποιημένη σχετική σημειογραφική χρονική αξία του γεγονότος $e$. Αν $n^{\mathrm{syl}}_{\ell}$ είναι ο αριθμός των κανονικών συλλαβών του προκαθορισμένου σταθερού κειμένου, η πρώτη περιγραφική κανονικοποιημένη ποσότητα είναι
\[
\eta^{\mathrm{txt}}_{w\ell}=\ln\!\left(\frac{Q^{\mathrm{txt}}_{w\ell}}{n^{\mathrm{syl}}_{\ell}}\right),
\]
η οποία αναφέρεται ως \emph{λογαριθμική σημειογραφική ανάπτυξη ανά συλλαβή}. Το $Q^{\mathrm{txt}}$ αφορά κειμενικό χωρίο και δεν επαναορίζει το $Q_x$, το οποίο αφορά πραγμάτωση μουσικής θέσης. Η $\eta^{\mathrm{txt}}$ υπολογίζεται μόνο με θετικό $Q^{\mathrm{txt}}$, θετικό $n^{\mathrm{syl}}_{\ell}$ και επαρκώς τεκμηριωμένη αντιστοίχιση γεγονότων. Οι μελοποιητικές επαναλήψεις παραμένουν στον αριθμητή, ενώ ο παρονομαστής διατηρεί τις συλλαβές του σταθερού κειμένου. Παύσεις ή χρονικά γεγονότα στα όρια χωρίων ακολουθούν κανόνα απόδοσης που έχει οριστεί πριν από την εξέταση των αποτελεσμάτων. Αν αναγνωρίζεται μόνο μέρος του χωρίου, το άθροισμα δηλώνεται ως μερική ανάπτυξη και δεν παρουσιάζεται ως συνολική τιμή. Αναφέρονται μαζί η ακατέργαστη ανάπτυξη, ο αριθμός γεγονότων, η κανονικοποιημένη ποσότητα και η πραγματική κάλυψη.

Η μονάδα ανάλυσης είναι το \textbf{έργο $\times$ μοναδικό κειμενικό χωρίο}. Οι θέσεις ενός έργου δεν αντιμετωπίζονται ως ανεξάρτητες παρατηρήσεις. Το πρώτο ερώτημα εξετάζει αν τα ίδια προκαθορισμένα χωρία παρουσιάζουν κοινές τάσεις ή ετερογένεια μεταξύ ανεξάρτητων μελοποιών και αναφέρεται αρχικά με πίνακες, γραφήματα και ενδοσυνθετικές αντιθέσεις. Σύνοψη ανά κειμενική κατηγορία επιχειρείται μόνο όταν υπάρχουν διαφορετικά χωρία ανά κατηγορία, ουσιαστική κοινή κάλυψη των δομικών και λειτουργικών συνθηκών και κατάλληλος πίνακας σχεδιασμού. Η κατηγορία είναι ιδιότητα του χωρίου: σταθερή επίδραση για κάθε ακριβές χωρίο απορροφά την κύρια επίδραση κατηγορίας, ενώ τυχαία επίδραση χωρίου εισάγει πρόσθετες παραδοχές και δεν θεραπεύει αυτομάτως την ταύτιση κατηγορίας, θέσης και λειτουργικού τμήματος. Όταν οι όροι δεν διαχωρίζονται, το αποτέλεσμα περιορίζεται στα συγκεκριμένα χωρία.

Το δεύτερο ερώτημα αφορά τη διαφοροποίηση του Νικολάου. Η σύγκριση αντιστοιχίζει ακριβώς τα ίδια χωρία και εξετάζει τη διαφορά της δικής του ενδοσυνθετικής αντίθεσης από τον ισοβαρή μέσο των αντίστοιχων αντιθέσεων των άλλων επιλέξιμων μελοποιών. Η κατανομή αναφοράς σε επίπεδο μελοποιού ακολουθεί τους υπάρχοντες αποκλεισμούς της M5 και το $K_{\min}$. Ο διαδοχικός αποκλεισμός κάθε μελοποιού από το σώμα αναφοράς σχηματίζει κατανομή αναφοράς και δεν ταυτίζεται με δοκιμασία αναδιατάξεων. Αν κατηγορία και λειτουργικό τμήμα συμπίπτουν, η ερμηνεία αφορά τη διαφορετική μεταχείριση των συγκεκριμένων χωρίων ή τμημάτων και όχι απομονωμένη σημασιολογική επίδραση.

Τα κρατήματα\footnote{Κράτημα: μη λεξικό, εκτενώς μελισματικό τμήμα της ψαλτικής που αναπτύσσεται με άσημες συλλαβές και όχι με κανονικό λεκτικό κείμενο. Βλ. Αναστασίου 2005.} διατηρούνται ως χωριστό πλαίσιο περιγραφής μη λεξικού υλικού. Για κάθε έργο και πηγή καταγράφονται χωριστά η παρουσία του κρατήματος, η πληρότητα ή έκταση του παραδιδόμενου τμήματος, τυχόν τεκμηριωμένη σύντμηση ή αντικατάσταση και η απόδοση σε μελοποιό ή μεταγενέστερη παρέμβαση. Η απουσία στην πηγή διακρίνεται από τεκμηριωμένη παράλειψη και η φθορά μάρτυρα από σύντμηση. Η απόδοση των σχετικών διαφορών στον μελοποιό επιτρέπεται μόνο στο εύρος που στηρίζεται από την προέλευση του συγκεκριμένου τμήματος. Το Αλληλούια παραμένει λεκτικό υλικό και καταγράφεται χωριστά ως λειτουργική επευφημία. Οι περιγραφές κρατημάτων μπορούν να στρωματώνονται ως προς ήχο, λειτουργικό τόπο, έκταση και μελισματική ανάπτυξη όπου υπάρχει επαρκής κοινή κάλυψη, χωρίς να μεταβάλλουν το κύριο πεδίο των υπαρχόντων εκτιμώμενων μεγεθών M2/M4/M5.

Η ανάλυση παραμένει προκαθορισμένη δευτερεύουσα μουσικολογική εργασία. Δεν προσθέτει νέα επιβεβαιωτική υπόθεση ή νέα οικογένεια διόρθωσης πολλαπλών ελέγχων. Οι μουσικές θέσεις δεν επαναοριοθετούνται με βάση τις νέες κειμενικές ετικέτες, οι UNKNOWN περιοχές διατηρούν τους υπάρχοντες κανόνες και όλες οι συγκρίσεις αναφέρουν το πραγματικό κοινό πεδίο και την αβεβαιότητά του. Οι νέες κειμενικές ετικέτες αποτελούν αποκλειστικά μεταδεδομένα της δευτερεύουσας ανάλυσης: δεν προστίθενται στο στρώμα $L$, δεν χρησιμοποιούνται ως είσοδοι ή εποπτεία στις M1/M2 και δεν επηρεάζουν την επιλογή των ερεθισμάτων της M3.


\subsection{Μελλοντικές κατευθύνσεις}

Τρεις προοπτικές παραμένουν ρητά εκτός του δεσμευτικού πρωτοκόλλου: (α) επέκταση της γνωσιακής μελέτης σε πληθυσμούς με διαφορετικά γνωσιακά ή νευροαναπτυξιακά χαρακτηριστικά, με αυτοτελές πρωτόκολλο δειγματοληψίας, δεοντολογίας και ισχύος· (β) πολυτροπική μελέτη φωνής και κίνησης σε νέα συλλογή εκτελέσεων, εφόσον προκύψει συνεργατική υποδομή· και (γ) φυσικά πληροφορημένη μοντελοποίηση φωνοπαραγωγής, μόνο μετά από ανεξάρτητη επικύρωση των απαιτούμενων φυσιολογικών μεταβλητών. Καμία από τις τρεις δεν απαιτείται για τα M1--M5 ή για τα επιβεβαιωτικά συμπεράσματα.


\section{Ανθρώπινη γνωσιακή επικύρωση και ιεραρχική μουσική πρόβλεψη}\label{app:human}

\subsection{Συμμετέχοντες και κοινή στρατολόγηση}\label{app:human-participants}

Η M3b και η M3c χρησιμοποιούν κοινό σύνολο διαθέσιμων συμμετεχόντων και κοινή δεοντολογική διαδικασία. Η βασική σύγκριση περιλαμβάνει ψάλτες και άτομα με συστηματική εκπαίδευση στη βυζαντινή μουσική και μουσικούς χωρίς βυζαντινή εκπαίδευση. Καταγράφονται ως συνεχείς μεταβλητές η γενική μουσική εμπειρία, η διάρκεια και ένταση βυζαντινής εκπαίδευσης, η ακουστική εξοικείωση με ψαλτικό ρεπερτόριο, η λειτουργική έκθεση και η εξοικείωση με τα συγκεκριμένα έργα. Η κύρια ερμηνεία της σύγκρισης ομάδων αφορά την εμπειρία στη συγκεκριμένη μουσική γλώσσα, ενώ οι συνεχείς δείκτες επιτρέπουν λεπτομερέστερη ανάλυση όταν υπάρχει επαρκής ανεξάρτητη διακύμανση.

Η συλλογή γίνεται σε εργαστήριο ή σε ελεγχόμενο διαδικτυακό περιβάλλον, με τρόπο που οριστικοποιείται μετά τον πιλοτικό έλεγχο της ακουστικής αναπαραγωγής, της προσήλωσης και της αξιοπιστίας των αποκρίσεων. Το τελικό μέγεθος δείγματος προκύπτει από προσομοίωση στατιστικής ισχύος για τα πραγματικά μικτά μοντέλα και τη δομή των ερεθισμάτων. Ο πιλοτικός έλεγχος εκτιμά τη διακύμανση των αποκρίσεων, την αξιοπιστία, τη διάρκεια των δοκιμών, την κόπωση και την κατανομή των κρίσιμων μεταβλητών πρόβλεψης.

\subsection{Κοινή τράπεζα ερεθισμάτων}\label{app:human-stimuli}

Η τράπεζα ερεθισμάτων κατασκευάζεται μετά την οριστικοποίηση του λεξιλογίου οικογενειών. Κάθε ερέθισμα διαθέτει αναγνωριστικό με ελεγχόμενη έκδοση και συνδέεται με οικογένεια, έργο, πηγή, εκτελεστή ή συνθετική απόδοση, ήχο, δομικό, κειμενικό και προσωδιακό περιβάλλον, διάρκεια, συμβολικές και ακουστικές αποστάσεις, $I_F$ και $I_{\mathrm{event}}$. Όπου υπάρχει φυσική αντιστοίχιση, διατηρούνται επίσης το αναγνωριστικό του προκαθορισμένου κειμενικού χωρίου και οι ανεξάρτητα ορισμένες κειμενικές ετικέτες ως μεταδεδομένα. Τα τελικά μουσικά χαρακτηριστικά του χωρίου μπορούν να προστεθούν αργότερα από οριστικοποιημένες αναλύσεις της M5 με κανόνα σύνοψης που έχει οριστεί εκ των προτέρων και χωρίς εξέταση των βιωματικών αποτελεσμάτων. Τα μεταδεδομένα αυτά δεν αλλάζουν τα κριτήρια επιλογής των ερεθισμάτων, την κάλυψη του $I_F\times I_{\mathrm{event}}$ ή την αντιστάθμιση του $R$. Η φυσικότητα των ερεθισμάτων αξιολογείται πριν από την κύρια συλλογή. Όπου χρησιμοποιούνται ελεγχόμενες συνθετικές αποδόσεις, η ίδια διαδικασία παραγωγής εφαρμόζεται σε όλες τις συνθήκες και η γενίκευση σε φυσικές εκτελέσεις εξετάζεται ως χωριστή ανάλυση.

\subsection{Λειτουργική οργάνωση του συμβολικού σώματος}\label{app:human-functional}

\textbf{Προκαθορισμένη δευτερεύουσα ανάλυση λειτουργικού τόπου.} Με αφετηρία την τεκμηριωμένη μορφολογική παρατήρηση για τη διαφοροποίηση της οργάνωσης πριν και μετά τη Μεγάλη Είσοδο (Καραγκούνης, 2005, σσ.~158--159), εξετάζεται στο συμβολικό σώμα αν η επανάληψη και η προβλεψιμότητα ακολουθιών οικογενειών διαφοροποιούνται μεταξύ τεκμηριωμένων τμημάτων πλήρων Χερουβικών πριν και μετά τη λειτουργική τομή της Μεγάλης Εισόδου\footnote{Μεγάλη Είσοδος: κεντρική τελετουργική στιγμή της Θείας Λειτουργίας, κατά την οποία μεταφέρονται τα Τίμια Δώρα. Στην ανάλυση χρησιμοποιείται ως προκαθορισμένο λατρευτικό/μορφολογικό σημείο αναφοράς και όχι ως μουσική μεταβλητή.}. Η σύγκριση είναι εντός έργου και αναφέρει χωριστά (α) δείκτη επανάληψης/επανεμφάνισης ακολουθιών και (β) log-loss σε υλικό εκτός εκπαίδευσης του ίδιου μοντέλου πρόβλεψης οικογενειών. Ελέγχονται μήκος τμήματος, αριθμός θέσεων, κειμενικός τόπος, εγγύτητα και πυκνότητα καταλήξεων, ήχος, μελοποιός και ιστορικό στρώμα. Η ανάλυση αφορά τη στατιστική οργάνωση του ρεπερτορίου και παραμένει διακριτή από το πειραματικό $R$, το οποίο μεταβάλλει μόνο τη μη μουσική ένδειξη λατρευτικού συμφραζομένου στον ανθρώπινο έλεγχο.

\subsection{M3c — Συνεδρία 1: μοντέλο πρόβλεψης σε επίπεδο οικογένειας}\label{app:human-m3c-model}

Το κύριο μοντέλο πρόβλεψης οικογενειών είναι PPM μεταβλητής τάξης με αποκλειστικά αριστερό συμφραζόμενο:
\[
P(F_t\mid F_{<t},M,L).
\]
Το $M$ δηλώνει την τροπική κατάσταση και το $L$ το δομικό, κειμενικό και προσωδιακό περιβάλλον. Οι είσοδοι $M,L$ κωδικοποιούν αντίστοιχα την τροπική κατάσταση και το δομικό, κειμενικό και προσωδιακό περιβάλλον του μουσικού υλικού και παραμένουν σταθερές για τον ίδιο μουσικό στόχο και το ίδιο πρόθεμα σε όλες τις λίστες. Το $R$ κωδικοποιεί χωριστά τη μη μουσική πειραματική ένδειξη στο ανθρώπινο στατιστικό μοντέλο. Το PPM συγκρίνεται με τη συνολική κατανομή $P(F_t)$ χωρίς συμφραζόμενο και με $P(F_t\mid M,L)$. Το κύριο κριτήριο επάρκειας του μοντέλου απαιτεί βελτίωση του log-loss με διασταυρωμένη επικύρωση έναντι του συγκριτή $P(F_t\mid M,L)$, ο οποίος χρησιμοποιεί τροπικό και δομικό συμφραζόμενο χωρίς την προηγούμενη διαδοχή οικογενειών σε έργα εκτός εκπαίδευσης. Ένας αιτιακός Transformer παρέχει συμπληρωματικό έλεγχο ανθεκτικότητας. Η αμφίδρομη συμπλήρωση χρησιμοποιείται μόνο ως βοηθητικός στόχος αναπαράστασης.

Το $I_{\mathrm{event}}$ υπολογίζεται από τον Transformer σε επίπεδο μελωδικού γεγονότος της M1 και συνοψίζεται με το μέσο πληροφοριακό περιεχόμενο της συνέχειας. Αθροιστικές και μέγιστες συνοπτικές ποσότητες αναφέρονται σε αναλύσεις ευαισθησίας.

\subsection{M3c — Συνεδρία 1: ανθρώπινη δοκιμασία και τετραμερής σχεδιασμός}\label{app:human-m3c-task}

Σε κάθε δοκιμή ο συμμετέχων ακούει το συμφραζόμενο και μία συνέχεια και αξιολογεί πόσο απρόσμενη του φαίνεται η συνέχεια σε προκαθορισμένη κλίμακα. Τα ερεθίσματα καλύπτουν τις τέσσερις περιοχές του $I_F\times I_{\mathrm{event}}$ και εξασφαλίζουν επαρκές πλήθος δοκιμών στους δύο αντιδιαγώνιους συνδυασμούς. Η κύρια επιβεβαιωτική παράμετρος είναι $\beta_F^{\mathrm{pred}}$ στο μοντέλο που προβλέπει πόσο απρόσμενη φαίνεται η συνέχεια.

Η σειρά των δοκιμών τυχαιοποιείται μέσα σε περιορισμούς ισορροπίας. Οι δοκιμές εξάσκησης χρησιμοποιούν διαφορετικό υλικό. Η εξοικείωση με το συγκεκριμένο απόσπασμα καταγράφεται μετά την κρίση ή σε χωριστή ομάδα δοκιμών ώστε να μην κατευθύνει την κρίση για το πόσο απρόσμενη είναι η συνέχεια.

\subsection{Ένδειξη λατρευτικού συμφραζομένου και αντιστάθμιση}\label{app:human-framing}

Η συνθήκη χρησιμοποιεί μη μουσική ένδειξη λατρευτικού συμφραζομένου. Η ένδειξη λατρευτικού πλαισίου μπορεί να είναι κειμενική, οπτική ή μη μελωδικά εκφωνούμενη και αντιστοιχίζεται με ουδέτερη ένδειξη ως προς διάρκεια και απαίτηση προσοχής. Το άμεσο μουσικό πρόθεμα και το ηχητικό απόσπασμα-στόχος είναι τα ίδια στις δύο εκδοχές της συνθήκης.

Κάθε μουσική συνέχεια παρουσιάζεται μία φορά ανά συμμετέχοντα. Οι πειραματικές λίστες κατανέμουν τις μουσικές συνέχειες στις δύο τιμές του $R$ με αντισταθμισμένο σχήμα λατινικού τετραγώνου. Η κύρια δευτερεύουσα ανάλυση συγκρίνει $R\times z(I_F)$ και $R\times z(I_{\mathrm{event}})$ μέσω του $\Delta_{\mathrm{ctx}}$. Ένα γενικό κύριο αποτέλεσμα του $R$ αναφέρεται χωριστά από τη σχετική στάθμιση των μεταβλητών πρόβλεψης των δύο επιπέδων.

\subsection{Βιωματική απόκριση}\label{app:human-affect}

Μετά από προκαθορισμένες ομάδες δοκιμών συλλέγονται σύντομες αξιολογήσεις βιωμένης απόκρισης. Οι τελικές κλίμακες επιλέγονται πριν από την κύρια συλλογή με κριτήριο εγκυρότητα, γλωσσική διαθεσιμότητα και χαμηλή επιβάρυνση των συμμετεχόντων. Η ανάλυση συσχετίζει ευαισθησία στο $I_F$, το $R$ και τις μεταβλητές εμπειρίας με τις βιωμένες διαστάσεις. Όπου οι ομάδες δοκιμών παρουσιάζουν πραγματική διακύμανση ως προς τα προκαθορισμένα κειμενικά χωρία και υπάρχουν διαθέσιμα συνοψισμένα μουσικά χαρακτηριστικά, εξετάζεται διερευνητικά η σχέση τους με τη βιωματική απόκριση στη μονάδα συμμετέχων επί ομάδα δοκιμών. Η ίδια αξιολόγηση ομάδας δεν επαναλαμβάνεται ως ανεξάρτητη απόκριση για κάθε ερέθισμα, και η κειμενική προέλευση ενός συνθετικού ή κατασκευασμένου ερεθίσματος δεν θεωρείται απόδειξη ότι ο συμμετέχων επεξεργάστηκε το νόημα του κειμένου. Η ανάλυση έχει διερευνητικό καθεστώς, εκτιμά συσχετίσεις και δεν ερμηνεύεται ως επίδραση της μουσικής έμφασης στο συναίσθημα.

\subsection{M3b — Συνεδρία 2: τριαδική επιλογή και κοινό πεδίο τιμών}\label{app:human-m3b-task}

Η M3b πραγματοποιείται μετά τη M3c. Οι μεικτές τριάδες επιλέγονται από την κοινή περιοχή συμβολικών και ακουστικών αποστάσεων, στην οποία εκπροσωπούνται και οι δύο τύποι σχέσεων σύμφωνα με τα προκαθορισμένα κριτήρια αντιστοίχισης, και οι ενδοοικογενειακές τριάδες κατασκευάζονται ώστε να καλύπτουν ουσιαστικό εύρος ενδοοικογενειακής παραλλαγής. Η απάντηση επιλογής του διαφορετικού μοντελοποιείται ως μία τριαδική επιλογή. Η πιθανοφάνεια διατηρεί την τριάδα ως μονάδα επιλογής.

\subsection{Ανθρώπινη αναπαράσταση, CCGP και ενδοοικογενειακή σύγκριση αναπαραστάσεων}\label{app:human-geometry}

Οι τριαδικές κρίσεις επιτρέπουν να εκτιμηθεί ένας χώρος ανθρώπινων αναπαραστάσεων από τις ίδιες τις επιλογές των ακροατών. Μια \textbf{μήτρα ανομοιότητας αναπαραστάσεων (Representational Dissimilarity Matrix, RDM)} καταγράφει τη σχετική απόσταση κάθε ερεθίσματος από τα υπόλοιπα. Η σύγκριση τέτοιων μητρών εξετάζει αν δύο αναπαραστάσεις οργανώνουν με παρόμοιο τρόπο τις σχετικές διαφορές μεταξύ των ίδιων ερεθισμάτων. Το \textbf{Cross-Condition Generalization Performance (CCGP)} ελέγχει διαφορετική ιδιότητα: αν μία οικογενειακή διάκριση που μαθαίνεται σε ορισμένες προκαθορισμένες συνθήκες μπορεί να μεταφερθεί σε άλλες συνθήκες που δεν χρησιμοποιήθηκαν για να μάθει ο αποκωδικοποιητής τη συγκεκριμένη διάκριση.

Ο χώρος των ανθρώπινων αναπαραστάσεων προσαρμόζεται από τις τριαδικές επιλογές χωρίς ετικέτες οικογένειας. Οι υπερπαράμετροι και η διάσταση του χώρου επιλέγονται με βάση την πρόβλεψη τριάδων σε υλικό ανάπτυξης που έχει κρατηθεί εκτός της προσαρμογής του μοντέλου. Η αξιοπιστία εκτιμάται με διαδικασίες διαίρεσης δείγματος και συμφωνίας μεταξύ κριτών κατάλληλες για τον σχεδιασμό και χρησιμοποιείται για την εκτίμηση του ανώτατου ορίου αξιοπιστίας.

Η ανάλυση ομοιότητας αναπαραστάσεων (Representational Similarity Analysis, RSA) συγκρίνει τις μήτρες ανομοιότητας του ανθρώπινου και του υπολογιστικού χώρου. Εδώ εφαρμόζεται χωριστά σε κάθε οικογένεια και ελέγχει τη συμφωνία τους αφού ληφθούν υπόψη οι συμβολικές και ακουστικές αποστάσεις. Το CCGP εκπαιδεύει αποκωδικοποιητή σε μία προκαθορισμένη διαμέριση συνθηκών και τον αξιολογεί σε διακριτές συνθήκες. Το ενδοοικογενειακό RSA υπολογίζει RDM ανά οικογένεια και συγκρίνει την υπολειμματική ευθυγράμμιση μοντέλου--ανθρώπου για EISE και B2 μετά τον έλεγχο των συμβολικών και ακουστικών RDM. Η τελική εξαγωγή συμπερασμάτων χρησιμοποιεί ιεραρχική συνάθροιση σε επίπεδο οικογένειας.

\subsection{Πολλαπλότητα, προεγγραφή και κριτήρια ενεργοποίησης}\label{app:human-gates}

Η οικογένεια επιβεβαιωτικών ανθρώπινων υποθέσεων περιέχει $H_{\mathrm{cat}}$ και $H_{\mathrm{pred}}$ και εφαρμόζει Holm. Το $\Delta_{\mathrm{ctx}}$, το CCGP, η ενδοοικογενειακή ευθυγράμμιση και οι αλληλεπιδράσεις εμπειρίας έχουν καθεστώς κύριων δευτερευόντων αποτελεσμάτων. Η κύρια συλλογή M3b προχωρά μετά την ικανοποίηση των κριτηρίων για κοινό πεδίο τιμών και επαρκή αξιοπιστία. Η M3c προχωρά μετά την οριστικοποίηση του λεξιλογίου οικογενειών, το κριτήριο επάρκειας του μοντέλου PPM, το κριτήριο επαρκούς κάλυψης των αντιδιαγώνιων συνδυασμών και τον τυφλό έλεγχο μουσικής φυσικότητας.

Όλα τα κατώφλια, κύριες συνοπτικές ποσότητες, κανόνες αποκλεισμού, τύποι μοντέλων και στόχοι στατιστικής ισχύος καταγράφονται πριν από την κύρια συλλογή. Αν ένα κριτήριο ενεργοποίησης δεν ικανοποιείται, δηλώνεται ότι το αντίστοιχο εκτιμώμενο μέγεθος δεν μπορεί να εκτιμηθεί αξιόπιστα. Η μη ικανοποίηση του κριτηρίου δεν οδηγεί σε εκ των υστέρων αλλαγή του κεντρικού ερωτήματος.

\subsection{M3ν — νευρική επέκταση}\label{app:human-neural}

Η νευρική επέκταση χρησιμοποιεί την ίδια τράπεζα ερεθισμάτων και τις ίδιες μεταβλητές πρόβλεψης. Το μοντέλο νευρικής κωδικοποίησης με χρονική ανάλυση συγκρίνει ακουστικούς δείκτες και δείκτες γεγονότων με μοντέλο που προσθέτει $I_F$, και το $R\times I_F$ εξετάζεται ως χωριστός όρος ευαίσθητος στο συμφραζόμενο. Επειδή κάθε συνέχεια αρχίζει σε όριο, το μοντέλο περιλαμβάνει προκαθορισμένες συμμεταβλητές για ακουστικά και δομικά χαρακτηριστικά του ορίου. Για το $I_{\mathrm{event}}$ εξετάζονται προκαθορισμένα μήκη άμεσου συμφραζομένου, με επιλογή αποκλειστικά σε δεδομένα ανάπτυξης, και αναφέρεται δεύτερο ευρύτερο χρονικό παράθυρο ως ανάλυση ευαισθησίας. Η ακριβής μέθοδος χρονικής νευρικής ανάλυσης οριστικοποιείται μετά από πιλοτικό έλεγχο και πριν από την κύρια συλλογή· η προσομοίωση ισχύος χρησιμοποιεί τον αναμενόμενο αριθμό αξιοποιήσιμων δοκιμών μετά τον καθαρισμό. Η M3ν ενεργοποιείται μόνο μετά από γραπτή επιβεβαίωση συνεργαζόμενου εργαστηρίου για διαθέσιμη EEG/MEG υποδομή, κατάλληλη εμπειρογνωμοσύνη ανάλυσης και ξεχωριστή δεοντολογική έγκριση. Η μη ενεργοποίησή της δεν επηρεάζει την ολοκλήρωση ή την επιβεβαιωτική ιεραρχία της έρευνας.

\section{Υλοποίηση και αναπαραγωγιμότητα}\label{app:impl}

Η αναπαραγωγιμότητα της υπολογιστικής αλυσίδας απαιτεί να μπορεί κάθε αποτέλεσμα να συνδεθεί με τα δεδομένα, τους κανόνες και την έκδοση του συστήματος που το παρήγαγαν. Για τον σκοπό αυτό καταγράφονται οι εκδόσεις των πηγών και των αναλυτών, τα αναγνωριστικά γεγονότων και χωρίων, οι κανόνες μετατροπής, οι διαχωρισμοί δεδομένων, οι οριστικοποιημένες ρυθμίσεις και οι τυχαιοποιήσεις όπου χρησιμοποιούνται. Οι δοκιμές ανακατασκευής ελέγχουν το ενιαίο συμβολικό επίπεδο αναφοράς και η διατήρηση των αναγνωριστικών επιτρέπει στις προβλέψεις, στις αποστάσεις και στα συνοπτικά μεγέθη να επιστρέφουν στα συγκεκριμένα χωρία των πηγών.

\subsection{Ευθυγράμμιση παρτιτούρας--εκτέλεσης}\label{app:impl-align}

Το σώμα διαθέτει \textbf{χειροκίνητα τοποθετημένα όρια συλλαβών}. Στο παπαδικό μέλος όμως μία εκτεταμένη συλλαβή περιέχει πολλαπλά μελωδικά γεγονότα. Επομένως:
Η ευθυγράμμιση στο επίπεδο μουσικού γεγονότος εξειδικεύει τη συλλαβική αντιστοίχιση με το ακόλουθο πρωτόκολλο.
Η μετάβαση από το ένα στο άλλο αποτελεί \textbf{το σημαντικότερο τεχνικό ρίσκο της M4} και ορίζεται ρητά.

\textbf{Πρωτόκολλο.} Τα χειροκίνητα όρια συλλαβών λειτουργούν ως \emph{σημεία αναφοράς}. Εντός κάθε συλλαβής, τα σημειογραφικά γεγονότα ευθυγραμμίζονται στον χρόνο με αλγόριθμο από προκαθορισμένη οικογένεια υποψηφίων: δυναμική στρέβλωση χρόνου με περιορισμούς μονοτονίας και τοπικού εύρους, ή αποκωδικοποίηση σε μοντέλο καταστάσεων καθοδηγούμενο από την παρτιτούρα. Η πιθανοφάνεια υπολογίζεται από την τροχιά θεμελιώδους συχνότητας, με το \textbf{ηλεκτρογλωττογράφημα} ως πρωτεύουσα πηγή επειδή δίνει ανθεκτικότερη εκτίμηση από τον ήχο μόνο του. Η επιλογή αλγορίθμου οριστικοποιείται στο πρώτο πακέτο εργασίας με βάση την επίδοση στα δοκιμαστικά έργα.

\textbf{Μεταβλητή ελέγχου ποιότητας.} Σφάλμα χρονισμού σε σημεία αναφοράς που έχουν επισημανθεί χειροκίνητα και έχουν διατηρηθεί εκτός της διαδικασίας ευθυγράμμισης, με αναφορά διαμέσου και ενενηκοστού πέμπτου εκατοστημορίου. Ένα υποσύνολο σημείων αναφοράς διατηρείται εκτός της ευθυγράμμισης για ανεξάρτητο έλεγχο ποιότητας.

\textbf{Κριτήριο ποιότητας και προσαρμογή της ανάλυσης στην ακρίβεια ευθυγράμμισης.} Ορίζεται ελάχιστο αποδεκτό επίπεδο ακρίβειας στο πρώτο πακέτο εργασίας. Η μετρημένη επίδοση της ευθυγράμμισης καθορίζει εκ των προτέρων τη χρονική κλίμακα της ανάλυσης: γεγονός όταν η ακρίβεια το επιτρέπει και συλλαβή ή νευματική ομάδα όταν απαιτείται μεγαλύτερη χρονική ανθεκτικότητα. Έτσι οι στατιστικές υποθέσεις αντιστοιχούν στην πραγματική μετρητική διακριτική ικανότητα των δεδομένων.

Κάθε ηχογράφηση φέρει επιπλέον \textbf{σημαία ποιότητας} που δηλώνει αν το τονικό ύψος είναι αξιόπιστα εξαγώγιμο. Ηχογραφήσεις με υποβαθμισμένη σημαία αποκλείονται από τις μετρικές που εξαρτώνται από το ύψος αλλά συμμετέχουν σε μετρήσεις χρονισμού και έκτασης.

\subsection{Μεταβλητές εκτελεστικής απόκλισης}\label{app:impl-vars}

Για κάθε εκτελεστή \(p\) και σημειογραφικό γεγονός \(e\) υπολογίζεται απόκλιση τονικού ύψους σε λεπτά της οκτάβας ως προς το αναμενόμενο:
\[
d_{pe}=1200\,\log_2\!\frac{F_{0,pe}}{F^{\text{αναμ}}_{0,e}}
\]
με κανονικοποίηση ως προς την τοπική βάση και τη φωνητική περιοχή του εκτελεστή. Επειδή το αναμενόμενο ύψος \(F^{\text{αναμ}}_{0,e}\) προϋποθέτει συγκεκριμένη παραμετροποίηση των βυζαντινών κλιμάκων, αναφέρεται \textbf{προκαθορισμένη ανάλυση ευαισθησίας} σε εναλλακτικές τεκμηριωμένες διαστηματικές παραμετροποιήσεις· ευρήματα που αντιστρέφονται μεταξύ τους δηλώνονται ως εξαρτώμενα από την παραδοχή κλίμακας. Παράλληλα υπολογίζονται κανονικοποιημένη διάρκεια, σχετικός χρονισμός, λόγος μελισματικής ανάπτυξης και μέτρα συμπεριφοράς κατάληξης.

Από τους πολλαπλούς εκτελεστές προκύπτει, για κάθε σημείο της παρτιτούρας, \textbf{προφίλ μεταβλητότητας μεταξύ εκτελεστών}. Το προφίλ αυτό αντιμετωπίζεται ως \textbf{δομημένος χώρος εκτελεστικών πραγματώσεων} της γραπτής μουσικής δομής: η διασπορά αποτελεί μεταβλητή προς εξήγηση και όχι υπόλοιπο που αφαιρείται από την ανάλυση. Η ερμηνεία αυτή συνδέει τη M4 με τη σύγχρονη υπολογιστική εθνομουσικολογία (Tzanetakis, 2014) και με τη μελέτη της μουσικής γνώσης όπως πραγματώνεται στην εκτελεστική πράξη. Συναφείς υπολογιστικές μελέτες φωνητικής εκτέλεσης σε προφορικές παραδόσεις δείχνουν επίσης την αξία της ερμηνείας ακουστικών χαρακτηριστικών σε συνάρτηση με τη γνώση της πρακτικής (Paschalidou, 2024). Η παρούσα έρευνα διατηρεί εδώ τις ήδη προβλεπόμενες ακουστικές και εκτελεστικές μετρήσεις της M4. Εφόσον το ΠΕ0 επιβεβαιώσει γραπτή πρόσβαση, τεκμηριωμένο συγχρονισμό και επαρκή ποιότητα του ηλεκτρογλωττογραφικού σήματος του DAMASKINOS, εξετάζεται επιπλέον διερευνητικά αν η διακύμανση ηλεκτρογλωττογραφικών δεικτών (πηλίκο επαφής και χαρακτηριστικά της παραγώγου του EGG· Herbst, 2020) παρουσιάζει συστηματική οργάνωση σε σχέση με τα όρια και την εσωτερική συντεταγμένη των θέσεων, με έλεγχο του σχετικού τονικού ύψους ως προς τη φωνητική περιοχή κάθε ψάλτη, της έντασης, του φωνήεντος, του εκτελεστή και του έργου. Η ερμηνεία περιορίζεται στα χαρακτηριστικά φωνητικής επαφής που τεκμηριώνει η επιλεγμένη μέθοδος και δεν αφορά σωματική προσπάθεια. Αποτυχία οποιουδήποτε κριτηρίου πρόσβασης ή ποιότητας οδηγεί σε μη ενεργοποίηση της ανάλυσης και δεν επηρεάζει κανένα επιβεβαιωτικό αποτέλεσμα.

\subsection{Οριστικοποιημένες ρυθμίσεις και πρωτόκολλο}\label{app:impl-config}

Στο τέλος του ΠΕ0 χρονοσημαίνεται η εφαρμογή των προκαθορισμένων κανόνων στα ευρήματα της απογραφής: έκδοση μητρώου και αναλυτών, σύνολο εκπαίδευσης, ορισμοί πεδίων, οικογενειακές σχέσεις, επαληθευμένες ακμές, διαχωρισμοί, κατώφλια, μετρική, σύνολα υποψηφίων, μορφή δομημένης εξόδου και επίπεδο εκτελεστικής ανάλυσης. Η \S\ref{sec:pilot-decision} διακρίνει τους κανόνες που έχουν οριστεί πριν από τον πιλοτικό έλεγχο από τις συγκεκριμένες επιλογές που προκύπτουν από την εφαρμογή τους. Οι υπερπαράμετροι μοντέλων ρυθμίζονται στη συνέχεια αποκλειστικά στα επιλέξιμα δεδομένα εκπαίδευσης και ανάπτυξης με τον κοινό υπολογιστικό προϋπολογισμό και οριστικοποιούνται πριν από κάθε τελική αξιολόγηση. Οι επιβεβαιωτικοί στόχοι διατηρούν την ανεξαρτησία τους σε όλη τη διαδικασία.

\subsection{Πακέτα εργασίας, πλήρες χρονοδιάγραμμα και παραδοτέα}\label{app:impl-wp}

\textbf{ΠΕ0 --- Απογραφή, υπογραφές τύπων, προέλευση, δικαιώματα και ερευνητική ετοιμότητα (μήνες 1--6).} Κατασκευή μητρώων έργων, πηγών, εμφανίσεων και ιστορικών σχέσεων. Συγκρότηση μικρού πιλοτικού λεξικού με πραγματικά θετικά και δύσκολα αρνητικά παραδείγματα. Ανάπτυξη των υπογραφών μόνο στο development material, πάγωμα έκδοσης και ανεξάρτητος έλεγχος αξιοπιστίας πριν από την παραγωγή κύριων B2/EISE pairs. Καταγραφή του \texttt{evidence\_role}, ώστε υλικό ορισμού να μη χρησιμοποιείται ως ανεξάρτητη επικύρωση. Απογραφή του \texttt{exegesis\_observability} και της κάλυψης $C_{\mathrm{ex}}(F)$ ανά οικογένεια. Έλεγχος της καρδιναλιότητας παλαιότερης--Νέας αντιστοίχισης πριν από τη δημιουργία ζευγών. Σύνταξη του εγχειριδίου επισημείωσης και συγκρότηση του πλήρως ελεγμένου υποσυνόλου ανάπτυξης. Έλεγχος ότι τα απαιτούμενα παρατηρήσιμα γνωρίσματα μπορούν να αναπαρασταθούν στα υπάρχοντα $N/T/M/L$ ή στα μεταδεδομένα, χωρίς νέα αρχιτεκτονική. Εφαρμογή του προκαθορισμένου κανόνα επιλογής της μορφής εξόδου πριν από τη M1. Οριστικοποίηση της ανοχής ορίων $\tau_b$ και γραπτή, χρονοσφραγισμένη δήλωση των $g_{\mathrm{use}}$ και $\delta_{\mathrm{use}}$ πριν από οποιαδήποτε επίδοση μοντέλου (§~\ref{app:stats-selection}). Προκαθορισμός του ιστορικού συνόλου υψηλής δυσκολίας, καταγραφή δικαιωμάτων και έλεγχος πρόσβασης στο ακατέργαστο ηχητικό και ηλεκτρογλωττογραφικό υλικό του DAMASKINOS. Η έκταση της εκτελεστικής ανάλυσης οριστικοποιείται μετά την απογραφή. Παράλληλα αρχίζει η δεοντολογική διαδικασία για M3b/M3c.

\textbf{ΠΕ1 --- Ενιαίο συμβολικό επίπεδο αναφοράς και πρώτη αλυσίδα EISE (μήνες 4--12).} Ορισμός ενιαίου επιπέδου αναφοράς για τα μελωδικά γεγονότα, αναγνωριστικά με καταγεγραμμένη προέλευση, προσαρμογείς και δοκιμές ανακατασκευής. Για τα λεκτικά γεγονότα ελέγχεται η ύπαρξη τεκμηριωμένης σύνδεσης με τη συλλαβή του σταθερού κειμένου ή ρητής άγνωστης τιμής, η δυνατότητα επιστροφής της γνωστής σύνδεσης στην πηγή, η συνέπεια των μελισματικών επαναλήψεων και η χωριστή αναφορά μη εφαρμοστέων και ελλιπών περιπτώσεων. Ακολουθούν η πρώτη λειτουργική αλυσίδα EISE σε υλικό ανάπτυξης, η υλοποίηση και πρώτη δοκιμή της μορφής εξόδου που έχει οριστικοποιηθεί στο ΠΕ0, καθώς και η οριστικοποίηση του κοινού πρωτοκόλλου της M1 και του λεξιλογίου ανάπτυξης· με το λεξιλόγιο αυτό αξιολογείται πιλοτικά το απλό οικογενειακό μοντέλο αναφοράς, υπολογίζεται μία φορά το $g_{\min}$ και κατόπιν οριστικοποιείται το πλήθος $n_{\mathrm{rep}}$ των αντιστοιχισμένων επαναλήψεων εκπαίδευσης, πριν από την κύρια εκπαίδευση της M1. Ο τελικός περιορισμένος έλεγχος της τμηματικής διαδικασίας πραγματοποιείται στη μετάβαση M1--M2, μετά την οριστικοποίηση της $R^\star$ και πριν από την πλήρη αντιστοιχισμένη εκπαίδευση. Αν έχει ενεργοποιηθεί η διαδρομή επικαλυπτόμενων τμημάτων, περιορίζονται πρώτα οι προαιρετικοί συγκριτές ώστε να διατηρηθεί ο ίδιος χρονικός προϋπολογισμός.

\textbf{ΠΕ2 / M1 --- Σύγκριση αναπαραστάσεων (μήνες 10--24).} $U\times F$, Transformer σε επίπεδο μελωδικού γεγονότος, επαγόμενη τμηματοποίηση, ανάκτηση και αναγνώριση ανοικτού συνόλου σε υλικό εκτός εκπαίδευσης, οριστικοποίηση της κύριας κλίμακας αναπαράστασης και κατόπιν μία ανεξάρτητη εφαρμογή του ιστορικού ελέγχου διακριτικής ικανότητας χωρίς ανατροφοδότηση προς την επιλογή.

\textbf{ΠΕ3 / M3a --- Μουσικολογική επικύρωση του λεξιλογίου οικογενειών (μήνες 12--24).} Το λεξιλόγιο ανάπτυξης που απαιτείται από τις M1--M2 έχει οριστικοποιηθεί νωρίτερα μόνο στο υλικό ανάπτυξης. Η M3a χρησιμοποιεί $C_{\mathrm{StA}}$, $C_{\mathrm{occ}}$ και ανεξάρτητο $C_{\mathrm{exp}}$ για εξωτερική μουσικολογική επικύρωση, αξιοπιστία, τροπική και ισοκρατηματική ανάλυση και δυνατότητα αξιόπιστης εκτίμησης του $A(z)$, χωρίς αναδρομική μεταβολή των αποτελεσμάτων M1--M2.

\textbf{M3c --- μοντελοποίηση πρόβλεψης (μήνες 20--28).} Πρωτότυπο PPM με προσωρινό λεξιλόγιο ανάπτυξης, τελικό PPM μετά την οριστικοποίηση του λεξιλογίου οικογενειών, σύγκριση με $P(F_t\mid M,L)$, έλεγχος ανθεκτικότητας με αιτιακό Transformer, τελικό $I_F$, $I_{\mathrm{event}}$ σε επίπεδο γεγονότος, τράπεζα υποψηφίων για τους αντιδιαγώνιους συνδυασμούς και τράπεζα ερεθισμάτων με ελεγχόμενες εκδόσεις.

\textbf{ΠΕ3β / M2 --- B2--EISE (μήνες 24--36).} Στην έναρξη, μετά την οριστικοποίηση της $R^\star$, πραγματοποιείται ο τελικός περιορισμένος έλεγχος της τμηματικής διαδικασίας μέσα στην ήδη επιλεγμένη μορφή εξόδου. Ακολουθούν αναδειγματοληψία αντιστοιχισμένων ζευγών, $\Delta_{\mathrm{hist}}$, μεταφορά σε υλικό εκτός εκπαίδευσης, αναλύσεις ανοικτού συνόλου, γεωμετρία τοπικής παραλλαγής και έλεγχοι ανθεκτικότητας.

\textbf{Πιλοτικοί έλεγχοι M3 --- μήνες 28--30.} Μουσική φυσικότητα, αξιοπιστία, κοινό πεδίο τιμών, κάλυψη των αντιδιαγώνιων συνδυασμών, διάρκεια εργασίας, διακύμανση αποκρίσεων και ισχύς βάσει προσομοίωσης. Η προεγγραφή οριστικοποιείται πριν από την κύρια συλλογή.

\textbf{Κύρια ανθρώπινη συλλογή M3 --- μήνες 30--40.} Συνεδρία 1: M3c και Συνεδρία 2: M3b στους ίδιους συμμετέχοντες. Κοινά μεταδεδομένα εμπειρίας, ελεγχόμενες εκδόσεις λιστών, τυφλή κωδικοποίηση ερεθισμάτων και προκαθορισμένα σενάρια ανάλυσης.

\textbf{M4 --- Εκτελεστική και ιστορική οργάνωση (μήνες 30--42).} Ευθυγράμμιση παρτιτούρας--ηχογράφησης, εκτελεστικές μεταβλητές, $\Gamma(z)$ και ιστορική–εκτελεστική σύνθεση στο επίπεδο που επιτρέπει ο κανόνας απόφασης GREEN/AMBER/RED.

\textbf{M5 / M3ν --- Ιστορική εφαρμογή και νευρική επέκταση (μήνες 36--44).} Νικόλαος: οικογένεια $\times$ έργο $\times$ λειτουργία και ενδοοικογενειακή πραγμάτωση. Η M3ν ενεργοποιείται μόνο με επιβεβαιωμένη υποδομή και χρησιμοποιεί τις οριστικοποιημένες μεταβλητές πρόβλεψης της M3c.

\textbf{ΠΕ7 --- Σύνθεση (μήνες 42--48).} Συνθετική ερμηνεία, έλεγχοι ανθεκτικότητας, δημοσιεύσεις, τεκμηρίωση και πακέτο αναπαραγωγιμότητας.

\subsection{Ιεραρχία επιστημονικής τεκμηρίωσης}\label{app:impl-evidence}

Η πρόταση οργανώνει τις προτεραιότητες σε πέντε επίπεδα. Ο \textbf{βασικός υπολογιστικός και μουσικολογικός πυρήνας} περιλαμβάνει M1, M2 με τις αντιθέσεις $\Delta_{\mathrm{hist}}$ και $\Delta_{\mathrm{learn}}$, M3a και το αναγκαίο ανεξάρτητο εκτελεστικό κριτήριο αξιολόγησης της M4. Η \textbf{ανθρώπινη επιβεβαιωτική οικογένεια} περιέχει ακριβώς $H_{\mathrm{cat}}$ και $H_{\mathrm{pred}}$. Η \textbf{M5-min} αποτελεί αυτοτελή ελάχιστη βαθμίδα εφαρμογής και διατηρεί τουλάχιστον τα $D_{\mathrm{fam}}$ και $D_{\mathrm{within}}$ σε συγκρίσιμο σώμα πλήρων έργων της Νέας Μεθόδου, με κατανομή αναφοράς σε επίπεδο μελοποιού μόνο όταν $K\ge K_{\min}$. Οι \textbf{δευτερεύουσες και ενισχυτικές αναλύσεις} περιλαμβάνουν $\Delta_{\mathrm{ctx}}$, CCGP, υπολειμματική ευθυγράμμιση μέσα στην ίδια οικογένεια, αλληλεπιδράσεις εμπειρίας, βιωματική απόκριση, ευρύτερη M4 και την πλήρη M5 με $D_{\mathrm{seq}}$, $D_{\mathrm{func}}$ και πρόσθετες ιστορικές ομάδες αναφοράς, με το ιδιαίτερο καθεστώς καθεμιάς. Οι \textbf{επεκτάσεις υπό προϋποθέσεις} είναι η νευρική M3ν και η διερευνητική ανάλυση δεικτών φωνητικής επαφής μέσα στη M4. Οι προϋποθέσεις ενεργοποίησης και το αν τα αντίστοιχα μεγέθη μπορούν να εκτιμηθούν αξιόπιστα καταγράφονται χωριστά. Η υλοποίηση της βαθμίδας 1 ακολουθεί την κλίμακα που τεκμηριώνει το ΠΕ0: η αξιολόγηση της M2 είναι επιβεβαιωτική σε GREEN/AMBER, διερευνητική σε RED και, χωρίς μετρήσιμο ανεξάρτητο κριτήριο αξιολόγησης, δηλώνεται ότι η $\Delta_{\mathrm{hist}}$ δεν μπορεί να εκτιμηθεί αξιόπιστα· για τις ανθρώπινες υποθέσεις και τη M5 αναφέρεται χωριστά αν υπάρχουν επαρκή δεδομένα για αξιόπιστη εκτίμηση (βλ. §~\ref{sec:core-hierarchy}).

Η σύνθεση καταγράφει τις εξαρτήσεις μεταξύ των διαφορετικών πηγών τεκμηρίωσης. Για κάθε αποτέλεσμα δηλώνονται οι πηγές δεδομένων που είναι κοινές με άλλα εκτιμώμενα μεγέθη. Η συνολική ερμηνεία στηρίζεται στη σύγκλιση ερωτημάτων που βασίζονται σε διαφορετικά δεδομένα και υπόκεινται σε διαφορετικούς μηχανισμούς σφάλματος.

\subsection{Ερευνητική ετοιμότητα και υποδομή}\label{app:impl-readiness}

Η ετοιμότητα παρακολουθείται με μητρώο με εκδόσεις. Για τις M1/M2 καταγράφονται η κάλυψη του σώματος, η προέλευση των πηγών, το πλήθος των οικογενειακών και εξηγητικών ακμών και ο υπολογιστικός προϋπολογισμός. Για M3a καταγράφονται στρατολόγηση ειδικών, ανεξαρτησία των ετικετών αξιολόγησης και αξιοπιστία. Για M3b καταγράφονται κοινό πεδίο τιμών, τράπεζα ερεθισμάτων, αξιοπιστία τριάδων και αξιοπιστία ανθρώπινων RDM. Για M3c καταγράφονται η οριστικοποίηση του λεξιλογίου οικογενειών, κριτήριο επάρκειας PPM, κάλυψη των αντιδιαγώνιων συνδυασμών, μουσική φυσικότητα, δεοντολογική έγκριση και στρατολόγηση συμμετεχόντων. Για M3ν καταγράφονται χωριστά νευρική υποδομή, εξειδικευμένη εμπειρία ανάλυσης και δεοντολογική έγκριση.

Η χρονική αλληλουχία οργανώνει τις εξαρτήσεις των εργασιών: οριστικοποίηση του λεξιλογίου ανάπτυξης πριν από την κύρια επιλογή M1/M2, οριστικοποίηση του επικυρωμένου λεξιλογίου οικογενειών έως τον μήνα 24, τελικό $I_F$ στους μήνες 24--28, πιλοτικοί έλεγχοι στους 28--30 και κύρια ανθρώπινη συλλογή στους 30--40. Η M3c συλλέγεται πριν από τη M3b. Κάθε κρίσιμη ρύθμιση, σπόρος τυχαιοποίησης, διαχωρισμός δεδομένων, λίστα ερεθισμάτων και παραγόμενο μοντέλο αποθηκεύεται με σταθερό αναγνωριστικό.


\begingroup
\setstretch{1.08}
\widowpenalty=10000
\section*{Βιβλιογραφία}
\addcontentsline{toc}{section}{Βιβλιογραφία}
\sloppy

\textbf{Ελληνόγλωσση Βιβλιογραφία}

\begin{enumerate}
\def\labelenumi{\arabic{enumi}.}
\tightlist

\item
 Αναστασίου, Γ. Γ. (2005). \emph{Τα Κρατήματα στην Ψαλτική Τέχνη}. Μελέται 12. Αθήνα: Ίδρυμα Βυζαντινής Μουσικολογίας.
\item
 Αποστολόπουλος, Θωμάς Κ. (2002). \emph{Ο Απόστολος Κώνστας ο Χίος και η συμβολή του στη θεωρία της μουσικής τέχνης: Μουσικολογική θεώρηση από έποψη ιστορική, κωδικογραφική, μελοποιητική και θεωρητική}. Μελέται 4. Αθήνα: Ίδρυμα Βυζαντινής Μουσικολογίας.
\item
 Αλυγιζάκης, Αντώνιος Ε. (1985). \emph{Η οκταηχία στην ελληνική λειτουργική υμνογραφία}. Θεσσαλονίκη: Π.~Σ.~Πουρναράς, 320 σ., ιδίως σσ.~81--82, 149--158, 198--201 και 235. Έντυπη έκδοση 1985· αρχειακή καταχώριση της διδακτορικής διατριβής στο ΕΑΔΔ με έτος 1986, \url{https://doi.org/10.12681/eadd/0578}. Οι ειδικές σελίδες παραμένουν προς αυτοψία του ίδιου του τόμου.
\item
  Αρβανίτης, Ιωάννης. (2010). \emph{Ο ρυθμός των εκκλησιαστικών μελών μέσα από την παλαιογραφική έρευνα και την εξήγηση της παλαιάς σημειογραφίας}. Διδακτορική διατριβή, 2 τόμοι, Τμήμα Μουσικών Σπουδών, Ιόνιο Πανεπιστήμιο. DOI: \url{https://doi.org/10.12681/eadd/21885}.
\item
 

Δεσπότης, Σ. Κ. (2007a). «Ερμηνευτικές προσεγγίσεις στο μουσικό υλικό της ελληνικής ψαλτικής τέχνης». \emph{Γρηγόριος ο Παλαμάς}, 818, 417--428.
\item
 Δεσπότης, Σ. Κ. (2007b). «Η ισοκρατηματική πρακτική της βυζαντινής εκκλησιαστικής μουσικής: ιστορική και μορφολογική προσέγγιση». \emph{Γρηγόριος ο Παλαμάς}, 816 (Μάρτιος-Απρίλιος), 143--175.
\item
 Ιωαννίδης, Χρυσοβαλάντης. (2022). «Η σύνθεση του Νικολάου Πρωτοψάλτου Σμύρνης στο Κοινωνικό ``Εἰς πᾶσαν τὴν γῆν'' σε ήχο πλ. Δ΄ και ο βαθμός εξάρτησής της από την προγενέστερη παράδοση της εκκλησιαστικής μελοποιίας». \emph{Θεολογία}, 93(4), 173-193.
\item
 Μποτωνάκης, Αντώνιος. (2014). «Η Σοφωτάτη και Αρίστη Παραλλαγή Ιωάννου Ιερέως του Πλουσιαδηνού». \emph{Επιστημονική Επετηρίδα Πατριαρχικής Ανώτατης Εκκλησιαστικής Ακαδημίας Κρήτης}, τόμ.~Γ΄, σσ.~337--373, ιδίως σ.~352. Δημοσιευμένη μεταγραφή της \emph{Ἑρμηνείας τῆς Παραλλαγῆς}, κατά τον κώδικα Διονυσίου~570· εκκρεμεί αυτοψία χειρογράφου.
\item
 Κανάκης, Σ. Δ. (2020). \emph{Θεωρητικές και διδακτικές προσεγγίσεις στην ερμηνεία και εκτέλεση ψαλτικών μελών}. Διδακτορική διατριβή, Πανεπιστήμιο Μακεδονίας.
\item
 Καραγκούνης, Κωνσταντῖνος Χαριλ. (2003). \emph{Η Παράδοση και Εξήγηση του Μέλους των Χερουβικών της Βυζαντινής και μετα-Βυζαντινής Μελοποιίας}. Μελέται 7. Αθήνα: Ίδρυμα Βυζαντινής Μουσικολογίας.
\item
 Καραγκούνης, Κωνσταντῖνος Χαριλ. (2005). \emph{Παραλειπόμενα περὶ τοῦ Χερουβικοῦ Ὕμνου}. Βόλος: Γραφικές Τέχνες Αλέξανδρος Ξουράφας.
\item
 Πολυκαρπίδης, Π. Γ. (2023). \emph{Υπολογιστική προσέγγιση της Βυζαντινής Μουσικής: Η περίπτωση του Ειρμολογίου του Θεοφάνη Καρύκη}. Διδακτορική διατριβή, Εθνικό και Καποδιστριακό Πανεπιστήμιο Αθηνών.
\item
 Σπανουδάκης, Δημοσθένης Π. (2017). \emph{Σύγχρονα μοντέλα ανάλυσης Βυζαντινής Μουσικής: Το φαινόμενο της μελισματικότητας στα κεκραγάρια}. Διδακτορική διατριβή, 3 τόμοι, Τμήμα Μουσικών Σπουδών, Αριστοτέλειο Πανεπιστήμιο Θεσσαλονίκης. \url{https://doi.org/10.12681/eadd/42193}.
\item
Φωτόπουλος, Κωνσταντίνος Ι. (2009). «Η μελοποιία ως αναπόσπαστο μέρος της βυζαντινής μουσικής και η προβληματική της στη σύγχρονη εποχή». Κείμενο ανακοίνωσης για το συνέδριο \emph{Current Problems in the Study of the Art of Chanting: Research and Practice}, Moscow P. I. Tchaikovsky Conservatoire, 12--16 Μαΐου 2009. Το διαθέσιμο αντίτυπο δεν τεκμηριώνει αυτοτελή έντυπη έκδοση.
\item
 Στάθης, Γρηγόριος Θ. (1982). «Συμπόσιον περί Βυζαντινής Μουσικής και ``Δειναί θέσεις'' και ``Εξήγησις''». \emph{Θεολογία}, 43(3), 749--782.
\item
 Στάθης, Γρηγόριος Θ. (1987). «Η εξήγηση της ψαλτικής τέχνης». \emph{Θεολογία}, 58(2), 337--371.
\item
 Στάθης, Γρηγόριος Θ. (1978). \emph{Η εξήγησις της Παλαιάς Βυζαντινής Σημειογραφίας}. Μελέται 2. Αθήνα: Ίδρυμα Βυζαντινής Μουσικολογίας.
\item
 Στάθης, Γρηγόριος Θ. (1993). \emph{Η εξήγησις της Παλαιάς Βυζαντινής Σημειογραφίας}. Μελέται 2, γ΄ έκδοση. Αθήνα: Ίδρυμα Βυζαντινής Μουσικολογίας.
\item
 Τόμα, Ι. Γ. (2018). \emph{Νικόλαος Γεωργίου, πρωτοψάλτης Σμύρνης. Βίος, ψαλτικό έργο, μουσικολογικές προσεγγίσεις}. Μεταπτυχιακή διπλωματική εργασία, Αριστοτέλειο Πανεπιστήμιο Θεσσαλονίκης.
\item
 Χαραλαμπίδης, Αρσένιος. (2024). \emph{Τα Κοινωνικά του Μπαλασίου Ιερέως σε Α΄ Ήχο: Μορφολογική Ανάλυση}. Πτυχιακή εργασία, Τμήμα Μουσικής Επιστήμης και Τέχνης, Πανεπιστήμιο Μακεδονίας, Θεσσαλονίκη.
\item
 Χαραλαμπίδης, Αρσένιος. (2026). \emph{Τα Κοινωνικά της Κυριακής σε Α΄ Ήχο: Μορφολογική και Τροπική Ανάλυση}. Μεταπτυχιακή διπλωματική εργασία, ΠΜΣ «Επιστήμες και Τέχνες της Μουσικής», Πανεπιστήμιο Μακεδονίας, Θεσσαλονίκη (επιβλέπουσα: Ε. Σπυράκου).
\item
 Χατζάρα, Α. (2015). \emph{Η μουσική ερμηνεία στην προφορική παράδοση της ελληνικής εκκλησιαστικής μουσικής}. Διδακτορική διατριβή, Τμήμα Μουσικών Σπουδών, Ιόνιο Πανεπιστήμιο, Κέρκυρα. \url{https://doi.org/10.12681/eadd/37215}.
\item
 Χαλδαιάκης, Αχιλλεύς Γ. (2014). \emph{Βυζαντινομουσικολογικά}, τ.~Γ΄: \emph{Μελοποιία}. Αθήνα: Εκδόσεις Άθως. Περιλαμβάνει τη μελέτη «Ο τροχός της οκταηχίας»· η ακριβής διατύπωση που παρατίθεται από τις σσ.~156--157 εκκρεμεί προς άμεση αυτοψία. \url{https://www.achilleaschaldaeakes.gr/}.
\item
 Χρύσανθος ὁ ἐκ Μαδύτων. (1832). \emph{Θεωρητικὸν Μέγα τῆς Μουσικῆς}. Τεργέστη.
\end{enumerate}

\textbf{Ξενόγλωσση Βιβλιογραφία}

\begin{enumerate}
\def\labelenumi{\arabic{enumi}.}

\tightlist

\item
 Alexandru, M., \& Troelsgård, C. (2020). ``The `Elements of the Papadikê' and Modality Features in Byzantine Chant.'' \emph{Series Musicologica Balcanica}, 1(1), 169-204.
\item
 Kanakis, Symeon. (2021). ``Addressing the Modality of Church Melodies, According to the Variety in the Choice of the Isokratema. Case Studies. \emph{Series Musicologica Balcanica}, 1(2), 233--261, ιδίως σ.~255. \url{https://doi.org/10.26262/smb.v1i2.7948}.
\item
 Alexandru, M., De Luca, E., \& Fujinaga, I. (2024). ``Brainstorming on the Encoding of Eastern Music Scripts: Middle-Byzantine Neumes as a Case Study.'' \emph{Series Musicologica Balcanica}, 5, 177-205.
\item
 Amargianakis, G. (1977). ``An Analysis of Stichera in the Deuteros Modes: The Stichera Idiomela for the Month of September in the Modes Deuteros, Plagal Deuteros, and Nenano.'' \emph{Cahiers de l'Institut du Moyen-Âge Grec et Latin}, 22-23, 1-263.
\item
 Anglada-Tort, M., Harrison, P. M. C., Lee, H., \& Jacoby, N. (2023). ``Large-Scale Iterated Singing Experiments Reveal Oral Transmission Mechanisms Underlying Music Evolution.'' \emph{Current Biology}, 33(8), 1472--1486.e12. \url{https://doi.org/10.1016/j.cub.2023.02.070}.
\item
 Antonopoulos, Spyridon. (2014 [submitted 2013]). \emph{The Life and Works of Manuel Chrysaphes the Lampadarios, and the Figure of Composer in Late Byzantium}. Doctoral dissertation, City University London, vol.~I, ιδίως σσ.~168--181 και 284--286. \url{https://openaccess.city.ac.uk/id/eprint/17439/}.
\item
 Anton, T. L., \& Trăușan-Matu, Ș. (2018). ``Byzantine Music Composition Using Markov Models.'' In \emph{Proceedings of RoCHI 2018}, 71-75.
\item
 Artstein, R., \& Poesio, M. (2008). ``Inter-Coder Agreement for Computational Linguistics.'' \emph{Computational Linguistics}, 34(4), 555-596.
\item
 Arvanitis, I. V. (2022). ``The Transcriptions of the Three Teachers: An Old or a New Sound?'' \emph{Epistēmēs Metron Logos}, 5, 143--149. \url{https://doi.org/10.12681/eml.28005}.
\item
 Ballen, G. A., Mühlová, K. H., \& Hajič Jr., J. (2026). ``What Did the Dove Sing to Pope Gregory? Ancestral Melody Reconstruction in Gregorian Chant Using Bayesian Phylogenetics.'' \emph{PLOS ONE}, 21(5), e0350014. \url{https://doi.org/10.1371/journal.pone.0350014}.
\item
 Bernardi, S., Benna, M. K., Rigotti, M., Munuera, J., Fusi, S., \& Salzman, C. D. (2020). ``The Geometry of Abstraction in Hippocampus and Prefrontal Cortex.'' \emph{Cell}, 183(4), 954--967.e21. \url{https://doi.org/10.1016/j.cell.2020.09.031}.
\item
 Bouris, N. G., Georgaki, A., \& Chaldaeakes, A. G. (2018). ``Encoding Byzantine Music Notation (ca. 1600-1814).'' \emph{Proceedings of SMC 2018}.
\item
 Cambouropoulos, E. (2001). ``The Local Boundary Detection Model (LBDM) and Its Application in the Study of Expressive Timing.'' \emph{Proceedings of ICMC}.
\item
 Byros, V. (2012). ``Meyer's Anvil: Revisiting the Schema Concept.'' \emph{Music Analysis}, 31(3), 273--346. \url{https://doi.org/10.1111/j.1468-2249.2012.00344.x}.
\item
 Cambouropoulos, E. (2006). ``Musical Parallelism and Melodic Segmentation: A Computational Approach.'' \emph{Music Perception}, 23(3), 249--268. \url{https://doi.org/10.1525/mp.2006.23.3.249}.
\item
 Cambouropoulos, E. (2009). ``How Similar Is Similar?'' \emph{Musicae Scientiae}, Discussion Forum 4B, 7--24. \url{https://jyx.jyu.fi/handle/123456789/22619}.
\item
 Chaldaeakes, A. G. (2010). ``Daniel the Protopsaltes (+1789): His Life and Work. A Preliminary Paper.'' \emph{Muzica}, 3, 39--53.
\item
 Chryssochoidis, G., Delviniotis, D., \& Kouroupetroglou, G. (2007). ``A Semi-Automated Tagging Methodology for Orthodox Ecclesiastic Chant Acoustic Corpora.'' \emph{Proceedings of the 4th Sound and Music Computing Conference (SMC 2007)}, Lefkada, Greece, 126--133.
\item
 Conklin, D., \& Anagnostopoulou, C. (2006). ``Segmental Pattern Discovery in Music.'' \emph{INFORMS Journal on Computing}, 18(3), 285-293.
\item
 Conomos, D. E. (1974). \emph{Byzantine Trisagia and Cheroubika of the Fourteenth and Fifteenth Centuries: A Study of Late Byzantine Liturgical Chant}. Thessaloniki: Patriarchal Institute for Patristic Studies.
\item
 Conomos, D. E. (1985a). \emph{The Late Byzantine and Slavonic Communion Cycle: Liturgy and Music}. Dumbarton Oaks Studies 21. Washington, DC: Dumbarton Oaks Research Library and Collection.
\item
 Conomos, D. E. (Ed.). (1985b). \emph{The Treatise of Manuel Chrysaphes, the Lampadarios: On the Theory of the Art of Chanting and on Certain Erroneous Views That Some Hold About It}. Monumenta Musicae Byzantinae, Corpus Scriptorum de Re Musica, 2. Vienna: Verlag der Österreichischen Akademie der Wissenschaften.
\item
 Cornelissen, B., Zuidema, W., \& Burgoyne, J. A. (2020a). ``Mode Classification and Natural Units in Plainchant.'' \emph{Proceedings of ISMIR 2020}.
\item
 Cornelissen, B., Zuidema, W., \& Burgoyne, J. A. (2020b). ``Studying Large Plainchant Corpora Using chant21.'' \emph{Proceedings of DLfM 2020}.
\item
 Dalitz, C., Michalakis, G. K., \& Pranzas, C. (2008). ``Optical Recognition of Psaltic Byzantine Chant Notation.'' \emph{International Journal on Document Analysis and Recognition}, 11, 143-158.
\item
 Davani, A. M., Díaz, M., \& Prabhakaran, V. (2022). ``Dealing with Disagreements: Looking Beyond the Majority Vote in Subjective Annotations.'' \emph{Transactions of the Association for Computational Linguistics}, 10, 92--110. \url{https://doi.org/10.1162/tacl_a_00449}.
\item
 Delviniotis, D., Kouroupetroglou, G., \& Theodoridis, S. (2008). ``Acoustic Analysis of Musical Intervals in Modern Byzantine Chant Scales.'' \emph{Journal of the Acoustical Society of America}, 124(4), EL262-EL269.
\item
 Delviniotis, D., \& Kouroupetroglou, G. (2011). ``DAMASKINOS: The Prototype Corpus of Greek Orthodox Ecclesiastical Chant Voices.'' \emph{Proceedings of Crossroads: Greece as an Intercultural Pole of Musical Thought and Creativity}, Thessaloniki, Greece.
\item
 Delviniotis, D. (2021). ``New Method of Byzantine Music Intervals' Measuring and Its Application in the Fourth Mode.'' \emph{Series Musicologica Balcanica}, 1(2), 209-232.
\item
 Di Liberto, G. M., Pelofi, C., Bianco, R., Patel, P., Mehta, A. D., Herrero, J. L., de Cheveign\'{e}, A., Shamma, S., \& Mesgarani, N. (2020). ``Cortical Encoding of Melodic Expectations in Human Temporal Cortex.'' \emph{eLife}, 9, e51784. \url{https://doi.org/10.7554/eLife.51784}.
\item
 Dvořáková, A., Eipert, T., Lacoste, D., \& Hajič Jr., J. (2026). ``Making Chant Computing Easy: CantusCorpus v1.0 and the PyCantus Library.'' arXiv:2603.11933. \url{https://arxiv.org/abs/2603.11933}.
\item
 Effland, T., \& Collins, M. (2021). ``Partially Supervised Named Entity Recognition via the Expected Entity Ratio Loss.'' \emph{Transactions of the Association for Computational Linguistics}, 9, 1320--1335. \url{https://doi.org/10.1162/tacl_a_00429}.
\item
 Finkensiep, C., Déguernel, K., Neuwirth, M., \& Rohrmeier, M. (2020). ``Voice-Leading Schema Recognition Using Rhythm and Pitch Features.'' \emph{Proceedings of the 21st International Society for Music Information Retrieval Conference (ISMIR 2020)}, 520--526. \url{https://doi.org/10.5281/zenodo.4245482}.
\item
 Fradet, N., Gutowski, N., Chhel, F., \& Briot, J.-P. (2023). ``Byte Pair Encoding for Symbolic Music.'' \emph{Proceedings of EMNLP 2023}.
\item
 Gezerlis, V. G., \& Theodoridis, S. (2002). ``Optical Character Recognition of the Orthodox Hellenic Byzantine Music Notation.'' \emph{Pattern Recognition}, 35(4), 895-914.
\item
 Gheorghiță, N. (2010). ``Observations on the Technique of Transcription (Exegesis) into the New Method of Analytical Music Notation of the Sunday Koinonikon of the 18th Century. \emph{Studii şi Cercetări de Istoria Artei. Teatru, Muzică, Cinematografie}, s.n., 4(48), 9--26. The same study is reprinted as chapter XIV of the author's 2010 volume, pp.~225--248.
\item
 Gheorghiță, N. (2010). \emph{Byzantine Chant between Constantinople and the Danubian Principalities: Studies in Byzantine Musicology}. Bucharest: Editura Sophia. Ιδίως κεφ. XI, σσ.~171--190· κεφ. XIII, σσ.~201--224· κεφ. XIV, σσ.~225--248.
\item
 Gingras, B., Pearce, M. T., Goodchild, M., Dean, R. T., Wiggins, G. A., \& McAdams, S. (2016). ``Linking Melodic Expectation to Expressive Performance Timing and Perceived Musical Tension.'' \emph{Journal of Experimental Psychology: Human Perception and Performance}, 42(4), 594--609. \url{https://doi.org/10.1037/xhp0000141}.
\item
 Gjerdingen, R. O. (2007). \emph{Music in the Galant Style}. Oxford: Oxford University Press. \url{https://doi.org/10.1093/oso/9780195313710.001.0001}.
\item
 Grigorjeva, V. (2023). ``Principles of Stratification of Melodic Formulas of Old Russian Out-of-Ichos Chants (with `Opekalovo' `Sepulchral' Chants).'' \emph{St. Tikhons' University Review. Series V: Christian Art}, 50, 9-23. \url{https://doi.org/10.15382/sturv202350.9-23}.
\item
 Gusev, V. D., Bakhmutova, I. V., Miroshnichenko, L. A., \& Titkova, T. N. (2020). ``Possible Approaches to Deciphering Russian Ancient Znamenny Chant.'' \emph{Journal of New Music Research}, 49(3), 298--306. \url{https://doi.org/10.1080/09298215.2020.1762666}.
\item
 Hajič Jr., J., Ballen, G. A., Mühlová, K. H., \& Vlhová-Wörner, H. (2023). ``Towards Building a Phylogeny of Gregorian Chant Melodies.'' \emph{Proceedings of the 24th International Society for Music Information Retrieval Conference (ISMIR 2023)}, 571-578.
\item
 Hajič Jr., J., Lanz, V., \& Ballen, G. A. (2025). ``Genome of Melody: Applying Bioinformatics to Study the Evolution of Gregorian Chant.'' \emph{Philosophical Transactions of the Royal Society B: Biological Sciences}, 380(1940), 20240274.
\item
 Hannick, C., \& Wolfram, G. (Eds.). (1985). \emph{Gabriel Hieromonachos: Abhandlung über den Kirchengesang}. Monumenta Musicae Byzantinae, Corpus Scriptorum de Re Musica, 1. Wien: Verlag der Österreichischen Akademie der Wissenschaften.
\item
 Hannick, C. (Ed.). (1991). \emph{Rhythm in Byzantine Chant: Acta of the Congress Held at Hernen Castle in November 1986}. Hernen, Holland: A.~A. Bredius Foundation.
\item
 Arvanitis, I. (2003). ``The Rhythmical and Metrical Structure of the Byzantine Heirmoi and Stichera as a Means to and as a Result of a New Rhythmical Interpretation of Byzantine Chant.'' \emph{Acta Musicae Byzantinae}, 6, 14--29.
\item
 Hansen, N. C., Kragness, H. E., Vuust, P., Trainor, L. J., \& Pearce, M. T. (2021). ``Predictive Uncertainty Underlies Auditory Boundary Perception.'' \emph{Psychological Science}, 32(9), 1416--1425. \url{https://doi.org/10.1177/0956797621997349}.
\item
 Hariharan, B., \& Girshick, R. (2017). ``Low-Shot Visual Recognition by Shrinking and Hallucinating Features.'' \emph{Proceedings of the IEEE International Conference on Computer Vision (ICCV)}. \url{https://openaccess.thecvf.com/content_iccv_2017/html/Hariharan_Low-Shot_Visual_Recognition_ICCV_2017_paper.html}.
\item
 Herbst, C. T. (2020). ``Electroglottography -- An Update.'' \emph{Journal of Voice}, 34(4), 503--526. \url{https://doi.org/10.1016/j.jvoice.2018.12.014}.
\item
 Herff, S. A., Harasim, D., Cecchetti, G., Finkensiep, C., \& Rohrmeier, M. A. (2021). ``Hierarchical Syntactic Structure Predicts Listeners' Sequence Completion in Music.'' \emph{Proceedings of the Annual Meeting of the Cognitive Science Society}, 43.
\item
 Holzapfel, A. (2010). \emph{Similarity Methods for Computational Ethnomusicology}. Doctoral dissertation, Department of Computer Science, University of Crete. \url{https://doi.org/10.12681/eadd/30917}.
\item


Janssen, B., Burgoyne, J. A., \& Honing, H. (2017). ``Predicting Variation of Folk Songs: A Corpus Analysis Study on the Memorability of Melodies.'' \emph{Frontiers in Psychology}, 8, 621. \url{https://doi.org/10.3389/fpsyg.2017.00621}.
\item
 Janssen, B., van Kranenburg, P., \& Volk, A. (2017). ``Finding Occurrences of Melodic Segments in Folk Songs Employing Symbolic Similarity Measures.'' \emph{Journal of New Music Research}, 46(2), 118--134. \url{https://doi.org/10.1080/09298215.2017.1316292}.
\item
 Kaliakatsos-Papakostas, M., Makris, D., Soiledis, K., Tsamis, K.-T., Katsouros, V., \& Cambouropoulos, E. (2025). ``HarmonyTok: Comparing Methods for Harmony Tokenization for Machine Learning.'' \emph{Information}, 16(9), 759. \url{https://doi.org/10.3390/info16090759}.
\item
 Kanakis, S. (2021). ``Addressing the Modality of Church Melodies, According to the Variety in the Choice of the Isokratema: Case Studies.'' \emph{Series Musicologica Balcanica}, 1(2), 233-261.
\item
 Karas, S. I. (1982). \emph{Μέθοδος της Ελληνικής Μουσικής: Θεωρητικόν}, Τόμοι Α'--Β'. Αθήνα: Σύλλογος προς Διάδοσιν της Εθνικής Μουσικής.
\item
 Karsdorp, F., van Kranenburg, P., \& Manjavacas, E. (2019). ``Learning Similarity Metrics for Melody Retrieval.'' \emph{Proceedings of the 20th International Society for Music Information Retrieval Conference (ISMIR 2019)}, 478--485.
\item
 Katsiavalos, A., Collins, T., \& Battey, B. (2019). ``An Initial Computational Model for Musical Schemata Theory.'' \emph{Proceedings of the 20th International Society for Music Information Retrieval Conference (ISMIR 2019)}, 166--172.
\item
 Kern, P., Heilbron, M., de Lange, F. P., \& Spaak, E. (2022). ``Cortical Activity During Naturalistic Music Listening Reflects Short-Range Predictions Based on Long-Term Experience.'' \emph{eLife}, 11, e80935. \url{https://doi.org/10.7554/eLife.80935}.
\item
 Koelsch, S., Rohrmeier, M., Torrecuso, R., \& Jentschke, S. (2013). ``Processing of Hierarchical Syntactic Structure in Music.'' \emph{Proceedings of the National Academy of Sciences}, 110(38), 15443--15448. \url{https://doi.org/10.1073/pnas.1300272110}.
\item
 Kokkinidis, K., Panagi, A., \& Manitsaris, A. (2016). ``Finding the Optimum Training Solution for Byzantine Music Recognition-A Max/Msp Approach.'' \emph{2016 5th International Conference on Modern Circuits and Systems Technologies (MOCAST)}. IEEE.
\item
 Kong, L., Dyer, C., \& Smith, N. A. (2016). ``Segmental Recurrent Neural Networks.'' \emph{4th International Conference on Learning Representations (ICLR 2016)}.
\item
 Lanz, V., \& Hajič Jr., J. (2023). ``Text Boundaries Do Not Provide a Better Segmentation of Gregorian Antiphons.'' \emph{Proceedings of DLfM 2023}.
\item
 Lanz, V., \& Hajič Jr., J. (2025). ``Gregorian Melody, Modality, and Memory: Segmenting Chant with Bayesian Nonparametrics.'' \emph{Proceedings of the 26th International Society for Music Information Retrieval Conference (ISMIR 2025)}, 638--646.
\item
 Lind, T. T. (2011). \emph{The Past Is Always Present: The Revival of the Byzantine Musical Tradition at Mount Athos}. Lanham, MD: Scarecrow Press. \url{https://doi.org/10.5771/9780810881488}.
\item
 Lingas, A. (2003). ``Performance Practice and the Politics of Transcribing Byzantine Chant.'' \emph{Acta Musicae Byzantinae}, 6, 56-76.
\item
 Mahajan, D., Tople, S., \& Sharma, A. (2021). ``Domain Generalization Using Causal Matching.'' \emph{Proceedings of the 38th International Conference on Machine Learning (ICML)}, PMLR 139. \url{https://proceedings.mlr.press/v139/mahajan21b.html}.
\item
 Mavromatis, P. (2005a). ``A Hidden Markov Model of Melody Production in Greek Church Chant.'' \emph{Computing in Musicology}, 14, 93-112.
\item
 Mavromatis, P. (2005b). \emph{The Echoi of Modern Greek Church Chant in Written and Oral Transmission: A Computational Model and Its Cognitive Implications}. Doctoral dissertation, Eastman School of Music, University of Rochester.
\item
 Mavromatis, P. (2009a). ``Minimum Description Length Modelling of Musical Structure.'' \emph{Journal of Mathematics and Music}, 3(3), 117-136.
\item
 Mavromatis, P. (2009b). ``HMM Analysis of Musical Structure: Identification of Latent Variables Through Topology-Sensitive Model Selection.'' \emph{Mathematics and Computation in Music 2009}.
\item
 Mavromatis, P. (2009c). ``A Multi-tiered Approach for Analyzing Expressive Timing in Music Performance.'' In T. Klouche \& T. Noll (Eds.), \emph{Mathematics and Computation in Music: MCM 2009}, 193--204. Berlin/Heidelberg: Springer. \url{https://doi.org/10.1007/978-3-642-02394-1_18}.
\item
 Mavromatis, P. (2019). ``Hierarchical Structural Patterns in Improvised Music: Implications for Cognition.'' Paper presented at the 42nd Annual Meeting of the Society for Music Theory, Columbus, Ohio, 7 November 2019. Published abstract in the official conference program, pp.~56--57.
\item
 Mongeau, M., \& Sankoff, D. (1990). ``Comparison of Musical Sequences.'' \emph{Computers and the Humanities}, 24(3), 161--175. \url{https://doi.org/10.1007/BF00117340}.
\item
 Neanes Project (2026). ``Byzantine Chant OCR.'' Open-source OCR toolset for Byzantine chant notation, GPL-3.0. \url{https://github.com/neanes/byzantine-chant-ocr}.
\item
 Neanes Project (2026). ``Neanes --- A Byzantine Chant Scorewriter.'' Open-source software, GPL-3.0. \url{https://github.com/neanes/neanes}.
\item
 Neuwirth, M., Finkensiep, C., \& Rohrmeier, M. (2023). ``Musical Schemata: Modelling Challenges and Pattern Finding (BachBeatles).'' In B. Schneider, B. Löffler, T. Mager, \& C. Hein (Eds.), \emph{Mixing Methods: Practical Insights from the Humanities in the Digital Age}, 147--164. Bielefeld: Bielefeld University Press. \url{https://doi.org/10.1515/9783839469132-016}.
\item
 Nuttall, T., Casado, M. G., Ferraro, A., Conklin, D., \& Caro Repetto, R. (2021). ``A Computational Exploration of Melodic Patterns in Arab-Andalusian Music.'' \emph{Journal of Mathematics and Music}, 15(2), 168-180. \url{https://doi.org/10.1080/17459737.2021.1917010}.
\item
 Nuttall, T., Plaja-Roglans, G., Pearson, L., \& Serra, X. (2022). ``In Search of Sañcāras: Tradition-Informed Repeated Melodic Pattern Recognition in Carnatic Music.'' \emph{Proceedings of the 23rd International Society for Music Information Retrieval Conference (ISMIR 2022)}, 337-344.
\item
 Park, D., Nam, J., \& Park, J. (2020). ``Novelty and Influence of Creative Works, and Quantifying Patterns of Advances Based on Probabilistic Reference Networks.'' \emph{EPJ Data Science}, 9, 2. \url{https://doi.org/10.1140/epjds/s13688-019-0214-8}.
\item
 Panteli, M., \& Purwins, H. (2013). ``A Quantitative Comparison of Chrysanthine Theory and Performance Practice of Scale Tuning, Steps, and Prominence of the Octoechos in Byzantine Chant.'' \emph{Journal of New Music Research}, 42(3), 205-221.
\item
 Papaioannou, C., Benetos, E., \& Potamianos, A. (2025). ``Universal Music Representations? Evaluating Foundation Models on World Music Corpora.'' \emph{Proceedings of the 26th International Society for Music Information Retrieval Conference (ISMIR 2025)}. \url{https://arxiv.org/abs/2506.17055}.
\item
 Papay, S., Klinger, R., \& Padó, S. (2022). ``Constraining Linear-chain CRFs to Regular Languages.'' \emph{International Conference on Learning Representations (ICLR)}. \url{https://arxiv.org/abs/2106.07306}.
\item
 Paschalidou, S. (2024). ``Effort in Hindustani Dhrupad Singing: Using Interviews to Develop Effort Inference Models of Acoustic and Movement Features.'' \emph{Yearbook for Traditional Music}, 56(1), 57--79. \url{https://doi.org/10.1017/ytm.2024.7}.
\item
 Pearce, M. T. (2018). ``Statistical Learning and Probabilistic Prediction in Music Cognition: Mechanisms of Stylistic Enculturation.'' \emph{Annals of the New York Academy of Sciences}, 1423(1), 378--395. \url{https://doi.org/10.1111/nyas.13654}.
\item
 Pearce, M. T., \& Wiggins, G. A. (2001). ``Towards a Framework for the Evaluation of Machine Compositions.'' In \emph{Proceedings of the AISB'01 Symposium on Artificial Intelligence and Creativity in the Arts and Sciences}, 22--32.
\item
 Pearce, M. T., Müllensiefen, D., \& Wiggins, G. A. (2010). ``The Role of Expectation and Probabilistic Learning in Auditory Boundary Perception: A Model Comparison.'' \emph{Perception}, 39(10), 1365--1389. \url{https://doi.org/10.1068/p6507}.
\item
 Plank, B. (2022). ``The `Problem' of Human Label Variation: On Ground Truth in Data, Modeling and Evaluation.'' In \emph{Proceedings of EMNLP 2022}, 10671-10682.
\item
 Polykarpidis, P., Kalofonos, D., Balageorgos, D., \& Anagnostopoulou, C. (2022). ``Three Related Corpora in Middle Byzantine Music Notation and a Preliminary Comparative Analysis.'' \emph{Proceedings of ISMIR 2022}, 306-313.
\item
 Popescu, T., \& Rohrmeier, M. (2024). ``Core Principles of Melodic Organisation Emerge from Transmission Chains with Random Melodies.'' \emph{Evolution and Human Behavior}, 45(6), 106619. \url{https://doi.org/10.1016/j.evolhumbehav.2024.106619}.
\item
 Rabinovitch, G., \& Carter-Enyi, A. (2024). ``Melodic Organization and Sequential Ordering of Galant Schemata: Implications for Eighteenth-Century European Musicianship.'' \emph{Music Theory Online}, 30(1). \url{https://doi.org/10.30535/mto.30.1.7}.
\item
 Raasted, J. (1966). \emph{Intonation Formulas and Modal Signatures in Byzantine Musical Manuscripts}. Monumenta Musicae Byzantinae, Subsidia 7. Copenhagen: Ejnar Munksgaard.
\item
 Ristad, E. S., \& Yianilos, P. N. (1998). ``Learning String-Edit Distance.'' \emph{IEEE Transactions on Pattern Analysis and Machine Intelligence}, 20(5), 522--532. \url{https://doi.org/10.1109/34.682181}.
\item
 Ryabinets, K. (2025). ``Melodic Formula / Constellation as a Structural Unit of the Sacred Monody on the Material of Ukrainian Notoline Irmologions of the XVI--XVIII Centuries.'' \emph{National Academy of Managerial Staff of Culture and Arts Herald}. \url{https://doi.org/10.32461/2226-3209.1.2025.328001}.
\item
 Sarawagi, S., \& Cohen, W. W. (2004). ``Semi-Markov Conditional Random Fields for Information Extraction.'' \emph{Advances in Neural Information Processing Systems 17}, 1185--1192.
\item
 Savage, P. E., Passmore, S., Chiba, G., Currie, T. E., Suzuki, H., \& Atkinson, Q. D. (2022). ``Sequence Alignment of Folk Song Melodies Reveals Cross-Cultural Regularities of Musical Evolution.'' \emph{Current Biology}, 32(6), 1395--1402.e8. \url{https://doi.org/10.1016/j.cub.2022.01.039}.
\item
 Schindler, A., Herdener, M., \& Bartels, A. (2013). ``Coding of Melodic Gestalt in Human Auditory Cortex.'' \emph{Cerebral Cortex}, 23(12), 2987--2993. \url{https://doi.org/10.1093/cercor/bhs289}.
\item
 Schiødt, N. (1971). ``A Computer-Aided Analysis of Thirty-Five Byzantine Hymns.'' In E. Wellesz \& M. Velimirović (Eds.), \emph{Studies in Eastern Chant}, Vol. 2, 129-154.
\item
 Serra, X. (2017). ``The Computational Study of a Musical Culture Through Its Digital Traces.'' \emph{Acta Musicologica}, 89(1), 24--44.
\item
 Smith, J. B. L., Burgoyne, J. A., Fujinaga, I., De Roure, D., \& Downie, J. S. (2011). ``Design and Creation of a Large-Scale Database of Structural Annotations.'' In \emph{Proceedings of ISMIR 2011}, 555-560.
\item
 St.~Anthony's Greek Orthodox Monastery. (n.d.-a). \emph{Byzantine Music Formulae}. Διαδικτυακή ειρμολογική/στιχηραρική συλλογή με περισσότερες από 10.000 δηλωμένες φόρμουλες, με αρχεία κατά ήχο και είδος, πρόσβαση Σεπτέμβριος 2026. \url{https://stanthonysmonastery.org/pages/byzantine-music-formulae}.
\item
 St.~Anthony's Greek Orthodox Monastery. (n.d.-b). \emph{Papadic Formulae}. Χωριστή διαδικτυακή συλλογή παπαδικών μελωδικών τύπων από 758 συνθέσεις, με μεταδεδομένα των αρχείων Word προέλευσης ανά φόρμουλα (τίτλος και συνθέτης), πρόσβαση Σεπτέμβριος 2026. \url{https://stanthonysmonastery.org/pages/papadic-formulae}.
\item
 Tofan, M. V. (2024). ``Stability and Flexibility in the Transmission of Melodic Archetypes in Greek Liturgical Chant.'' \emph{Bulletin of the Transilvania University of Braşov, Series VIII: Performing Arts}, 17(66), 193-202.
\item
 Troelsgård, C. (2011). \emph{Byzantine Neumes: A New Introduction to the Middle Byzantine Musical Notation}. MMB Subsidia 9.
\item
 Tsuboi, Y., Kashima, H., Mori, S., Oda, H., \& Matsumoto, Y. (2008). ``Training Conditional Random Fields Using Incomplete Annotations.'' \emph{Proceedings of the 22nd International Conference on Computational Linguistics (COLING 2008)}, 897--904. \url{https://aclanthology.org/C08-1113/}.
\item
 Tzanetakis, G. (2014). ``Computational Ethnomusicology: A Music Information Retrieval Perspective.'' In A. Georgaki \& G. Kouroupetroglou (Eds.), \emph{Proceedings of ICMC|SMC|2014}, Athens, Greece.
\item
 Uma, A., Fornaciari, T., Hovy, D., Paun, S., Plank, B., \& Poesio, M. (2021). ``Learning from Disagreement: A Survey.'' \emph{Journal of Artificial Intelligence Research}, 72, 1385-1470.
\item
 Urbano, J., Schedl, M., \& Serra, X. (2013). ``Evaluation in Music Information Retrieval.'' \emph{Journal of Intelligent Information Systems}, 41(3), 345-369.
\item
 Volk, A., \& van Kranenburg, P. (2012). ``Melodic Similarity among Folk Songs: An Annotation Study on Similarity-Based Categorization in Music.'' \emph{Musicae Scientiae}, 16(3), 317--339. \url{https://doi.org/10.1177/1029864912448329}.
\item
 Yan, Y., Cwitkowitz, F., \& Duan, Z. (2021). ``Skipping the Frame-Level: Event-Based Piano Transcription With Neural Semi-CRFs.'' \emph{Advances in Neural Information Processing Systems}, 34, 20583--20595.
\item
 Zaratiana, U., Tomeh, N., El Khbir, N., Holat, P., \& Charnois, T. (2023). ``Filtered Semi-Markov CRF.'' \emph{Findings of the Association for Computational Linguistics: EMNLP 2023}, 222--235. \url{https://doi.org/10.18653/v1/2023.findings-emnlp.17}.
\end{enumerate}


\endgroup
\end{document}
